\documentclass[11pt,a4paper]{book}

% Document setup and language
\usepackage[english]{babel}
\usepackage[utf8]{inputenc}
\usepackage[T1]{fontenc}
\usepackage{subfiles}
\setsecnumdepth{subsection}

% Mathematics
\usepackage{amsfonts,amsmath,amssymb,amsthm}
\usepackage{dsfont}
\usepackage{mathtools}

% Typography
\usepackage{newtxtext,newtxmath}
\usepackage{csquotes}
\usepackage{microtype}
\setlength{\parskip}{\smallskipamount}
\setlength{\parindent}{0pt}

% Diagrams and tables
\usepackage{tikz-cd,tikz}
\usepackage{graphicx}
\usepackage{booktabs}

% Lists
\usepackage{enumitem}
\setlist[enumerate,1]{label=(\arabic*),leftmargin=*}
\setlist[itemize]{leftmargin=*}

% Links and references
\usepackage[hyphens]{url}
\urlstyle{same}
\usepackage{bookmark}
\usepackage{hyperref}
\hypersetup{
  colorlinks=true,
  linkcolor=black,
  filecolor=black,
  urlcolor=black,
  citecolor=black,
  bookmarksdepth=subsubsection,
  pdftitle={Derived manifolds: MIGOS seminar notes}
}

% Bibliography
\usepackage[
  backend=biber,
  natbib,
  style=alphabetic,
  maxcitenames=4,
  maxalphanames=5,
  maxbibnames=99
]{biblatex}
\renewbibmacro{in:}{}
\addbibresource{Bibliography.bib}

% Basic notation. Add project-specific notation only when it is first used.
\newcommand{\Cinfty}{C^\infty}
\newcommand{\R}{\mathbb{R}}
\newcommand{\N}{\mathbb{N}}
\newcommand{\Z}{\mathbb{Z}}
\DeclareMathOperator{\An}{An}
\newcommand{\Cat}{\mathrm{Cat}}
\newcommand{\catop}{^{\mathrm{op}}}
\newcommand{\unit}{\mathds{1}}
\newcommand{\iso}{\xrightarrow{\;\smash{\raisebox{-0.5ex}{\ensuremath{\scriptstyle\sim}}}\;}}
\DeclareMathOperator{\Hom}{Hom}
\DeclareMathOperator{\Fun}{Fun}
\DeclareMathOperator{\Spec}{Spec}
\DeclareMathOperator{\Sec}{Sec}
\DeclareMathOperator{\Sol}{Sol}

% Theorem environments
\theoremstyle{plain}
\newtheorem{theorem}{Theorem}[chapter]
\newtheorem{conjecture}[theorem]{Conjecture}
\newtheorem{corollary}[theorem]{Corollary}
\newtheorem{lemma}[theorem]{Lemma}
\newtheorem{proposition}[theorem]{Proposition}

\theoremstyle{definition}
\newtheorem{claim}[theorem]{Claim}
\newtheorem{construction}[theorem]{Construction}
\newtheorem{convention}[theorem]{Convention}
\newtheorem{definition}[theorem]{Definition}
\newtheorem{discussion}[theorem]{Discussion}
\newtheorem{example}[theorem]{Example}
\newtheorem{exercise}[theorem]{Exercise}
\newtheorem{fact}[theorem]{Fact}
\newtheorem{notation}[theorem]{Notation}
\newtheorem{observation}[theorem]{Observation}
\newtheorem{question}[theorem]{Question}
\newtheorem{remark}[theorem]{Remark}
\newtheorem{terminology}[theorem]{Terminology}
\newtheorem{warning}[theorem]{Warning}

\newtheorem*{claim*}{Claim}
\newtheorem*{convention*}{Convention}
\newtheorem*{definition*}{Definition}
\newtheorem*{example*}{Example}
\newtheorem*{fact*}{Fact}
\newtheorem*{question*}{Question}
\newtheorem*{remark*}{Remark}
\newtheorem*{theorem*}{Theorem}
\newtheorem*{warning*}{Warning}

% Pullback and pushout markers for tikz-cd.
\tikzcdset{pullback/.style={phantom,"\lrcorner"{very near start}}}
\tikzcdset{pushout/.style={phantom,"{\rotatebox{180}{$\lrcorner$}}"{very near end}}}

\usepackage[capitalise]{cleveref}
\crefname{claim}{claim}{claims}
\crefname{conjecture}{conjecture}{conjectures}
\crefname{construction}{construction}{constructions}
\crefname{convention}{convention}{conventions}
\crefname{corollary}{corollary}{corollaries}
\crefname{definition}{definition}{definitions}
\crefname{discussion}{discussion}{discussions}
\crefname{example}{example}{examples}
\crefname{exercise}{exercise}{exercises}
\crefname{fact}{fact}{facts}
\crefname{lemma}{lemma}{lemmas}
\crefname{notation}{notation}{notations}
\crefname{observation}{observation}{observations}
\crefname{proposition}{proposition}{propositions}
\crefname{question}{question}{questions}
\crefname{remark}{remark}{remarks}
\crefname{terminology}{terminology}{terminologies}
\crefname{theorem}{theorem}{theorems}
\crefname{warning}{warning}{warnings}

\Crefname{claim}{Claim}{Claims}
\Crefname{conjecture}{Conjecture}{Conjectures}
\Crefname{construction}{Construction}{Constructions}
\Crefname{convention}{Convention}{Conventions}
\Crefname{corollary}{Corollary}{Corollaries}
\Crefname{definition}{Definition}{Definitions}
\Crefname{discussion}{Discussion}{Discussions}
\Crefname{example}{Example}{Examples}
\Crefname{exercise}{Exercise}{Exercises}
\Crefname{fact}{Fact}{Facts}
\Crefname{lemma}{Lemma}{Lemmas}
\Crefname{notation}{Notation}{Notations}
\Crefname{observation}{Observation}{Observations}
\Crefname{proposition}{Proposition}{Propositions}
\Crefname{question}{Question}{Questions}
\Crefname{remark}{Remark}{Remarks}
\Crefname{terminology}{Terminology}{Terminologies}
\Crefname{theorem}{Theorem}{Theorems}
\Crefname{warning}{Warning}{Warnings}


\begin{document}

\frontmatter

\title{Derived manifolds\\[10pt]
\Large MIGOS seminar notes, University of Regensburg, 2026/27}
\author{\large Organizers: Bastiaan Cnossen and Sil Linskens}

\maketitle

\begin{KeepFromToc}
  \tableofcontents
\end{KeepFromToc}

\chapter*{About these notes}
\addcontentsline{toc}{chapter}{About these notes}

These notes are AI-generated via Codex, model GPT-6 Astra.

These notes accompany a semester of MIGOS on derived manifolds. They are
written for first-year PhD students with varied backgrounds in homotopy
theory, algebraic geometry, differential geometry, and analysis. The chapters
follow the seminar program, from smooth function rings and derived
intersections to Steffens's representability theorem for elliptic moduli
problems and its application to pseudo-holomorphic curves.

The exposition develops the constructions through examples and proofs of
their principal formal properties. Substantial background theorems are
stated with references. Sections called ``Further reading'' provide
material for reading beyond the corresponding seminar talk.

The main route has three stages.
\begin{enumerate}
  \item Chapters 1--3 pass from geometric intersections to smooth function
    rings. The regular-equation calculation explains why transverse
    pullbacks are preserved; Chapter 3 collects its detailed globalization
    proof after the main conclusions.
  \item Chapters 4--9 construct derived intersections, their intrinsic
    tangent complexes, and local charts. Derived bordism gives a first
    geometric application, and stacks provide the gluing criterion.
  \item Chapters 10--14 formulate elliptic solution stacks, prove their
    relative representability using coordinates, augmentation, and analysis,
    and apply this argument to pseudo-holomorphic maps.
\end{enumerate}
The chapter openings identify the geometric question and the construction
that answers it. The final sections contain further proofs and examples;
they can be read separately from the main route. In particular, the
comparison and existence theorems for the derived formalism are quoted
where needed, while their consequences for equations are calculated.

Manifolds are finite-dimensional, Hausdorff, and second countable, and have
no boundary unless stated otherwise. A closed manifold is compact and has
no boundary. Throughout our derived $\infty$-categories we use chain complexes
with homological grading: their differentials lower degree by one.
The dg-algebra models for animated smooth rings are concentrated in
nonnegative degrees, with $\pi_i=H_i$.

\mainmatter

% The title and chapter boundary remain provisional. This chapter is the
% complete working manuscript for the first seminar talk.

\chapter{Why derived manifolds?}\label[chapter]{chap:Why_Derived_Manifolds}

\section{Transverse fiber products}
\label[section]{sec:Transverse_Fiber_Products}

Consider the following two intersections of curves in $\R^2$.

\begin{figure}[ht]
  \centering
  \begin{tikzpicture}[scale=1.05]
    \begin{scope}[xshift=-3.3cm]
      \draw[gray!45] (-1.7,0)--(1.7,0);
      \draw[gray!45] (0,-1.35)--(0,1.35);
      \draw[very thick,blue] (-1.55,0)--(1.55,0);
      \draw[very thick,red] (0,-1.2)--(0,1.2);
      \fill (0,0) circle (1.6pt);
      \node at (0,-1.7) {Transverse axes};
    \end{scope}
    \begin{scope}[xshift=3.3cm]
      \draw[gray!45] (-1.7,0)--(1.7,0);
      \draw[gray!45] (0,-.3)--(0,1.45);
      \draw[very thick,blue] (-1.55,0)--(1.55,0);
      \draw[very thick,red,domain=-1.35:1.35,samples=80]
        plot (\x,{0.72*\x*\x});
      \fill (0,0) circle (1.6pt);
      \node at (0,-.55) {A tangent line and parabola};
    \end{scope}
  \end{tikzpicture}
  \caption{Two intersections with the same underlying set.}
  \label[figure]{fig:Two_Intersections}
\end{figure}

Set-theoretically, the answer is the same in both panels of
\Cref{fig:Two_Intersections}: the intersection consists only of the origin.
Geometrically, the two answers behave very differently. The transverse
intersection is stable locally under sufficiently small perturbations. For the
tangent intersection, replace the parabola by
\[
  y=x^2+\varepsilon.
\]
Its intersection with the $x$-axis consists of two points when
$\varepsilon<0$, one point when $\varepsilon=0$, and no points when
$\varepsilon>0$. The singleton at $\varepsilon=0$ therefore conceals both an
infinitesimal direction and the
obstruction which prevents that direction from extending to an actual family
of solutions.

\begin{question*}
  What geometric object should replace a nontransverse intersection so that it
  remembers the defining equations, agrees with the ordinary intersection in
  the transverse case, and behaves well in families?
\end{question*}

This question is the point of entry to derived manifolds. We will not define
them in this chapter. Instead, we first recall exactly what transversality
provides, then calculate what is lost when it fails, and finally see the same
deformation--obstruction pattern in elliptic moduli problems. The calculations
will give us a short list of requirements that later definitions must satisfy.
The same intersection-theoretic problem is the starting point of Spivak's
construction of derived smooth manifolds
\cite[Sections~1--3]{Spivak_Derived_Smooth_Manifolds}.

Let $f\colon X\to Z$ and $g\colon Y\to Z$ be smooth maps between smooth
manifolds. Their fiber product as sets is
\[
  X\times_ZY
  =\{(x,y)\in X\times Y\mid f(x)=g(y)\}.
\]
Without an additional hypothesis, this subset of $X\times Y$ need not be a
smooth manifold of the expected dimension. Transversality is the hypothesis
which rules out this defect.

\begin{definition}[Transverse maps]\label[definition]{def:Transverse_Maps}
  Let $f\colon X\to Z$ and $g\colon Y\to Z$ be smooth maps. For
  $(x,y)\in X\times_ZY$, with common image $z=f(x)=g(y)$, define the
  \emph{linearized matching map}
  \[
    \delta_{x,y}
    =D_xf-D_yg
    \colon T_xX\oplus T_yY\longrightarrow T_zZ.
  \]
  We say that $f$ and $g$ are \emph{transverse} if $\delta_{x,y}$ is
  surjective for every $(x,y)\in X\times_ZY$.
\end{definition}

Equivalently, $f$ and $g$ are transverse if
\[
  \operatorname{im}(D_xf)+\operatorname{im}(D_yg)=T_zZ
\]
whenever $f(x)=g(y)=z$. In particular, if $A,B\subseteq Z$ are embedded
submanifolds and $f$ and $g$ are their inclusions, this condition becomes
\[
  T_zA+T_zB=T_zZ
\]
at every $z\in A\cap B$.

We will use the following standard consequence of the regular value theorem.

\begin{proposition}[Transverse preimages]
  \label[proposition]{prop:Transverse_Preimages}
  Let $h\colon M\to N$ be a smooth map and let $S\subseteq N$ be an embedded
  submanifold. Suppose that
  \[
    \operatorname{im}(D_xh)+T_{h(x)}S=T_{h(x)}N
  \]
  for every $x\in h^{-1}(S)$. Then $h^{-1}(S)$ is an embedded submanifold of
  $M$ of codimension $\operatorname{codim}_N(S)$, and
  \[
    T_xh^{-1}(S)=(D_xh)^{-1}(T_{h(x)}S)
  \]
  for every $x\in h^{-1}(S)$.
\end{proposition}

This result is also called the transverse preimage theorem; see
\cite[Chapter~8]{lee2000smooth}. We take it as recalled differential topology
and use it to prove the principal result of this subsection.

\begin{theorem}[Transverse fiber product theorem]
  \label[theorem]{thm:Transverse_Fiber_Product}
  Let $f\colon X\to Z$ and $g\colon Y\to Z$ be transverse smooth maps. Then
  $X\times_ZY$ is a smooth manifold. If it is nonempty, its dimension is
  \[
    \dim(X\times_ZY)=\dim X+\dim Y-\dim Z.
  \]
  For every $(x,y)\in X\times_ZY$, there is a natural isomorphism
  \[
    T_{(x,y)}(X\times_ZY)
    \cong
    \ker\bigl(\delta_{x,y}\colon T_xX\oplus T_yY\to T_zZ\bigr).
  \]
\end{theorem}

\begin{proof}
  Consider the product map
  \[
    h=f\times g\colon X\times Y\longrightarrow Z\times Z,
    \qquad
    h(x,y)=(f(x),g(y)).
  \]
  If $\Delta_Z\subseteq Z\times Z$ denotes the diagonal, then
  \[
    X\times_ZY=h^{-1}(\Delta_Z).
  \]
  At a point $(z,z)\in\Delta_Z$, its tangent space is
  \[
    T_{(z,z)}\Delta_Z=\{(w,w)\mid w\in T_zZ\}.
  \]
  The linear map
  \[
    \pi_z\colon T_zZ\oplus T_zZ\longrightarrow T_zZ,
    \qquad
    \pi_z(a,b)=a-b,
  \]
  is surjective and has kernel $T_{(z,z)}\Delta_Z$. Consequently, $h$ is
  transverse to $\Delta_Z$ at $(x,y)$ if and only if the composite
  \[
    T_xX\oplus T_yY
    \xrightarrow{D_{(x,y)}h}
    T_zZ\oplus T_zZ
    \xrightarrow{\pi_z}
    T_zZ
  \]
  is surjective. Here we use the elementary fact that a linear map to a vector
  space $W$ is transverse to a subspace $U\subseteq W$ if and only if its
  composite with $W\to W/U$ is surjective. The displayed composite is
  precisely
  \[
    (u,v)\longmapsto D_xf(u)-D_yg(v)=\delta_{x,y}(u,v).
  \]
  Thus the transversality of $f$ and $g$ is equivalent to the transversality
  of $h$ to $\Delta_Z$.

  By \Cref{prop:Transverse_Preimages}, the inverse image
  $h^{-1}(\Delta_Z)$ is a smooth submanifold of $X\times Y$. Since the diagonal
  has codimension $\dim Z$ in $Z\times Z$, this proves the dimension formula.

  The tangent-space statement in \Cref{prop:Transverse_Preimages} gives
  \begin{align*}
    T_{(x,y)}h^{-1}(\Delta_Z)
    &\cong
      \{(u,v)\in T_xX\oplus T_yY
        \mid (D_xf(u),D_yg(v))\in T_{(z,z)}\Delta_Z\} \\
    &=\{(u,v)\in T_xX\oplus T_yY
        \mid D_xf(u)=D_yg(v)\} \\
    &=\ker(\delta_{x,y}),
  \end{align*}
  as asserted.
\end{proof}

In particular, under the hypotheses of
\Cref{thm:Transverse_Fiber_Product}, there is a short exact sequence
\[
  0\longrightarrow T_{(x,y)}(X\times_ZY)
  \longrightarrow T_xX\oplus T_yY
  \xrightarrow{\delta_{x,y}}T_zZ
  \longrightarrow0.
\]
The question motivating the following examples is what remains meaningful when
the last map is no longer surjective.

\section{Nontransverse fiber products}
\label[section]{sec:Nontransverse_Fiber_Products}

\begin{notation}\label[notation]{not:Linearized_Matching_Complex}
  For arbitrary smooth maps $f\colon X\to Z$ and $g\colon Y\to Z$, and a point
  $(x,y)\in X\times_ZY$, write
  \[
    \mathbb{T}_{x,y}(f,g)
    =\bigl[T_xX\oplus T_yY\xrightarrow{\delta_{x,y}}T_zZ\bigr],
  \]
  with the source in cohomological degree $0$ and the target in degree $1$.
  We call this the \emph{linearized matching complex}. Thus
  \[
    H^0\mathbb{T}_{x,y}(f,g)=\ker(\delta_{x,y}),
    \qquad
    H^1\mathbb{T}_{x,y}(f,g)=\operatorname{coker}(\delta_{x,y}).
  \]
\end{notation}

The notation is deliberately provisional. Later, after derived intersections
have been defined, this complex will be identified with their tangent complex.
For now it is simply a two-term complex extracted from the derivatives of the
two maps. No homological algebra beyond the displayed kernel and cokernel
calculations is needed in this chapter.

\subsection{What the underlying set forgets}
\label[subsection]{sec:Underlying_Set_Forgets}

We now return to the two intersections pictured in
\Cref{fig:Two_Intersections}. The following two examples calculate their
linearized matching complexes. In both cases the set-theoretic intersection
consists only of the origin, but the complexes are different.

\begin{example}[Two transverse axes]\label[example]{ex:Transverse_Axes}
  Parametrize the two coordinate axes by
  \[
    i\colon\R\longrightarrow\R^2,
    \quad i(s)=(s,0),
    \qquad
    k\colon\R\longrightarrow\R^2,
    \quad k(t)=(0,t).
  \]
  Their intersection consists of the origin. The linearized matching map is
  \[
    \delta_{0,0}\colon\R^2\longrightarrow\R^2,
    \qquad
    \delta_{0,0}(u,v)=(u,-v).
  \]
  It is an isomorphism, so the maps are transverse and
  \[
    H^0\mathbb{T}_{0,0}(i,k)=0,
    \qquad
    H^1\mathbb{T}_{0,0}(i,k)=0.
  \]
\end{example}

\begin{example}[A tangent line and parabola]
  \label[example]{ex:Tangent_Line_And_Parabola}
  Parametrize the $x$-axis and the parabola $y=x^2$ by
  \[
    i\colon\R\longrightarrow\R^2,
    \quad i(s)=(s,0),
    \qquad
    j\colon\R\longrightarrow\R^2,
    \quad j(t)=(t,t^2).
  \]
  A point $(s,t)$ lies in the fiber product precisely when
  \[
    s=t,
    \qquad
    t^2=0.
  \]
  Thus the underlying fiber product is again a single point. At this point,
  however, the linearized matching map is
  \[
    \delta_{0,0}\colon\R^2\longrightarrow\R^2,
    \qquad
    \delta_{0,0}(u,v)=(u-v,0).
  \]
  It follows that
  \[
    H^0\mathbb{T}_{0,0}(i,j)
    =\{(u,u)\mid u\in\R\}\cong\R
  \]
  and
  \[
    H^1\mathbb{T}_{0,0}(i,j)
    =\R^2/(\R\times\{0\})\cong\R.
  \]
  We refer provisionally to these as the infinitesimal deformation space and
  the obstruction space. Their alternating dimension is zero, which agrees
  with the expected dimension of the intersection.
\end{example}

The nonzero class in degree $0$ has a direct interpretation. The vector
$(1,1)$ solves the linearization of the equations $s=t$ and $t^2=0$. It does
not arise as the velocity of an actual curve of solutions. More precisely, a
curve with this initial velocity would have expansions
\[
  s(\tau)=\tau+O(\tau^2),
  \qquad
  t(\tau)=\tau+O(\tau^2).
\]
The second equation would then give
\[
  t(\tau)^2=\tau^2+O(\tau^3),
\]
whose quadratic coefficient is nonzero. This coefficient lies in the second
coordinate of the target, which represents the cokernel of
$\delta_{0,0}(u,v)=(u-v,0)$. Thus it is an explicit second-order obstruction
to extending the infinitesimal solution.

In particular, there cannot be an actual curve of solutions with velocity
$(1,1)$. Indeed, the identity $t(\tau)^2=0$ for every $\tau$ forces
$t(\tau)=0$ and then $s(\tau)=0$ for every $\tau$. This is an elementary
instance of obstruction theory. A formal treatment will require more structure
than this calculation alone.

The same example can be presented as a zero locus in one variable. Eliminating
$s$ gives the smooth function
\[
  q\colon\R\longrightarrow\R,
  \qquad
  q(t)=t^2.
\]
Its linearized complex at the origin is
\[
  [\R\xrightarrow{D_0q}\R]
  =
  [\R\xrightarrow{0}\R].
\]
This is not merely an analogy with the matching complex in
\Cref{ex:Tangent_Line_And_Parabola}. Let
\[
  a=(1,1),
  \qquad
  b=(1,0)
\]
be a basis of the source of $\delta_{0,0}$, and let $e_1,e_2$ be the standard
basis of its target. Then
\[
  \delta_{0,0}(a)=0,
  \qquad
  \delta_{0,0}(b)=e_1.
\]
Consequently, the matching complex splits as
\[
  \mathbb{T}_{0,0}(i,j)
  \cong
  [\R b\iso\R e_1]
  \oplus
  [\R a\xrightarrow{0}\R e_2].
\]
The first summand is acyclic, while the second is isomorphic to the linearized
complex of $q$.

\begin{example}[The same point from different equations]
  \label[example]{ex:Same_Point_Different_Equations}
  Compare
  \[
    r,q\colon\R\longrightarrow\R,
    \qquad
    r(t)=t,
    \qquad
    q(t)=t^2.
  \]
  Both zero sets are the singleton $\{0\}$. Their linearized complexes at the
  origin are respectively
  \[
    [\R\xrightarrow{1}\R]
    \qquad\text{and}\qquad
    [\R\xrightarrow{0}\R].
  \]
  The first complex is acyclic, while the second has a one-dimensional
  cohomology group in each of degrees $0$ and $1$. Thus the passage from an
  equation to its underlying set of solutions forgets information already
  visible in the derivative.
\end{example}

The linearized complex does not remember the entire equation. For example,
$t^2$ and $t^3$ have the same derivative at the origin. Their different orders
of vanishing are higher-order information. Thus a derived intersection should
retain the nonlinear equation itself, while its tangent complex should extract
the first-order deformation and obstruction spaces.

\subsection{Complexes in families}
\label[subsection]{sec:Complexes_In_Families}

The previous examples concern a single equation. A one-parameter family shows
why it is useful to retain the entire two-term complex rather than only its
kernel or cokernel.

\begin{example}[A family with jumping kernels]
  \label[example]{ex:Family_With_Jumping_Kernels}
  For $t\in\R$, consider the linear map
  \[
    L_t\colon\R\longrightarrow\R,
    \qquad
    L_t(u)=tu.
  \]
  Its kernel and cokernel are
  \[
    \ker L_t=
    \begin{cases}
      0,&t\neq0,\\
      \R,&t=0,
    \end{cases}
    \qquad
    \operatorname{coker}L_t=
    \begin{cases}
      0,&t\neq0,\\
      \R,&t=0.
    \end{cases}
  \]
  Neither the kernels nor the cokernels form a vector bundle of constant rank
  over the parameter line.

  By contrast, let $E^0=E^1=\R\times\R$ be trivial line bundles over the first
  factor and define a bundle map
  \[
    d\colon E^0\longrightarrow E^1,
    \qquad
    d(t,u)=(t,tu).
  \]
  This gives a smooth two-term complex of vector bundles
  \[
    [E^0\xrightarrow{d}E^1]
  \]
  whose fiber over $t$ is $[\R\xrightarrow{L_t}\R]$. The complex varies
  smoothly even though the dimensions of its cohomology groups jump. The
  alternating dimension
  \[
    \dim\ker L_t-\dim\operatorname{coker}L_t=0
  \]
  is independent of $t$.

  The total set of solutions is
  \[
    \{(t,u)\in\R^2\mid tu=0\},
  \]
  which is the union of the two coordinate axes. The additional branch over
  $t=0$ records the jump in $\ker L_t$, and the two branches meet at the
  singular origin.
\end{example}

The point of \Cref{ex:Family_With_Jumping_Kernels} is not that the complex has
constant cohomology. It does not. The point is that its terms are vector
bundles and its differential varies smoothly. This is the stable datum from
which the jumping kernels and cokernels can both be recovered.

\subsection{Requirements on a derived replacement}
\label[subsection]{sec:Requirements_On_A_Derived_Replacement}

We can now collect what the examples require from a replacement of singular
fiber products.

\begin{enumerate}
  \item It should admit arbitrary pullbacks and recover the ordinary smooth
    fiber product when the maps are transverse.
  \item It should retain the nonlinear defining equations in a
    presentation-independent way, not merely their set of solutions or their
    first derivatives.
  \item It should organize deformation and obstruction data in families, even
    when the dimensions of the individual deformation and obstruction spaces
    jump.
\end{enumerate}

The second requirement contains two lessons from the preceding calculations.
The tangent line--parabola intersection and the zero locus of $t^2$ are two
presentations of the same problem, so the desired object should not depend on
the redundant variable. On the other hand, the zero set and the linearized
complex are both too coarse: the first cannot distinguish $t$ from $t^2$, and
the second cannot distinguish $t^2$ from $t^3$.

These are requirements, not a definition. They describe the finite-dimensional
objects we seek. The next question is why the same objects should govern moduli
problems arising from infinite-dimensional differential equations.

\section{Elliptic equations and finite-dimensional reduction}
\label[section]{sec:Elliptic_Moduli}

This section is a quoted preview of the analytic endpoint. In the live talk we
use only the semilinear example, the linearized two-term complex, and the
statement of finite-dimensional reduction. The bundle-theoretic setup and the
analytic hypotheses are retained here so that the preview is mathematically
precise, but they are not prerequisites for the next chapters.

Why should singular zero loci matter in geometry? Many moduli problems in
differential geometry, symplectic geometry, and gauge theory are spaces of
solutions to nonlinear elliptic partial differential equations. Such problems
usually also involve symmetries and compactifications. We suppress those
additional layers for now and study the equation itself.

Let $M$ be a compact smooth manifold, let $\pi\colon V\to M$ be a smooth fiber
bundle, and let $F\to M$ be a smooth vector bundle. Suppose that
\[
  P\colon\Gamma(V)\longrightarrow\Gamma(F)
\]
is a nonlinear elliptic differential operator. We write
\[
  \operatorname{Sol}(P)=P^{-1}(0).
\]
A point of $\operatorname{Sol}(P)$ is a section satisfying the geometric
equation encoded by $P$. Thus $\operatorname{Sol}(P)$ is our simplified moduli
set of solutions. When a symmetry group acts, the correct moduli problem is
typically a quotient stack rather than this set, but the local elliptic
mechanism discussed below remains the starting point.

For a concrete model, take $V=F=M\times\R$ over a closed connected
Riemannian manifold. The semilinear equation
\[
  P(u)=\Delta u+u^2=0
\]
defines a subset $\operatorname{Sol}(P)\subseteq C^\infty(M)$. This example is
not our eventual geometric application. Its purpose is to make clear that
$\operatorname{Sol}(P)$ is simply the locus of solutions to a nonlinear
elliptic equation. The zero function is a solution and its linearization is
$D_0P=\Delta$. Since $M$ is closed and connected, its kernel consists of the
constant functions, and self-adjointness identifies its cokernel with the same
one-dimensional space. We will return to this example when we carry out a
finite-dimensional reduction later in the seminar.

There is an apparent mismatch with the preceding sections. The spaces
$\Gamma(V)$ and $\Gamma(F)$ are infinite-dimensional, whereas derived
manifolds will be built locally from finite-dimensional smooth data. The role
of ellipticity is to bridge this gap.

Write
\[
  T^{\mathrm{vert}}V
  =\ker(D\pi\colon TV\longrightarrow TM)
\]
for the vertical tangent bundle. At a solution $u$, the tangent space to the
space of sections is
\[
  T_u\Gamma(V)\cong\Gamma(u^*T^{\mathrm{vert}}V).
\]
If $V$ is a vector bundle, then $T^{\mathrm{vert}}V\cong\pi^*V$, and this
reduces to the familiar identification $T_u\Gamma(V)\cong\Gamma(V)$.
The linearization of $P$ at $u$ is therefore the elliptic operator
\[
  D_uP\colon
  \Gamma(u^*T^{\mathrm{vert}}V)
  \longrightarrow
  \Gamma(F).
\]
We therefore obtain the two-term complex
\[
  \mathbb{T}_uP
  =\bigl[
      \Gamma(u^*T^{\mathrm{vert}}V)
      \xrightarrow{D_uP}
      \Gamma(F)
    \bigr]
\]
in cohomological degrees $0$ and $1$. Its cohomology is the infinitesimal
deformation space and obstruction space of the equation at $u$:
\[
  H^0(\mathbb{T}_uP)=\ker(D_uP),
  \qquad
  H^1(\mathbb{T}_uP)=\operatorname{coker}(D_uP).
\]
Ellipticity and compactness imply that these spaces are finite-dimensional.
As $u$ varies, their dimensions can jump, so the complex $\mathbb{T}_uP$ is
again more stable than its cohomology groups separately. This is the
infinite-dimensional version of \Cref{ex:Family_With_Jumping_Kernels}.\footnote{
  This is the motivating linear picture in
  \cite[Sections~1--2]{Pardon_Representability_Elliptic_Fredholm}.}

Finite-dimensional deformation and obstruction spaces do not yet make
$\operatorname{Sol}(P)$ a finite-dimensional geometric object. The remaining
analytic input says that all other directions can be solved away locally. We
use the following statement only as a quoted analytic principle. In addition
to the compactness and ellipticity assumptions above, one needs the standard
hypotheses ensuring that, after choosing Sobolev completions of sufficiently
high regularity, $P$ is a smooth map of Banach manifolds, its linearization is
Fredholm, the Banach inverse function theorem applies, and elliptic regularity
identifies the resulting nearby solutions with smooth sections.

\begin{fact}[Local finite-dimensional reduction]
  \label[fact]{fact:Finite_Dimensional_Reduction}
  Let $P$ satisfy the analytic hypotheses just stated, and let
  $u\in\operatorname{Sol}(P)$. There are
  finite-dimensional vector spaces $K$ and $C$, an open neighborhood $U$ of
  $0$ in $K$, and a smooth map
  \[
    \kappa\colon U\longrightarrow C,
    \qquad
    \kappa(0)=0,
  \]
  such that a neighborhood of $u$ in $\operatorname{Sol}(P)$ is identified
  with $\kappa^{-1}(0)$. Moreover, at every $v\in\kappa^{-1}(0)$, the
  finite-dimensional complex
  \[
    [K\xrightarrow{D_v\kappa}C]
  \]
  is quasi-isomorphic to the complex $\mathbb{T}_{u_v}P$ at the corresponding
  solution $u_v$.
\end{fact}

For this talk, ``quasi-isomorphic'' means that the complexes can be connected
by a zigzag of chain maps which induce isomorphisms on $H^0$ and $H^1$. In
particular, they have the same deformation and obstruction spaces. No
derived-category formalism is needed here.

The mechanism behind the fact is worth keeping in mind. After passing to
Sobolev completions, $D_uP$ is a Fredholm operator. One chooses splittings
which separate its kernel, image, and cokernel. The inverse function theorem
solves the equation in the complementary directions, leaving a
finite-dimensional obstruction map $\kappa$. Elliptic bootstrapping then shows
that the resulting solutions are smooth. We quote this analytic argument; the
precise conclusion above is
\cite[Fact~1.0.0.1]{Steffens_Derived_Differential_Geometry}.

\Cref{fact:Finite_Dimensional_Reduction} has two distinct consequences. First,
it solves the apparent dimension problem: near any solution, the
infinite-dimensional equation $P=0$ is controlled by a finite-dimensional
zero locus $\kappa^{-1}(0)$. Second, it preserves the linearized
deformation--obstruction theory. The examples from the preceding sections are
therefore not merely toys. They describe the local form of genuine elliptic
moduli problems.

The nonlinear map $\kappa$, and not only its derivative, remains essential.
As with $t\mapsto t^2$, its zero set can hide higher-order obstruction data.

\section{From local zero loci to coherent moduli problems}
\label[section]{sec:Local_Zero_Loci_Not_Enough}

\Cref{fact:Finite_Dimensional_Reduction} solves the infinite-dimensionality
problem, but it does not solve the invariance problem. The construction of
$\kappa$ depends on choices of Sobolev completions, splittings,
finite-dimensional subspaces, and neighborhoods. Different choices can
produce zero-locus presentations
\[
  \kappa\colon U\longrightarrow C
  \qquad\text{and}\qquad
  \kappa'\colon U'\longrightarrow C'
\]
of different dimensions. An identification of their underlying zero sets says
too little, just as equality of the zero sets of $t$ and $t^2$ says too little.

There are therefore two global problems.

\begin{enumerate}
  \item One must specify when two local finite-dimensional reductions present
    the same geometric object.
  \item One must glue these local objects on overlaps, together with
    compatibilities between the comparison maps and all higher
    compatibilities.
\end{enumerate}

The second point is where higher-categorical language becomes useful. It keeps
track not only of comparison maps, but also of homotopies between comparisons
and their iterated compatibilities. No knowledge of a particular model of
$\infty$-categories is needed for the present conclusion: ordinary equality
is too rigid, while an unstructured collection of coordinate changes is not
coherent enough.

There is also a conceptual alternative to writing down all coordinate changes
by hand. One can define the moduli problem as a functor of families and ask
whether it is represented by a derived smooth object. A representing object,
if it exists, is unique. This separates two questions that should not be
conflated: first represent the moduli problem, then use additional structures
  to construct virtual fundamental classes or enumerative invariants. Pardon's
  proceedings article advocates precisely this representability viewpoint, but
  its implementation belongs later in the seminar
  \cite[Section~1]{Pardon_Representability_Elliptic_Fredholm}.

The route from this local picture to the seminar's main theorem has three
stages. First, we construct derived zero loci and their tangent complexes, then
pass from affine calculations to locally affine derived smooth geometry.
Second, we formulate moduli problems intrinsically as functors of families and
construct their solution stacks. Third, Fredholm reduction and locality of
representability identify those stacks with quasi-smooth derived
$C^\infty$-schemes. Pseudo-holomorphic curves will then provide the principal
geometric application.

\section{Summary and supplementary materials}
\label[section]{sec:Why_Derived_Summary_And_Further_Reading}

\subsection{The live route}
\label[subsection]{sec:Why_Derived_Live_Route}

The default live route occupies 70 minutes. Its engine is the transverse
fiber-product theorem and the precise point where its proof uses
surjectivity. The tangent line--parabola intersection and the family
$[\R\xrightarrow{t}\R]$ then motivate deformation-obstruction complexes.
The elliptic section is an eight-minute application: it states
\Cref{fact:Finite_Dimensional_Reduction} and explains why the preceding
finite-dimensional examples occur in genuine moduli problems, without
developing the analytic proof. The talk ends with the invariance and coherence
problems for the resulting local zero-locus presentations.

\subsection{Takeaways and next step}
\label[subsection]{sec:Takeaways_And_Next_Step}

The argument of this chapter can be compressed into four statements.

\begin{enumerate}
  \item Transversality makes an ordinary fiber product smooth, with tangent
    space given by the kernel of the linearized matching map.
  \item When transversality fails, the same linearization produces both a
    deformation space and an obstruction space, but the nonlinear equation
    contains further information.
  \item In families, the two-term complex varies more robustly than its kernel
    or cokernel separately.
  \item Elliptic moduli problems are locally finite-dimensional zero loci, but
    those local presentations are nonunique and must be compared coherently.
\end{enumerate}

Derived manifolds are designed to package these requirements into a category
with well-behaved pullbacks. To construct them, we first need a language in
which smooth equations can be treated algebraically. The next chapter begins
with $C^\infty$-rings and explains how an ordinary manifold is encoded by its
algebra of smooth functions.

\subsection{Related literature}
\label[subsection]{sec:Why_Derived_Related_Literature}

The differential-topological input, including transverse preimages and fiber
products, is standard; Lee gives a convenient reference
\cite[Chapter~8]{lee2000smooth}. Spivak develops derived manifolds explicitly
as a solution to the failure of ordinary intersections and proves their basic
intersection-theoretic properties
\cite[Sections~1--3 and~9]{Spivak_Derived_Smooth_Manifolds}.

For the passage from elliptic equations to finite-dimensional obstruction
maps, the direct source used here is the opening chapter of Steffens's thesis
\cite[Chapter~1]{Steffens_Derived_Differential_Geometry}. Pardon's completed
proceedings article gives a complementary discussion of complexes of elliptic
operators, Kuranishi charts, and the representability viewpoint
\cite[Sections~1--5]{Pardon_Representability_Elliptic_Fredholm}. Its principal
representability argument is explicitly presented as a proof sketch, a status
that later chapters will keep separate from Steffens's completed theorem.

% The title and chapter boundary remain provisional. This chapter is the
% complete working manuscript for the second seminar talk. Sections marked as
% supplements are intended for the written notes, not for the default live
% route through the talk.

\chapter{Smooth geometry through \texorpdfstring{$\Cinfty$}{C-infinity}-rings}
\label[chapter]{chap:Smooth_Geometry_Through_Cinfty_Rings}

\section{Smooth functional calculus}
\label[section]{sec:Smooth_Functional_Calculus}

The preceding chapter ended with a basic loss of information. The equations
\[
  t=0
  \qquad\text{and}\qquad
  t^2=0
\]
have the same set of solutions in $\R$, but their linearizations at the origin
are different. Passing immediately to the zero set forgets the equation which
cut it out.

Ordinary algebraic geometry suggests retaining a quotient of a function ring.
In the smooth setting, the two candidates are
\[
  \Cinfty(\R)/(t)
  \qquad\text{and}\qquad
  \Cinfty(\R)/(t^2).
\]
These expressions lead to two questions.

\begin{enumerate}
  \item What structure on $\Cinfty(\R)$ remembers substitution by arbitrary
    smooth functions, rather than only addition and multiplication?
  \item Why does that structure descend to a quotient by an ordinary ideal?
\end{enumerate}

The answer to the first question is the notion of a $\Cinfty$-ring. The answer
to the second is Hadamard's lemma. Once both are in place, we will compute the
two displayed quotients and see exactly what the second remembers beyond its
single real point.

There is also a more global test for the language. If a manifold is replaced
by its ring of smooth functions, can we still recover its smooth maps? The
principal geometric result of the chapter says that we can. In fact, the same
argument works for arbitrary closed subsets of Euclidean spaces equipped with
their smooth functions.

Let us begin with the motivating example. If $M$ is a manifold, a smooth map
$f\colon\R^n\to\R$ can be applied pointwise to an $n$-tuple of smooth
functions on $M$:
\[
  (a_1,\ldots,a_n)
  \longmapsto
  f(a_1,\ldots,a_n),
  \qquad
  f(a_1,\ldots,a_n)(x)
  =f(a_1(x),\ldots,a_n(x)).
\]
A $\Cinfty$-ring is a set on which all these smooth operations make sense,
subject to the relations imposed by composition of smooth maps.
The operations formulation below follows Moerdijk--Reyes
\cite[Chapter~I, Section~1]{Moerdijk_Reyes_Smooth_Infinitesimal_Analysis}; a
modern account, including the functor formulation, is given by Joyce
\cite[Sections~2.1--2.2]{Joyce_Introduction_Cinfty_Schemes}.

\begin{definition}[$\Cinfty$-ring]
  \label[definition]{def:Cinfty_Ring_Operations}
  A \emph{$\Cinfty$-ring} is a set $A$ equipped, for every $n\geq 0$ and every
  smooth map $f\colon\R^n\to\R$, with an operation
  \[
    \Phi_f\colon A^n\longrightarrow A.
  \]
  These operations satisfy the following two conditions.

  \begin{enumerate}
    \item If $\pi_i\colon\R^n\to\R$ is the $i$th projection, then
      $\Phi_{\pi_i}(a_1,\ldots,a_n)=a_i$.
    \item If $f_1,\ldots,f_m\colon\R^n\to\R$ and
      $g\colon\R^m\to\R$ are smooth, then
      \[
        \Phi_{g\circ(f_1,\ldots,f_m)}
        =\Phi_g\circ(\Phi_{f_1},\ldots,\Phi_{f_m}).
      \]
  \end{enumerate}
\end{definition}

The case $n=0$ supplies an element of $A$ for every constant in $\R$. The
second axiom says that evaluating a composite smooth function gives the same
answer as evaluating its pieces successively.

\begin{example}[Smooth functions]
  \label[example]{ex:Smooth_Functions_Cinfty_Ring}
  For every manifold $M$, the set $\Cinfty(M)$ is a $\Cinfty$-ring under
  pointwise smooth functional calculus. More generally, the same construction
  works whenever a notion of smooth real-valued function has been specified
  and is closed under smooth substitution.
\end{example}

\begin{definition}[$\Cinfty$-homomorphism]
  \label[definition]{def:Cinfty_Homomorphism}
  A homomorphism of $\Cinfty$-rings is a map $\varphi\colon A\to B$ satisfying
  \[
    \varphi\bigl(\Phi_f(a_1,\ldots,a_n)\bigr)
    =\Phi_f\bigl(\varphi(a_1),\ldots,\varphi(a_n)\bigr)
  \]
  for every smooth $f\colon\R^n\to\R$ and every
  $(a_1,\ldots,a_n)\in A^n$.
\end{definition}

If $u\colon M\to N$ is smooth, then pullback gives a homomorphism
\[
  u^*\colon\Cinfty(N)\longrightarrow\Cinfty(M),
  \qquad
  h\longmapsto h\circ u.
\]
Thus smooth geometry enters the algebra contravariantly. This reversal of
arrows will remain in force throughout the next several chapters.

\begin{observation}[The underlying commutative algebra]
  \label[observation]{obs:Underlying_Commutative_Algebra}
  Every $\Cinfty$-ring has an underlying commutative $\R$-algebra. Addition,
  multiplication, scalar multiplication, zero, and one are obtained from the
  corresponding smooth maps between Euclidean spaces. Every
  $\Cinfty$-homomorphism is consequently a homomorphism of commutative
  $\R$-algebras.
\end{observation}

The converse is false. The extra smooth operations impose genuine algebraic
conditions. For example, in every $\Cinfty$-ring and for every $a\in A$, the
element $1+a^2$ is invertible. Its inverse is
\[
  \Phi_g(a),
  \qquad
  g(x)=\frac{1}{1+x^2},
\]
because the identity $(1+x^2)g(x)=1$ is an identity of smooth functions.
Therefore the usual commutative algebra $\R[t]$ cannot be the underlying
algebra of a $\Cinfty$-ring whose distinguished element $t$ is its polynomial
generator: the polynomial $1+t^2$ has no inverse in $\R[t]$.

There is a compact categorical formulation of the same definition. Let
$\mathrm{Euc}$ be the category whose objects are the Euclidean spaces
$\R^n$, for $n\geq 0$, and whose morphisms are smooth maps. Its finite
products are the usual products of Euclidean spaces.

\begin{proposition}[Functor formulation]
  \label[proposition]{prop:Cinfty_Ring_Functor_Formulation}
  To give a $\Cinfty$-ring is equivalently to give a finite-product-preserving
  functor
  \[
    A\colon\mathrm{Euc}\longrightarrow\mathrm{Set}.
  \]
  Under this correspondence, the underlying set of the ring is $A(\R)$.
\end{proposition}

\begin{proof}
  Suppose first that the operations of
  \Cref{def:Cinfty_Ring_Operations} are given. Set
  $A(\R^n)=A^n$. If
  $f=(f_1,\ldots,f_m)\colon\R^n\to\R^m$, define
  \[
    A(f)=(\Phi_{f_1},\ldots,\Phi_{f_m})\colon A^n\longrightarrow A^m.
  \]
  The projection and composition axioms say precisely that this is a functor,
  and the construction visibly preserves finite products.

  Conversely, let $A\colon\mathrm{Euc}\to\mathrm{Set}$ preserve finite
  products and put $A=A(\R)$. The product comparison identifies
  $A(\R^n)$ with $A^n$. For a smooth map $f\colon\R^n\to\R$, the map $A(f)$
  then defines $\Phi_f\colon A^n\to A$. Functoriality gives the two axioms in
  \Cref{def:Cinfty_Ring_Operations}.
\end{proof}

This is an instance of a Lawvere theory, but no general theory of Lawvere
theories is needed here. The useful point is simply that all smooth maps
between Euclidean spaces act functorially.

\section{Free smooth algebras and singular equations}
\label[section]{sec:Free_Smooth_Algebras_And_Singular_Equations}

\subsection{Free smooth algebras}
\label[subsection]{sec:Free_Smooth_Algebras}

Let $x_1,\ldots,x_n\in\Cinfty(\R^n)$ be the coordinate functions. Every
smooth function $f\in\Cinfty(\R^n)$ is obtained by applying the operation
$\Phi_f$ to these coordinates:
\[
  f=\Phi_f(x_1,\ldots,x_n).
\]
This tautological identity is the key to the universal property.

\begin{theorem}[Free $\Cinfty$-rings]
  \label[theorem]{thm:Free_Cinfty_Rings}
  For every $\Cinfty$-ring $A$, evaluation on the coordinate functions induces
  a bijection
  \[
    \Hom_{\Cinfty\text{-}\mathrm{Ring}}
    \bigl(\Cinfty(\R^n),A\bigr)
    \longrightarrow A^n,
    \qquad
    \varphi\longmapsto
    \bigl(\varphi(x_1),\ldots,\varphi(x_n)\bigr).
  \]
  Thus $\Cinfty(\R^n)$ is the free $\Cinfty$-ring on
  $x_1,\ldots,x_n$.
\end{theorem}

This is Moerdijk--Reyes, Proposition~1.1
\cite[Chapter~I, Proposition~1.1]{Moerdijk_Reyes_Smooth_Infinitesimal_Analysis}.

\begin{proof}
  Given $(a_1,\ldots,a_n)\in A^n$, define
  \[
    \operatorname{ev}_{a_1,\ldots,a_n}
    \colon\Cinfty(\R^n)\longrightarrow A,
    \qquad
    f\longmapsto\Phi_f(a_1,\ldots,a_n).
  \]
  The composition axiom for smooth functional calculus says exactly that this
  map preserves every smooth operation, so it is a $\Cinfty$-homomorphism. It
  sends $x_i$ to $a_i$.

  Conversely, if $\varphi\colon\Cinfty(\R^n)\to A$ is a
  $\Cinfty$-homomorphism, then
  \[
    \varphi(f)
    =\varphi\bigl(\Phi_f(x_1,\ldots,x_n)\bigr)
    =\Phi_f\bigl(\varphi(x_1),\ldots,\varphi(x_n)\bigr).
  \]
  Hence $\varphi$ is uniquely determined by the images of the coordinate
  functions, and the two constructions are inverse.
\end{proof}

Smooth functions on $\R^n$ therefore play the role that polynomials in $n$
variables play in ordinary commutative algebra. The analogy is useful, but it
should not be taken to say that a smooth function is determined by a finite
algebraic expression. The allowed operations already include every smooth map
$\R^n\to\R$.

\subsection{Quotients and smooth equations}
\label[subsection]{sec:Quotients_And_Smooth_Equations}

To use quotient rings, we must show that an ordinary ideal is respected by all
smooth operations. The required fact is a multivariable form of the
fundamental theorem of calculus.

\begin{lemma}[Hadamard's lemma]
  \label[lemma]{lem:Hadamard}
  For every smooth map $f\colon\R^n\to\R$, there are smooth maps
  $g_1,\ldots,g_n\colon\R^n\times\R^n\to\R$ such that
  \[
    f(y)-f(x)=\sum_{i=1}^n(y_i-x_i)g_i(x,y).
  \]
\end{lemma}

\begin{proof}
  Put
  \[
    g_i(x,y)
    =\int_0^1
      \frac{\partial f}{\partial x_i}\bigl(x+s(y-x)\bigr)\,ds.
  \]
  Differentiation under the integral sign shows that $g_i$ is smooth. The
  fundamental theorem of calculus applied to
  $s\mapsto f(x+s(y-x))$ gives
  \[
    f(y)-f(x)
    =\int_0^1\frac{d}{ds}f(x+s(y-x))\,ds
    =\sum_{i=1}^n(y_i-x_i)g_i(x,y).
  \]
\end{proof}

\begin{proposition}[Quotient $\Cinfty$-rings]
  \label[proposition]{prop:Quotient_Cinfty_Rings}
  Let $A$ be a $\Cinfty$-ring, and let $I\subseteq A$ be an ideal in its
  underlying commutative algebra. There is a unique $\Cinfty$-ring structure
  on $A/I$ for which the quotient map $A\to A/I$ is a
  $\Cinfty$-homomorphism.
\end{proposition}

This is Moerdijk--Reyes, Proposition~1.2
\cite[Chapter~I, Proposition~1.2]{Moerdijk_Reyes_Smooth_Infinitesimal_Analysis}.

\begin{proof}
  The only possible definition is
  \[
    \Phi_f(\overline{a}_1,\ldots,\overline{a}_n)
    =\overline{\Phi_f(a_1,\ldots,a_n)}.
  \]
  We check that it is independent of the representatives. Suppose that
  $b_i-a_i\in I$ for every $i$. Apply \Cref{lem:Hadamard} inside $A$ to obtain
  \[
    \Phi_f(b_1,\ldots,b_n)-\Phi_f(a_1,\ldots,a_n)
    =\sum_{i=1}^n(b_i-a_i)
      \Phi_{g_i}(a_1,\ldots,a_n,b_1,\ldots,b_n).
  \]
  The right-hand side lies in $I$. Thus the operations are well defined. Their
  projection and composition axioms descend from those in $A$. Uniqueness is
  forced by the requirement that the quotient map preserve every operation.
\end{proof}

This result is special to smooth functional calculus. It says that every
ordinary ideal in the underlying algebra is already compatible with all the
smooth operations.

\begin{definition}[Finite generation and finite presentation]
  \label[definition]{def:Cinfty_Finiteness}
  A $\Cinfty$-ring $A$ is \emph{finitely generated} if there are
  $a_1,\ldots,a_n\in A$ such that no proper $\Cinfty$-subring contains all the
  $a_i$. It is \emph{finitely presented} if it is isomorphic to
  \[
    \Cinfty(\R^n)/(f_1,\ldots,f_k)
  \]
  for some finite collections of generators and relations.
\end{definition}

\begin{proposition}
  \label[proposition]{prop:Finitely_Generated_Cinfty_Quotient}
  A $\Cinfty$-ring is finitely generated if and only if it is isomorphic to
  $\Cinfty(\R^n)/I$ for some $n$ and some ideal
  $I\subseteq\Cinfty(\R^n)$.
\end{proposition}

\begin{proof}
  If $a_1,\ldots,a_n$ generate $A$, the universal property in
  \Cref{thm:Free_Cinfty_Rings} gives a homomorphism
  $\Cinfty(\R^n)\to A$ sending $x_i$ to $a_i$. Its image is a
  $\Cinfty$-subring containing the generators, so it is surjective. By
  \Cref{prop:Quotient_Cinfty_Rings}, the induced isomorphism of commutative
  algebras $\Cinfty(\R^n)/I\to A$ is an isomorphism of $\Cinfty$-rings. The
  converse follows because the images of the coordinate functions generate
  every such quotient.
\end{proof}

Finite generation here means generation under every smooth operation. It does
not imply that the underlying commutative $\R$-algebra is finitely generated.
Likewise, finite presentation is a genuine additional condition: smooth
function rings are not Noetherian. The next chapter will return to finite
presentation for manifolds.

\subsection{One point, two equations}
\label[subsection]{sec:One_Point_Two_Equations}

We can now compute the two quotients from the opening. In one variable,
\Cref{lem:Hadamard} gives
\[
  f(t)=f(0)+t g(t)
\]
for some smooth $g$. Evaluation at the origin therefore induces an
isomorphism
\[
  \Cinfty(\R)/(t)\cong\R.
\]

For the second quotient, Taylor's formula with smooth remainder gives
\[
  f(t)=f(0)+f'(0)t+t^2h(t).
\]
Consequently, the map
\[
  \Cinfty(\R)\longrightarrow\R[\varepsilon]/(\varepsilon^2),
  \qquad
  f\longmapsto f(0)+f'(0)\varepsilon
\]
is a surjective homomorphism of commutative $\R$-algebras with kernel $(t^2)$.
It identifies
\[
  \Cinfty(\R)/(t^2)
  \cong
  \R[\varepsilon]/(\varepsilon^2).
\]
The ring on the right is called the ring of \emph{dual numbers}. Its
$\Cinfty$-structure is not additional arbitrary data. It is the quotient
structure from \Cref{prop:Quotient_Cinfty_Rings}, and its operations are given
by the first-order Taylor formula
\[
  f(a+b\varepsilon)=f(a)+b f'(a)\varepsilon.
\]

To compare these quotients with their ordinary solution sets, we isolate the
appropriate notion of point.

\begin{definition}[Real point]
  \label[definition]{def:Real_Point}
  A \emph{real point} of a $\Cinfty$-ring $A$ is a
  $\Cinfty$-homomorphism $A\to\R$, where $\R$ carries its evident smooth
  functional calculus.
\end{definition}

\begin{proposition}[Real points of a quotient]
  \label[proposition]{prop:Real_Points_Of_Quotient}
  For every ideal $I\subseteq\Cinfty(\R^n)$, evaluation induces a bijection
  \[
    \Hom_{\Cinfty\text{-}\mathrm{Ring}}
    \bigl(\Cinfty(\R^n)/I,\R\bigr)
    \cong
    \{p\in\R^n\mid f(p)=0\text{ for every }f\in I\}.
  \]
\end{proposition}

\begin{proof}
  By \Cref{thm:Free_Cinfty_Rings}, every homomorphism
  $\Cinfty(\R^n)\to\R$ is evaluation at a unique point
  $p=(p_1,\ldots,p_n)$. It factors through the quotient precisely when every
  element of $I$ vanishes at $p$.
\end{proof}

Both $\Cinfty(\R)/(t)$ and $\Cinfty(\R)/(t^2)$ therefore have exactly one real
point. They are nevertheless not isomorphic: the class of $t$ is zero in the
first and is a nonzero square-zero element in the second. The quotient by
$(t^2)$ remembers first-order information which the zero set forgets.

\begin{warning}
  \label[warning]{warn:Ordinary_Cinfty_Not_Derived}
  This example does not yet construct a derived zero locus. An ordinary
  quotient $\Cinfty$-ring can retain nilpotents and dependence on the chosen
  ideal, but it does not by itself encode homotopy-coherent equations or
  derived pullbacks. It is the classical algebraic input from which the later
  derived theory will be built.
\end{warning}

\section{Recovering smooth spaces from functions}
\label[section]{sec:Recovering_Smooth_Spaces_From_Functions}

\subsection{Recovering closed smooth subsets}
\label[subsection]{sec:Recovering_Closed_Smooth_Subsets}

We now turn from equations to geometry. Let $X\subseteq\R^n$ be an arbitrary
closed subset. It need not be a manifold.

\begin{definition}[Smooth functions and maps on subsets]
  \label[definition]{def:Smooth_Functions_On_Subsets}
  A function $a\colon X\to\R$ is \emph{smooth} if it extends to a smooth
  function on an open neighborhood of $X$ in $\R^n$. We write $\Cinfty(X)$
  for the resulting $\Cinfty$-ring.

  A map $u\colon X\to Y\subseteq\R^m$ is \emph{smooth} if its component
  functions are smooth on $X$. Equivalently, it extends to a smooth map into
  $\R^m$ on an open neighborhood of $X$.
\end{definition}

The equivalence in the second paragraph uses only that there are finitely many
component functions: intersect their extension neighborhoods and combine the
extensions.

\begin{lemma}[Smooth extension from a closed subset]
  \label[lemma]{lem:Smooth_Extension_Closed_Subset}
  If $X\subseteq\R^n$ is closed, every element of $\Cinfty(X)$ is the
  restriction of a globally defined smooth function on $\R^n$.
\end{lemma}

\begin{proof}
  Let $a\in\Cinfty(X)$, and choose a smooth extension $\widetilde a$ on an open
  neighborhood $U$ of $X$. A smooth partition of unity subordinate to the
  cover
  \[
    \{U,\R^n\setminus X\}
  \]
  supplies a smooth function $\chi\colon\R^n\to[0,1]$ whose support is
  contained in $U$ and whose restriction to $X$ is one. Define
  \[
    \overline a(x)=
    \begin{cases}
      \chi(x)\widetilde a(x),&x\in U,\\
      0,&x\notin U.
    \end{cases}
  \]
  Since the support of $\chi$ is contained in $U$, this formula defines a
  smooth function on all of $\R^n$. Its restriction to $X$ is $a$.
\end{proof}

Let
\[
  \mathfrak m_X
  =\{h\in\Cinfty(\R^n)\mid h|_X=0\}.
\]
The restriction homomorphism is surjective by
\Cref{lem:Smooth_Extension_Closed_Subset}, with kernel $\mathfrak m_X$.
Hence
\[
  \Cinfty(X)\cong\Cinfty(\R^n)/\mathfrak m_X.
\]
In particular, $\Cinfty(X)$ is a finitely generated $\Cinfty$-ring, even when
$X$ is singular or has complicated topology.

The following theorem is the main geometric payoff of the chapter.

\begin{theorem}[Full faithfulness for closed subsets]
  \label[theorem]{thm:Full_Faithfulness_Closed_Subsets}
  Let $X\subseteq\R^n$ and $Y\subseteq\R^m$ be closed subsets. Pullback induces
  a bijection
  \[
    \{\text{smooth maps }Y\to X\}
    \longrightarrow
    \Hom_{\Cinfty\text{-}\mathrm{Ring}}
    \bigl(\Cinfty(X),\Cinfty(Y)\bigr).
  \]
  Consequently, the functor
  \[
    \Cinfty(-)\colon
    (\text{closed subsets of Euclidean spaces})\catop
    \longrightarrow
    \Cinfty\text{-}\mathrm{Ring}
  \]
  is fully faithful.
\end{theorem}

\begin{proof}
  Let
  \[
    \Psi\colon\Cinfty(X)\longrightarrow\Cinfty(Y)
  \]
  be a homomorphism. Write $x_1,\ldots,x_n$ for the coordinate functions on
  $\R^n$, and define
  \[
    g_i=\Psi(x_i|_X)\in\Cinfty(Y),
    \qquad
    g=(g_1,\ldots,g_n)\colon Y\longrightarrow\R^n.
  \]
  The map $g$ is smooth by \Cref{def:Smooth_Functions_On_Subsets}. For every
  $h\in\Cinfty(\R^n)$, preservation of smooth functional calculus gives
  \[
    \Psi(h|_X)
    =\Psi\bigl(\Phi_h(x_1|_X,\ldots,x_n|_X)\bigr)
    =\Phi_h(g_1,\ldots,g_n)
    =h\circ g.
  \]

  We claim that $g$ lands in $X$. If $g(y)=p\notin X$ for some $y\in Y$,
  choose a smooth bump function $h\colon\R^n\to\R$ supported in
  $\R^n\setminus X$ and satisfying $h(p)=1$. Then $h|_X=0$, whereas the
  displayed identity gives
  \[
    0=\Psi(h|_X)(y)=h(g(y))=1,
  \]
  a contradiction.

  It remains to identify $\Psi$. Given $a\in\Cinfty(X)$, choose a global
  smooth extension $h\in\Cinfty(\R^n)$ using
  \Cref{lem:Smooth_Extension_Closed_Subset}. Then
  \[
    \Psi(a)=\Psi(h|_X)=h\circ g=a\circ g.
  \]
  Thus $\Psi=g^*$. The map $g$ is unique because its component functions are
  the images of the restricted coordinates.

  Conversely, pullback along a smooth map $g\colon Y\to X$ preserves every
  smooth operation pointwise, and hence defines a $\Cinfty$-homomorphism. The
  two constructions are inverse.
\end{proof}

The corresponding result for all locally closed subsets is
Moerdijk--Reyes, Proposition~1.5
\cite[Chapter~I, Proposition~1.5]{Moerdijk_Reyes_Smooth_Infinitesimal_Analysis}.

For the usual class of finite-dimensional, second-countable smooth manifolds,
Whitney's embedding theorem may be strengthened to give a proper Euclidean
embedding \cite{lee2000smooth}. Its image is closed. The intrinsic smooth
functions and smooth maps agree with those in
\Cref{def:Smooth_Functions_On_Subsets}. We therefore obtain the case which
will be used later.

\begin{corollary}[Recovering manifolds]
  \label[corollary]{cor:Recovering_Manifolds}
  The contravariant functor
  \[
    \Cinfty(-)\colon
    (\text{smooth manifolds})\catop
    \longrightarrow
    \Cinfty\text{-}\mathrm{Ring}
  \]
  is fully faithful.
\end{corollary}

The conclusion is stronger than a reconstruction of points. The entire smooth
map $Y\to X$ is recovered from the homomorphism on function rings by looking
at the images of finitely many ambient coordinate functions.

\section{Summary and supplementary materials}
\label[section]{sec:Cinfty_Summary_And_Further_Reading}

\subsection{The live route}
\label[subsection]{sec:Cinfty_Live_Route}

The default live route occupies 70 minutes. Its engine is the free
$\Cinfty$-ring theorem together with Hadamard's lemma and the quotient
construction. The dual numbers are the principal calculation. Full
faithfulness for closed subsets is the geometric payoff and the handoff to
Talk 3. The locally closed extension and Borel's theorem remain available as
supplementary reading, but no later talk depends on their proofs.

\subsection{Takeaways and next step}
\label[subsection]{sec:Cinfty_Takeaways_And_Next_Step}

The main route through the chapter can be compressed into four statements.

\begin{enumerate}
  \item A $\Cinfty$-ring is an algebraic object in which every smooth map
    $\R^n\to\R$ acts as an operation.
  \item The ring $\Cinfty(\R^n)$ is free on its coordinate functions, and
    Hadamard's lemma implies that every ordinary ideal defines a quotient
    $\Cinfty$-ring.
  \item The quotients by $(t)$ and $(t^2)$ have the same real point but
    different algebraic structures. The latter retains a first-order
    thickening which the zero set forgets.
  \item Passing contravariantly from a manifold to its $\Cinfty$-ring of
    smooth functions loses neither the manifold nor its smooth maps.
\end{enumerate}

We have therefore found a classical algebraic language for smooth equations.
Two important tasks remain before deriving it. First, we need to understand
which $\Cinfty$-rings actually arise from manifolds and how geometric products
and transverse pullbacks appear algebraically. Second, ordinary quotients do
not yet have the homotopy-invariant universal property required for
nontransverse intersections. The next chapter addresses the first task and
prepares the second.

\subsection{Supplement: infinite jets and Borel's theorem}
\label[subsection]{sec:How_Large_Smooth_Algebra_Can_Be}

Borel's theorem is a useful calibration of the size of smooth algebra, but no
later chapter depends on it. The preceding examples were finite-dimensional
as ordinary $\R$-algebras. The next theorem shows how misleading that special
feature can be. Let
$\mathfrak m_0^\infty\subseteq\Cinfty(\R^n)$ be the ideal of functions
\emph{flat at the origin}, meaning that every partial derivative of every
order vanishes at the origin. The Taylor homomorphism is
\[
  T_0\colon\Cinfty(\R^n)\longrightarrow\R[[x_1,\ldots,x_n]],
  \qquad
  f\longmapsto
  \sum_{\alpha\in\N^n}\frac{D^\alpha f(0)}{\alpha!}x^\alpha.
\]

\begin{theorem}[Borel's theorem]
  \label[theorem]{thm:Borel}
  The Taylor homomorphism $T_0$ is surjective and has kernel
  $\mathfrak m_0^\infty$. It therefore induces an isomorphism of commutative
  $\R$-algebras
  \[
    \Cinfty(\R^n)/\mathfrak m_0^\infty
    \cong
    \R[[x_1,\ldots,x_n]].
  \]
  In particular, the formal power series ring inherits a $\Cinfty$-ring
  structure from the quotient on the left.
\end{theorem}

This is the form of Borel's theorem recorded as Moerdijk--Reyes,
Theorem~1.3
\cite[Chapter~I, Theorem~1.3]{Moerdijk_Reyes_Smooth_Infinitesimal_Analysis}.

The assertion about the kernel is immediate from the definition. Surjectivity
is the content of the theorem: every formal Taylor series is the Taylor series
of some smooth function. We record the standard construction in one variable
as a sketch, both to indicate why the result is plausible and to keep the
logical status of the argument clear.

\begin{discussion}[Proof sketch in one variable]
  \label[discussion]{disc:Borel_Proof_Sketch}
  Choose a compactly supported smooth function $\chi$ which is identically one
  near the origin. Given prescribed derivatives $(a_k)_{k\geq0}$, consider
  \[
    \sum_{k\geq0}\frac{a_k}{k!}t^k\chi(\lambda_k t).
  \]
  The numbers $\lambda_k$ are chosen recursively and sufficiently large. For
  every $j<k$, the $C^j$-norm of the $k$th summand can then be made at most
  $2^{-k}$. It follows that the series and each of its derivatives converge
  uniformly on $\R$. Near the origin the cutoff in the $k$th summand is
  constant, so this summand contributes $a_k$ to the $k$th derivative at zero
  and contributes nothing to the other Taylor coefficients. The resulting
  smooth function has the prescribed Taylor series. The multivariable proof
  uses the same rapidly shrinking cutoff construction, organized by total
  degree.
\end{discussion}

The quotient in \Cref{thm:Borel} is generated as a $\Cinfty$-ring by the
images of $x_1,\ldots,x_n$, but its underlying commutative algebra contains
arbitrary formal power series. This is the promised warning about finiteness:
generation under all smooth operations is a much weaker condition than
generation of the underlying commutative algebra.

\subsection{Supplement: locally closed subsets}
\label[subsection]{sec:Supplement_Locally_Closed_Subsets}

The live route through the talk can stop after
\Cref{cor:Recovering_Manifolds}. For completeness, this subsection proves the
extension used by Moerdijk and Reyes: the same full-faithfulness statement
holds for locally closed subsets. The only additional task is to replace an
open subset by a closed hypersurface in one higher-dimensional Euclidean
space. The construction follows their Lemma~1.4 and Proposition~1.5
\cite[Chapter~I, Lemma~1.4 and Proposition~1.5]{Moerdijk_Reyes_Smooth_Infinitesimal_Analysis}.

\begin{lemma}[Open subsets are smooth cozero loci]
  \label[lemma]{lem:Open_Subset_Cozero_Locus}
  For every open subset $U\subseteq\R^n$, there is a smooth function
  $f\colon\R^n\to[0,\infty)$ such that
  \[
    U=\{x\in\R^n\mid f(x)\neq0\}.
  \]
\end{lemma}

\begin{proof}
  Choose a countable collection of open balls $B_k$ which covers $U$ and whose
  closures lie in $U$. For each $k$, choose a nonnegative smooth function
  $\rho_k$ supported in $U$ and strictly positive on $B_k$. Choose constants
  $\varepsilon_k>0$ so small that
  \[
    \varepsilon_k
    \sup_{x\in\R^n}|D^\alpha\rho_k(x)|
    \leq 2^{-k}
  \]
  for every multi-index $\alpha$ with $|\alpha|\leq k$. Then
  \[
    f=\sum_{k\geq1}\varepsilon_k\rho_k
  \]
  converges together with all its derivatives and hence defines a smooth
  function. It vanishes outside $U$. At every point of $U$, at least one
  $\rho_k$ is positive, so $f$ is positive there.
\end{proof}

\begin{proposition}[Closed realization]
  \label[proposition]{prop:Locally_Closed_Closed_Realization}
  Every locally closed subset $X\subseteq\R^n$ is diffeomorphic, in the sense
  of \Cref{def:Smooth_Functions_On_Subsets}, to a closed subset of
  $\R^{n+1}$.
\end{proposition}

\begin{proof}
  First let $U\subseteq\R^n$ be open, and choose $f$ as in
  \Cref{lem:Open_Subset_Cozero_Locus}. The graph
  \[
    H_f
    =\{(x,y)\in\R^n\times\R\mid y f(x)=1\}
  \]
  is closed, since it is the zero locus of the continuous function
  $(x,y)\mapsto y f(x)-1$. The map
  \[
    U\longrightarrow H_f,
    \qquad
    x\longmapsto\bigl(x,1/f(x)\bigr)
  \]
  is a diffeomorphism with inverse the first projection.

  Now write the locally closed subset as $X=F\cap U$, with $F$ closed in
  $\R^n$ and $U$ open. Under the displayed diffeomorphism, $X$ corresponds to
  \[
    H_f\cap(F\times\R),
  \]
  which is closed in $\R^{n+1}$.
\end{proof}

\begin{corollary}[Moerdijk--Reyes]
  \label[corollary]{cor:Full_Faithfulness_Locally_Closed}
  The contravariant functor from locally closed subsets of Euclidean spaces and
  smooth maps to finitely generated $\Cinfty$-rings is fully faithful.
\end{corollary}

\begin{proof}
  Replace the source and target by the closed realizations from
  \Cref{prop:Locally_Closed_Closed_Realization}, and apply
  \Cref{thm:Full_Faithfulness_Closed_Subsets}. Their function rings are
  finitely generated because each closed realization is a quotient of a free
  $\Cinfty$-ring.
\end{proof}

\subsection{Related literature}
\label[subsection]{sec:Cinfty_Related_Literature}

The primary source for this chapter is Moerdijk--Reyes, Chapter~I, Section~1
\cite[Chapter~I, Section~1]{Moerdijk_Reyes_Smooth_Infinitesimal_Analysis}.
It contains the free-ring theorem, quotient construction, Borel's theorem, the
cozero-locus lemma, and full faithfulness for locally closed subsets. Our proof
of the closed-subset case makes the ambient-coordinate and bump-function
argument more explicit than the source.

Joyce gives a concise modern treatment of $C^\infty$-rings, finite generation
and presentation, Weil algebras, and the passage to $C^\infty$-schemes
\cite[Sections~2--3]{Joyce_Introduction_Cinfty_Schemes}. The latter passage is
deliberately postponed to Talk 7, after the universal construction, the first
affine derived calculations, and their stable linearization. Spivak's
construction shows how ordinary
$C^\infty$-rings are replaced by their homotopical counterparts in one model
of derived smooth manifolds
\cite[Sections~5--7]{Spivak_Derived_Smooth_Manifolds}.

% !TEX root = ../main.tex

\chapter{Geometric constructions in \texorpdfstring{$\Cinfty$}{C-infinity}-rings}
\label[chapter]{chap:Manifolds_As_Cinfty_Rings}

The preceding chapter showed that a smooth manifold and all of its smooth maps
are encoded by its ring of smooth functions. We now ask whether geometric
constructions are encoded as well. Since the functor
\[
  \Cinfty(-)\colon
  (\text{smooth manifolds})\catop
  \longrightarrow
  \Cinfty\text{-}\mathrm{Ring}
\]
is contravariant, the first test is whether a pullback of manifolds becomes a
pushout of $\Cinfty$-rings.

\begin{question}
  \label[question]{question:When_Pushout_Is_Manifold}
  When is the algebraic pushout of smooth function rings again the function
  ring of the geometric pullback?
\end{question}

Transversality is a sufficient condition. To understand why, we first show
that independent smooth equations generate the full ideal of functions
vanishing on their common zero locus. This lets us translate geometric
constructions into presentations of smooth rings.

Two intermediate results prepare the general pullback theorem. Every
manifold has a finitely presented function ring, and products and open
restrictions have the expected algebraic descriptions. We then prove the
transverse-pullback theorem locally and explain its consequences. The final
proof section supplies the globalization argument. Together with the full
faithfulness proved in the preceding chapter, these results identify
manifolds inside smooth algebra and explain which intersections that
embedding preserves.

\section{Equations and zero loci}
\label[section]{sec:Equations_And_Coproducts}

Return to the intersections in \Cref{fig:Two_Intersections}. Both the
transverse axes and the tangent line and parabola intersect in a single
point. Their equations nevertheless give different algebraic pushouts.
We begin by making this calculation precise, then ask when a quotient by
equations is the function ring of their zero locus.

\subsection{Coproducts impose simultaneous equations}
\label[subsection]{sec:Coproducts_Impose_Equations}

We first fix the algebraic notation. The category of $\Cinfty$-rings has all
small colimits. Following Moerdijk and Reyes, we write
\[
  B\otimes_\infty C
\]
for the coproduct of two $\Cinfty$-rings. This notation records the
categorical role of the construction. It is not, in general, the tensor
product of the underlying commutative $\R$-algebras.

\begin{proposition}[Coproducts of free rings and quotients]
  \label[proposition]{prop:Cinfty_Coproduct_Formulas}
  The following formulas hold.
  \begin{enumerate}
    \item There is a natural isomorphism
      \[
        \Cinfty(\R^n)\otimes_\infty\Cinfty(\R^m)
        \cong \Cinfty(\R^{n+m}).
      \]
    \item If $I\subseteq B$ and $J\subseteq C$ are ideals, then
      \[
        (B/I)\otimes_\infty(C/J)
        \cong
        (B\otimes_\infty C)/(I,J),
      \]
      where $(I,J)$ denotes the ideal generated by the images of $I$ and $J$.
    \item If $I,J\subseteq B$, then
      \[
        (B/I)\mathbin{\amalg_B}(B/J)
        \cong B/(I+J).
      \]
  \end{enumerate}
\end{proposition}

\begin{proof}
  A homomorphism from $\Cinfty(\R^n)\otimes_\infty\Cinfty(\R^m)$ to a
  $\Cinfty$-ring $D$ is a pair of homomorphisms from the two factors. By
  \Cref{thm:Free_Cinfty_Rings}, this is the same as a pair consisting of an
  $n$-tuple and an $m$-tuple of elements of $D$, hence the same as an
  $(n+m)$-tuple. The free-ring theorem identifies the representing object as
  $\Cinfty(\R^{n+m})$.

  A homomorphism from the coproduct in the second formula to $D$ is a pair of
  homomorphisms $B\to D$ and $C\to D$ which kill $I$ and $J$, respectively.
  By the universal properties of the coproduct and the quotient, this is the
  same as a homomorphism from $(B\otimes_\infty C)/(I,J)$ to $D$. The third
  formula is proved in the same way: a compatible pair of maps from $B/I$ and
  $B/J$ is a map from $B$ killing both ideals.
\end{proof}

Set
\[
  A=\Cinfty(\R^2).
\]
For the two coordinate axes, the function rings are $A/(x)$ and $A/(y)$.
Imposing both equations gives
\[
  A/(x)\mathbin{\amalg_A}A/(y)
  \cong A/(x,y)
  \cong \R.
\]
This is the function ring of their intersection point. For the $x$-axis and
the parabola $y=x^2$, the analogous calculation begins
\[
  A/(y)\mathbin{\amalg_A}A/(y-x^2)
  \cong A/(y,y-x^2).
\]
Eliminating $y$ gives $\Cinfty(\R)/(x^2)$, the dual-number ring.

The first calculation recovers the function ring of the geometric
intersection; the second retains a nonzero nilpotent element. The
regular-equation theorem below explains why independent equations
recover the function ring of the zero locus.

\subsection{Regular equations}
\label[subsection]{sec:Regular_Equations}

Let $M$ be a smooth manifold and let $g_1,\ldots,g_q\in\Cinfty(M)$. Write
\[
  Z(g_1,\ldots,g_q)
  =\{p\in M\mid g_1(p)=\cdots=g_q(p)=0\}
\]
and
\[
  \mathfrak m_Z
  =\{h\in\Cinfty(M)\mid h|_Z=0\}.
\]
The inclusion $(g_1,\ldots,g_q)\subseteq\mathfrak m_Z$ is automatic. Equality
is a geometric assertion about the equations.

We first isolate the globalization argument which will occur several times.

\begin{lemma}[Local ideal membership]
  \label[lemma]{lem:Local_Ideal_Membership}
  Let $h,g_1,\ldots,g_q\in\Cinfty(M)$. If every point of $M$ has an open
  neighborhood $U$ such that
  \[
    h|_U\in(g_1|_U,\ldots,g_q|_U),
  \]
  then $h\in(g_1,\ldots,g_q)$ in $\Cinfty(M)$.
\end{lemma}

\begin{proof}
  Choose a locally finite open cover $\{U_\alpha\}$ on which
  \[
    h|_{U_\alpha}
    =\sum_{i=1}^q a_{\alpha i}g_i|_{U_\alpha},
  \]
  and choose a smooth partition of unity $\{\rho_\alpha\}$ subordinate to
  this cover. Since $\operatorname{supp}(\rho_\alpha)\subseteq U_\alpha$, the
  product $\rho_\alpha a_{\alpha i}$ extends by zero to a smooth function on
  $M$. Local finiteness then gives smooth functions
  \[
    a_i=\sum_\alpha\rho_\alpha a_{\alpha i}\in\Cinfty(M).
  \]
  Since $\sum_\alpha\rho_\alpha=1$, we obtain
  \[
    h
    =\sum_\alpha\rho_\alpha h
    =\sum_{i=1}^q a_i g_i.
  \]
\end{proof}

\begin{theorem}[Regular equations]
  \label[theorem]{thm:Regular_Equations}
  Let $g=(g_1,\ldots,g_q)\colon M\to\R^q$, and let $Z=g^{-1}(0)$. Suppose
  that
  \[
    D_pg\colon T_pM\longrightarrow\R^q
  \]
  is surjective for every $p\in Z$. Then
  \[
    \mathfrak m_Z=(g_1,\ldots,g_q).
  \]
\end{theorem}

This is \citet[Chapter~I, Lemma~2.1]{Moerdijk_Reyes_Smooth_Infinitesimal_Analysis}.

\begin{proof}
  Let $h\in\mathfrak m_Z$. We prove ideal membership locally and apply
  \Cref{lem:Local_Ideal_Membership}.

  Suppose first that $p\notin Z$. Some $g_i(p)$ is nonzero, so after shrinking
  to a neighborhood $U$ of $p$, the function $g_i|_U$ is invertible. Hence
  \[
    h|_U=(h/g_i)g_i|_U.
  \]

  Now suppose that $p\in Z$. The submersion theorem gives local coordinates
  \[
    (x_1,\ldots,x_q,y_1,\ldots,y_{d-q})
  \]
  centered at $p$ in which $g_i=x_i$ for $1\leq i\leq q$. After shrinking,
  we may take the coordinate neighborhood to be a product which is convex in
  the $x$-variables. In these coordinates, $h(0,y)=0$. The parameterized form
  of Hadamard's lemma gives
  \[
    h(x,y)
    =h(0,y)+\sum_{i=1}^q x_i
      \int_0^1\frac{\partial h}{\partial x_i}(tx,y)\,dt
    =\sum_{i=1}^q x_i h_i(x,y)
  \]
  for smooth functions $h_i$. Thus $h$ belongs locally to the ideal generated
  by the $g_i$. The local ideal-membership lemma finishes the proof.
\end{proof}

The theorem says more than the regular-value theorem. The latter identifies
$Z$ as a submanifold; the former identifies the entire ideal of functions
vanishing on $Z$.

\begin{corollary}[Function ring of a regular zero locus]
  \label[corollary]{cor:Function_Ring_Regular_Zero_Locus}
  Under the hypotheses of \Cref{thm:Regular_Equations}, restriction induces an
  isomorphism
  \[
    \Cinfty(M)/(g_1,\ldots,g_q)\cong\Cinfty(Z).
  \]
\end{corollary}

\begin{proof}
  The zero locus $Z$ is a closed embedded submanifold of $M$. Every smooth
  function on $Z$ extends to $M$: extend it along a tubular-neighborhood
  retraction, multiply by a cutoff supported in that neighborhood, and extend
  by zero. Hence the restriction map $\Cinfty(M)\to\Cinfty(Z)$ is surjective.
  Its kernel is $\mathfrak m_Z$, which is $(g_1,\ldots,g_q)$ by
  \Cref{thm:Regular_Equations}.
\end{proof}

The regularity hypothesis is essential. For $g(t)=t^2$, the zero locus is the
origin, but
\[
  (t^2)\subsetneq(t)=\mathfrak m_{\{0\}}.
\]
Consequently,
\[
  \Cinfty(\R)/(t^2)\not\cong\Cinfty(\{0\}).
\]
This is exactly the defect seen in the tangent
line--parabola intersection.

\section{Manifold function rings are finitely presented}
\label[section]{sec:Finite_Presentation_And_Products}

Our first application is a finiteness theorem: every manifold function ring
can be described by finitely many generators and smooth relations. Here
finite presentation is understood in the sense of
\Cref{def:Cinfty_Finiteness}.

\begin{example}[The circle]
  \label[example]{ex:Circle_Finite_Presentation}
  The circle has the finite presentation
  \[
    \Cinfty(S^1)\cong\Cinfty(\R^2)/(x^2+y^2-1).
  \]
  The gradient $(2x,2y)$ is nonzero on the circle, so this follows from
  the regular-equation theorem.
\end{example}

For a general manifold, a short global list of regular equations need not
be available. We will instead use a tubular neighborhood and a retraction.

The first step is to present the function ring of an arbitrary open subset
of Euclidean space. One extra generator and one relation suffice,
regardless of the shape of the open subset.

\begin{proposition}[One-relation presentation of an open subset]
  \label[proposition]{prop:Open_Subset_One_Relation}
  Let $U\subseteq\R^n$ be open. Choose a smooth function
  $a\colon\R^n\to[0,\infty)$ such that $U=\{a\neq0\}$. Then
  \[
    \Cinfty(U)
    \cong
    \Cinfty(\R^{n+1})/(ya(x)-1).
  \]
\end{proposition}

This is \citet[Chapter~I, Corollary~2.2]{Moerdijk_Reyes_Smooth_Infinitesimal_Analysis}.

\begin{proof}
  The map
  \[
    U\longrightarrow\R^{n+1},
    \qquad
    x\longmapsto\bigl(x,1/a(x)\bigr)
  \]
  is a diffeomorphism onto the hypersurface
  \[
    H_a=\{(x,y)\mid ya(x)-1=0\}.
  \]
  Along $H_a$, the partial derivative of $ya(x)-1$ with respect to $y$ is
  $a(x)$, which is nonzero. Thus the defining equation is regular, and
  \Cref{cor:Function_Ring_Regular_Zero_Locus} gives
  \[
    \Cinfty(U)\cong\Cinfty(H_a)
    \cong\Cinfty(\R^{n+1})/(ya(x)-1).
  \]
\end{proof}

A properly embedded manifold is a retract of an open tubular neighborhood
in Euclidean space. On function rings, it is therefore enough to know that
finite presentation passes to retracts.

\begin{lemma}[Retracts of finitely presented rings]
  \label[lemma]{lem:Retracts_Finitely_Presented_Cinfty_Rings}
  A retract of a finitely presented $\Cinfty$-ring is finitely presented.
\end{lemma}

\begin{proof}
  Let $A$ be a retract of
  \[
    P=\Cinfty(\R^n)/(r_1,\ldots,r_k),
  \]
  with maps $s\colon A\to P$ and $p\colon P\to A$ satisfying
  $p\circ s=\operatorname{id}_A$. The composite $e=s\circ p$ is an
  idempotent endomorphism of $P$ whose image is $s(A)$. Choose smooth
  functions $e_i\in\Cinfty(\R^n)$ representing the images under $e$ of the
  coordinate generators $x_i$ of $P$. Then $A$ is the coequalizer of
  $\operatorname{id}_P$ and $e$, and hence
  \[
    A\cong
    \Cinfty(\R^n)/
    (r_1,\ldots,r_k,x_1-e_1,\ldots,x_n-e_n).
  \]
  Indeed, a homomorphism $u\colon P\to D$ factors through $p$ precisely when
  $u=u\circ e$, and it is enough to check this equality on the coordinate
  generators. The displayed quotient therefore has the universal property of
  $A$.
\end{proof}

\begin{theorem}[Finite presentation for manifolds]
  \label[theorem]{thm:Manifold_Function_Ring_Finitely_Presented}
  If $M$ is a finite-dimensional, second-countable smooth manifold without
  boundary, then $\Cinfty(M)$ is a finitely presented $\Cinfty$-ring.
\end{theorem}

This is \citet[Chapter~I, Theorem~2.3]{Moerdijk_Reyes_Smooth_Infinitesimal_Analysis}.

\begin{proof}
  Choose a proper smooth embedding $i\colon M\hookrightarrow\R^N$. Its image
  is closed. The tubular-neighborhood theorem supplies an open neighborhood
  $U\subseteq\R^N$ of $M$ and a smooth retraction $r\colon U\to M$. Since
  $r\circ i=\operatorname{id}_M$, the induced maps
  \[
    \Cinfty(M)
    \xrightarrow{r^*}
    \Cinfty(U)
    \xrightarrow{i^*}
    \Cinfty(M)
  \]
  have composite the identity. By \Cref{prop:Open_Subset_One_Relation}, the
  middle ring is finitely presented. The result follows from
  \Cref{lem:Retracts_Finitely_Presented_Cinfty_Rings}.
\end{proof}

The theorem does not say that the underlying commutative $\R$-algebra
$\Cinfty(M)$ is finitely presented. All finiteness statements take place in
the algebraic theory whose operations are arbitrary smooth functions.

\section{Products and open restrictions}
\label[section]{sec:Products_Become_Coproducts}

The same passage from open subsets to manifolds also computes products.
Once products are understood, the equation $ya=1$ identifies restriction
to the open subset $\{a\neq0\}$ with adjoining an inverse to $a$. These
two constructions will let us reduce transverse pullbacks to local
equations and then compare the calculations on overlapping charts.

For manifolds $M$ and $P$, pullback along the two projections determines a
natural comparison homomorphism
\[
  \theta_{M,P}\colon
  \Cinfty(M)\otimes_\infty\Cinfty(P)
  \longrightarrow
  \Cinfty(M\times P).
\]

\begin{lemma}[Products of open subsets]
  \label[lemma]{lem:Products_Open_Subsets}
  If $U\subseteq\R^n$ and $V\subseteq\R^m$ are open, then
  $\theta_{U,V}$ is an isomorphism.
\end{lemma}

This is \citet[Chapter~I, Lemma~2.4]{Moerdijk_Reyes_Smooth_Infinitesimal_Analysis}.

\begin{proof}
  Choose smooth functions $a$ and $b$ with
  $U=\{a\neq0\}$ and $V=\{b\neq0\}$. By
  \Cref{prop:Open_Subset_One_Relation,prop:Cinfty_Coproduct_Formulas},
  \begin{align*}
    \Cinfty(U)\otimes_\infty\Cinfty(V)
    &\cong
    \frac{\Cinfty(\R^{n+m+2})}
      {(ya(x)-1,zb(w)-1)}.
  \end{align*}
  The two displayed equations have independent differentials along their
  common zero locus, since the coefficients of $dy$ and $dz$ are $a(x)$ and
  $b(w)$ there. Their common zero locus is
  \[
    H_a\times H_b\cong U\times V.
  \]
  The regular-equation theorem therefore identifies the displayed quotient
  with $\Cinfty(U\times V)$. Under these identifications the isomorphism is
  the natural comparison map $\theta_{U,V}$.
\end{proof}

\begin{proposition}[Products of manifolds]
  \label[proposition]{prop:Products_Manifolds_Coproducts}
  For smooth manifolds $M$ and $P$, the natural map
  \[
    \theta_{M,P}\colon
    \Cinfty(M)\otimes_\infty\Cinfty(P)
    \iso
    \Cinfty(M\times P)
  \]
  is an isomorphism.
\end{proposition}

This is \citet[Chapter~I, Proposition~2.5]{Moerdijk_Reyes_Smooth_Infinitesimal_Analysis}.

\begin{proof}
  Choose proper embeddings of $M$ and $P$, open tubular neighborhoods
  $U\supseteq M$ and $V\supseteq P$, and retractions $U\to M$ and $V\to P$.
  Thus $(M,P)$ is a retract of $(U,V)$. Naturality of $\theta$ gives a
  corresponding retraction diagram in the arrow category of $\Cinfty$-rings:
  $\theta_{M,P}$ is a retract of $\theta_{U,V}$. The latter is an isomorphism
  by \Cref{lem:Products_Open_Subsets}. A retract of an isomorphism is an
  isomorphism, so the result follows.
\end{proof}

This result is the first complete answer to
\Cref{question:When_Pushout_Is_Manifold}: a product is a pullback over a point,
and its function ring is always the corresponding coproduct.

Every open subset of a manifold can be written as $\{a\neq0\}$ for some
smooth function $a$. We now extend the one-relation presentation of an
open subset of Euclidean space to this setting.

\begin{proposition}[Smooth localization]
\label[proposition]{prop:Smooth_Localization}
Let $M$ be a manifold, let $a\in\Cinfty(M)$, and put $U=\{a\neq0\}$.
Restriction exhibits $\Cinfty(U)$ as the universal $\Cinfty$-ring under
$\Cinfty(M)$ in which $a$ becomes invertible.
\end{proposition}

\begin{proof}
By the product theorem, adjoining one free smooth generator $y$ gives
\[
 \Cinfty(M)\otimes_\infty\Cinfty(\R)\cong\Cinfty(M\times\R).
\]
The quotient by $ya-1$ therefore universally makes $a$ invertible.
The function $(m,y)\mapsto ya(m)-1$ is regular along its zero set, since its
derivative in $y$ is the nonzero number $a(m)$. That zero manifold is the
graph $m\mapsto(m,1/a(m))$ over $U$. The regular-equation theorem identifies
the quotient with $\Cinfty(U)$, and its structure map with restriction.
\end{proof}

For $M=\R^n$, this is
\citet[Chapter~I, Proposition~1.6 and Corollary~2.2]
{Moerdijk_Reyes_Smooth_Infinitesimal_Analysis}.
The argument extends it to manifolds using the product calculation in
\Cref{prop:Products_Manifolds_Coproducts}.

\section{Transverse pullbacks}
\label[section]{sec:Fiber_Products_And_The_Boundary_Of_Transversality}

We now return to \Cref{question:When_Pushout_Is_Manifold}. The essential
case is the inverse image of an embedded submanifold: transversality
ensures that its local defining equations remain independent after
pullback. General transverse fiber products reduce to this case by
using the diagonal.

\begin{theorem}[Transverse inverse images become pushouts]
  \label[theorem]{thm:Transverse_Inverse_Images_Pushouts}
  Let $f\colon M\to N$ be transverse to an embedded submanifold
  $i\colon Z\hookrightarrow N$, and put $P=f^{-1}(Z)$. Then
  \[
    \begin{tikzcd}
      \Cinfty(N) \rar["f^*"] \dar["i^*"'] \drar[pushout]
        & \Cinfty(M) \dar \\
      \Cinfty(Z) \rar
        & \Cinfty(P)
    \end{tikzcd}
  \]
  is a pushout of $\Cinfty$-rings.
\end{theorem}

This is \citet[Chapter~I, Proposition~2.6]{Moerdijk_Reyes_Smooth_Infinitesimal_Analysis}.

The local calculation contains the geometric content. Suppose
$U\subseteq N$ is an open set on which $Z$ is defined by functions
$g_1,\ldots,g_q$ with independent differentials along $U\cap Z$. Set
$V=f^{-1}(U)$. Transversality implies that
\[
  g_1\circ f,\ldots,g_q\circ f
\]
have independent differentials along $V\cap P$. Hence
\Cref{cor:Function_Ring_Regular_Zero_Locus} gives
\[
  \Cinfty(U\cap Z)
  \cong\Cinfty(U)/(g_1,\ldots,g_q)
\]
and
\[
  \Cinfty(V\cap P)
  \cong
  \Cinfty(V)/(g_1\circ f,\ldots,g_q\circ f).
\]
The local function-ring square is therefore isomorphic to
\[
  \begin{tikzcd}
    \Cinfty(U) \rar["f^*"] \dar \drar[pushout]
      & \Cinfty(V) \dar \\
    \Cinfty(U)/(g_1,\ldots,g_q) \rar
      & \Cinfty(V)/(g_1\circ f,\ldots,g_q\circ f).
  \end{tikzcd}
\]
It is a pushout by \Cref{prop:Cinfty_Coproduct_Formulas}: the lower-right ring
is obtained by imposing on $\Cinfty(V)$ the images of the same equations.

This calculation also identifies the exact use of transversality. It ensures
that the pulled-back equations remain regular, so that their algebraic
quotient is the function ring of the geometric inverse image.

The local calculation proves the theorem on coordinate neighbourhoods.
Globalization uses partitions of unity and descent for functions modulo
finitely many equations. We give that argument in
\Cref{sec:Further_Transverse_Proof}; first we explain the consequence
for arbitrary transverse fibre products.

We can now answer \Cref{question:When_Pushout_Is_Manifold} for general fiber
products. Let $f\colon M\to N$ and $g\colon P\to N$ be transverse. Then
\[
  f\times g\colon M\times P\longrightarrow N\times N
\]
is transverse to the diagonal $\Delta_N$, and
\[
  (f\times g)^{-1}(\Delta_N)=M\times_NP.
\]

\begin{theorem}[Transverse pullbacks become pushouts]
  \label[theorem]{thm:Transverse_Pullbacks_Pushouts}
  For transverse smooth maps $f\colon M\to N$ and $g\colon P\to N$, the two
  squares
  \[
    \begin{tikzcd}
      M\times_NP \rar \dar \drar[pullback]
        & P \dar["g"] \\
      M \rar["f"'] & N
    \end{tikzcd}
    \qquad
    \begin{tikzcd}
      \Cinfty(N) \rar["g^*"] \dar["f^*"'] \drar[pushout]
        & \Cinfty(P) \dar \\
      \Cinfty(M) \rar
        & \Cinfty(M\times_NP)
    \end{tikzcd}
  \]
  are respectively a pullback of manifolds and a pushout of $\Cinfty$-rings.
\end{theorem}

\begin{proof}
  Apply \Cref{thm:Transverse_Inverse_Images_Pushouts} to
  $f\times g$ and $\Delta_N\subseteq N\times N$. Then use
  \Cref{prop:Products_Manifolds_Coproducts} to identify the function rings of
  $M\times P$ and $N\times N$ with the corresponding coproducts. The resulting
  universal property says precisely that a map out of
  $\Cinfty(M\times_NP)$ is a pair of maps out of $\Cinfty(M)$ and
  $\Cinfty(P)$ whose restrictions along $f^*$ and $g^*$ agree. This is the
  pushout property of the square displayed above.
\end{proof}

Combining \Cref{cor:Recovering_Manifolds,thm:Manifold_Function_Ring_Finitely_Presented,thm:Transverse_Pullbacks_Pushouts}
gives the classical structural theorem.

\begin{theorem}[Manifolds inside smooth algebra]
  \label[theorem]{thm:Manifolds_Inside_Smooth_Algebra}
  The contravariant functor
  \[
    \Cinfty(-)\colon
    (\text{smooth manifolds})\catop
    \longrightarrow
    \Cinfty\text{-}\mathrm{Ring}
  \]
  is fully faithful, takes values in finitely presented $\Cinfty$-rings, and
  sends transverse pullback squares to pushout squares.
\end{theorem}

This is \citet[Chapter~I, Theorem~2.8]{Moerdijk_Reyes_Smooth_Infinitesimal_Analysis}.

The tangency calculation in \Cref{sec:Coproducts_Impose_Equations} shows
why the transversality hypothesis cannot simply be dropped. The geometric
intersection is a point, with function ring $\R$, while the pushout is
$\Cinfty(\R)/(x^2)\cong\R[\varepsilon]/(\varepsilon^2)$.
The comparison homomorphism kills $\varepsilon$: it forgets the
infinitesimal information retained by the equations.

This pushout shows that ordinary smooth algebra already contains
intersections which are not manifolds. It does not retain all intersection
data, however: the ordinary pushout for the self-intersection
$*\to\R\leftarrow *$ is just $\R$. The excess direction seen in
\Cref{ex:Excess_Self_Intersection}
requires a derived construction. The final exercise asks us to compare
these two descriptions.
The transverse-pullback theorem proved here is the compatibility
condition for both the ordinary and the derived universal constructions
in the next chapter.

\section{Further proofs: globalizing the pushout}
\label[section]{sec:Further_Transverse_Proof}

We supply the globalization used in
\Cref{thm:Transverse_Inverse_Images_Pushouts}. The local regular-equation
calculation is already established. The remaining task is to glue maps
into a test ring, which may contain infinitesimal elements invisible
on its real points.

\subsection{Test rings and descent}
\label[subsection]{sec:Finite_Presentation_And_Test_Rings}

To verify a pushout, we must test maps into arbitrary smooth rings. The
following criterion will let us restrict to finitely presented test rings,
where we can glue functions modulo a finite list of equations.

\begin{proposition}[Filtered-colimit criterion]
  \label[proposition]{prop:Filtered_Colimit_Criterion}
  A $\Cinfty$-ring $A$ is finitely presented if and only if
  \[
    \Hom_{\Cinfty\text{-}\mathrm{Ring}}(A,-)
  \]
  preserves filtered colimits.
\end{proposition}

\begin{proof}
  Suppose first that
  $A=\Cinfty(\R^n)/(r_1,\ldots,r_k)$. For every $\Cinfty$-ring $B$,
  \[
    \Hom(A,B)
    \cong
    \bigl\{b\in B^n\mid
      \Phi_{r_1}(b)=\cdots=\Phi_{r_k}(b)=0\bigr\}.
  \]
  Filtered colimits of $\Cinfty$-rings are computed on underlying sets, and
  filtered colimits of sets commute with finite limits. The displayed functor
  therefore preserves filtered colimits.

  Conversely, every $\Cinfty$-ring is a filtered colimit of finitely presented
  ones mapping to it. Concretely, finite sets of elements give maps from
  finitely generated free rings, finite sets of relations give finitely
  presented quotients, two such stages admit a common refinement by combining
  their generators and relations, and parallel maps become equal after adding
  finitely many relations on a finite set of generators. Apply this canonical
  filtered presentation to $A$. If $\Hom(A,-)$ preserves its colimit, then
  $\operatorname{id}_A$ factors through one finitely presented stage
  $P\to A$. Thus $A$ is a retract of $P$, and
  \Cref{lem:Retracts_Finitely_Presented_Cinfty_Rings} completes the proof.
\end{proof}

For a finitely presented test ring, the following descent statement
lets us glue the local maps constructed by imposing regular equations.

\begin{lemma}[Cech descent]
  \label[lemma]{lem:Cech_Descent_Finitely_Presented_Cinfty_Ring}
  Let
  \[
    B=\Cinfty(\R^k)/(h_1,\ldots,h_p),
    \qquad
    X=Z(h_1,\ldots,h_p),
  \]
  and let $\{W_\alpha\}$ be a cover of $X$ by open subsets of $\R^k$. For a
  nonempty finite set $S$ of indices, put
  \[
    W_S=\bigcap_{\alpha\in S}W_\alpha,
    \qquad
    B_S=
    \Cinfty(W_S)/(h_1|_{W_S},\ldots,h_p|_{W_S}).
  \]
  Then restriction induces an isomorphism
  \[
    B\iso\varprojlim_{S}B_S,
  \]
  where the limit is taken over the Cech diagram of nonempty finite
  intersections.
\end{lemma}

This is \citet[Chapter~I, Lemma~2.7]{Moerdijk_Reyes_Smooth_Infinitesimal_Analysis}.

\begin{proof}
  It is enough to prove the usual gluing statement. Suppose that we have
  $g_\alpha\in\Cinfty(W_\alpha)$ such that on every pairwise intersection
  \[
    g_\alpha-g_\beta
    \in(h_1,\ldots,h_p).
  \]
  We must construct a unique class $[g]\in B$ whose restriction is
  $[g_\alpha]$ for every $\alpha$.

  Add $\R^k\setminus X$ to the cover and choose a locally finite refinement
  $\{V_\lambda\}$ of the resulting cover of $\R^k$. If $V_\lambda$ lies in
  some $W_\alpha$, let $k_\lambda=g_\alpha|_{V_\lambda}$, choosing one such
  $\alpha$. If $V_\lambda$ lies in $\R^k\setminus X$, let $k_\lambda=0$.
  The choices are compatible modulo $(h_1,\ldots,h_p)$. Indeed, this follows
  from the original compatibility on overlaps of the $W_\alpha$, while on an
  open set disjoint from $X$ the function $h_1^2+\cdots+h_p^2$ is invertible,
  so $(h_1,\ldots,h_p)$ is the unit ideal.

  Let $\{\rho_\lambda\}$ be a partition of unity subordinate to the refined
  cover and define
  \[
    g=\sum_\lambda\rho_\lambda k_\lambda,
  \]
  extending each summand by zero. This is a globally defined smooth function.
  Fix an original index $\beta$. Near any point of $W_\beta$, only finitely
  many $\rho_\lambda$ are nonzero, so
  \[
    g-g_\beta
    =\sum_\lambda\rho_\lambda(k_\lambda-g_\beta)
  \]
  belongs locally to $(h_1,\ldots,h_p)$. By
  \Cref{lem:Local_Ideal_Membership}, it belongs to that ideal on
  $W_\beta$. Thus $[g]$ restricts to $[g_\beta]$.

  If $g'$ is another choice, then $g-g'$ belongs locally to
  $(h_1,\ldots,h_p)$ on a cover of $\R^k$. Another application of
  \Cref{lem:Local_Ideal_Membership} shows that $g-g'$ belongs to the global
  ideal. Hence $[g]=[g']$ in $B$. The construction respects every smooth
  operation, since those operations are computed componentwise in the limit.
\end{proof}

We also record how a homomorphism into a finitely presented test ring restricts
to an open cover. This is the bridge between the local geometric calculation
and \Cref{lem:Cech_Descent_Finitely_Presented_Cinfty_Ring}.

\begin{lemma}[Restricting a test-ring map]
  \label[lemma]{lem:Restricting_Test_Ring_Map}
  Let $\eta\colon\Cinfty(M)\to B$ be a homomorphism, where
  $B=\Cinfty(\R^k)/(h_1,\ldots,h_p)$, and let $\{V_\alpha\}$ be an open cover
  of $M$. There is an open cover $\{W_\alpha\}$ of
  $X=Z(h_1,\ldots,h_p)$ such that $\eta$ induces compatible homomorphisms
  \[
    \eta_S\colon
    \Cinfty(V_S)\longrightarrow B_S
  \]
  on every nonempty finite intersection.
\end{lemma}

\begin{proof}
  Choose $a_\alpha\in\Cinfty(M)$ with
  $V_\alpha=\{a_\alpha\neq0\}$. Represent $\eta(a_\alpha)$ by a smooth
  function $\widetilde a_\alpha$ on $\R^k$, and set
  \[
    W_\alpha
    =\{\widetilde a_\alpha\neq0\}.
  \]
  These open subsets of $\R^k$ cover $X$. Indeed, evaluation at $x\in X$ gives a real point
  $B\to\R$, and its composite with $\eta$ is evaluation at a point
  $m_x\in M$. Since the $V_\alpha$ cover $M$, some $a_\alpha(m_x)$ is
  nonzero, and this value equals $\widetilde a_\alpha(x)$.

  In $B_S$, each image of $\eta(a_\alpha)$ for $\alpha\in S$ is invertible.
  The universal property of adjoining smooth inverses in
  \Cref{prop:Smooth_Localization} extends
  $\eta$ uniquely to $\Cinfty(V_S)\to B_S$. Uniqueness makes these extensions
  compatible under restriction.
\end{proof}

\subsection{Globalizing the local calculation}
\label[subsection]{sec:Transverse_Inverse_Image_Theorem}

\begin{proof}[Proof of \Cref{thm:Transverse_Inverse_Images_Pushouts}]
  We first reduce to a closed embedding. An embedded submanifold $Z$
  is locally closed in $N$, so it is closed in some open neighbourhood
  $N'\subseteq N$. The inclusions $N'\to N$ and $f^{-1}(N')\to M$
  give smooth localizations. Their function-ring square is a pushout
  by the universal property of localization. Pasting with the square
  over $N'$ therefore reduces the assertion to the case where $Z$
  is closed in $N$.


  For a $\Cinfty$-ring $B$, consider the natural map
  \begin{align*}
    \Hom(\Cinfty(P),B)
    \longrightarrow
    \Hom(\Cinfty(M),B)
    \mathbin{\times_{\Hom(\Cinfty(N),B)}}
    \Hom(\Cinfty(Z),B).
  \end{align*}
  All four manifold function rings are finitely presented by
  \Cref{thm:Manifold_Function_Ring_Finitely_Presented}. Both sides therefore
  preserve filtered colimits in $B$, by
  \Cref{prop:Filtered_Colimit_Criterion} and the fact that filtered colimits of
  sets commute with finite limits. Since every $\Cinfty$-ring is a filtered
  colimit of finitely presented ones, it is enough to prove that the displayed
  map is a bijection when
  \[
    B=\Cinfty(\R^k)/(h_1,\ldots,h_p).
  \]

  Let homomorphisms
  \[
    \varphi\colon\Cinfty(M)\to B,
    \qquad
    \psi\colon\Cinfty(Z)\to B
  \]
  agree after restriction from $\Cinfty(N)$. Choose an open cover
  $\{U_\alpha\}$ of $N$ such that $U_\alpha\cap Z$ is cut out in $U_\alpha$
  by independent functions. Put $V_\alpha=f^{-1}(U_\alpha)$. The local
  calculation in \Cref{sec:Fiber_Products_And_The_Boundary_Of_Transversality} shows that the
  function-ring square over every nonempty finite intersection $U_S$ is a
  pushout.

  Choose $b_\alpha\in\Cinfty(N)$ with
  $U_\alpha=\{b_\alpha\neq0\}$. Then $V_\alpha$ is the cozero locus of
  $b_\alpha\circ f$, while $U_\alpha\cap Z$ is the cozero locus of
  $b_\alpha|_Z$. The compatibility of $\varphi$ and $\psi$ says that these two
  elements have the same image in $B$. Consequently,
  \Cref{lem:Restricting_Test_Ring_Map} produces one open cover
  $\{W_\alpha\}$ of $X=Z(h_1,\ldots,h_p)$ on which both maps restrict. We
  obtain compatible
  local maps
  \[
    \Cinfty(V_S)\longrightarrow B_S,
    \qquad
    \Cinfty(U_S\cap Z)\longrightarrow B_S.
  \]
  The local pushout property gives a unique homomorphism
  \[
    \chi_S\colon\Cinfty(V_S\cap P)\longrightarrow B_S.
  \]
  These homomorphisms are compatible on intersections, by their uniqueness.

  For $a\in\Cinfty(P)$, apply $\chi_S$ to the restriction of $a$ to
  $V_S\cap P$. The resulting elements form a compatible family in the Cech
  diagram $\{B_S\}$. By
  \Cref{lem:Cech_Descent_Finitely_Presented_Cinfty_Ring}, they determine a
  unique element $\chi(a)\in B$. This construction commutes with every smooth
  operation, so it defines a $\Cinfty$-homomorphism
  \[
    \chi\colon\Cinfty(P)\longrightarrow B.
  \]
  Locally, $\chi$ has the required composites with the restriction maps from
  $M$ and $Z$; descent makes these equalities global. The same descent argument
  shows uniqueness. Thus the displayed map of Hom-sets is a bijection.
\end{proof}

\section{Further calculations}
\label[section]{sec:Manifolds_As_Cinfty_Summary_And_Further_Reading}

\begin{example}[Smooth coproducts versus algebraic tensor products]
The difference between a smooth coproduct and an algebraic tensor product
is visible in one function. The image of
\[
 \Cinfty(\R)\otimes_\R\Cinfty(\R)\longrightarrow\Cinfty(\R^2)
\]
consists of finite sums $\sum_jf_j(x)g_j(y)$. The function $e^{xy}$
cannot be written in this form.

Indeed, if it were a sum of $r$ such terms, choose $r+1$ distinct real
numbers $x_0,\ldots,x_r$. All functions $e^{x_i y}$ would lie in the
span of $g_1,\ldots,g_r$ and hence be linearly dependent.
But differentiating a linear relation at $y=0$ in orders $0,\ldots,r$
gives a Vandermonde system with invertible matrix $(x_i^j)$.
The functions are therefore linearly independent, a contradiction.
The smooth coproduct includes $e^{xy}$ because it permits smooth
substitution involving variables from both factors.
\end{example}

\begin{exercise}
Give a finite presentation of $\Cinfty(\R\setminus\{0\})$ and use
it to calculate its coproduct with $\Cinfty(S^1)$.
\end{exercise}

\begin{exercise}
Show that the pushout for the self-intersection $*\to\R\leftarrow*$
in ordinary smooth rings is $\R$. Compare the result with the
linearized matching complex $[0\to\R]$ of \Cref{chap:Why_Derived_Manifolds}.
\end{exercise}

The proofs follow
\cite[Chapter~I, Section~2]{Moerdijk_Reyes_Smooth_Infinitesimal_Analysis}.
The differential-topological ingredients are standard tubular
neighbourhood, partition-of-unity, and submersion results; see
\cite{lee2000smooth}.
The transverse compatibility becomes part of the universal property in
\cite{Carchedi_Steffens_Universal_Derived_Manifolds}.

% !TEX root = ../main.tex

\chapter{Derived manifolds}
\label[chapter]{chap:From_Smooth_Algebra_To_Derived_Manifolds}

The preceding chapter established the structural theorem
\[
  \Cinfty(-)\colon
  \mathrm{Mfd}\catop
  \longrightarrow
  \Cinfty\text{-}\mathrm{Ring}^{\mathrm{fp}}.
\]
This functor is fully faithful, and it sends transverse pullback squares of
manifolds to pushout squares of $\Cinfty$-rings. Thus smooth algebra already
contains every manifold and every transverse intersection. It contains more:
the nontransverse line--parabola intersection produced the finitely presented
ring
\[
  \Cinfty(\R)/(x^2),
\]
which is not the function ring of a manifold.

These examples suggest a universal enlargement of the category of manifolds.

\begin{question}
  \label[question]{question:Universal_Finite_Limit_Completion}
  What is the universal way to enlarge manifolds so that all finite limits
  exist, while preserving transverse pullbacks and the terminal object?
\end{question}

The answer depends on the categories in which we ask the question.
Ordinary smooth rings retain nilpotent thickenings, but their pushouts
collapse the excess self-intersections seen in
\Cref{ex:Excess_Self_Intersection}. To retain that information, we pass
to $\infty$-categories, where morphisms carry identifications between them,
identifications between those identifications, and so on. Pullbacks can
then retain this higher information.

We first describe the universal construction for ordinary smooth rings
and locate its limitation in a self-intersection calculation. We then
introduce the pullback formalism in $\infty$-categories and extend smooth
algebra to this setting. The final step relates the resulting animated smooth rings
to a universal category of derived manifolds. We quote the algebraic
and geometric extension theorems, explaining their content through
presentations and examples. The next chapter will calculate the resulting
derived intersections.

\section{The ordinary finite-limit envelope}
\label[section]{sec:Ordinary_Finite_Limit_Envelope}

We begin with Euclidean spaces and freely adjoin finite limits while
preserving their finite products. The
result is the opposite category of finitely presented smooth rings:
generators specify coordinates, and relations specify equations.
Its universal property says that a finite-limit-preserving functor out
of this category is determined by an interpretation of smooth operations
in its target.

Recall that $\mathrm{Euc}$ consists of Euclidean spaces and smooth
maps. In \Cref{prop:Cinfty_Ring_Functor_Formulation}, an ordinary
$\Cinfty$-ring was a finite-product-preserving functor
$\mathrm{Euc}\to\mathrm{Set}$. The same definition works in any ordinary
category with finite products.

\begin{definition}[$\Cinfty$-ring object]
  \label[definition]{def:Internal_Cinfty_Ring}
  Let $C$ be an ordinary category with finite products. A \emph{$\Cinfty$-ring object in
  $C$} is a finite-product-preserving functor
  \[
    R\colon\mathrm{Euc}\longrightarrow C.
  \]
  We write $\operatorname{Alg}_{\Cinfty}(C)$ for the category of
  $\Cinfty$-ring objects in $C$.
\end{definition}

The object $R(\R)$ carries an operation in $C$ for every smooth map between
Euclidean spaces. Since $R(\R^n)\cong R(\R)^n$, the functor is determined by
this one object together with its smooth operations. For $C=\mathrm{Set}$ we
recover ordinary $\Cinfty$-rings.

A finite-limit-preserving functor transports such a ring object,
because it preserves finite products and the identities between the
smooth operations. We will use this to characterize an extension
from Euclidean spaces by its value on the smooth line.

To describe the universal target for these operations, set
\[
  \operatorname{Aff}_{\Cinfty}^{\mathrm{fp}}
  =
  (\Cinfty\text{-}\mathrm{Ring}^{\mathrm{fp}})\catop.
\]
An object of this opposite category is an affine smooth object described
by its function ring. Its spectrum and structure sheaf will be constructed
in \Cref{chap:Local_Derived_Geometry_Kuranishi_Charts}.

There is a functor
\[
  j\colon\mathrm{Euc}
  \longrightarrow
  \operatorname{Aff}_{\Cinfty}^{\mathrm{fp}}
\]
which sends $\R^n$ to the free ring $\Cinfty(\R^n)$, regarded as an object of
the opposite category. A smooth map $f\colon\R^n\to\R^m$ induces
\[
  f^*\colon\Cinfty(\R^m)\longrightarrow\Cinfty(\R^n),
\]
which becomes a morphism in the same direction as $f$ after taking opposites.
The functor $j$ preserves finite products.

The category of formal affine smooth objects has the following universal
property.

\begin{theorem}[Ordinary finite-limit envelope]
  \label[theorem]{thm:Ordinary_Cinfty_Finite_Limit_Envelope}
  Let $C$ be an ordinary category with finite limits. Restriction along $j$
  induces an equivalence of categories
  \[
    \Fun^{\mathrm{lex}}
    \bigl(\operatorname{Aff}_{\Cinfty}^{\mathrm{fp}},C\bigr)
    \simeq
    \operatorname{Alg}_{\Cinfty}(C),
  \]
  where the left side consists of finite-limit-preserving functors.
\end{theorem}

This is the $0$-truncated case of the finite-limit-envelope theorem for an
algebraic theory
\cite[Theorem~3.22 and Remark~3.26]{Carchedi_Steffens_Universal_Derived_Manifolds}.

The universal property has an explicit consequence: a functor preserving
finite limits must send a presentation by equations to the corresponding
zero locus in its target. A finitely presented $\Cinfty$-ring can be written
\[
  A=
  \frac{\Cinfty(\R^n)}{(f_1,\ldots,f_k)}.
\]
Equivalently, $A$ is the coequalizer
\[
  \begin{tikzcd}[column sep=large]
    \Cinfty(\R^k)
      \ar[r,shift left=.6ex,"{(f_1,\ldots,f_k)^*}"]
      \ar[r,shift right=.6ex,"0^*"']
      & \Cinfty(\R^n) \rar & A,
  \end{tikzcd}
\]
where the $i$th coordinate function on $\R^k$ is sent respectively to $f_i$
and to zero. In the opposite category, this becomes an equalizer diagram
\[
  \begin{tikzcd}[column sep=large]
    A \rar & \Cinfty(\R^n)
      \ar[r,shift left=.6ex]
      \ar[r,shift right=.6ex]
      & \Cinfty(\R^k).
  \end{tikzcd}
\]
Consequently, if $F$ preserves finite limits and
\[
  R=F\circ j\colon\mathrm{Euc}\longrightarrow C,
\]
then $F(A)$ is forced to be the internal zero locus
\[
  F(A)
  \cong
  \operatorname{Eq}
  \left(
    \begin{tikzcd}[column sep=large]
      R(\R)^n
        \ar[r,shift left=.6ex,"{(f_1,\ldots,f_k)}"]
        \ar[r,shift right=.6ex,"0"']
        & R(\R)^k
    \end{tikzcd}
  \right).
\]
Thus the chosen $\Cinfty$-ring object determines the value of $F$ on every
finite presentation.

This argument explains uniqueness, but a theorem is still required. One must
show that the construction is independent of the chosen presentation,
functorial in $A$, and compatible with natural transformations. The general
proof observes that finitely presented algebras are generated from finitely
generated free algebras by finite colimits and retracts. Taking opposites
turns those operations into finite limits and retracts. The same generation
statements underlie the animated version of the theorem. Extending from
Euclidean spaces to all manifolds requires a geometric theorem, which we
will state in \Cref{sec:Universal_Property_Derived_Manifolds}.

\section{Why ordinary intersections are not enough}
\label[section]{sec:What_Ordinary_Envelope_Forgets}

The ordinary envelope gives a formal zero locus for every smooth map,
even when the geometric zero set is singular. Quotients retain the
equations and can therefore distinguish nilpotent thickenings. They do
not, however, distinguish a closed subobject from its self-intersection.
This calculation will explain what needs to change in the universal
problem.

Let $M$ be a manifold and let
\[
  f=(f_1,\ldots,f_k)\colon M\longrightarrow\R^k
\]
be smooth. The diagram of function rings is
\[
\begin{tikzcd}
  \Cinfty(\R^k) \rar["f^*"] \dar["0^*"'] \drar[pushout]
    & \Cinfty(M) \dar \\
  \R \rar & A_f.
\end{tikzcd}
\]
By the quotient formula from \Cref{prop:Cinfty_Coproduct_Formulas},
\[
  A_f
  \cong
  \frac{\Cinfty(M)}{(f_1,\ldots,f_k)}.
\]
After taking opposites, the corresponding formal affine object is the
pullback of
\[
  M\xrightarrow{f}\R^k\xleftarrow{0}*.
\]
This construction makes sense without a spectrum or a structure sheaf.

\begin{example}[A transverse zero]
  \label[example]{ex:Ordinary_Envelope_Transverse_Zero}
  The quotient $\Cinfty(\R)/(x)\cong\R$ says that the formal zero locus
  of $x$ is the ordinary point, as required by transverse compatibility.
\end{example}

\begin{example}[A singular zero]
  \label[example]{ex:Ordinary_Envelope_Singular_Zero}
  The dual-number quotient $\Cinfty(\R)/(x^2)$ from the preceding
  chapters is the formal zero locus of $x^2$. Thus the ordinary envelope
  already retains its nilpotent thickening.
\end{example}

Now consider the self-intersection of a closed subobject. Let $A$ be a
$\Cinfty$-ring, let $I\subseteq A$ be an ideal, and put
\[
  B=A/I.
\]
The self-intersection of the corresponding formal closed subobject is computed
by
\[
  B\amalg_A B.
\]

\begin{proposition}[Ordinary self-intersections]
  \label[proposition]{prop:Ordinary_Self_Intersections_Collapse}
  There is a natural isomorphism
  \[
    (A/I)\amalg_A(A/I)
    \cong
    A/I.
  \]
\end{proposition}

\begin{proof}
  Apply the quotient pushout formula:
  \[
    (A/I)\amalg_A(A/I)
    \cong
    A/(I+I)
    =A/I.
  \]
\end{proof}

For example, let $A=\Cinfty(\R^2)$ and $I=(y)$. Then $B$ represents the
$x$-axis. Its ordinary self-intersection inside the plane is again represented
by $B$. The quotient remembers that $y$ vanishes, but it has no record that
the relation $y=0$ was imposed twice.

The collapse reflects how pullbacks work in an ordinary category.
An ordinary pullback is characterized by an isomorphism of sets
\[
  \Hom_C(T,X\times_ZY)
  \cong
  \Hom_C(T,X)
  \times_{\Hom_C(T,Z)}
  \Hom_C(T,Y).
\]
An element on the right is a pair of maps whose composites to $Z$ are equal.
In a set, equality is a property. Once the equality holds, there is no further
datum recording how it holds. In particular, imposing the same relation twice
cannot create new equality data.

\section{Pullbacks with higher identifications}
\label[section]{sec:Equality_As_Structure}

We want a pullback to record how two composites are identified. For this,
we first introduce animae, whose elements admit identifications and
higher identifications. The based loops in a circle will illustrate
what a self-pullback can retain. We then formulate pullbacks in
$\infty$-categories before returning to smooth algebra.

\begin{definition}[Anima, informal description]
\label[definition]{def:Anima_Conceptual}
An \emph{anima} behaves like a set in which the identifications between
any two elements themselves form an anima. Identifications can be composed
and inverted, with the usual identities holding through coherent higher
identifications.
\end{definition}

For the purposes of this seminar, we take the notion of an anima as a
foundational concept, much as we take the notion of a set in ordinary
mathematics. It is possible to construct explicit models using, for example,
simplicial sets. Such constructions require substantial setup and do not
by themselves convey the intended intuition. We will instead work with
the heuristic picture above: an anima behaves like a set, except that its
identifications themselves form animae.

The following examples give different meanings to these identifications.
\begin{itemize}
  \item Every set $X$ determines an anima. For $x,y\in X$, the anima of
    identifications between $x$ and $y$ is a point if $x=y$, and empty
    otherwise. These are the $0$-truncated animae.
  \item Every ordinary category $C$ determines an anima $C^{\simeq}$,
    its \emph{groupoid core}. Its points are the objects of $C$, and its
    identifications are isomorphisms in $C$. For example, there is an
    anima of manifolds whose identifications are diffeomorphisms.
  \item Every topological space $X$ determines an anima $\Pi_\infty(X)$.
    Its points are points of $X$, and its identifications are paths in $X$.
    Here the anima of identifications can itself be interesting:
    identifications between paths are homotopies relative to their
    endpoints, and continuing in this way gives iterated homotopies.
\end{itemize}

An ordinary category can be described by a set of objects and a set of
morphisms between each pair of objects. In an $\infty$-category $C$, we
instead have an anima of objects and, for every two objects $X,Y$, an anima
\[
  \Hom_C(X,Y)
\]
of morphisms. An identification in $\Hom_C(X,Y)$ is called a
\emph{homotopy} between morphisms. Composition is associative and unital
with coherent higher identifications. Universal properties that are
ordinarily phrased using equalities between morphisms are now phrased
using homotopies between morphisms, together with their higher coherence.

Ordinary categories are the special case in which every mapping anima is
a set, so there are no nontrivial homotopies between morphisms. When an
ordinary category $C$ is viewed in this way, its anima of objects is its
core $C^{\simeq}$; isomorphisms between objects are still retained.

We can now write $\An$ for the $\infty$-category of animae. We will use
products, pullbacks, functor categories, and equivalences in this setting
without choosing a model.

The effect of retaining identifications is already visible in a
self-pullback of animae. Let $X$ be an anima and let $x\colon *\to X$
be a point. Form the pullback
\[
\begin{tikzcd}
  *\times_X* \rar \dar \drar[pullback]
    & * \dar["x"] \\
  * \rar["x"'] & X.
\end{tikzcd}
\]
An element of this pullback consists of the two chosen points together with
an identification between their images in $X$. Both images are $x$.
Writing $X(x,x)$ for the anima of identifications of $x$ with itself,
including all higher identifications, we obtain an isomorphism
\[
  *\times_X*
  \simeq
  X(x,x).
\]

In sets, the right side is a singleton. In animae, it can be highly
nontrivial. If
\[
  X=\Pi_\infty(S^1)
\]
is the underlying anima of the circle with base point $x$, an identification
of $x$ with itself is represented by a continuous path
\[
  \gamma\colon[0,1]\longrightarrow S^1,
  \qquad \gamma(0)=\gamma(1)=x.
\]
The endpoint condition is a strict equality of points, but the datum in the
pullback is the entire path $\gamma$. It may travel around the circle any
number of times. Paths between these choices are homotopies of loops that
keep the base point fixed. Thus
\[
  X(x,x)
  \simeq
  \Omega\Pi_\infty(S^1)
  \simeq
  \Z.
\]
The integers record the winding numbers of loops. Each winding-number
component is contractible, which gives the last isomorphism with the
discrete anima $\Z$.

The loop anima is a prototype for the information a homotopy pullback
can retain. A geometric derived self-intersection uses the same
pullback formalism in a different category.

The same description of pullbacks applies in any $\infty$-category.
Let
\[
  X\xrightarrow{f}Z\xleftarrow{g}Y
\]
be a diagram in an $\infty$-category $C$. Its pullback is characterized by
natural isomorphisms of animae
\[
  \Hom_C(T,X\times_ZY)
  \simeq
  \Hom_C(T,X)
  \times_{\Hom_C(T,Z)}
  \Hom_C(T,Y)
\]
for all $T\in C$.

Informally, a generalized $T$-point consists of
\[
  a\colon T\to X,
  \qquad
  b\colon T\to Y,
  \qquad
  \alpha\colon f\circ a\cong g\circ b.
\]
The equality $\alpha$ is part of the datum. Equalities between choices of
$\alpha$, and all subsequent coherence, are contained in the pullback of hom
animae.

This is the only pullback notion needed in the sequel. In older sources the
phrase \emph{homotopy pullback} is often used to distinguish a construction in
a chosen presentation from an ordinary strict pullback. We will work
synthetically and simply say \emph{pullback} in an $\infty$-category.

The passage from hom sets to hom animae now gives a precise reason to reconsider
\Cref{question:Universal_Finite_Limit_Completion}. We should ask for a
universal finite-limit envelope among $\infty$-categories, not merely among
ordinary categories.

\section{Animated smooth rings and finite limits}
\label[section]{sec:Animated_Finite_Limit_Envelope}

We now apply the preceding pullback formalism to smooth algebra. The
operations are still indexed by smooth maps between Euclidean spaces,
but their values are animae. The finite-limit-envelope theorem has a
corresponding animated form, with compactness replacing ordinary finite
presentation and with retracts included explicitly.

The definition in \Cref{def:Internal_Cinfty_Ring} extends to an $\infty$-category $C$
with finite products: a smooth ring object is a finite-product-preserving
functor $\mathrm{Euc}\to C$, with $\mathrm{Euc}$ regarded as an
$\infty$-category. We retain the notation $\operatorname{Alg}_{\Cinfty}(C)$.
Taking $C=\An$ gives the following definition.

\begin{definition}[Animated $\Cinfty$-ring]
  \label[definition]{def:Animated_Cinfty_Ring}
  An \emph{animated $\Cinfty$-ring} is a finite-product-preserving functor
  \[
    A\colon\mathrm{Euc}\longrightarrow\An.
  \]
  We denote their $\infty$-category by
  \[
    \operatorname{Alg}_{\Cinfty}(\An)
    =
    \Fun^\times(\mathrm{Euc},\An).
  \]
\end{definition}

An ordinary $\Cinfty$-ring is the special case in which every value $A(\R^n)$
is a set, regarded as a $0$-truncated anima. Animated rings allow the smooth
operations themselves to interact with higher equality data.

For ordinary smooth rings, \Cref{prop:Filtered_Colimit_Criterion} expresses
finite presentation by preservation of filtered colimits. We use the same
criterion here, applied to the whole mapping anima.

\begin{definition}[Homotopical finite presentation]
  \label[definition]{def:Homotopically_Finitely_Presented_Cinfty_Ring}
  An animated $\Cinfty$-ring $A$ is \emph{homotopically finitely presented} if
  it is compact in $\operatorname{Alg}_{\Cinfty}(\An)$: the functor
  \[
    \Hom_{\operatorname{Alg}_{\Cinfty}(\An)}(A,-)
    \colon
    \operatorname{Alg}_{\Cinfty}(\An)
    \longrightarrow
    \An
  \]
  preserves filtered colimits.
\end{definition}

We use the subscript $\mathrm{fp}$ for the full subcategory of such objects.
The word \emph{homotopically} matters: finite presentation is now a statement
about the entire hom anima, not only its set of connected components.

The next theorem identifies the universal target for Euclidean smooth
operations when finite limits and retracts are required. It applies to
every algebraic theory; we state only its $\Cinfty$ instance.

\begin{theorem}[Animated finite-limit envelope]
  \label[theorem]{thm:Animated_Cinfty_Finite_Limit_Envelope}
  Let $C$ be an idempotent-complete $\infty$-category with finite limits.
  Restriction along the finitely generated free animated $\Cinfty$-rings
  induces an equivalence of $\infty$-categories
  \[
    \Fun^{\mathrm{lex}}
    \left(
      \operatorname{Alg}_{\Cinfty}(\An)_{\mathrm{fp}}\catop,
      C
    \right)
    \simeq
    \operatorname{Alg}_{\Cinfty}(C).
  \]
\end{theorem}

This is \citet[Theorem~3.22]{Carchedi_Steffens_Universal_Derived_Manifolds}.
The ordinary \Cref{thm:Ordinary_Cinfty_Finite_Limit_Envelope} is its
$0$-truncated counterpart.

Recall that an $\infty$-category is idempotent-complete if every retract
diagram which exists formally is represented by an actual object. Ordinary
categories with finite limits are automatically idempotent-complete. This is
not automatic for $\infty$-categories, which is why the condition appears in
the theorem.

We sketch the ingredients of the quoted theorem. The starting point is
compact generation by free algebras.

\begin{enumerate}
  \item The free $\Cinfty$-ring on $n$ generators is compact. Mapping out of
    it is evaluation on an $n$-tuple, and evaluation commutes with filtered
    colimits.
  \item Every animated $\Cinfty$-ring is assembled as a colimit of free
    rings. This is the homotopical form of presenting an algebra by generators
    and relations.
  \item A compact animated $\Cinfty$-ring is a retract of a finite colimit of
    finitely generated free rings.
  \item After taking opposites, the preceding statement says that every object
    of
    $\operatorname{Alg}_{\Cinfty}(\An)_{\mathrm{fp}}\catop$
    is a retract of a finite limit of Euclidean objects. This describes the values forced on any finite-limit-preserving
    extension from $\mathrm{Euc}$.
\end{enumerate}

The complete argument constructs the extension by right Kan extension and
checks that it takes values in the chosen target $C$. The compactness and
idempotent-completeness hypotheses are precisely what make that last step
work. Finite colimits and retracts of free algebras explain the form of the
universal property; the Kan-extension argument supplies existence and
compatibility with all presentations.

Replacing sets by animae changes both the limits in the universal
problem and the finiteness condition on its algebraic objects.
The new pullbacks retain homotopies and their coherence. Compactness
must likewise control the whole mapping anima.

\begin{warning}
  \label[warning]{warn:Ordinary_Fp_Not_Necessarily_Animated_Fp}
  An ordinary finitely presented $\Cinfty$-ring need not remain compact when it
  is regarded as a discrete animated $\Cinfty$-ring. Thus the ordinary
  envelope is not simply a full subcategory of the animated envelope obtained
  by declaring every ordinary object discrete.
\end{warning}

The analogous warning for an arbitrary algebraic theory is recorded in
\citet[Warning~3.28]{Carchedi_Steffens_Universal_Derived_Manifolds}.
The truncation functor $\pi_0$ does send homotopically finitely presented
animated rings to finitely presented ordinary rings, but the reverse passage
can require retaining higher relation data.

In particular, the animated pushout associated with a self-intersection
can retain higher homotopy groups beyond its ordinary quotient.

\section{The universal property of derived manifolds}
\label[section]{sec:Universal_Property_Derived_Manifolds}

The animated envelope was constructed from Euclidean spaces. To answer
the chapter's question, we must relate its universal property to functors
on all manifolds that preserve transverse pullbacks. A geometric extension
theorem provides this link. We first state that theorem, then define the
universal category of derived manifolds and identify it with the opposite
category of compact animated smooth rings.

\Cref{chap:Manifolds_As_Cinfty_Rings} showed that the function-ring functor
preserves transverse pullbacks. We need a stronger statement: any
interpretation of Euclidean smooth operations extends uniquely to
manifolds while preserving transverse pullbacks and the terminal object.
We quote this geometric extension theorem.

Let
\[
  q\colon\mathrm{Euc}\longrightarrow\mathrm{Mfd}
\]
be the inclusion. For every
idempotent-complete $\infty$-category $C$ with finite limits, restriction along
$q$ induces an equivalence
\[
  \Fun^{\pitchfork}(\mathrm{Mfd},C)
  \simeq
  \Fun^\times(\mathrm{Euc},C)
  =
  \operatorname{Alg}_{\Cinfty}(C).
\]
Here $\Fun^\pitchfork$ denotes the functors preserving transverse
pullbacks and the terminal object. This is
\citet[Lemma~5.1 and Corollary~5.2]{Carchedi_Steffens_Universal_Derived_Manifolds}.

The geometric idea comes from the presentations in
\Cref{chap:Manifolds_As_Cinfty_Rings}. An open Euclidean subset is a
transverse zero locus of $ya(x)-1$, and every manifold is a retract of
an open tubular neighborhood in Euclidean space. A product-preserving
functor on $\mathrm{Euc}$ can therefore assign an object to each such
presentation using finite limits and retracts. This explains the candidate
extension. Proving its independence of choices, functoriality, and
preservation of transverse pullbacks is the substantive part of the
theorem; the proof uses right Kan extension.

The earlier transverse-pullback theorem concerned the actual function-ring
functor. This extension theorem applies to any internal $\Cinfty$-ring:
it determines a unique transverse-pullback-preserving functor on manifolds.

For an ordinary category $C$ with finite limits, combining this extension
theorem with \Cref{thm:Ordinary_Cinfty_Finite_Limit_Envelope} gives
\[
  \Fun^{\mathrm{lex}}(\operatorname{Aff}_{\Cinfty}^{\mathrm{fp}},C)
  \simeq\Fun^\pitchfork(\mathrm{Mfd},C).
\]
Thus ordinary finitely presented smooth rings solve the ordinary version
of the chapter's universal problem. We now formulate the version for
$\infty$-categories.

\begin{definition}[Universal $\infty$-category of derived manifolds]
  \label[definition]{def:Universal_Infinity_Category_Derived_Manifolds}
  The $\infty$-category $\mathrm{DMfd}$ of \emph{derived manifolds} is equipped
  with a functor
  \[
    i\colon\mathrm{Mfd}\longrightarrow\mathrm{DMfd}
  \]
  and is characterized by the following properties.
  \begin{enumerate}
    \item The $\infty$-category $\mathrm{DMfd}$ has finite limits and is
      idempotent-complete.
    \item The functor $i$ preserves transverse pullbacks and the terminal
      object.
    \item For every idempotent-complete $\infty$-category $C$ with finite
      limits, restriction along $i$ induces an equivalence
      \[
        \Fun^{\mathrm{lex}}(\mathrm{DMfd},C)
        \simeq
        \Fun^{\pitchfork}(\mathrm{Mfd},C),
      \]
      where the right side consists of functors preserving transverse
      pullbacks and the terminal object.
  \end{enumerate}
\end{definition}

This is \citet[Definition~2.1]{Carchedi_Steffens_Universal_Derived_Manifolds}.
The universal property determines $\mathrm{DMfd}$ uniquely.
It fixes the existing transverse pullbacks, adjoins all finite
homotopy limits, and includes their retracts.

The geometric extension theorem turns the defining universal property
under manifolds into a universal property in terms of smooth ring objects.
Comparing this with the animated envelope will both establish existence
and give an algebraic description of $\mathrm{DMfd}$.

\begin{theorem}[Universal property through $\Cinfty$-rings]
  \label[theorem]{thm:Universal_Property_Derived_Manifolds_Cinfty_Rings}
  For every idempotent-complete $\infty$-category $C$ with finite limits,
  there is a natural equivalence
  \[
    \Fun^{\mathrm{lex}}(\mathrm{DMfd},C)
    \simeq
    \operatorname{Alg}_{\Cinfty}(C).
  \]
\end{theorem}

\begin{proof}
  By the universal property under manifolds,
  \[
    \Fun^{\mathrm{lex}}(\mathrm{DMfd},C)
    \simeq
    \Fun^{\pitchfork}(\mathrm{Mfd},C).
  \]
  Restriction to Euclidean spaces gives
  \[
    \Fun^{\pitchfork}(\mathrm{Mfd},C)
    \simeq
    \Fun^\times(\mathrm{Euc},C)
    =
    \operatorname{Alg}_{\Cinfty}(C).
  \]
  Composing these equivalences proves the claim.
\end{proof}

This is \citet[Theorem~5.3]{Carchedi_Steffens_Universal_Derived_Manifolds}.

The animated finite-limit envelope has exactly the same universal property.
We therefore obtain the promised algebraic presentation.

\begin{corollary}[Derived manifolds as formal animated affines]
  \label[corollary]{cor:Derived_Manifolds_Animated_Cinfty_Rings}
  There is an equivalence of $\infty$-categories
  \[
    \mathrm{DMfd}
    \simeq
    \operatorname{Alg}_{\Cinfty}(\An)_{\mathrm{fp}}\catop.
  \]
\end{corollary}

\begin{proof}
  By
  \Cref{thm:Animated_Cinfty_Finite_Limit_Envelope,thm:Universal_Property_Derived_Manifolds_Cinfty_Rings},
  both sides are idempotent-complete $\infty$-categories with finite limits
  which universally carry a $\Cinfty$-ring object. Universal objects are
  unique, so the two $\infty$-categories are equivalent.
\end{proof}

This is \citet[Corollary~5.4]{Carchedi_Steffens_Universal_Derived_Manifolds}.
This identifies the universal category with an algebraic construction.
In particular, manifolds embed fully faithfully, transverse pullbacks
are preserved, and all finite pullbacks exist.

For a smooth map $f\colon U\to\R^r$, the function ring of its derived
zero locus is consequently the pushout
\[
 \Cinfty(U)\otimes^{\mathbb L}_{\infty,\Cinfty(\R^r)}\R
\]
in animated smooth rings. \Cref{chap:Derived_Zero_Loci_And_Tangent_Complexes} computes it by a differential
graded resolution. The category of derived schemes introduced in
\Cref{chap:Local_Derived_Geometry_Kuranishi_Charts} allows these affine objects to be described locally and glued.

\section{Further details: retracts and coherence}
\label[section]{sec:From_Smooth_Algebra_Summary_Supplementary_Materials}

The universal properties above require idempotent-complete targets.
We explain more precisely what this means in an $\infty$-category and
why finite limits alone do not supply it. The exercises revisit this
point and the self-pullback example.

Informally, an idempotent looks like a morphism $e\colon X\to X$ together with
a homotopy from $e\circ e$ to $e$. Formally, an idempotent in an
$\infty$-category is a coherent idempotent diagram: the displayed homotopy is
part of the data, together with its higher compatibilities. Splitting it means
finding an object $Y$ and maps
\[
  Y\xrightarrow{i}X\xrightarrow{r}Y
\]
together with coherent identifications of $r\circ i$ with
$\operatorname{id}_Y$ and of $i\circ r$ with the given idempotent. In an
ordinary category with finite limits, an idempotent $e\colon X\to X$
splits by taking the equalizer $i\colon Y\to X$ of $e$ and the identity.
The map $e$ factors as $ir$, and $ri=\operatorname{id}_Y$ follows
from the equalizer property. For $\infty$-categories
this is not forced by finite limits alone. The compact algebraic objects are
closed under retracts, so their opposite category is naturally
idempotent-complete. The universal property of derived manifolds must record
this closure explicitly.

\begin{exercise}
Prove the equalizer construction of a splitting just described.
Compare it with the explicit finite presentation for a retract ring
in \Cref{chap:Manifolds_As_Cinfty_Rings}.
\end{exercise}

\begin{exercise}
Use the universal property of a homotopy pullback to identify
$*\times_{BG}*$ with a discrete group $G$, where $BG$ is its
classifying anima and both maps select the base point.
Compare this with the self-pullback over a set.
\end{exercise}

The primary source is
\cite[Definition~2.1, Theorem~3.22, Remark~3.26,
and Section~5]{Carchedi_Steffens_Universal_Derived_Manifolds}.
It treats both universal properties, their geometric comparison, and
the relation with Spivak's construction.

% !TEX root = ../main.tex

\chapter{Calculating derived intersections}
\label[chapter]{chap:Derived_Zero_Loci_And_Tangent_Complexes}

A smooth map $f\colon U\to\R^r$ determines a derived manifold
\[
Z^{\mathrm{der}}(f)=U\times_{\R^r}\{0\}.
\]
The preceding chapter guarantees that this pullback exists. We now ask
what its functions remember about the equations defining it. To calculate
them, we need an algebraic presentation of the pullback.
The resulting algebra has one generator for each equation, together with a
differential recording that equation. Products of the generators keep track
of relations among the equations.

Write $f=(f_1,\ldots,f_r)$. Under the
contravariant description of derived manifolds by animated $\Cinfty$-rings,
the defining pullback becomes the pushout
\[
\begin{tikzcd}
 \Cinfty(\R^r) \rar["f^*"] \dar["0^*"'] \drar[pushout]
   & \Cinfty(U) \dar\\
 \R \rar & \mathcal O(Z^{\mathrm{der}}(f)).
\end{tikzcd}
\]
Thus
\begin{equation}\label{eq:Derived_Zero_Locus_Animated_Pushout}
 \mathcal O(Z^{\mathrm{der}}(f))
 \cong \Cinfty(U)\mathbin{\otimes^{\mathbb L}_{\infty,\Cinfty(\R^r)}}\R.
\end{equation}
Here the tensor notation denotes a pushout of animated smooth rings.
The corresponding ordinary pushout is
$\Cinfty(U)/(f_1,\ldots,f_r)$. Computing that quotient first would discard
the higher homotopy of the animated pushout.

There are already two different classical objects in this discussion.
The ring $\Cinfty(U)/(f_1,\ldots,f_r)$ records the equations, whereas the
set $f^{-1}(0)$ only records their common real solutions. For example,
$x$ and $x^2$ have the same zero set but give different quotients.
The derived construction adds a third level: homotopy groups beyond
$\pi_0$.

We first introduce smooth differential graded algebras as a tool for
computing animated smooth rings, then construct the Koszul model of a
zero locus. Comparing a transverse equation, a multiple zero, and a
repeated equation will make these three levels visible. The same
construction will then apply to families and to sections of vector bundles.

\section{Smooth differential graded algebras}
\label[section]{sec:Smooth_Dold_Kan_Bridge}

Smooth differential graded algebras make the higher homotopy of an
animated smooth ring accessible through homology. We use homological
grading, as in Steffens: the differential lowers degree by one. Our
connective dg-algebra models are concentrated in nonnegative degrees.
After describing their smooth structure, we state the comparison theorem
that justifies using them to compute derived pushouts.

\begin{definition}[Smooth dg-algebra]
\label[definition]{def:Smooth_Dg_Algebra}
A \emph{connective $\Cinfty$-dg-algebra} is a commutative differential
graded $\R$-algebra
\[
 \cdots\longrightarrow A_2\xrightarrow{d}A_1
 \xrightarrow{d}A_0
\]
together with a $\Cinfty$-ring structure on $A_0$ whose underlying
$\R$-algebra structure is the given one. A morphism preserves the differential,
the multiplication, and the smooth operations on degree zero.
\end{definition}

For homogeneous elements, commutativity and the differential rule mean
\[
 ab=(-1)^{|a||b|}ba,\qquad
 d(ab)=d(a)b+(-1)^{|a|}a\,d(b).
\]
In particular, an odd-degree generator has square zero over $\R$.
The smooth operations are supplied in degree zero. There is no operation
$f(a_1,\ldots,a_n)$ for arbitrary nonhomogeneous elements in this definition.

This is the convention of the foundations paper \cite[Definition~3.3.1]
{Steffens_Derived_Cinfty_Geometry_I}. Because $d$ vanishes on $A_0$, there
is no additional chain-rule condition to impose on its degree-zero elements.

\begin{example}
An ordinary $\Cinfty$-ring is a smooth dg-algebra concentrated in degree
zero. More generally, for an ordinary smooth ring $A$ and elements
$f_1,\ldots,f_r\in A$, form
\[
 A[\varepsilon_1,\ldots,\varepsilon_r],
 \qquad |\varepsilon_i|=1,\qquad d\varepsilon_i=f_i.
\]
The notation denotes a free graded-commutative extension: as a graded
$A$-module it is $A\otimes_\R\Lambda^\bullet(\R^r)$, with
$\Lambda^q$ placed in degree $q$. Its differential is determined by the
displayed equations and the Leibniz rule. In particular,
\[
 d(\varepsilon_i\varepsilon_j)
 =f_i\varepsilon_j-f_j\varepsilon_i.
\]
The equality $d^2=0$ follows on generators and hence on the whole algebra.
\end{example}

The image of $d\colon A_1\to A_0$ is an ideal: for $a\in A_0$ and
$b\in A_1$, $a\,d(b)=d(ab)$. Therefore $H_0(A)$ is an ordinary
$\Cinfty$-ring. To see that all smooth operations descend, apply Hadamard's
lemma to a smooth function $\varphi\colon\R^n\to\R$:
\[
 \varphi(x)-\varphi(y)=\sum_i(x_i-y_i)\psi_i(x,y).
\]
Replacing each $x_i$ by an element congruent to $y_i$ modulo the ideal
changes the output by an element of the same ideal. The groups $H_i(A)$
are modules over $H_0(A)$, and their products form a graded algebra.

The passage from dg-algebras to animated smooth rings identifies maps
that induce isomorphisms on homology.

\begin{definition}
\label[definition]{def:Smooth_Dga_Quasi_Isomorphism}
A morphism of smooth dg-algebras is a \emph{quasi-isomorphism} if it induces
an isomorphism on every homology group.
\end{definition}

\begin{theorem}[Smooth Dold--Kan correspondence]
\label[theorem]{fact:Smooth_Dold_Kan_Comparison}
The localization of connective $\Cinfty$-dg-algebras at quasi-isomorphisms,
taken as an $\infty$-categorical localization, is equivalent to the
$\infty$-category of animated $\Cinfty$-rings. If $A_{\mathrm{an}}$
corresponds to a dg-algebra $A$, then
\[
 \pi_i(A_{\mathrm{an}})\cong H_i(A),\qquad i\geq0.
\]
\end{theorem}

We use this as a comparison theorem
\cite[Theorem~3.3.12]{Steffens_Derived_Cinfty_Geometry_I}.
It includes multiplicative and smooth structure. The ordinary Dold--Kan
correspondence for vector spaces alone does not supply that structure.

One way to describe the animated object attached to $A$ is through its
values on free smooth rings. In the localized category $D$ of smooth
dg-algebras, set
\[
 A_{\mathrm{an}}(\R^n)=\Hom_D(\Cinfty(\R^n),A).
\]
The functor is covariant in Euclidean spaces because smooth functions are
contravariant. Free smooth rings turn finite products of Euclidean spaces
into coproducts, so this functor preserves finite products. It is therefore
an animated $\Cinfty$-ring. The mapping objects in this formula are animae;
using only strict dg-algebra maps would lose the higher homotopies.

The theorem allows computations using dg-algebras provided we compute
their \emph{derived} pushouts. An arbitrary strict pushout of dg-algebras
need not have the required universal property. In the zero-locus
calculation, we will replace evaluation at the origin by a free resolution
before taking the strict pushout.

\section{The Koszul model of a zero locus}
\label[section]{sec:Koszul_Presentation_Derived_Zero_Locus}

We compute the pushout in
\Cref{eq:Derived_Zero_Locus_Animated_Pushout} in two stages. First we
resolve evaluation at the origin of $\R^r$. We then substitute the
functions defining our zero locus into this resolution.

Start with one coordinate $t$ and put $B=\Cinfty(\R)$. Consider the map
\[
 B[\varepsilon]\longrightarrow\R,\qquad
 |\varepsilon|=1,\quad d\varepsilon=t,
\]
sending $t$ and $\varepsilon$ to zero. Its underlying complex is
\[
 B\xrightarrow{\ t\ }B
\]
in degrees $1,0$. Multiplication by $t$ is injective: a smooth function
annihilated by $t$ vanishes away from the origin, hence everywhere.
Its cokernel is $\R$, since a smooth function vanishing at zero is divisible
by $t$. Thus this map is a quasi-isomorphism.

For $B=\Cinfty(\R^r)$, the same construction gives
\[
 K_B(t_1,\ldots,t_r)
 =\bigl(B[\varepsilon_1,\ldots,\varepsilon_r],
          d\varepsilon_i=t_i\bigr)\longrightarrow\R.
\]
The coordinate functions form a regular sequence. Indeed, quotienting by
the first $j$ coordinates gives the smooth functions on the remaining
coordinate subspace, and the next coordinate is a non-zero-divisor there.
Successively applying the one-variable calculation proves that the
augmentation is a quasi-isomorphism.

The map $B\to K_B(t_1,\ldots,t_r)$ is built by adjoining generators in
positive degree with prescribed differential. Together with the
augmentation to $\R$, it is a relative free resolution of evaluation.
Taking its strict pushout computes derived base change. The criterion
used here is \citet[Corollary~3.3.15]{Steffens_Derived_Cinfty_Geometry_I}:
this map is an isomorphism in
degree zero, is surjective on $H_0$, and is a relative cell extension
of underlying commutative dg-algebras. The resulting zero-locus
calculation is given in
\citet[Example~3.3.16]
{Steffens_Derived_Cinfty_Geometry_I}.

Now base change along $f^*\colon B\to A=\Cinfty(U)$. Each $t_i$ becomes
$f_i$, so the pushout is
\begin{equation}\label{eq:Koszul_Algebra_Function}
 K_A(f)=\bigl(A[\varepsilon_1,\ldots,\varepsilon_r],
               d\varepsilon_i=f_i\bigr).
\end{equation}

\begin{theorem}[Koszul model]
\label[theorem]{thm:Koszul_Presentation_Derived_Zero_Locus}
The smooth dg-algebra $K_A(f)$ presents the animated function ring of
$Z^{\mathrm{der}}(f)$. Consequently,
\[
 \pi_i\mathcal O(Z^{\mathrm{der}}(f))\cong H_i(K_A(f)).
\]
\end{theorem}

\begin{proof}
Replace the evaluation map $B\to\R$ in
\Cref{eq:Derived_Zero_Locus_Animated_Pushout} by the relative free
resolution $B\to K_B(t)$. Its strict pushout along $B\to A$ computes
the derived pushout and is exactly \Cref{eq:Koszul_Algebra_Function}.
Apply the smooth Dold--Kan correspondence.
\end{proof}

The construction resolves the coordinates on $\R^r$, which are a regular
sequence, before substituting the possibly singular functions $f_i$.
No regular-sequence assumption on the $f_i$ is needed. The
positive-degree homology can retain additional relations among them.

In particular,
\[
 H_0(K_A(f))=A/(f_1,\ldots,f_r).
\]
In degree $1$, a cycle $\sum_i a_i\varepsilon_i$ is a relation
$\sum_i a_if_i=0$. Boundaries include the relations
$f_i\varepsilon_j-f_j\varepsilon_i$. Thus higher homology measures
relations that survive after these elementary ones have been identified.
The algebra also retains the multiplication between such classes.

\section{Equations with the same zero set}

The preceding chapters computed ordinary quotient rings. We now return
to those equations to ask what survives in the positive degrees of the
Koszul model. In particular, a multiple zero need not have higher
homotopy of functions, whereas an excess equation can produce it.

\begin{example}[A transverse equation]
For $f(x)=x$, the model is
\[
 \Cinfty(\R)[\varepsilon],\qquad d\varepsilon=x.
\]
The resolution calculation shows that it is quasi-isomorphic to $\R$.
Thus the derived zero locus is the ordinary point.

More generally, if $f\colon U\to\R^r$ is transverse to zero, its components
can be completed locally to coordinates along the zero set. The same
Koszul calculation shows that the derived zero locus is the ordinary
manifold $f^{-1}(0)$. Away from the zero set, one of the functions is
invertible, and the Koszul complex is contractible. These local statements
recover the transverse-pullback property of derived manifolds.
\end{example}

\begin{example}[A multiple zero]
For $f(x)=x^2$, multiplication by $x^2$ is still injective on $\Cinfty(\R)$.
Hence the Koszul algebra has no positive homology, but
\[
 H_0(K(x^2))=\Cinfty(\R)/(x^2)\cong\R[\delta]/(\delta^2).
\]
The last identification follows from Taylor's formula:
\[
 g(x)=g(0)+xg'(0)+x^2r(x),
\]
with $r$ smooth. The smooth operations on the quotient are given by the
first-order Taylor formula. For instance,
$\varphi(a+b\delta)=\varphi(a)+\varphi'(a)b\delta$.

Thus $Z^{\mathrm{der}}(x^2)$ is not the ordinary point, even though its
animated ring is discrete. Nilpotents in $\pi_0$ can already distinguish
it. The same argument gives
$\Cinfty(\R)/(x^m)\cong\R[\delta]/(\delta^m)$ for $m\geq1$.
\end{example}

\begin{example}[A zero equation]
For the zero function on $\R$, the differential vanishes:
\[
 K(0)=\Cinfty(\R)[\varepsilon],\qquad d\varepsilon=0.
\]
Its degree-zero homology is $\Cinfty(\R)$, and its degree $1$
homology is another copy of $\Cinfty(\R)$. The classical truncation is
the line, while the derived object remembers that an equation was imposed
which vanishes identically.

Similarly, the self-intersection of $\{0\}\subset\R$ has function algebra
$\R[\varepsilon]$ with $d\varepsilon=0$. Its underlying topological space
is a point and its classical function ring is $\R$, but $\pi_1=\R$.
This is a local instance of the excess-intersection phenomena in
\citet[Example~2.7]{Spivak_Derived_Smooth_Manifolds}.
\end{example}

\begin{example}[A repeated equation]
The maps $x\mapsto x$ and $x\mapsto(x,x)$ have the same zero set and
the same ordinary quotient ring. For the second map the Koszul algebra is
\[
 \Cinfty(\R)[\varepsilon_1,\varepsilon_2],
 \qquad d\varepsilon_1=d\varepsilon_2=x.
\]
Set $\eta=\varepsilon_2-\varepsilon_1$. Then $d\eta=0$, and the algebra
becomes $K(x)\otimes_\R\Lambda(\eta)$. It is quasi-isomorphic to
$\R[\eta]$ with zero differential. The repeated equation produces a
nonzero $\pi_1$.
\end{example}

These examples separate three comparisons. A point and a multiple point
differ in their classical function rings. An ordinary point and a derived
self-intersection differ in higher homotopy. Different lists generating
the same ordinary ideal can define different derived zero loci, because
the derived construction remembers the chosen map to the space of equations.

\section{Families of equations}

The model also lets us follow an equation as its parameters vary.
Even when the total zero locus is a manifold, its fibres can acquire
multiplicity or higher homotopy.

Let $S$ be a parameter manifold and let
$f\colon S\times U\to\R^r$ be smooth. The relative equation defines a
derived object over $S$:
\[
 X=(S\times U)\times_{S\times\R^r}(S\times\{0\}).
\]
Its Koszul model has coefficients in $\Cinfty(S\times U)$ and
differential $d\varepsilon_i=f_i(s,x)$.

For a smooth map $g\colon T\to S$, pullback gives the same construction
with $f(g(t),x)$. This follows by pasting the defining pullback squares.
In particular, the fibre at $s\in S$ is the derived zero locus of $f_s$.
The fibre is a derived pullback even when the total object is an ordinary
manifold.

For example, consider $f(s,x)=x^2-s$. The total zero locus is the ordinary
parabola $s=x^2$, since the derivative in the $s$-direction is $-1$.
Over $s>0$ its fibre consists of two points; over $s<0$ it is empty.
The derived fibre at $s=0$ has function ring $\R[\delta]/(\delta^2)$.
The parameter map from the parabola to $S$ has a critical point there.

A second family, $f(s,x)=sx$, illustrates the appearance of higher
homotopy in a fibre. For $s\neq0$, the fibre is an ordinary point.
For $s=0$, its Koszul model is
$\Cinfty(\R)[\varepsilon]$ with zero differential. Although multiplication
by $sx$ on $\Cinfty(\R^2)$ is injective, substitution $s=0$ destroys
injectivity. Derived base change retains this change in the equation.

\section{Zero loci of vector bundle sections}
\label[section]{sec:Supplement_Vector_Bundle_Koszul_Formula}

The same model works without a chosen list of global equations. Let
$p\colon E\to U$ be a vector bundle of rank $r$, and let $s$ be a section.
Its derived zero locus is the pullback of $s$ and the zero section as maps
from $U$ to $E$. Define a sheaf of smooth dg-algebras on $U$ by
\[
 \mathcal K(s)_q
   =\Gamma(-,\Lambda^qE^\vee),\qquad d=\iota_s.
\]
Here $\iota_s$ contracts an alternating form with $s$. In a local frame,
this is the Koszul algebra already calculated. Changes of frame respect
contraction, so the local algebras glue.

\begin{proposition}
\label[proposition]{prop:Vector_Bundle_Koszul_Presentation}
The sheaf $\mathcal K(s)$ presents the functions of the derived zero locus
of $s$, supported on its ordinary zero set. In particular, for the zero
section its differential vanishes and its homotopy sheaves are
$\Lambda^qE^\vee$ in degree $q$.
\end{proposition}

\begin{proof}
Trivialize $E$ on an open cover and apply the Koszul model theorem.
On the complement of the zero set, contraction has a local contracting
homotopy: choose $\lambda\in E^\vee$ with $\lambda(s)=1$ and use exterior
multiplication by $\lambda$. The frame-independent contraction formula
identifies the models on overlaps, giving the stated sheaf.
\end{proof}

This is the bundle form of the calculation in
\cite[Example~3.3.17]{Steffens_Derived_Cinfty_Geometry_I}.
It explains how the local excess directions of a self-intersection fit
into a vector bundle rather than a fixed vector space.

\section{Further reading}

The following observations explain the role of regular sequences in the
examples and the effect of allowing generators in higher degrees.
The exercises give further comparisons between presentations of zero loci.

A regular sequence $f_1,\ldots,f_r$ in an ordinary
smooth ring has Koszul homology concentrated in degree zero. One proves
this inductively: adjoining the last generator forms the cone of
multiplication by $f_r$ on the preceding Koszul complex. If the preceding
homology is $A/(f_1,\ldots,f_{r-1})$ and $f_r$ is a non-zero-divisor
there, the cone has only the desired quotient as its degree-zero
homology. This explains both the transverse case and the multiple-zero
example. Concentration in degree zero is weaker than smoothness of the
resulting ordinary ring.

The smooth dg-algebra model also permits generators in degrees above $1$.
For example, a free generator $z$ in degree $2$ with $dz=0$ gives
$\R[z]$. Since $z$ has even degree, its powers need not vanish:
$H_{2j}(\R[z])=\R z^j$ for every $j\geq0$.
A finite list of dg-generators therefore does not imply that the
function algebra has only finitely many nonzero homotopy groups.
The boundedness of a tangent complex is a separate issue, treated next.

\begin{exercise}
Compute the Koszul homology of $(x,0)\colon\R\to\R^2$ and compare it
with the repeated-equation example.
\end{exercise}

\begin{exercise}
For $f=(x_1,\ldots,x_q,0,\ldots,0)\colon\R^n\to\R^r$, show that the
Koszul algebra is quasi-isomorphic to
$\Cinfty(\R^{n-q})[\eta_1,\ldots,\eta_{r-q}]$ with zero differential.
Interpret the surviving generators as excess equations.
\end{exercise}

\begin{exercise}
Show directly that adjoining a new coordinate $y$ and the equation $y=0$
does not change the derived zero locus. Compare this with adjoining the
identically zero equation.
\end{exercise}

% This chapter is the complete working manuscript for the sixth seminar talk.
% The final section contains supplementary material outside the default live
% route.

\chapter{Derived bordism and fundamental classes}
\label[chapter]{chap:Derived_Bordism_And_Fundamental_Classes}

The preceding chapter constructed a derived zero locus
\[
  Z^{\mathrm{der}}(s)
  =
  M\mathbin{\times_E^{\mathrm{der}}}M
\]
for every section $s\colon M\to E$ of a finite-rank vector bundle, without
requiring $s$ to be transverse to the zero section. It then calculated this
object through a Koszul presentation. We now ask a different question. If $M$
is compact and $s'$ is a transverse perturbation of $s$, how are the derived
zero locus $Z^{\mathrm{der}}(s)$ and the ordinary manifold $Z(s')$ related?

They are generally not equivalent. They need not even have homeomorphic
underlying spaces. The relation between them is deliberately coarser:

\begin{quote}
  A compact derived zero locus and a transverse perturbation represent the
  same bordism class.
\end{quote}

This statement gives derived intersections their first topological payoff.
The unperturbed intersection exists as a geometric object, while its bordism
class is the classical class obtained by perturbation. The same statement also
shows how much bordism forgets. A derived object may have nontrivial animated
structure even when its bordism class vanishes.

There is one scope convention which will remain in force throughout the
chapter. Spivak's bordism theorem concerns compact \emph{quasi-smooth} derived
manifolds, meaning the derived manifolds locally presented as zero loci of
maps between ordinary manifolds. All derived intersections and vector-bundle
zero loci used in this chapter are of this kind. We do not claim the theorem
for arbitrary higher-amplitude objects in the full $\infty$-category
$\mathrm{DMfd}$. Carchedi--Steffens compare Spivak's framework with the modern
one used in this manuscript, so the quasi-smooth objects considered below may
be regarded as objects of the same $\infty$-category as in the preceding
chapters \cite[Theorem~5.10]{Carchedi_Steffens_Universal_Derived_Manifolds}.

\section{From perturbations to derived cobordisms}
\label[section]{sec:From_Perturbations_To_Derived_Cobordisms}

\subsection{A homotopy of sections}
\label[subsection]{sec:Homotopy_Of_Sections_Cobordism}

Let $M$ be a compact manifold, let $q\colon E\to M$ be a rank-$r$ vector
bundle, and let $s\colon M\to E$ be a section. Recall that its derived zero
locus is defined by the pullback square
\[
\begin{tikzcd}
  Z^{\mathrm{der}}(s) \rar \dar \drar[pullback]
    & M \dar["0"] \\
  M \rar["s"'] & E .
\end{tikzcd}
\]
Its underlying set is the ordinary zero set of $s$, but its structure sheaf
retains the equations defining that set.

Choose another section $s'\colon M\to E$. We first assume that $s'$ is
transverse to the zero section. Choose a smooth function
\[
  \beta\colon\R\longrightarrow[0,1]
\]
which is identically zero on a neighborhood of $0$ and identically one on a
neighborhood of $1$. Fiberwise affine combination defines a section
\[
  H\colon\R\times M\longrightarrow\operatorname{pr}_M^*E,
  \qquad
  H(u,m)=(1-\beta(u))s(m)+\beta(u)s'(m).
\]
Its derived zero locus is taken over $\R\times M$:
\[
\begin{tikzcd}
  W \rar \dar \drar[pullback]
    & \R\times M \dar["0"] \\
  \R\times M \rar["H"'] & \operatorname{pr}_M^*E .
\end{tikzcd}
\]
Projection to the first factor induces a map
\[
  p\colon W=Z^{\mathrm{der}}(H)\longrightarrow\R.
\]
The important point is that the single derived object $W$ contains both zero
loci as fibers.

\begin{proposition}[A homotopy of sections gives a cobordism]
  \label[proposition]{prop:Homotopy_Of_Sections_Derived_Cobordism}
  With notation as above, the map $p\colon W\to\R$ is proper and is a product
  near its fibers over $0$ and $1$. Moreover,
  \[
    W_{p=0}\simeq Z^{\mathrm{der}}(s),
    \qquad
    W_{p=1}\simeq Z(s').
  \]
\end{proposition}

\begin{proof}
  Pulling the defining square of $W$ back along $\{0\}\to\R$ replaces $H$
  by $s$, so
  \[
    W_{p=0}\simeq Z^{\mathrm{der}}(s).
  \]
  The same argument at $1$ gives
  $W_{p=1}\simeq Z^{\mathrm{der}}(s')$. Since $s'$ is transverse to the zero
  section, its derived zero locus is the ordinary zero manifold $Z(s')$.

  Choose $\varepsilon>0$ so that $\beta$ is zero on
  $(-\varepsilon,\varepsilon)$ and one on
  $(1-\varepsilon,1+\varepsilon)$. Over these two intervals the section $H$
  is constant in the parameter. Consequently,
  \[
    W|_{(-\varepsilon,\varepsilon)}
    \simeq
    Z^{\mathrm{der}}(s)\times(-\varepsilon,\varepsilon)
  \]
  and
  \[
    W|_{(1-\varepsilon,1+\varepsilon)}
    \simeq
    Z(s')\times(1-\varepsilon,1+\varepsilon).
  \]

  It remains to check properness. Let $K\subset\R$ be compact. The
  underlying space of $p^{-1}(K)$ is the ordinary zero set of $H$ inside
  $K\times M$. It is closed because the zero section of the vector bundle is
  closed. Since $K\times M$ is compact, the space underlying $p^{-1}(K)$ is
  compact. Thus $p$ is proper.
\end{proof}

The proof separates the two roles of the hypotheses. The constancy of the
homotopy near $0$ and $1$ supplies product neighborhoods of the endpoint
fibers. Compactness of $M$ prevents zeros from escaping to infinity while the
section is perturbed.

\subsection{The definition of derived cobordism}
\label[subsection]{sec:Definition_Derived_Cobordism}

The preceding construction contains every clause in the general definition.

\begin{definition}[Collared map]
  \label[definition]{def:Collared_Map_Derived_Manifolds}
  Let $p\colon W\to\R$ be a map of derived manifolds and let $i\in\R$. We say
  that $p$ is \emph{$i$-collared} if there is an $\varepsilon>0$ and an
  equivalence over $(i-\varepsilon,i+\varepsilon)$
  \[
    W|_{|p-i|<\varepsilon}
    \simeq
    W_{p=i}\times(-\varepsilon,\varepsilon).
  \]
  On the right, the interval is translated so that its origin maps to $i$.
\end{definition}

\begin{definition}[Derived cobordism]
  \label[definition]{def:Derived_Cobordism}
  Let $X_0$ and $X_1$ be compact quasi-smooth derived manifolds. A
  \emph{derived cobordism} from $X_0$ to $X_1$ consists of a quasi-smooth
  derived manifold $W$ and a proper map
  \[
    p\colon W\longrightarrow\R
  \]
  such that:
  \begin{enumerate}
    \item The map $p$ is $0$-collared and $1$-collared.
    \item There are equivalences
      \[
        W_{p=0}\simeq X_0,
        \qquad
        W_{p=1}\simeq X_1.
      \]
  \end{enumerate}
  If the objects carry maps to a manifold $T$, a derived cobordism
  \emph{over $T$} also carries a continuous map from the underlying space of
  $W$ to $T$ which restricts to the endpoint maps.
\end{definition}

This is Spivak's boundaryless formulation of cobordism
\cite[Definitions~3.5--3.6]{Spivak_Derived_Smooth_Manifolds}. Instead of
introducing derived manifolds with boundary, one uses a manifold without
boundary equipped with a proper map to $\R$. Properness is the compactness
condition appropriate to such a family. The collars permit two cobordisms to
be glued along a common endpoint fiber.

\begin{remark}[Cobordism is not equivalence]
  \label[remark]{rem:Cobordism_Not_Equivalence}
  A homotopy of sections does not generally produce an equivalence between
  their derived zero loci. It produces a larger derived object over a
  parameter line. Passing from objects to bordism classes discards the
  difference between the two endpoint fibers.
\end{remark}

\begin{corollary}[Homotopic sections]
  \label[corollary]{cor:Homotopic_Sections_Derived_Cobordant}
  Let $M$ be compact and let $s_0,s_1$ be sections of the same vector bundle
  over $M$. Then their derived zero loci are derived cobordant. If $s_1$ is
  transverse, then $Z^{\mathrm{der}}(s_0)$ is derived cobordant to the
  ordinary manifold $Z(s_1)$.
\end{corollary}

\begin{proof}
  Apply the construction above to a homotopy of sections which is constant
  near its endpoints. Transversality is needed only to identify the second
  endpoint with an ordinary manifold.
\end{proof}

\section{Classical and derived bordism}
\label[section]{sec:Classical_And_Derived_Bordism}

\subsection{The comparison map}
\label[subsection]{sec:Derived_Bordism_Comparison_Map}

Let $T$ be a manifold. Write $\Omega_d(T)$ for the ordinary unoriented
bordism group of proper maps from $d$-manifolds to $T$. Equivalently,
$\Omega_*(T)$ is the homology theory represented by $MO$. Define
$\Omega_d^{\mathrm{der}}(T)$ in the same way using compact quasi-smooth
derived manifolds of virtual dimension $d$ and derived cobordisms over $T$.
Disjoint union gives the group operation.

The inclusion of ordinary manifolds into derived manifolds induces a map
\[
  i_*
  \colon
  \Omega_*(T)\longrightarrow\Omega_*^{\mathrm{der}}(T).
\]
At first sight the target might be larger. Derived manifolds provide many new
representatives, including singular zero loci and nontransverse
intersections. The comparison theorem says that these representatives do not
produce new bordism classes.

\begin{theorem}[Spivak's bordism comparison]
  \label[theorem]{thm:Spivak_Derived_Bordism_Comparison}
  For every manifold $T$, inclusion induces an isomorphism
  \[
    i_*
    \colon
    \Omega_*(T)
    \iso
    \Omega_*^{\mathrm{der}}(T).
  \]
\end{theorem}

The absolute assertion is Spivak's Theorem 3.12, and the relative assertion is
Corollary 3.13
\cite[Theorem~3.12 and Corollary~3.13]{Spivak_Derived_Smooth_Manifolds}.
The theorem says that every compact quasi-smooth derived manifold is
\emph{derived cobordant} to an ordinary manifold. It does not say that it is
equivalent to one.

\subsection{Global equations for a compact derived manifold}
\label[subsection]{sec:Global_Equations_Compact_Derived_Manifold}

The proof reduces the general case to the construction from the first
section. Two geometric inputs make this possible.

\begin{theorem}[Embedding and global zero-locus presentation]
  \label[theorem]{thm:Compact_Derived_Manifold_Global_Zero_Locus}
  Let $X$ be a compact quasi-smooth derived manifold in Spivak's framework.
  Then there are:
  \begin{enumerate}
    \item An embedding $j\colon X\to\R^N$.
    \item An open neighborhood $U\subseteq\R^N$ of the underlying space of
      $X$.
    \item A finite-rank vector bundle $E\to U$.
    \item A section $s\colon U\to E$.
  \end{enumerate}
  With this data, there is an equivalence over $U$
  \[
    X\simeq Z^{\mathrm{der}}(s).
  \]
\end{theorem}

\begin{proof}[Source and proof architecture]
  The embedding is Spivak's Proposition 3.3. Once $X$ is embedded, the
  normal-bundle axiom, Definition 2.1(7), supplies the neighborhood $U$, the
  vector bundle, and the defining section
  \cite[Definition~2.1(7) and Proposition~3.3]
  {Spivak_Derived_Smooth_Manifolds}. The proof of the embedding theorem is a
  derived version of the familiar partition-of-unity proof for compact
  manifolds. We use the result as a quoted geometric input.
\end{proof}

There is a small but essential difference from the opening construction. The
ambient manifold $U$ need not be compact. Compactness of $X$ nevertheless
allows the perturbation to be supported near its zero locus.

Choose a compact subset $K\subset U$ whose interior contains the underlying
space of $X$. Since $s$ has no zeros outside the interior of $K$, relative
transversality gives a transverse section $s'$ which agrees with $s$ outside
$K$. The straight-line homotopy, made constant near its endpoints, is also
fixed outside $K$. Every zero of every section in the homotopy therefore lies
in $K$, which restores the properness argument.

\subsection{Proof of the comparison theorem}
\label[subsection]{sec:Proof_Derived_Bordism_Comparison}

\begin{proof}[Proof of \Cref{thm:Spivak_Derived_Bordism_Comparison}]
  We first prove surjectivity in the absolute case. Let $X$ be a compact
  quasi-smooth derived manifold. Choose a global zero-locus presentation
  \[
    X\simeq Z^{\mathrm{der}}(s)
  \]
  as in \Cref{thm:Compact_Derived_Manifold_Global_Zero_Locus}. Perturb $s$ to
  a transverse section $s'$ while keeping the perturbation fixed outside a
  compact neighborhood of the zero locus. The homotopy-zero-locus
  construction of \Cref{prop:Homotopy_Of_Sections_Derived_Cobordism} gives a
  derived cobordism
  \[
    X\sim_{\mathrm{bord}}Z(s').
  \]
  The endpoint $Z(s')$ is an ordinary compact manifold. Hence every derived
  bordism class is in the image of $i_*$.

  For injectivity, suppose that ordinary compact manifolds $M_0$ and $M_1$
  are derived cobordant. Let
  \[
    p\colon W\longrightarrow\R
  \]
  be a derived cobordism between them. Restrict to the part of $W$ over an
  open interval containing $[0,1]$. Spivak's relative embedding theorem
  embeds this restricted cobordism over the parameter into
  $\R^N\times\R$ \cite[Corollary~3.4]{Spivak_Derived_Smooth_Manifolds}.
  The normal-bundle axiom presents it as the derived zero locus of a section.
  Because $p$ is collared, this presentation is product-like near $M_0$ and
  $M_1$. Relative transversality permits us to perturb the defining section
  while keeping those collared regions fixed. The perturbed zero locus is an
  ordinary smooth cobordism from $M_0$ to $M_1$. Thus two ordinary manifolds
  which become derived cobordant were already classically cobordant, and
  $i_*$ is injective.

  For the relative theorem over $T$, one performs the same constructions
  while retaining the map to $T$. The only additional point is that a map
  from the compact zero locus to the manifold $T$ extends to a neighborhood
  in the ambient Euclidean space. The perturbation and cobordism can then be
  carried out over $T$.
\end{proof}

The two halves of the proof have distinct meanings. Surjectivity says that
derived representatives can be regularized without changing their bordism
class. Injectivity says that allowing derived cobordisms does not identify two
classical bordism classes which were previously distinct.

\begin{example}[A singular fiber of a Morse function]
  \label[example]{ex:Figure_Eight_Derived_Bordism}
  Let $f\colon T^2\to\R$ be a Morse function and let $c$ be a critical value
  of index one. The underlying space of the derived fiber
  \[
    T^2\mathbin{\times_{\R}^{\mathrm{der}}}\{c\}
  \]
  is a figure eight. It is not the underlying space of a one-dimensional
  manifold. Moving $c$ slightly to either side produces a regular fiber,
  either one circle or two circles depending on the side. The parameter
  interval gives derived cobordisms from the singular fiber to both regular
  fibers. Thus the two regular fibers are classically cobordant, while the
  singular derived fiber remains a distinct geometric object. This is one of
  Spivak's examples \cite[Example~2.7]{Spivak_Derived_Smooth_Manifolds}.
\end{example}

\section{Intersection products and Euler classes}
\label[section]{sec:Intersection_Products_And_Euler_Classes}

\subsection{The nontransverse cup-product formula}
\label[subsection]{sec:Nontransverse_Cup_Product_Formula}

Let $M$ be a closed manifold, and let $A,B\subseteq M$ be closed
submanifolds. We use cohomological notation for their Poincare-dual
unoriented bordism classes:
\[
  [A]\in MO^{\operatorname{codim}(A)}(M),
  \qquad
  [B]\in MO^{\operatorname{codim}(B)}(M).
\]
Equivalently, one may work homologically with the proper maps $A\to M$ and
$B\to M$ and then apply bordism Poincare duality. This convention prevents a
common ambiguity: the intersection formula is a cup-product formula in
cohomology, while the geometric representatives are manifolds over $M$.

If $A$ and $B$ are transverse, their ordinary intersection represents the
cup product. Derived geometry removes the transversality condition from the
geometric representative.

\begin{theorem}[General cup-product formula]
  \label[theorem]{thm:General_Cup_Product_Derived_Intersection}
  Let $A$ and $B$ be closed submanifolds of a closed manifold $M$. Then
  \[
    [A]\smile[B]
    =
    [A\mathbin{\times_M^{\mathrm{der}}}B]
  \]
  in the unoriented bordism cohomology of $M$.
\end{theorem}

\begin{proof}
  Perturb the inclusion $A\to M$ through smooth maps until it is transverse
  to $B\to M$. Denote the perturbed map by $A'\to M$. Pulling the homotopy
  back along $B\to M$ gives a derived cobordism over $M$ from
  \[
    A\mathbin{\times_M^{\mathrm{der}}}B
  \]
  to the ordinary transverse fiber product $A'\pitchfork B$. Hence the two
  objects determine the same class over $M$. The classical transverse
  intersection formula gives
  \[
    [A]\smile[B]=[A'\pitchfork B].
  \]
  Combining the two equalities proves the result.
\end{proof}

This is Spivak's Corollary 3.13
\cite[Corollary~3.13]{Spivak_Derived_Smooth_Manifolds}. Notice where
transversality remains in the argument. It is not needed to form the derived
intersection. It is used only to identify the bordism class of that canonical
object with the familiar classical product.

\subsection{The self-intersection of a zero section}
\label[subsection]{sec:Zero_Section_Self_Intersection}

Let $M$ be a closed $m$-manifold and let $q\colon E\to M$ be a rank-$r$
vector bundle. Write $z\colon M\to E$ for the zero section. Its derived
self-intersection is
\[
  X_E
  =
  M\mathbin{\times_E^{\mathrm{der}}}M.
\]
The two maps to $E$ in this expression are both $z$. The underlying space of
$X_E$ is all of $M$, but the derived object remembers that the same normal
equation has been imposed twice.

Choose a transverse section $s'\colon M\to E$. Its graph is a perturbation of
one copy of the zero section. The transverse intersection of $z(M)$ and
$s'(M)$ is precisely the zero manifold $Z(s')$. The homotopy from $z$ to
$s'$ therefore gives a derived cobordism over $M$
\[
  X_E
  \sim_{\mathrm{bord}}
  Z(s').
\]

The bordism Euler class
\[
  e_{MO}(E)\in MO^r(M)
\]
is represented by the zero locus of a transverse section. We have proved the
following self-intersection formula.

\begin{corollary}[Zero-section self-intersection]
  \label[corollary]{cor:Zero_Section_Self_Intersection_Euler_Class}
  The cohomological bordism class represented by the derived
  self-intersection of the zero section is the Euler class:
  \[
    [M\mathbin{\times_E^{\mathrm{der}}}M]
    =
    e_{MO}(E).
  \]
  In homological notation, the corresponding class over $M$ is
  \[
    e_{MO}(E)\cap[M]
    =
    [Z(s')\longrightarrow M]
    \in MO_{m-r}(M).
  \]
\end{corollary}

The derived self-intersection is canonical, whereas the transverse section
$s'$ is a choice. Any two choices are joined by a homotopy of sections, so
their zero manifolds are cobordant. Derived geometry packages the class before
the choice of perturbation is made.

\begin{example}[The trivial line bundle]
  \label[example]{ex:Trivial_Line_Derived_Self_Intersection}
  Let $E=M\times\R$ be the trivial line bundle. Its derived zero-section
  self-intersection is
  \[
    X_{\underline{\R}}
    =
    M\mathbin{\times_{M\times\R}^{\mathrm{der}}}M.
  \]
  Locally, its algebra of functions is the Koszul algebra
  \[
    \bigl(\Cinfty(U)\otimes\Lambda(\eta),\partial\eta=0\bigr).
  \]
  Thus its underlying space and classical truncation are both $M$, while it
  has a nonzero degree-one class generated by $\eta$.

  On the other hand, the zero section can be perturbed to the nowhere-zero
  constant section $1$. Its zero locus is empty. Hence
  \[
    [X_{\underline{\R}}]=0
  \]
  in bordism. The derived object is visibly nonempty and nonclassical, but its
  bordism class vanishes. This is the cleanest example of the information
  which bordism deliberately forgets.
\end{example}

\section{Fundamental classes and their limits}
\label[section]{sec:Fundamental_Classes_And_Their_Limits}

\subsection{From bordism to generalized homology}
\label[subsection]{sec:From_Bordism_To_Generalized_Homology}

Let $X$ be a compact quasi-smooth derived manifold of virtual dimension $d$.
The comparison theorem supplies a class
\[
  [X]\in\Omega_d^{\mathrm{der}}
  \cong
  \Omega_d
  =MO_d.
\]
More generally, a proper map $X\to T$ supplies a class in $MO_d(T)$.

Suppose that a generalized homology theory $\mathbb E_*$ is equipped with a
morphism
\[
  MO\longrightarrow\mathbb E
\]
between the stable objects which represent the two theories. This induces a
natural transformation
\[
  MO_*(T)\longrightarrow\mathbb E_*(T).
\]

\begin{definition}[Generalized fundamental class]
  \label[definition]{def:Generalized_Fundamental_Class_Derived_Manifold}
  The image of $[X\to T]$ under this natural transformation is the
  $\mathbb E$-homology fundamental class of $X$ over $T$:
  \[
    [X\to T]_{\mathbb E}
    \in
    \mathbb E_d(T).
  \]
\end{definition}

For the live talk, the displayed morphism should be treated as the datum that
makes ordinary unoriented manifold fundamental classes available in
$\mathbb E$-homology. The point of the comparison theorem is then immediate:
the same datum assigns compatible classes to compact quasi-smooth derived
manifolds. This is the observation made in Spivak's introduction
\cite[Section~1]{Spivak_Derived_Smooth_Manifolds}.

\subsection{Orientations}
\label[subsection]{sec:Orientations_Derived_Bordism}

Everything above is unoriented. This is not a cosmetic convention. To obtain
an integral or otherwise oriented class, one must supply orientation data.
Spivak formulates an orientation using a normal bundle. A compact derived
manifold is embedded in some $\R^N$, and the normal-bundle axiom supplies a
normal vector bundle. After adding trivial summands, this normal bundle is
stable under changing the embedding. An orientation of the stable normal
bundle is the input for oriented derived bordism
\cite[Remark~3.8]{Spivak_Derived_Smooth_Manifolds}.

In the oriented case, $MO$ is replaced by the appropriate Thom object, for
example $MSO$ for oriented bordism. A morphism from that object to
$\mathbb E$ transports the oriented bordism class to $\mathbb E$-homology.
Nothing in the existence of a derived zero locus automatically provides such
an orientation.

This brief normal-bundle discussion also points forward. Bordism only retains
the stable normal information needed for a class. The next chapter will
construct a finer stable linear object intrinsic to the derived manifold: its
tangent and cotangent complexes.

\subsection{What has not been constructed}
\label[subsection]{sec:What_Derived_Bordism_Does_Not_Construct}

The phrase \enquote{fundamental class} can obscure several separate
requirements. The construction in this chapter requires a compact
quasi-smooth derived manifold, or more generally a proper map to the chosen
target. It does not prove that a moduli problem is compact. It does not add
limiting objects to a noncompact moduli space. It does not orient a determinant
line. It also does not treat isotropy or produce a virtual integration theory
for derived orbifolds or stacks.

These distinctions matter for the elliptic moduli problems later in the
seminar. A representability theorem has the form
\[
  \text{moduli functor}
  \simeq
  \text{quasi-smooth derived geometric object}.
\]
It does not by itself imply
\[
  \text{compactness},
  \qquad
  \text{orientation},
  \qquad
  \text{virtual fundamental class}.
\]
If the represented object is later shown to be compact and is equipped with
the required normal orientation, derived bordism becomes one possible route
to a class. In applications with automorphisms, compactifications, or more
general Kuranishi structures, additional theory is needed.

We can now formulate the handoff to the next chapter. The trivial-line
self-intersection in \Cref{ex:Trivial_Line_Derived_Self_Intersection} has zero
bordism class, but its derived structure has not disappeared. What intrinsic
first-order invariant retains that structure, and how does it record
deformations and obstructions? The answer will be the stable tangent complex.

\section{Summary and supplementary materials}
\label[section]{sec:Derived_Bordism_Summary_And_Supplements}

\subsection{Takeaways}
\label[subsection]{sec:Derived_Bordism_Takeaways}

The live route through the chapter can be compressed into seven statements.

\begin{enumerate}
  \item A homotopy of vector-bundle sections gives a derived family of their
    zero loci.
  \item Making the homotopy constant near its endpoints gives collars, while
    compactness or compact support gives properness.
  \item Every compact quasi-smooth derived manifold has a global zero-locus
    presentation in Spivak's framework.
  \item Inclusion induces an isomorphism from classical unoriented bordism to
    derived bordism.
  \item The theorem compares bordism classes, not derived manifolds up to
    equivalence.
  \item A nontransverse derived intersection represents the classical
    cup product, and a zero-section self-intersection represents the bordism
    Euler class.
  \item A map from $MO$ to a generalized homology theory transports the
    bordism class to a fundamental class, while compactness and orientation
    remain additional inputs.
\end{enumerate}

The remainder of the chapter is supplementary. It supplies proofs and
examples which are useful for reference but are not part of the default
70-minute route.

\subsection{Supplement: derived cobordism is an equivalence relation}
\label[subsection]{sec:Supplement_Derived_Cobordism_Equivalence_Relation}

\begin{proposition}
  \label[proposition]{prop:Derived_Cobordism_Equivalence_Relation}
  Derived cobordism is an equivalence relation on compact quasi-smooth
  derived manifolds.
\end{proposition}

\begin{proof}
  Reflexivity is represented by the projection
  \[
    X\times\R\longrightarrow\R.
  \]
  This map is proper when $X$ is compact, is collared everywhere, and has
  fiber $X$ over both $0$ and $1$.

  Symmetry follows by reversing the parameter. If $p\colon W\to\R$ is a
  cobordism from $X_0$ to $X_1$, then $1-p$ is a cobordism from $X_1$ to
  $X_0$.

  For transitivity, let $W_{01}$ be a cobordism from $X_0$ to $X_1$ and let
  $W_{12}$ be a cobordism from $X_1$ to $X_2$. Their collars identify open
  neighborhoods of the common endpoint with
  $X_1\times(-\varepsilon,\varepsilon)$. After translating the parameters,
  glue the two derived manifolds along these equivalent open subobjects. The
  product descriptions ensure that the parameter maps agree on the overlap.
  The resulting map is proper and collared at the remaining endpoints, hence
  is a derived cobordism from $X_0$ to $X_2$.
\end{proof}

This is Spivak's Proposition 3.10
\cite[Proposition~3.10]{Spivak_Derived_Smooth_Manifolds}. The collar condition
is exactly what makes the gluing step formal.

\subsection{Supplement: the relative injectivity argument}
\label[subsection]{sec:Supplement_Relative_Injectivity_Derived_Bordism}

The injectivity proof hides two relative statements. First, a derived
cobordism $p\colon W\to\R$ can be embedded over the parameter, after
restricting to a neighborhood of $[0,1]$:
\[
\begin{tikzcd}
  W' \rar[hook,"j"] \dar["p"']
    & \R^N\times(-1,2) \dar["\operatorname{pr}_2"] \\
  (-1,2) \rar[equal] & (-1,2).
\end{tikzcd}
\]
The map to the parameter remains proper. This is the role of Spivak's
Corollary 3.4.

Second, the defining section for $W'$ may be perturbed relative to the
collared regions. In those regions the endpoint fibers are already ordinary
manifolds, and the section is already transverse in the normal directions.
Relative transversality changes the section only away from those regions. Its
ordinary zero locus is therefore a smooth cobordism with the same endpoints.

For a cobordism over $T$, the map from the compact part of the zero locus to
$T$ is extended to a neighborhood before perturbing. All subsequent
constructions retain this extended map. This is why the absolute comparison
upgrades to an isomorphism of bordism homology theories, rather than merely an
isomorphism of coefficient groups.

\subsection{Supplement: the Pontryagin--Thom picture}
\label[subsection]{sec:Supplement_Pontryagin_Thom_Derived_Bordism}

For an ordinary compact manifold $N$ embedded in $\R^k$, the
Pontryagin--Thom collapse produces a stable map from a sphere to the Thom
space of the normal bundle. Stabilizing over all embeddings yields the class
of $N$ in $MO_*$. This construction explains why the normal bundle, rather
than a chosen tangent bundle, is the natural input for unoriented bordism.

For a compact quasi-smooth derived manifold $X$, the embedding and
normal-bundle theorem again supplies stable normal data. A transverse
perturbation of a defining section produces an ordinary manifold $N$ derived
cobordant to $X$. Define the Pontryagin--Thom class of $X$ to be the class of
$N$. If a second perturbation produces $N'$, a homotopy of perturbations gives
a derived cobordism from $N$ to $N'$. Injectivity in
\Cref{thm:Spivak_Derived_Bordism_Comparison} turns this into an ordinary
cobordism, so the stable class is independent of the perturbation.

This picture is useful conceptually, but it should not replace the
homotopy-of-sections construction in the live talk. The latter makes the
geometric content visible without requiring an introduction to stable
homotopy theory.

\subsection{Supplement: the Möbius line bundle}
\label[subsection]{sec:Supplement_Mobius_Euler_Class}

Let $L\to S^1$ be the Möbius line bundle. A section can be described by a
smooth function $f\colon\R\to\R$ satisfying
\[
  f(t+1)=-f(t).
\]
The function $f(t)=\sin(\pi t)$ defines a transverse section. Modulo the
identification $t\sim t+1$, it has one zero. Therefore the derived
self-intersection of the zero section is derived cobordant to one point over
$S^1$. Its mod-$2$ Euler number is nonzero.

This complements the trivial-line example. For the trivial line, a
nowhere-zero section makes the self-intersection class vanish. For the Möbius
line, every section has a zero and the transverse zero set detects the
nontrivial unoriented Euler class.

\subsection{Related literature}
\label[subsection]{sec:Derived_Bordism_Related_Literature}

The definitions of collared derived cobordism, the embedding theorem, the
comparison with ordinary bordism, and the general cup-product formula are all
from Spivak's \emph{Derived smooth manifolds}
\cite[Sections~1--3]{Spivak_Derived_Smooth_Manifolds}. The comparison which
places Spivak's model inside the universal modern framework is
Carchedi--Steffens, especially their comparison with Spivak in Section 5.2
\cite[Section~5.2]{Carchedi_Steffens_Universal_Derived_Manifolds}.

Joyce's theory of d-manifolds and d-orbifolds develops oriented bordism and
virtual classes in a framework which also accommodates finite isotropy. Those
results are relevant to later moduli problems, but they require definitions
and hypotheses not developed in this chapter. In particular, they should not
be conflated with the bordism comparison for compact quasi-smooth derived
manifolds proved here.

% !TEX root = ../main.tex

\chapter{Quasi-smooth derived geometry and Kuranishi charts}
\label[chapter]{chap:Local_Derived_Geometry_Kuranishi_Charts}

The preceding chapters described the functions and tangent complexes of
derived zero loci. We now ask when a derived object admits such a
presentation locally. The answer is detected by its cotangent complex:
the quasi-smooth objects are precisely those locally presented as zero
loci of sections of finite-rank vector bundles.

To make this local statement precise, we first pass from animated smooth
rings to their spectra and structure sheaves. Smooth localization supplies
the restriction maps, and gluing affine objects leads to derived
$\Cinfty$-schemes. We then state the local zero-locus theorem and study
its presentations, called Kuranishi charts. Stabilization and changes of
bundle frame illustrate how different equations can describe the same
derived object. These comparisons will be used later in the finite-dimensional
reduction of elliptic equations.

\section{Smooth localization}
\label[section]{sec:Localization_To_Derived_Spectra}

For a smooth function $a$ on a manifold $M$, put
$D(a)=\{x\in M:a(x)\neq0\}$. The smooth localization satisfies
\[
 \Cinfty(M)[a^{-1}]\cong\Cinfty(D(a)).
\]
One way to see this is to identify $D(a)$ with the graph of $1/a$ in
$M\times\R$. This graph is the transverse zero locus of
$ya(x)-1$, so adjoining a smooth inverse to $a$ gives precisely its
function ring.

Smooth localization is larger than algebraic localization of the underlying
ring. For instance, $\Cinfty(\R)[x^{-1}]$ contains the smooth function
$\exp(1/x^2)$ on the punctured line. It is not a fraction $g(x)/x^n$
with $g$ smooth on all of $\R$: multiplication by any power of $x$
leaves its divergence at zero. Smooth functional calculus is part of
the localization.

Let now $A$ be an animated $\Cinfty$-ring and $a\in\pi_0A$.

\begin{definition}
\label[definition]{def:Derived_Localization}
The smooth localization $A[a^{-1}]$ is initial among animated smooth
rings under $A$ in which the image of $a$ is invertible in $\pi_0$.
\end{definition}

Choose a representative of $a$, hence a map $\Cinfty(\R)\to A$.
The localization is the pushout
\[
 \begin{tikzcd}
 \Cinfty(\R)\rar\dar\drar[pushout]&A\dar\\
 \Cinfty(\R\setminus\{0\})\rar&A[a^{-1}].
 \end{tikzcd}
\]
The universal property makes the result independent of this choice.

\begin{theorem}[Flatness of smooth localization]
\label[theorem]{thm:Homotopy_Groups_Smooth_Localization}
There are natural isomorphisms
\[
 \pi_0(A[a^{-1}])\cong(\pi_0A)[a^{-1}],\qquad
 \pi_i(A[a^{-1}])
 \cong\pi_i(A)\otimes_{\pi_0A}(\pi_0A)[a^{-1}].
\]
The smooth localization of the ordinary ring $\pi_0A$ is flat as
a module over its underlying ring.
\end{theorem}

We use this result without proof
\cite[Proposition~4.1.3.13]{Steffens_Derived_Differential_Geometry}.
The tensor product in its second formula is an ordinary module tensor
product. Flatness explains why restricting functions creates no
additional higher homotopy from a derived tensor product.
It restricts each existing homotopy module.

The corresponding relative cotangent complex vanishes:
\[
 \mathbb L_{A[a^{-1}]/A}\cong0.
\]
By transitivity,
\begin{equation}\label{eq:Cotangent_Complex_Localizes}
 \mathbb L_{A[a^{-1}]}\cong\mathbb L_A\otimes_A A[a^{-1}].
\end{equation}
This is the compatibility needed to regard cotangent complexes as local
geometric objects.

\section{The spectrum and its structure sheaf}

Assume $A$ is finitely presented as an animated smooth ring.
Its spectrum has underlying set
\[
 |\Spec A|=\Hom_{\Cinfty\text{-rings}}(\pi_0A,\R).
\]
These are real points, not arbitrary prime ideals of an underlying
commutative ring. A point determines a maximal ideal with residue field
$\R$, namely the kernel of its evaluation map. The evaluation map
is the useful datum in smooth geometry.

For $a\in\pi_0A$, the subsets
\[
 D(a)=\{x:x(a)\neq0\}
\]
form a basis for a topology, since $D(a)\cap D(b)=D(ab)$.
If $\pi_0A=\Cinfty(\R^n)/I$, evaluation identifies this topological
space with the common zero set $Z(I)\subseteq\R^n$, with its subspace
topology. To check the topology, bump functions give basic
neighbourhoods inside any prescribed Euclidean open set.

\begin{construction}[Derived spectrum]
\label[construction]{con:Derived_Spectrum}
The derived spectrum is
\[
 \Spec A=(|\Spec A|,\mathcal O_{\Spec A}),
\]
where the sheaf of animated smooth rings is obtained from smooth
localizations. In the finite-presentation setting,
\[
 \mathcal O_{\Spec A}(D(a))\cong A[a^{-1}].
\]
The local ring at $x$ is the filtered colimit of these localizations
over basic neighbourhoods of $x$.
\end{construction}

The assertion that these assignments have the required sheaf and
reconstruction properties is a theorem about smooth rings.
It is part of the affine comparison
\cite[Section~4.1.3, especially Corollary~4.1.3.34]
{Steffens_Derived_Differential_Geometry}.
For unrestricted smooth rings one must take care with sheafification
and reconstruction; finite presentation is the class used here.

The term \emph{locally animated $\Cinfty$-ringed space} means that
$\pi_0$ of every stalk is a local smooth ring. At $x$ its residue map
is evaluation in $\R$. The higher homotopy modules remain in the stalk.
A stalk belongs to a point; the values on basic opens are the
localizations before taking that colimit.

For an ordinary manifold, the construction recovers its familiar
ringed space:
\[
 \Spec\Cinfty(M)\cong(M,\Cinfty_M).
\]
Real homomorphisms are evaluations at points, the basic opens give
the manifold topology, and the localization formula identifies the
structure sheaf with smooth functions.

The examples from \Cref{chap:Derived_Zero_Loci_And_Tangent_Complexes} now have an explicit topological meaning.
The spectra of $\R$, $\R[\delta]/(\delta^2)$, and
$\R[\varepsilon]$ with $|\varepsilon|=1$, $d=0$, all have one point.
Their structure sheaves differ. The second has nilpotents in degree zero;
the third has a nonzero homotopy group in degree one.

\section{Derived schemes and quasi-smoothness}

The spectrum construction supplies affine objects and their open subsets.
We can therefore define locally affine objects, then use the cotangent
complex to single out those that should admit local presentations by
vector bundle sections.

An open subset $V\subseteq|X|$ of a locally ringed object defines
an open subobject $(V,\mathcal O_X|_V)$. In particular,
\[
 \Spec A[a^{-1}]
 \cong(D(a),\mathcal O_{\Spec A}|_{D(a)}).
\]
Thus algebraic localization becomes geometric open restriction.

\begin{definition}
\label[definition]{def:Derived_Cinfty_Scheme}
A derived $\Cinfty$-scheme locally of finite presentation is a locally
animated $\Cinfty$-ringed space with an open cover by spectra of
finitely presented animated $\Cinfty$-rings.
\end{definition}

The affine objects in this definition are exactly the derived manifolds
of the universal-property convention in \Cref{chap:From_Smooth_Algebra_To_Derived_Manifolds}:
\[
 \mathrm{DMfd}\simeq
 \bigl(\operatorname{Alg}_{\Cinfty}(\An)^{\mathrm{fp}}\bigr)^{\mathrm{op}}.
\]
In particular, an ordinary manifold is affine in this sense.
Affineness does not mean that it is a single Euclidean coordinate chart.
For example, $S^1=\Spec\Cinfty(S^1)$ is affine.
Allowing locally affine objects supplies the category in which
open gluing can be performed.

The classical truncation is
\[
 t_0X=(|X|,\pi_0\mathcal O_X).
\]
It keeps the topology and removes higher homotopy from the structure
sheaf. It can still contain nilpotents. Thus classical truncation and
the underlying zero set are different operations.

Cotangent complexes restrict by
\Cref{eq:Cotangent_Complex_Localizes}, so they glue over affine opens.
A module over an animated ring $A$ is \emph{perfect} if it is a retract
of a finite construction by cofibers from shifts of finite sums of $A$.
For an ordinary ring this is equivalent to having a bounded complex
of finitely generated projective modules as a presentation. Perfection
on a derived scheme is checked on its affine opens. A module has Tor-amplitude
in $[a,b]$ if tensoring with any discrete module produces homology
only in degrees $a,\ldots,b$.
This condition concerns derived tensor products; it is stronger than
merely locating the homology of the original complex.

\begin{definition}[Quasi-smoothness]
\label[definition]{def:Quasi_Smooth_Derived_Cinfty_Scheme}
A derived $\Cinfty$-scheme locally of finite presentation is
\emph{quasi-smooth} if its cotangent complex is perfect of
Tor-amplitude contained in $[0,1]$.
For a morphism, the same definition uses the relative cotangent complex.
\end{definition}

Perfection and Tor-amplitude can be checked locally.
A two-term complex of finite-rank vector bundles in degrees $1,0$
has this amplitude. The cotangent calculation of \Cref{chap:Stable_Linearization_And_Tangent_Complexes} therefore
proves that every derived zero locus of a section of a finite-rank
bundle is quasi-smooth.

For a perfect complex the amplitude condition can also be tested on
residue fields. Consequently, at each real point of a quasi-smooth
object its tangent complex has homology only in degrees $0,-1$.
The bundle presentation of the cotangent complex, rather than just
a dimension count at individual points, is the local structure behind
this assertion.

\section{Local presentations by Kuranishi charts}
\label[section]{sec:Kuranishi_Charts_As_Presentations}

The cotangent calculation already shows that a vector bundle zero locus
is quasi-smooth. The local normal-form theorem supplies the converse.
We first specify what data is included in a chart.

\begin{definition}[Kuranishi chart]
\label[definition]{def:Kuranishi_Chart}
A Kuranishi presentation is a triple $(U,E,s)$ consisting of a
finite-dimensional manifold, a finite-rank vector bundle on it,
and a smooth section. Its associated derived object is
$Z^{\mathrm{der}}(s)$.
A Kuranishi chart for $X$ also includes an identification of this
derived zero locus with an open subobject of $X$.
\end{definition}

This usage concerns charts without isotropy, boundary, or corners.
It is the local presentation needed in the representability proof.
A choice of charts alone does not specify any of the more elaborate
global notions called Kuranishi structures in the literature.

\begin{theorem}[Local zero-locus form]
\label[theorem]{thm:Local_Zero_Locus_Form}
A derived $\Cinfty$-scheme locally of finite presentation is
quasi-smooth if and only if it admits an open cover by Kuranishi charts.
\end{theorem}

The forward implication is the substantive normal-form theorem.
The affine statement is proved in
\citet[Proposition~5.1.1.11 and Corollary~5.1.1.13]
{Steffens_Derived_Differential_Geometry}.
In that statement, a finitely presented quasi-smooth affine object
can be expressed as the zero locus of a bundle section over an open
subset of a Euclidean space. Applying it on affine opens proves the
displayed local form. The reverse implication is the cotangent calculation.

The proof begins by choosing generators for the ring, giving an
ambient smooth affine object. The relative cotangent complex describes
its equations. The amplitude assumption implies that the shifted
relative cotangent complex is a finite projective module, which can
be extended to a vector bundle on an ambient neighbourhood.
Lifting the relations then produces a section of that bundle.
A cotangent criterion identifies its derived zero locus with the
original object. This explains the role of the amplitude assumption;
the extension and comparison theorems in this argument are quoted.

At a point of a chart, the tangent complex is
$[T_xU\to E_x]$, and the virtual dimension is
$\dim U-\operatorname{rank}E$. Quasi-smoothness permits both singular
classical truncations and nonzero higher homotopy. The zero locus of
$x^2$ illustrates the first, and a zero equation illustrates the second.

\section{Changing a presentation}

The quotient rings, Koszul algebras, and tangent complexes calculated
earlier distinguish different geometric objects. We now use those
calculations to recognize changes of equations that preserve an object.

Let $f\colon U\to\R^r$, and define
\[
 f'\colon U\times\R^q\longrightarrow\R^r\times\R^q,
 \qquad f'(x,y)=(f(x),y).
\]
The additional equation forces the additional variable to zero.

\begin{proposition}[Stabilization]
\label[proposition]{prop:Stabilization_Kuranishi_Chart}
Projection onto $U$ induces an isomorphism
$Z^{\mathrm{der}}(f')\cong Z^{\mathrm{der}}(f)$.
\end{proposition}

\begin{proof}
The second factor is the derived zero locus of the identity of $\R^q$.
That map is transverse to zero, so its zero locus is the ordinary point.
Pasting pullbacks gives the stated isomorphism.
Equivalently, the new Koszul variables satisfy $d\eta_j=y_j$.
Evaluation at $y=\eta=0$ removes an acyclic Koszul resolution.
\end{proof}

For example, $x^2$ and $(x^2,y)$ give the same derived object.
Their tangent presentations differ by
$[\R\xrightarrow1\R]$, an acyclic summand.
Both ambient dimension and the number of equations increase by one,
leaving virtual dimension unchanged.

Changing a bundle frame likewise changes the list of equations
without changing the object. If $g(x)$ is an invertible matrix of
smooth functions, then $f$ and $g f$ present isomorphic derived zero loci.
Explicitly, the map from the Koszul algebra for $gf$ to that for $f$
sends its generator $\varepsilon'_i$ to $\sum_jg_{ij}\varepsilon_j$.
The inverse matrix gives the inverse dg-algebra map.
The intrinsic bundle model with differential $\iota_s$ makes this
independence immediate.

In contrast, $x$ and $(x,0)$ have the same ordinary zero set but
different derived zero loci. The extra zero equation contributes a
degree-minus-one tangent direction and a degree-one homotopy class of functions.
Nor does equality of tangent complexes suffice: $x^2$ and $x^3$
have the same tangent complex at zero but distinct classical
function rings. Comparisons of charts must respect their equations.

\section{Further reading}

Two refinements of the chart description will recur later: removing
transverse directions makes a chart minimal at a point, while a
nontrivial obstruction bundle shows what local frames can conceal.

A useful refinement makes a chart minimal at a chosen point.
Suppose $f\colon U\subseteq\R^n\to\R^r$ vanishes at $x$ and $d_xf$
has rank $q$. After linear changes of coordinates, $q$ components
of $f$ have independent differentials. The inverse function theorem
allows these components to be used as coordinates $y\in\R^q$.
Writing the remaining coordinates as $z$, the map becomes
\[
 (y,z)\longmapsto(y,g(y,z)).
\]
Hadamard's lemma gives a smooth matrix $H$ with
$g(y,z)=g(0,z)+H(y,z)y$.
An invertible change of equation generators subtracts $H(y,z)y$,
giving $(y,g(0,z))$. Stabilization now removes the coordinates $y$.

The resulting local chart is the zero locus of
\[
 \kappa\colon\R^{n-q}\supseteq V\longrightarrow\R^{r-q},
 \qquad\kappa(z)=g(0,z),
\]
with $d\kappa$ zero at the chosen point.
Its source tangent and target identify with
$\ker d_xf$ and $\operatorname{coker}d_xf$.
These identifications use choices; the derived object they present is
the original open neighbourhood. The analytic reduction of an elliptic
equation will produce a chart of this kind from an initially
infinite-dimensional problem.

\begin{example}[An obstruction bundle with no global frame]
Let $L\to S^1$ be the M\"obius line bundle and form
$X=S^1\times_L S^1$ using two zero sections.
The intrinsic Koszul algebra is the exterior algebra of $L^\vee$
with zero differential. Therefore
\[
 \pi_0\mathcal O_X=\Cinfty_{S^1},\qquad
 \pi_1\mathcal O_X=\Gamma(-,L^\vee).
\]
Each trivializing open gives a single zero equation. On overlaps the
generator changes by the transition function of $L^\vee$.
There is no globally nonvanishing generator, although the derived
self-intersection is defined globally.

After restriction to the classical circle, the tangent complex is
$[TS^1\xrightarrow0 L]$. Its obstruction sheaf remembers the
nontrivial line bundle. This shows why vector bundles, rather than only
fixed vector spaces of equations, naturally occur in Kuranishi charts.
\end{example}

\begin{exercise}
Show that the derived zero locus of $(x,xy)\colon\R^2\to\R^2$
is locally, and in this case globally, the product of a line and
the derived zero locus of the zero map from a point to $\R$.
\end{exercise}

\begin{exercise}
For $f(x,y)=x^2+y^2$, calculate the tangent complex at the origin
and compare its virtual dimension with the dimension of its zero set.
Explain why the zero set alone cannot carry this information.
\end{exercise}

For ordinary smooth spectra and their structure sheaves, see
\cite[Sections~2.2 and~3.1]{Joyce_Introduction_Cinfty_Schemes}.
The localization and normal-form statements used here are in
\cite[Sections~4.1.3 and~5.1.1]{Steffens_Derived_Differential_Geometry}.
They provide the passage from the affine universal theory to the
local charts used in geometric moduli problems.

% !TEX root = ../main.tex

% This chapter is the complete working manuscript for the eighth seminar
% talk. The final section contains supplementary material outside the default
% live route.

\chapter{Local derived geometry and Kuranishi charts}
\label[chapter]{chap:Local_Derived_Geometry_Kuranishi_Charts}

Talk 7 ended with the local normal-form slogan
\[
  \text{quasi-smooth}
  \quad\Longleftrightarrow\quad
  \text{locally a derived zero locus}.
\]
The right-hand side is geometric, but the word ``locally'' has not yet been
explained. Our affine description by animated $\Cinfty$-rings provides
pullbacks, Koszul presentations, and cotangent complexes. It does not visibly
display an underlying topological space, open subsets, or restriction of
functions.

The first task of this chapter is therefore to make affine derived objects
local. The basic operation is smooth localization. For a manifold $M$ and a
function $a\in\Cinfty(M)$, inverting $a$ restricts $M$ to the open set where
$a$ does not vanish. For an animated $\Cinfty$-ring $A$, the spectrum
construction organizes all such localizations into a topological space with
a sheaf of animated functions. Its guiding slogan is
\begin{quote}
  The topology of $\Spec A$ is controlled by $\pi_0A$, while its structure
  sheaf retains the full animated algebra $A$.
\end{quote}

The second task is to understand the nonuniqueness of the local zero-locus
presentations produced in this way. A triple
\[
  (U,E,s),
\]
consisting of a manifold, a vector bundle, and a section, is a presentation
of a derived zero locus. It is not the intrinsic derived object itself. We
will see this concretely by comparing
\[
  x\longmapsto x^2
  \qquad\text{and}\qquad
  (x,y)\longmapsto(x^2,y).
\]
The second presentation has one additional variable and one additional
equation, but presents the same derived object. This operation is called
stabilization.

Here $\Spec A$ is a geometric spectrum. It is unrelated to the stable
homotopy-theoretic coefficient objects mentioned in Talk 7. Stable modules
linearize an animated algebra; the geometric spectrum equips that algebra
with points, open restrictions, and a structure sheaf.

\section{From localization to derived spectra}
\label[section]{sec:Localization_To_Derived_Spectra}

\subsection{Open restriction in smooth algebra}
\label[subsection]{sec:Open_Restriction_Smooth_Algebra_Talk_Eight}

Let $M$ be a manifold, let $a\in\Cinfty(M)$, and write
\[
  D(a)=\{x\in M\mid a(x)\neq0\}.
\]
Talk 3 proved that restriction to $D(a)$ is the smooth localization at $a$.

\begin{proposition}[Smooth localization]
  \label[proposition]{prop:Smooth_Localization_Recall_Talk_Eight}
  Restriction induces an isomorphism of $\Cinfty$-rings
  \[
    \Cinfty(M)[a^{-1}]
    \iso
    \Cinfty(D(a)).
  \]
\end{proposition}

The graph
\[
  D(a)\longrightarrow M\times\R,
  \qquad
  x\longmapsto(x,a(x)^{-1})
\]
identifies $D(a)$ with the regular zero locus of
\[
  (x,y)\longmapsto ya(x)-1.
\]
This is the geometric reason that adjoining an inverse to $a$ produces
precisely the functions on $D(a)$.

It matters that this is localization in $\Cinfty$-rings. Algebraically
inverting the coordinate $x$ in the underlying commutative ring of
$\Cinfty(\R)$ produces only fractions $f(x)/x^n$. The smooth localization is
$\Cinfty(\R\setminus\{0\})$ and contains arbitrary smooth functions on the
punctured line, including functions whose growth near zero is not bounded by
any power of $x^{-1}$. The full smooth functional calculus is part of the
localization.

Every open subset of a manifold is a nonvanishing locus. If $V\subseteq M$
is open, there is a smooth function $a\colon M\to[0,\infty)$ such that
$V=D(a)$. Thus ordinary open restriction is completely visible in smooth
algebra.

\subsection{Derived localization}
\label[subsection]{sec:Derived_Localization_Talk_Eight}

Let $A$ be an animated $\Cinfty$-ring and let $a\in\pi_0A$. A morphism
$A\to B$ carries $a$ to a well-defined element of $\pi_0B$, so it makes
sense to require that image to be invertible.

\begin{definition}[Derived localization]
  \label[definition]{def:Derived_Localization_Talk_Eight}
  The \emph{localization of $A$ at $a$} is a morphism
  \[
    A\longrightarrow A[a^{-1}]
  \]
  which is initial among morphisms from $A$ carrying $a$ to an invertible
  element of $\pi_0$.
\end{definition}

This initial object is unique in the homotopical sense: its universal
property includes the homotopies and all their higher compatibilities. It is
computed by the pushout
\[
\begin{tikzcd}
  \Cinfty(\R) \rar \dar \drar[pushout]
    & A \dar \\
  \Cinfty(\R\setminus\{0\}) \rar
    & A[a^{-1}],
\end{tikzcd}
\]
where the upper morphism classifies $a$.

The decisive calculation says that this construction only restricts the
derived information already present.

\begin{theorem}[Homotopy groups of a smooth localization]
  \label[theorem]{thm:Homotopy_Groups_Smooth_Localization_Talk_Eight}
  Let $A$ be an animated $\Cinfty$-ring and let $a\in\pi_0A$. Then
  \[
    \pi_0(A[a^{-1}])
    \cong
    (\pi_0A)[a^{-1}],
  \]
  and for every $n\geq0$ the canonical morphism
  \[
    \pi_n(A)\otimes_{\pi_0A}(\pi_0A)[a^{-1}]
    \longrightarrow
    \pi_n(A[a^{-1}])
  \]
  is an isomorphism.
\end{theorem}

The theorem is Steffens, Proposition~4.1.3.13
\cite[Proposition~4.1.3.13]{Steffens_Derived_Differential_Geometry}. We quote
its proof, which uses flatness of smooth open restriction. The formula says
that localization is \emph{strong}: it creates no new higher homotopy groups.
Each existing homotopy group is simply restricted to the chosen open set.

There is also a linear version of the same assertion. Smooth localization is
formally etale, so
\[
  \mathbb L_{A[a^{-1}]/A}\simeq0.
\]
Transitivity therefore gives
\begin{equation}
  \label{eq:Cotangent_Complex_Localizes_Talk_Eight}
  \mathbb L_{A[a^{-1}]}
  \simeq
  \mathbb L_A\otimes_AA[a^{-1}].
\end{equation}
This is the bridge back to Talk 7: the stable linearization of an open
restriction is the restriction of the stable linearization.

If $A$ is homotopically finitely presented, then so is $A[a^{-1}]$. This
follows either from the displayed pushout or from Steffens,
Remark~4.1.3.14
\cite[Remark~4.1.3.14]{Steffens_Derived_Differential_Geometry}.

\subsection{Points, basic opens, and animated functions}
\label[subsection]{sec:Derived_Spectrum_Slogan_Talk_Eight}

We now organize the localizations of $A$ into a geometric object. Assume for
the moment that $A$ is homotopically finitely presented. Its set of real
points is
\[
  |\Spec A|
  =
  \Hom_{\Cinfty\text{-rings}}(\pi_0A,\R).
\]
For $a\in\pi_0A$, define the basic open
\[
  D(a)
  =
  \{x\colon\pi_0A\to\R\mid x(a)\neq0\}.
\]
These opens generate a topology. On $D(a)$, place the full animated
localization $A[a^{-1}]$ rather than only its discrete truncation.

\begin{construction}[Derived spectrum]
  \label[construction]{con:Derived_Spectrum_Talk_Eight}
  The \emph{derived spectrum} of $A$ is the locally animated
  $\Cinfty$-ringed space
  \[
    \Spec A
    =
    \bigl(|\Spec A|,\mathcal O_{\Spec A}\bigr),
  \]
  whose structure sheaf is obtained from the assignments
  \[
    \mathcal O_{\Spec A}(D(a))
    \simeq
    A[a^{-1}].
  \]
\end{construction}

For the finite-presentation objects used in the live argument, these
assignments already have the expected locality properties. The ordinary
structured-space construction and its extension to animated rings are
discussed in \Cref{sec:Supplement_Real_Spectrum_Talk_Eight}.

Truncation has an especially simple description:
\[
  t_0(\Spec A)
  =
  \bigl(|\Spec A|,\pi_0\mathcal O_{\Spec A}\bigr)
  \cong
  \Spec(\pi_0A).
\]
The topological space does not change. Truncation only discards the higher
homotopy sheaves of functions.

A one-point example makes the distinction visible. Let
\[
  K(0)=\bigl(\R\otimes\Lambda(e),\partial e=0\bigr)
\]
be the animated function algebra of the derived self-intersection
$*\times_{\R}*$. Then
\[
  \pi_0K(0)=\R,
  \qquad
  \pi_1K(0)=\R.
\]
The following three objects have the same one-point underlying space:
\[
\begin{array}{@{}lll@{}}
  \text{object} & \pi_0\text{ of the stalk} & \pi_1\text{ of the stalk}\\
  \hline
  \text{ordinary point} & \R & 0\\
  \text{dual-number point} & \R[\varepsilon]/(\varepsilon^2) & 0\\
  \text{derived self-intersection} & \R & \R.
\end{array}
\]
Topology distinguishes none of them. The ordinary and higher parts of the
structure sheaf distinguish all three.

\subsection{The affine comparison theorem}
\label[subsection]{sec:Affine_Comparison_Talk_Eight}

There is always a natural morphism
\[
  A\longrightarrow\Gamma(\Spec A,\mathcal O_{\Spec A}).
\]
For arbitrary animated $\Cinfty$-rings it need not be an equivalence. In the
finite-presentation range of derived manifolds, it is.

\begin{definition}[Affine derived $\Cinfty$-scheme]
  \label[definition]{def:Affine_Derived_Cinfty_Scheme_Talk_Eight}
  An \emph{affine derived $\Cinfty$-scheme of finite presentation} is a
  locally animated $\Cinfty$-ringed space equivalent to $\Spec A$ for a
  homotopically finitely presented animated $\Cinfty$-ring $A$.
\end{definition}

\begin{theorem}[Affine comparison in finite presentation]
  \label[theorem]{thm:Affine_Comparison_Finite_Presentation_Talk_Eight}
  Spectrum and global sections induce an equivalence of $\infty$-categories
  \[
    \operatorname{Alg}_{\Cinfty}(\An)_{\mathrm{fp}}\catop
    \simeq
    \mathrm{d}\Cinfty\mathrm{Aff}_{\mathrm{fp}}.
  \]
  In particular, for every homotopically finitely presented animated
  $\Cinfty$-ring $A$, the unit induces an isomorphism
  \[
    A
    \iso
    \Gamma(\Spec A,\mathcal O_{\Spec A}).
  \]
\end{theorem}

The proof passes through the class of fair animated $\Cinfty$-rings. Steffens
first proves spectrum and global sections to be inverse on fair objects and
then proves that almost finitely presented animated $\Cinfty$-rings are fair.
We use only the resulting finite-presentation statement. The relevant results
are Theorem~4.1.3.22, Proposition~4.1.3.33, and Corollary~4.1.3.34
\cite[Theorem~4.1.3.22, Proposition~4.1.3.33, and
Corollary~4.1.3.34]{Steffens_Derived_Differential_Geometry}.

The theorem does not introduce a second affine theory. It equips the affine
objects of Talks 4 to 7 with an explicit topology, open restrictions, and a
structure sheaf. Combining it with the algebraic presentation from Talk 4
gives
\[
  \mathrm{DMfd}
  \simeq
  \mathrm{d}\Cinfty\mathrm{Aff}_{\mathrm{fp}}.
\]

\section{Local derived smooth geometry}
\label[section]{sec:Local_Derived_Smooth_Geometry_Talk_Eight}

The affine comparison theorem gives a geometric meaning to a single animated
algebra. We now allow these affine objects to be restricted and glued.

\subsection{Open subobjects and locally affine schemes}
\label[subsection]{sec:Open_Subobjects_Locally_Affine_Talk_Eight}

Let $(X,\mathcal O_X)$ be a locally animated $\Cinfty$-ringed space and let
$V\subseteq|X|$ be open. Restriction gives another such space
\[
  (V,\mathcal O_X|_V)
\]
and a morphism to $X$. We call this an \emph{open subobject}.

On an affine derived spectrum, principal open restriction is precisely the
geometric form of localization:
\begin{equation}
  \label{eq:Principal_Open_Derived_Spectrum}
  \Spec A[a^{-1}]
  \simeq
  \bigl(D(a),\mathcal O_{\Spec A}|_{D(a)}\bigr).
\end{equation}
The underlying topological spaces agree by construction. On every smaller
principal open, the structure sheaves agree by the universal properties of
the corresponding iterated localizations.

\begin{definition}[Derived $\Cinfty$-scheme]
  \label[definition]{def:Locally_Affine_Derived_Cinfty_Scheme_Talk_Eight}
  A \emph{derived $\Cinfty$-scheme locally of finite presentation} is a
  locally animated $\Cinfty$-ringed space which admits an open cover by
  affine derived $\Cinfty$-schemes of finite presentation.
\end{definition}

This is a genuine enlargement of the affine category. The universal
$\infty$-category $\mathrm{DMfd}$ used so far is equivalent to the affine
finite-presentation category. A locally affine derived $\Cinfty$-scheme may
instead require several affine opens.

The sheaf condition is what permits these affine restrictions to be glued.
For the live route, it is enough to remember its familiar meaning: animated
functions defined locally and agreeing coherently on all intersections arise
from one function on the union. The Cech formulation and the distinction
between sheaf descent here and moduli descent in Talk 9 are recorded in the
supplement.

\subsection{Truncation and quasi-smoothness are local}
\label[subsection]{sec:Truncation_Quasi_Smoothness_Local_Talk_Eight}

The ordinary truncation of a derived $\Cinfty$-scheme is
\[
  t_0X
  =
  \bigl(|X|,\pi_0\mathcal O_X\bigr).
\]
On an affine open it restricts to
\[
  t_0(\Spec A)=\Spec(\pi_0A).
\]
Again the topological space is unchanged. Only the animated structure sheaf
has been truncated.

The cotangent complex restricts to opens by
\Cref{eq:Cotangent_Complex_Localizes_Talk_Eight}. This lets us globalize the
finiteness condition from Talk 7.

\begin{definition}[Quasi-smooth derived $\Cinfty$-scheme]
  \label[definition]{def:Quasi_Smooth_Derived_Cinfty_Scheme_Talk_Eight}
  A derived $\Cinfty$-scheme locally of finite presentation is
  \emph{quasi-smooth} if, on an affine open cover $\Spec A_i$, every
  cotangent complex $\mathbb L_{A_i}$ is perfect of Tor-amplitude contained
  in $[-1,0]$.
\end{definition}

\begin{proposition}[Locality of quasi-smoothness]
  \label[proposition]{prop:Locality_Quasi_Smoothness_Talk_Eight}
  The condition in
  \Cref{def:Quasi_Smooth_Derived_Cinfty_Scheme_Talk_Eight} is independent of
  the chosen affine open cover and may be checked after any affine open
  refinement.
\end{proposition}

\begin{proof}
  \Cref{eq:Cotangent_Complex_Localizes_Talk_Eight} shows that
  perfection and the stated Tor-amplitude persist under restriction.
  Conversely, perfection and Tor-amplitude are local properties of stable
  modules on an affine spectrum. If they hold on an open cover, the local
  duals and amplitude bounds agree on overlaps and give the corresponding
  global properties.
\end{proof}

We can now restate the normal-form theorem from Talk 7 without leaving the
word ``locally'' undefined.

\begin{theorem}[Local zero-locus form]
  \label[theorem]{thm:Local_Zero_Locus_Form_Talk_Eight}
  Let $X$ be a derived $\Cinfty$-scheme locally of finite presentation. Then
  $X$ is quasi-smooth if and only if every point of $t_0X$ has an open
  neighborhood equivalent to the derived zero locus of a section of a
  finite-rank vector bundle over a manifold.
\end{theorem}

One direction was calculated in Talk 7: the cotangent complex of a
finite-rank zero locus is perfect in degrees $-1$ and $0$. The converse is a
derived inverse-function and normal-form theorem. We quote it from Steffens,
Proposition~5.1.1.11 and Corollary~5.1.1.13
\cite[Proposition~5.1.1.11 and
Corollary~5.1.1.13]{Steffens_Derived_Differential_Geometry}.

The converse is a substantial input. The tangent-complex calculation alone
does not produce local coordinates or equations. It only verifies the
amplitude of an object already presented as a zero locus.

\subsection{One global zero locus and its local equations}
\label[subsection]{sec:Global_Zero_Locus_Local_Charts_Talk_Eight}

Let $q\colon E\to M$ be a finite-rank vector bundle and let
$s\colon M\to E$ be a section. The global derived zero locus
\[
  Z=Z^{\mathrm{der}}(s)
\]
was constructed before any trivialization of $E$ was chosen:
\[
\begin{tikzcd}
  Z \rar \dar \drar[pullback]
    & M \dar["s"] \\
  M \rar["0"'] & E.
\end{tikzcd}
\]

Choose a bundle-trivializing open cover $M=\bigcup_iU_i$ and identify
\[
  E|_{U_i}\cong U_i\times\R^r.
\]
In this frame the section is a smooth map
\[
  s_i=(s_{i,1},\ldots,s_{i,r})
  \colon
  U_i\longrightarrow\R^r.
\]
Let $Z_i\subseteq Z$ be the inverse image of $U_i$ along the projection
$Z\to M$.

\begin{proposition}[Local Koszul presentations]
  \label[proposition]{prop:Local_Koszul_Presentations_Talk_Eight}
  There are natural equivalences
  \[
    Z_i
    \simeq
    \Spec K(s_i),
  \]
  where
  \[
    K(s_i)
    =
    \bigl(
      \Cinfty(U_i)\otimes\Lambda(e_1,\ldots,e_r),
      \partial e_j=s_{i,j}
    \bigr).
  \]
\end{proposition}

\begin{proof}
  Open restriction commutes with the pullback defining $Z$, so
  \[
    Z_i
    \simeq
    U_i\times_{E|_{U_i}}U_i.
  \]
  After the chosen trivialization, this is the derived zero locus of
  $s_i\colon U_i\to\R^r$. The Koszul calculation from Talk 5 identifies its
  animated function algebra with $K(s_i)$.
\end{proof}

On an overlap $U_i\cap U_j$, the two local formulas are related by the frame
change of $E$. More intrinsically, both formulas present the restriction of
the same global pullback $Z$. Their comparison is therefore not additional
structure imposed on $Z$. It is induced by restriction from the one object
which was present before either frame was chosen.

The tangent complexes behave in the same way. Globally along the classical
zero set,
\[
  \mathbb T_Z
  \simeq
  [TM|_{Z(s)}\xrightarrow{Ds}E|_{Z(s)}]
\]
in cohomological degrees $0$ and $1$. In a frame on $U_i$, this restricts to
\[
  [TU_i|_{Z(s_i)}\xrightarrow{Ds_i}U_i\times\R^r].
\]
The change of frame acts simultaneously on the obstruction bundle and on the
derivative. The local complexes are presentations of one global complex.

This example joins the two halves of the chapter. The spectrum construction
explains how one global derived zero locus acquires local equations. We now
ask how much of one such system of equations belongs to the object itself.

\section{Kuranishi charts as presentations}
\label[section]{sec:Kuranishi_Charts_As_Presentations_Talk_Eight}

\subsection{Presentation, object, and chart}
\label[subsection]{sec:Presentation_Object_Chart_Talk_Eight}

We work throughout without isotropy groups, boundaries, or corners. This is
the finite-dimensional local situation needed for the later elliptic
applications.

\begin{definition}[Kuranishi presentation]
  \label[definition]{def:Kuranishi_Presentation_Talk_Eight}
  A \emph{Kuranishi presentation without isotropy} is a triple
  \[
    K=(U,E,s),
  \]
  where $U$ is a manifold, $E\to U$ is a finite-rank vector bundle, and
  $s\colon U\to E$ is a section. It presents the derived zero locus
  \[
    \mathbf Z^{\mathrm{der}}(K)
    :=
    Z^{\mathrm{der}}(s).
  \]
  Its ordinary zero locus is
  \[
    Z(K)=\{x\in U\mid s(x)=0_x\}.
  \]
\end{definition}

When a presentation is used as a chart for a fixed derived
$\Cinfty$-scheme $X$, the data also include an equivalence
\[
  X|_V
  \simeq
  \mathbf Z^{\mathrm{der}}(K)
\]
for some open $V\subseteq|X|$. On ordinary truncations this gives a
homeomorphism from $V$ to the zero set $Z(K)$, often called the footprint
map. In informal usage the triple itself is frequently called a Kuranishi
chart. The distinction between presentation and chart relative to $X$ is
useful whenever choices are being compared.

There are therefore three levels:

\begin{enumerate}
  \item The \emph{presentation} $(U,E,s)$ consists of auxiliary
    finite-dimensional data.
  \item The \emph{presented object} is the intrinsic derived zero locus
    $\mathbf Z^{\mathrm{der}}(K)$.
  \item A \emph{comparison} is a morphism of presentations which induces a
    morphism between the presented derived objects.
\end{enumerate}

At $x\in Z(K)$, Talk 7 gives
\begin{equation}
  \label{eq:Tangent_Complex_Kuranishi_Presentation_Talk_Eight}
  \mathbb T_x\mathbf Z^{\mathrm{der}}(K)
  \simeq
  [T_xU\xrightarrow{D_xs}E_x]
\end{equation}
in cohomological degrees $0$ and $1$. Thus
\[
  \operatorname{Def}_x(K)=\ker(D_xs),
  \qquad
  \operatorname{Obs}_x(K)=\operatorname{coker}(D_xs),
\]
and
\[
  \operatorname{vdim}K
  =
  \dim U-\operatorname{rank}E.
\]

The triple contains more information than the object it presents. It records
one ambient manifold, one obstruction bundle, and one system of equations.
The remainder of the chapter exhibits operations which change these choices.

\subsection{The tangent complex is not a complete invariant}
\label[subsection]{sec:Tangent_Complex_Not_Complete_Talk_Eight}

Before comparing presentations, we must rule out a tempting shortcut.
Consider
\[
  f(x)=x^2
  \qquad\text{and}\qquad
  g(x)=x^3.
\]
Both derivatives vanish at the origin, so their tangent complexes agree:
\[
  \mathbb T_0Z^{\mathrm{der}}(f)
  \simeq
  [\R\xrightarrow{0}\R]
  \simeq
  \mathbb T_0Z^{\mathrm{der}}(g).
\]
Their ordinary truncations are nevertheless different. Talk 5 gives
\[
  \pi_0\mathcal O\bigl(Z^{\mathrm{der}}(f)\bigr)
  \cong
  \frac{\Cinfty(\R)}{(x^2)},
  \qquad
  \pi_0\mathcal O\bigl(Z^{\mathrm{der}}(g)\bigr)
  \cong
  \frac{\Cinfty(\R)}{(x^3)}.
\]
Repeated use of Hadamard's lemma identifies the corresponding local rings at
the origin with
\[
  \R[\varepsilon]/(\varepsilon^2)
  \qquad\text{and}\qquad
  \R[\varepsilon]/(\varepsilon^3).
\]
They are not isomorphic.

\begin{warning}[Linear data do not classify the nonlinear object]
  \label[warning]{warn:Tangent_Complex_Not_Classifier_Talk_Eight}
  Isomorphic deformation and obstruction spaces, or even isomorphic abstract
  tangent complexes, do not imply that two derived zero loci are equivalent.
  A comparison of nonlinear presentations requires a morphism between them.
\end{warning}

This is analogous to the ordinary inverse-function theorem. The derivative
of a specified smooth map may prove that the map is locally invertible. An
accidental isomorphism between two tangent spaces does not produce a map of
manifolds. Later we will quote a derived version of this principle for a
specified morphism of Kuranishi presentations.

\subsection{Shrinking and strict coordinate changes}
\label[subsection]{sec:Elementary_Changes_Kuranishi_Talk_Eight}

Two changes of presentation are immediate.

First, let $W\subseteq U$ be open. Restriction gives
\[
  K|_W=(W,E|_W,s|_W).
\]
Its derived zero locus is the corresponding open restriction:
\begin{equation}
  \label{eq:Kuranishi_Shrinking_Open_Restriction}
  \mathbf Z^{\mathrm{der}}(K|_W)
  \simeq
  \mathbf Z^{\mathrm{der}}(K)|_{W\cap Z(s)}.
\end{equation}
If $W$ contains all of $Z(s)$, the ambient presentation has been shrunk
without changing the complete presented object. If $W$ contains only part of
$Z(s)$, the result presents a proper open subobject. The word ``shrinking''
is used for both operations, so the footprint should always be stated.

Second, suppose that $f_b\colon U\to V$ is a diffeomorphism and
$f_v\colon E\to F$ is a vector-bundle isomorphism over $f_b$ satisfying
\[
  f_v\circ s=t\circ f_b.
\]
Then $(f_b,f_v)$ is a strict coordinate and frame change. It carries the
pullback square defining $Z^{\mathrm{der}}(s)$ isomorphically to the pullback
square defining $Z^{\mathrm{der}}(t)$.

These operations change the visible equations only mildly. Stabilization is
the first operation which changes both the dimension of the ambient manifold
and the rank of the obstruction bundle while preserving the derived object.

\section{Stabilization and invariance}
\label[section]{sec:Stabilization_And_Invariance_Talk_Eight}

\subsection{Adding a variable and a transverse equation}
\label[subsection]{sec:Definition_Stabilization_Talk_Eight}

Let $K=(U,E,s)$ be a Kuranishi presentation and let $p\colon G\to U$ be a
finite-rank vector bundle. On its total space, the pullback bundle $p^*G$ has
a tautological section
\[
  \tau_G\colon\operatorname{Tot}(G)\longrightarrow p^*G,
  \qquad
  \tau_G(g)=g\in G_{p(g)}.
\]

\begin{definition}[Stabilization]
  \label[definition]{def:Stabilization_Kuranishi_Talk_Eight}
  The \emph{stabilization of $K$ by $G$} is the Kuranishi presentation
  \[
    K^G
    =
    \left(
      \operatorname{Tot}(G),
      p^*E\oplus p^*G,
      (p^*s,\tau_G)
    \right).
  \]
\end{definition}

The new equation $\tau_G=0$ forces the new ambient variable to lie on the
zero section $U\subseteq\operatorname{Tot}(G)$. Since this equation is
transverse, it should contribute no residual derived structure. The original
equation $s=0$ then remains.

Projection $p\colon\operatorname{Tot}(G)\to U$ together with projection
onto $p^*E$ defines a morphism of presentations
\begin{equation}
  \label{eq:Stabilization_Morphism_Talk_Eight}
  \sigma_G\colon K^G\longrightarrow K.
\end{equation}
Indeed, projection carries $(p^*s,\tau_G)$ to $s\circ p$.

\subsection{The stabilized square equation}
\label[subsection]{sec:Stabilized_Square_Equation_Talk_Eight}

Begin with
\[
  K
  =
  (\R,\R\times\R,x\mapsto x^2).
\]
Stabilization by the trivial line gives
\[
  K^{\R}
  =
  (\R^2,\R^2,(x,y)\mapsto(x^2,y)).
\]
We compare the two presentations at three levels.

At the level of ordinary zero loci, both equations have only the origin as a
solution. Projection $(x,y)\mapsto x$ identifies the two zero sets.

At the level of stable linearization, the original tangent complex is
\[
  \mathbb T_0\mathbf Z^{\mathrm{der}}(K)
  \simeq
  [\R\xrightarrow{0}\R].
\]
The stabilized tangent complex is
\[
  \mathbb T_0\mathbf Z^{\mathrm{der}}(K^{\R})
  \simeq
  \left[
    \R^2
    \xrightarrow{
      \left(
        \begin{smallmatrix}
          0&0\\
          0&1
        \end{smallmatrix}
      \right)}
    \R^2
  \right].
\]
It decomposes as
\begin{equation}
  \label{eq:Square_Stabilization_Tangent_Decomposition}
  \left[
    \R^2
    \xrightarrow{
      \left(
        \begin{smallmatrix}
          0&0\\
          0&1
        \end{smallmatrix}
      \right)}
    \R^2
  \right]
  \cong
  [\R\xrightarrow{0}\R]
  \oplus
  [\R\xrightarrow{1}\R].
\end{equation}
The second summand is zero in the stable module category. Projection induces
a quasi-isomorphism from the stabilized complex to the original complex.

The complete nonlinear calculation takes place on animated function
algebras. The two Koszul presentations are
\[
  \mathcal K(K)
  =
  \bigl(
    \Cinfty(\R)\otimes\Lambda(e),
    \partial e=x^2
  \bigr)
\]
and
\[
  \mathcal K(K^{\R})
  =
  \bigl(
    \Cinfty(\R^2)\otimes\Lambda(e,f),
    \partial e=x^2,
    \partial f=y
  \bigr).
\]
The pair $(y,f)$ is the Koszul resolution of the transverse equation $y=0$.
Eliminating it restricts smooth functions to the $x$-axis and induces an
equivalence
\begin{equation}
  \label{eq:Square_Stabilization_Koszul_Equivalence}
  \mathcal K(K^{\R})
  \simeq
  \mathcal K(K).
\end{equation}

We have checked stabilization at three levels:

\begin{enumerate}
  \item The ordinary zero loci are identified.
  \item The tangent complexes differ by an acyclic summand.
  \item The full animated function algebras are equivalent.
\end{enumerate}

The third statement is the intrinsic conclusion. The first two are its
ordinary and infinitesimal shadows. This is different from the comparison of
$x^2$ and $x^3$: there we had an accidental agreement of tangent complexes
but no equivalence of the nonlinear algebras.

\subsection{Why stabilization preserves the derived object}
\label[subsection]{sec:General_Stabilization_Proof_Talk_Eight}

The preceding calculation globalizes without choosing a frame for $G$.

\begin{theorem}[Invariance under stabilization]
  \label[theorem]{thm:Stabilization_Preserves_Derived_Object_Talk_Eight}
  For every Kuranishi presentation $K=(U,E,s)$ and every finite-rank vector
  bundle $G\to U$, projection induces an equivalence
  \[
    \mathbf Z^{\mathrm{der}}(K^G)
    \simeq
    \mathbf Z^{\mathrm{der}}(K).
  \]
\end{theorem}

\begin{proof}
  The tautological section $\tau_G$ is transverse to the zero section of
  $p^*G$. Its ordinary zero locus is the zero section
  \[
    U\longrightarrow\operatorname{Tot}(G).
  \]
  By transverse exactness, its derived zero locus is this ordinary manifold
  $U$, with no additional derived structure.

  The derived zero locus of the direct-sum section
  $(p^*s,\tau_G)$ can be formed in two stages. First impose
  $\tau_G=0$, and then impose the restriction of $p^*s=0$. Finite pullbacks
  may be pasted in either order, so this iterated zero locus is equivalent to
  the zero locus of the pair. After the first stage the ambient manifold is
  the zero section $U$, the bundle $p^*E$ restricts to $E$, and the section
  $p^*s$ restricts to $s$. Hence the second stage is precisely
  $Z^{\mathrm{der}}(s)$.
\end{proof}

The tangent calculation is compatible with this equivalence. Along the zero
section there is a canonical splitting
\[
  T_{0_x}\operatorname{Tot}(G)
  \cong
  T_xU\oplus G_x.
\]
The derivative of the stabilized section is block diagonal:
\[
  \left[
    T_xU\oplus G_x
    \xrightarrow{
      \left(
        \begin{smallmatrix}
          D_xs&0\\
          0&1
        \end{smallmatrix}
      \right)}
    E_x\oplus G_x
  \right].
\]
Consequently,
\begin{equation}
  \label{eq:General_Stabilization_Tangent_Decomposition}
  \mathbb T_x\mathbf Z^{\mathrm{der}}(K^G)
  \simeq
  \mathbb T_x\mathbf Z^{\mathrm{der}}(K)
  \oplus
  [G_x\xrightarrow{1}G_x].
\end{equation}
The virtual dimension is unchanged:
\[
  \dim\operatorname{Tot}(G)
  -\operatorname{rank}(p^*E\oplus p^*G)
  =
  \dim U-\operatorname{rank}E.
\]

\subsection{Recognizing safe changes of presentation}
\label[subsection]{sec:Recognizing_Safe_Changes_Talk_Eight}

Stabilization was proved directly. There is also a general practical
criterion for a specified comparison map.

\begin{definition}[Morphisms of Kuranishi presentations]
  \label[definition]{def:Morphism_Kuranishi_Presentation_Talk_Eight}
  Let $K=(U,E,s)$ and $L=(V,F,t)$. A \emph{morphism of presentations}
  \[
    f=(f_b,f_v)\colon K\longrightarrow L
  \]
  consists of a smooth map $f_b\colon U\to V$ and a fiberwise-linear map
  $f_v\colon E\to F$ over $f_b$ such that
  \[
    f_v\circ s=t\circ f_b.
  \]
\end{definition}

The compatibility carries $Z(s)$ into $Z(t)$ and induces a morphism
\[
  \mathbf Z^{\mathrm{der}}(f)
  \colon
  \mathbf Z^{\mathrm{der}}(K)
  \longrightarrow
  \mathbf Z^{\mathrm{der}}(L).
\]
At a zero $x$, differentiation gives a map of tangent complexes
\begin{equation}
  \label{eq:Tangent_Map_Kuranishi_Morphism_Talk_Eight}
  \begin{tikzcd}
    T_xU \rar["D_xs"] \dar["D_xf_b"']
      & E_x \dar["(f_v)_x"] \\
    T_{f_b(x)}V \rar["D_{f_b(x)}t"']
      & F_{f_b(x)}.
  \end{tikzcd}
\end{equation}

\begin{definition}[Weak equivalence of presentations]
  \label[definition]{def:Weak_Equivalence_Kuranishi_Talk_Eight}
  A morphism $f\colon K\to L$ is a \emph{weak equivalence} if it induces a
  bijection $Z(K)\to Z(L)$ and the map in
  \Cref{eq:Tangent_Map_Kuranishi_Morphism_Talk_Eight} is a
  quasi-isomorphism at every zero.
\end{definition}

\begin{theorem}[Presentation-invariance criterion]
  \label[theorem]{thm:Weak_Equivalence_Presents_Equivalence_Talk_Eight}
  A morphism of finite-dimensional Kuranishi presentations is a weak
  equivalence if and only if its induced morphism of derived zero loci is an
  equivalence.
\end{theorem}

The direction from derived equivalence to weak equivalence follows by taking
ordinary truncations and tangent complexes. The converse is a derived
inverse-function theorem. Pointwise quasi-isomorphism makes the relative
cotangent complex vanish at every zero. Finite presentation and derived
Nakayama extend this vanishing to open neighborhoods, where the morphism is
etale. The bijection of zero loci then globalizes these local equivalences.
A fuller source-based proof is recorded in
\Cref{sec:Supplement_Weak_Equivalence_Kuranishi}.

The theorem must not be used without its input map. The comparison of $x^2$
and $x^3$ supplies isomorphic abstract tangent complexes but no morphism whose
linearization is being tested. The theorem recognizes equivalences; it does
not construct nonlinear comparisons from linear data.

\section{Summary and supplementary materials}
\label[section]{sec:Summary_And_Supplementary_Materials_Talk_Eight}

\subsection{Takeaways and the coherence problem}
\label[subsection]{sec:Takeaways_Coherence_Problem_Talk_Eight}

Let us collect the route of the chapter.
\begin{enumerate}
  \item Smooth localization is open restriction written on functions.
  \item The topology of $\Spec A$ is controlled by $\pi_0A$, while its
    structure sheaf retains the full animated $\Cinfty$-ring $A$.
  \item In finite presentation, affine derived spectra and animated
    $\Cinfty$-rings contain the same information.
  \item Quasi-smoothness is local because cotangent complexes restrict under
    smooth localization.
  \item Quasi-smooth derived $\Cinfty$-schemes are locally derived zero loci.
  \item A Kuranishi presentation is auxiliary data presenting such a local
    derived object.
  \item Stabilization changes the presentation by an acyclic variable-equation
    pair and preserves the presented derived object.
  \item Tangent complexes test a specified comparison map, but do not classify
    nonlinear derived objects by themselves.
\end{enumerate}

There is one important problem which this chapter has deliberately not
solved. Suppose that a moduli problem is covered by local Kuranishi
presentations $K_i$. On a double overlap one wants a comparison between
$K_i$ and $K_j$. On a triple overlap, the two composites of comparisons need
not agree strictly, so one needs a homotopy between them. On a quadruple
overlap, those homotopies must themselves be compatible. The pattern
continues.

Pairwise weak equivalences therefore do not constitute a global Kuranishi
object. The next chapter replaces a list of chosen charts by an intrinsic
moduli functor and packages all overlap data in an $\infty$-categorical
descent condition. Its main representability input will be the following
result.

\begin{theorem}[Open local representability, preview]
  \label[theorem]{thm:Open_Local_Representability_Preview}
  Let $\mathcal X$ be a derived $\Cinfty$-stack with an effective open atlas
  by affine derived $\Cinfty$-schemes belonging to a class stable under open
  restriction. Then $\mathcal X$ is represented by a derived
  $\Cinfty$-scheme locally belonging to that class.
\end{theorem}

The definitions and proof belong to Talk~9. For now, the point is only the
shape of the answer: local presentations glue when their comparison data are
the restrictions of one object satisfying descent. This is why the local
geometry developed here is an input to, rather than a substitute for, the
coherence theorem.

\subsection{Supplement: the real spectrum and its structure sheaf}
\label[subsection]{sec:Supplement_Real_Spectrum_Talk_Eight}

The live talk used the derived spectrum at slogan level. We now record the
ordinary construction more carefully. It explains where the topology and
local rings come from.

\begin{definition}[Real spectrum]
  \label[definition]{def:Real_Spectrum_Talk_Eight}
  Let $A$ be a $\Cinfty$-ring. Its \emph{real spectrum} is the set
  \[
    \Spec_{\R}A
    =\Hom_{\Cinfty\text{-}\mathrm{Ring}}(A,\R).
  \]
  For $a\in A$, set
  \[
    D(a)=\{x\in\Spec_{\R}A\mid x(a)\neq0\}.
  \]
  The subsets $D(a)$ form a basis for a topology on $\Spec_{\R}A$.
\end{definition}

Indeed,
\[
  D(1)=\Spec_{\R}A,
  \qquad
  D(a)\cap D(b)=D(ab).
\]
Every $a\in A$ defines an evaluation function
\[
  a^*\colon\Spec_{\R}A\longrightarrow\R,
  \qquad
  a^*(x)=x(a).
\]
The spectrum topology is the weakest topology for which all evaluation
functions are continuous. To see the nontrivial direction, choose, for an
open $V\subseteq\R$, a smooth function $h$ with
$V=\{t\mid h(t)\neq0\}$. Then
\[
  (a^*)^{-1}(V)=D(\Phi_h(a)).
\]
It is also Hausdorff: two distinct homomorphisms differ on some $a\in A$, and
disjoint intervals around the two values pull back to disjoint
neighborhoods.

If $A=\Cinfty(\R^n)/I$, evaluation identifies the real spectrum with
\[
  Z(I)=\{p\in\R^n\mid f(p)=0\text{ for all }f\in I\}
\]
with its ordinary subspace topology. Thus the real spectrum recovers the
expected point set, but it cannot see nilpotent or higher derived structure.

For $x\in\Spec_{\R}A$, define the local $\Cinfty$-ring
\[
  A_x=A[(A\setminus x^{-1}(0))^{-1}].
\]
Its elements are germs of elements of $A$ near $x$. The structure sheaf may
be described without making a choice of basis presentation.

\begin{construction}[Structure sheaf]
  \label[construction]{con:Structure_Sheaf_Talk_Eight}
  For $U\subseteq\Spec_{\R}A$ open, let $\mathcal O_A(U)$ be the set of
  families
  \[
    \sigma(x)\in A_x,
    \qquad x\in U,
  \]
  which are locally represented by elements of $A$. Pointwise smooth
  operations and restriction make $\mathcal O_A$ a sheaf of
  $\Cinfty$-rings. Define
  \[
    \Spec A=(\Spec_{\R}A,\mathcal O_A).
  \]
\end{construction}

By construction, $\mathcal O_{A,x}\cong A_x$. Moreover, localization induces
an isomorphism of local $\Cinfty$-ringed spaces
\begin{equation}
  \label{eq:Principal_Open_Ordinary_Spectrum_Talk_Eight}
  \Spec A[a^{-1}]
  \iso
  (D(a),\mathcal O_A|_{D(a)}).
\end{equation}
For a finitely presented $\Cinfty$-ring, the basis assignment
$D(a)\mapsto A[a^{-1}]$ already gives the same sheaf after sheafification.
For arbitrary rings, the local-germ construction avoids subtleties caused by
different localization maps presenting the same real open subset.

\begin{theorem}[Reconstruction of manifolds]
  \label[theorem]{thm:Reconstruction_Manifolds_Talk_Eight}
  For a manifold $M$, evaluation induces an isomorphism of local
  $\Cinfty$-ringed spaces
  \[
    \Spec\Cinfty(M)\iso(M,\Cinfty_M).
  \]
\end{theorem}

\begin{proof}
  Every homomorphism $\Cinfty(M)\to\R$ is evaluation at a unique point of
  $M$. The basic open $D(a)$ is the nonvanishing locus of the smooth function
  $a$. Conversely, if $x\in U\subseteq M$ with $U$ open, a bump function
  supported in $U$ and nonzero at $x$ produces a basic neighborhood
  $x\in D(a)\subseteq U$. Hence the spectrum topology is the manifold
  topology.

  Smooth localization identifies the stalk at $x$ with the ring
  $\Cinfty_{M,x}$ of smooth germs. Finally, a section of $\mathcal O_A$ is
  locally represented by global smooth functions and is therefore smooth.
  Conversely, cutoff functions show that every smooth function on an open
  subset is locally the restriction of a global smooth function.
\end{proof}

For an animated $\Cinfty$-ring, the same construction is applied to
$\pi_0A$ to obtain the point set and topology, while the localizations
$A[a^{-1}]$ provide the animated structure sheaf. The homotopy groups of
that sheaf are the sheafifications of
\[
  D(a)\longmapsto
  \pi_n(A)\otimes_{\pi_0A}(\pi_0A)[a^{-1}].
\]
This makes precise the slogan that the topology sees $\pi_0A$, whereas the
structure sheaf sees all of $A$.

\subsection{Supplement: the Möbius self-intersection}
\label[subsection]{sec:Supplement_Mobius_Self_Intersection_Talk_Eight}

The zero section of a nontrivial bundle shows why local equations need not
admit a global choice of generators. Let $L\to S^1$ be the Möbius line bundle
and form the derived self-intersection of its zero section:
\[
  X=S^1\times_LS^1.
\]
Its classical truncation is $S^1$.

Choose an open cover $S^1=U_0\cup U_1$ on which $L$ is trivial. On $U_i$,
the zero section is represented by the zero function, and hence
\[
  \mathcal O_X|_{U_i}
  \simeq
  \bigl(\Cinfty_{U_i}\otimes\Lambda(e_i),\partial e_i=0\bigr).
\]
Thus $\pi_0\mathcal O_X|_{U_i}\cong\Cinfty_{U_i}$ and
$\pi_1\mathcal O_X|_{U_i}\cong\Cinfty_{U_i}e_i$.

Let $g_{01}\colon U_0\cap U_1\to\R^\times$ be the transition function of
$L$. The degree-one Koszul generator transforms as a local frame of
$L^\vee$:
\[
  e_1=g_{01}^{-1}e_0.
\]
Consequently,
\[
  \pi_1\mathcal O_X
  \cong
  \bigl(U\longmapsto\Gamma(U,L^\vee)\bigr).
\]
There is no global nowhere-vanishing degree-one generator. Dually, the
tangent complex is
\[
  \mathbb T_X\simeq[TS^1\xrightarrow{0}L],
\]
so its obstruction sheaf is the Möbius bundle itself. The derived object is
global, but every elementary equation displaying it is local and depends on
a frame.

\subsection{Supplement: weak equivalences of presentations}
\label[subsection]{sec:Supplement_Weak_Equivalence_Kuranishi}

We explain the source of
\Cref{thm:Weak_Equivalence_Presents_Equivalence_Talk_Eight}. Let
$f\colon K\to L$ be a morphism of the finite-dimensional presentations
considered in this chapter, and write
\[
  F=\mathbf Z^{\mathrm{der}}(f)
  \colon X\longrightarrow Y
\]
for the induced morphism of finite-presentation affine derived
$\Cinfty$-schemes.

If $F$ is an equivalence, its truncation is an isomorphism and its tangent
map is a quasi-isomorphism at every point. Thus $f$ is a weak equivalence.
For the converse, suppose that $f$ induces a bijection on zero loci and a
quasi-isomorphism on tangent complexes. The fiber of the relative cotangent
complex $\mathbb L_{X/Y}$ at a point $x$ is dual to the cofiber of the map of
tangent complexes at $x$. It therefore vanishes. Since $\mathbb L_{X/Y}$ is
perfect in the present finite-presentation setting, derived Nakayama implies
that it vanishes on a neighborhood of every point. Hence $F$ is etale near
every point.

The derived inverse-function theorem then identifies a neighborhood of $x$
with a neighborhood of $F(x)$. Because the map on underlying zero loci is a
bijection, these local equivalences cover both source and target and assemble
to an equivalence $X\simeq Y$. The finite-presentation and geometricity
hypotheses are essential here. The statement is not a principle saying that
arbitrary derived objects with isomorphic tangent complexes are equivalent.

There is also an algebraic way to recognize what is happening. After
trivializing the bundles, the presentations produce Koszul algebras. A
morphism of presentations induces a morphism between these algebras. Weak
equivalence says that the induced map has the same classical points and is a
quasi-isomorphism on the linearized complexes at those points. The argument
above promotes that pointwise linear information, through the relative
cotangent complex and finite presentation, to a local equivalence of the
full nonlinear Koszul algebras.

\subsection{Supplement: terminology and truncations}
\label[subsection]{sec:Supplement_Terminology_Truncations_Talk_Eight}

The terminology surrounding Kuranishi geometry is not uniform. In this
chapter a \emph{Kuranishi presentation} means only a finite-dimensional
triple $(U,E,s)$ without isotropy, boundary, or corners. A \emph{Kuranishi
chart for $X$} additionally includes an identification of its derived zero
locus with an open subobject of $X$. This usage isolates exactly the local
calculation needed later.

Other theories package different amounts of information. Joyce's
d-manifolds and d-orbifolds are 2-categorical geometric structures, while
Kuranishi spaces are built from charts and coordinate changes with their own
specified coherence. Fully derived manifolds retain higher mapping animae
which cannot in general be recovered from a 2-categorical truncation. These
frameworks are closely related, but they should not be identified without a
comparison theorem and its hypotheses. In particular, the local
stabilization calculation in this chapter does not by itself compare their
global notions of atlas.

\subsection{Related literature}
\label[subsection]{sec:Related_Literature_Talk_Eight}

The derived localization theorem, affine comparison, local zero-locus normal
form, and weak-equivalence criterion used in the chapter come from Steffens
\cite[Proposition~4.1.3.13, Corollary~4.1.3.34, Corollary~5.1.1.13, and
Remarks~1.0.0.4 and~1.0.0.6]{Steffens_Derived_Differential_Geometry}.
Joyce gives detailed accounts of ordinary $\Cinfty$-spectra, affine
$\Cinfty$-schemes, and the reconstruction of manifolds
\cite[Sections~2.2 and~3.1]{Joyce_Introduction_Cinfty_Schemes}
\cite[Sections~4.4 and~4.5]{Joyce_Algebraic_Geometry_Cinfty_Rings}.
The universal derived-manifold viewpoint and its relation to local
zero-locus constructions are developed by Carchedi and Steffens
\cite{Carchedi_Steffens_Universal_Derived_Manifolds}. The representability
theorem previewed at the end of the live route is
\cite[Proposition~2.1.48]{Steffens_Representability_Elliptic_Moduli}.

% !TEX root = ../main.tex

\chapter{Stacks and local representability}
\label[chapter]{chap:Coherence_Descent_Local_Representability}

A moduli problem assigns families to parameter objects. For example, a
family of zeros of $f\colon U\to\R^r$ parametrized by $T$ is a map
$T\to U$ satisfying the equation. In derived geometry the equation
includes a specified homotopy to zero. Pulling the family back along
$T'\to T$ gives a family over $T'$.

This viewpoint defines a moduli problem before a geometric representing
object has been constructed. Descent expresses that families and their
isomorphisms can be glued. The main result of this chapter gives a way
to recognize when the resulting stack is a derived scheme: construct
an open cover by representable stacks.

This returns to the coherence problem posed in
\Cref{sec:Local_Zero_Loci_Not_Enough}. Starting with a globally defined
stack supplies the overlaps and their higher compatibilities; the work
is to construct open charts of that functor. We develop descent and
atlases with the vector bundle zero locus as a running example, then
formulate representability relative to a parameter stack. The distinction
between gluing families and gluing charts will be essential throughout.

\section{The open topology and descent}

Descent concerns a cover of a parameter object: it reconstructs an
anima of global families from coherent local families. We specify the
test objects and covers first, then express this reconstruction as a limit.

We use the category $\mathrm{DMfd}$ of derived manifolds as test objects.
By the affine comparison, these are spectra of finitely presented animated
smooth rings. An open embedding is an inclusion of an open subobject
with its restricted structure sheaf. An \emph{open covering family}
$\{T_i\to T\}$ is a family of open embeddings whose underlying images
cover $|T|$. Smooth localization describes these embeddings algebraically.
They are preserved by pullback.

\begin{definition}[Prestack]
\label[definition]{def:Prestack_Derived_Cinfty}
A derived smooth prestack is a functor
$\mathcal F\colon\mathrm{DMfd}^{\mathrm{op}}\to\An$.
\end{definition}

The anima $\mathcal F(T)$ contains the $T$-families, their isomorphisms,
and higher identifications between those isomorphisms.
Replacing it by its set of components loses automorphisms as well as
higher coherence. Pullback along a map of parameter objects is part
of the functor.

For an open cover, write $T_{i_0\cdots i_n}$ for the iterated intersection.
Restriction gives a cosimplicial diagram with degree-$n$ term
\[
 C^n=\prod_{i_0,\ldots,i_n}\mathcal F(T_{i_0\cdots i_n}).
\]
Its face maps restrict to smaller intersections, and its degeneracy maps
repeat indices. The following formula for its limit avoids any
need for an infinite disjoint union to be a test object.

\begin{definition}[Stack]
\label[definition]{def:Stack_Derived_Cinfty}
A prestack is a stack if, for every open covering family, the natural map
\begin{equation}\label{eq:Stack_Descent_Parameter_Cech}
 \mathcal F(T)\longrightarrow
 \lim_{[n]\in\Delta}
 \prod_{i_0,\ldots,i_n}\mathcal F(T_{i_0\cdots i_n})
\end{equation}
is an isomorphism of animae.
\end{definition}

For a set-valued functor this reduces to the usual sheaf condition:
local elements must agree on double overlaps.
For groupoid-valued functors one instead chooses isomorphisms on double
overlaps and imposes the cocycle condition on triple overlaps.
For general animae the cocycle condition itself is supplied by a homotopy,
followed by higher compatibility data. The limit in
\Cref{eq:Stack_Descent_Parameter_Cech} packages all of it.

Descent applies to maps between families too. It says that the anima of
global families is the anima of coherent local data, rather than merely
that some global family exists.

The resulting $\infty$-category of stacks is denoted
$\mathrm{dSt}_{\Cinfty}^{\mathrm{fp}}$ here.
The superscript records our site of finitely presented test objects.
Steffens embeds this category fully faithfully into his category using
all affine derived smooth schemes, and calls its objects stacks locally
of finite presentation
\cite[Proposition~2.1.40 and Definition~2.1.41]
{Steffens_Representability_Elliptic_Moduli}.
Our constructions take place in this subcategory.

\section{Represented stacks and the zero-locus example}

Representability asks whether all families arise as maps to a geometric
object. The universal property of a derived pullback answers this
question for the zero-locus functor and gives a model for the later
elliptic moduli problem.

For a derived scheme $X$ locally of finite presentation, its functor
of points is
\[
 h_X(T)=\Hom(T,X).
\]
It satisfies descent because morphisms of locally ringed objects can
be glued over an open cover of their source.
This is the subcanonical property of the open topology.
Moreover, the functor $X\mapsto h_X$ is fully faithful.

\begin{definition}[Representability]
\label[definition]{def:Represented_Stack}
A stack $\mathcal F$ is represented by a derived $\Cinfty$-scheme $X$
if there is a natural isomorphism
$\mathcal F(T)\cong\Hom(T,X)$ for every test object $T$.
\end{definition}

The representing object is unique. If $X$ is not affine, the notation
$\mathcal F(X)$ means the value obtained by descent from an affine
open cover. Under representability, the identity of $X$ corresponds
to a universal family in $\mathcal F(X)$. Every other family is its
pullback, by the Yoneda lemma.

Let $f\colon U\to\R^r$ be smooth. Define
\begin{equation}\label{eq:Derived_Zero_Locus_Prestack}
 \mathcal Z_f(T)=\Hom(T,U)\times_{\Hom(T,\R^r)}\Hom(T,*),
\end{equation}
with pullback in animae.
A point is a map $u\colon T\to U$ together with a homotopy from $f\circ u$
to zero. Identifications in this pullback give the isomorphisms between
such data and their higher compatibilities. Limits of stacks are computed on each test object, so
$\mathcal Z_f$ is a stack.

By the universal property of a pullback,
\[
 \mathcal Z_f(T)\cong\Hom(T,U\times_{\R^r}^{\mathrm{der}}*).
\]
Thus it is represented by the derived zero locus computed in \Cref{chap:Derived_Zero_Loci_And_Tangent_Complexes}.
This is a simple instance of the representability question that will
later arise for differential equations.

For a section $s\colon U\to E$ of a vector bundle, the analogous formula
is most conveniently written in the category over $U$:
\[
 \mathcal Z_s=U\times_E^{\mathrm{der}}U,
\]
using $s$ and the zero section. Restricting to a trivializing open
$U_i$ gives the Koszul presentation for its list of local components.
These local representatives will form an open atlas of $\mathcal Z_s$.

\section{Smooth stacks as derived stacks}

A smooth stack is a sheaf of animae on ordinary manifolds with their
open-cover topology. The inclusion of manifolds into derived manifolds
induces a fully faithful functor
\[
 j_!\colon\mathrm{Shv}(\mathrm{Mfd})
 \longrightarrow\mathrm{dSt}_{\Cinfty}^{\mathrm{fp}}.
\]
It is obtained by left Kan extension and sheafification, and sends
an ordinary manifold to its derived functor of points.
We can therefore regard smooth stacks as particular derived stacks.

The derived values of $j_!\mathcal F$ are not obtained just by
evaluating $\mathcal F$ on an underlying point set.
They are determined by this extension construction.
Conversely, restriction to ordinary test manifolds does not recover
every derived stack. For example, a multiple point and an ordinary
point have the same families over ordinary manifolds, while maps
from infinitesimal derived tests distinguish them.

There is also an essential distinction between the two categories of
pullbacks. In smooth stacks,
\[
 *\times_\R *=*
\]
for the two origins. After passing to derived stacks, the pullback
is the derived self-intersection, with algebra $\R[\varepsilon]$,
where $|\varepsilon|=1$ and $d\varepsilon=0$.
Thus $j_!$ does not preserve arbitrary pullbacks.
The pullbacks needed in the analytic argument will be compared using
submersiveness in \Cref{chap:Elliptic_Representability_Smooth_Derived_Bridge}.

\section{Open atlases and local representability}

We now turn from gluing families over a parameter cover to presenting
the moduli stack itself by charts. Here a \v{C}ech diagram reconstructs
the stack as a colimit. Effectiveness expresses that the charts cover
the functor; openness makes it possible to glue their representing
objects into a derived scheme.

Let $p\colon\mathcal U\to\mathcal F$ be a morphism of stacks.
Its \v{C}ech nerve is the simplicial stack
\[
 \mathcal U_n=
 \underbrace{\mathcal U\times_{\mathcal F}\cdots
 \times_{\mathcal F}\mathcal U}_{n+1\text{ factors}}.
\]
The augmentation maps its geometric realization to $\mathcal F$.
Here geometric realization means a colimit in stacks, computed by
taking the presheaf colimit and then sheafifying.

\begin{definition}[Effective epimorphism]
\label[definition]{def:Effective_Epimorphism}
The map $p$ is an effective epimorphism if
$|\mathcal U_\bullet|\to\mathcal F$ is an isomorphism.
\end{definition}

On the open site there is a useful test: every family in $\mathcal F(T)$
must, after an open cover of $T$, lift to $\mathcal U$, with an
identification of its image with the given family.
Equivalently, the induced map of sheaves of components is surjective.
This characterizes effective epimorphisms in the $\infty$-topos of stacks.

Starting with a globally defined stack supplies the overlap data.
If $U_i\to\mathcal F$ and $U_j\to\mathcal F$ are two charts, their
overlap is $U_i\times_{\mathcal F}U_j$. Its associativity and higher
compatibilities come from pullbacks in the ambient $\infty$-category.
The following condition makes these overlaps geometric open subsets.

\begin{definition}[Open atlas]
\label[definition]{def:Open_Atlas}
A morphism $\mathcal V\to\mathcal F$ is an open inclusion if every
pullback along $T\to\mathcal F$ is represented by an open subobject of $T$.
An open atlas is a family of open inclusions $h_{U_i}\to\mathcal F$
from derived schemes such that
$\coprod_i h_{U_i}\to\mathcal F$ is an effective epimorphism.
\end{definition}

The coproduct here is in stacks. Each map is a condition on families:
after choosing a family over $T$, belonging to the indicated chart
cuts out an open subset of $T$.

The word \emph{open} is crucial for representing a stack by a scheme.
For example, a nontrivial finite group $G$ has a classifying stack $BG$.
The map $*\to BG$ is an effective epimorphism, expressing local
triviality of principal $G$-bundles. It is not an open inclusion:
its pullback along itself is $G\to*$.
The automorphisms of the trivial bundle also show that $BG$ cannot be
represented by a derived scheme, whose real points form a set.
More general notions of atlas lead to orbifolds and other geometric
stacks, but an arbitrary effective epimorphism is insufficient for
scheme representability.

An open atlas contains exactly the data needed to reconstruct a
derived scheme from its local presentations.

\begin{theorem}[Open local representability]
\label[theorem]{thm:Open_Local_Representability}
A derived smooth stack is represented by a derived $\Cinfty$-scheme
locally of finite presentation if and only if it has an open atlas
by affine derived schemes of finite presentation.
One may equivalently allow derived schemes locally of finite presentation
as the atlas objects.
\end{theorem}

This is \citet[Proposition~2.1.48]{Steffens_Representability_Elliptic_Moduli}.

\begin{proof}[Gluing argument]
A represented stack has such an atlas by the definition of a derived
scheme. Conversely, set $U_{ij}=U_i\times_{\mathcal F}U_j$.
Since each chart map is open, $U_{ij}$ is represented by an open
subscheme of both $U_i$ and $U_j$.
The two maps identify these opens. On triple and higher overlaps
the identifications have the compatibility supplied by the \v{C}ech nerve.

Derived schemes glue along this open descent datum, producing a scheme
$X$. On the underlying topological level this is the usual gluing of
open subsets; the animated structure sheaves glue by descent.
The functor of points of $X$ realizes the same \v{C}ech diagram.
Since the atlas is an effective epimorphism,
$h_X\cong|\mathcal U_\bullet|\cong\mathcal F$.
Allowing non-affine charts gives the same criterion by refining each
one with an affine open cover.
\end{proof}

In the bundle-section example, the charts are the zero loci over a
trivializing cover of $U$. Their overlaps are the zero loci over the
corresponding intersections. The transition maps of the bundle identify
the Koszul models, and the theorem recovers the global derived zero locus.

The same argument preserves a local property such as quasi-smoothness.
If the charts are quasi-smooth, the cotangent amplitude condition
holds on the resulting scheme.

\section{Relative representability and base change}

Moduli problems often have a parameter stack $S$.
For example, the metric, the domain manifold, or a coefficient in an
equation may vary. The correct representability statement is then
relative to $S$.

\begin{definition}
\label[definition]{def:Relative_Representability}
A map $\mathcal F\to S$ is representable by derived $\Cinfty$-schemes
locally of finite presentation if, for each affine derived test
$T\to S$, the pullback $\mathcal F\times_S T$ is represented by such
a scheme. It is relatively quasi-smooth if these representing
morphisms to $T$ are quasi-smooth.
\end{definition}

This asserts representability of every pulled-back family.
It does not assert that the total stack $\mathcal F$ is itself a scheme
when $S$ has automorphisms.

\begin{corollary}[Locality on the base]
\label[corollary]{cor:Local_Representability_Relative}
Representability by derived schemes locally of finite presentation
can be checked after an effective epimorphic covering family of the base.
The same holds with the additional requirement of relative
quasi-smoothness.
\end{corollary}

\begin{proof}[Argument]
Pull the covering family back along a test object $T\to S$.
The local lifting property of an effective epimorphism gives an open
cover $\{T_j\to T\}$ on which these maps lift to members of the cover.
By assumption, the restrictions of $\mathcal F\times_S T$ to the
$T_j$ are represented. They form an open atlas of
$\mathcal F\times_S T$, so the theorem glues them.
Relative cotangent complexes commute with base change, and their
perfection and amplitude are local, giving the last assertion.
\end{proof}

The source formulation is
\cite[Definition~2.1.42 and Corollary~2.1.49]
{Steffens_Representability_Elliptic_Moduli}.
This permits a proof over a smooth stack to be carried out on manifold
charts of the base, provided the moduli construction is compatible
with pullback.

\section{Further reading}

The local representability criterion does not make the construction of
a chart automatic.
One must identify the entire derived functor on an open neighbourhood,
including its values on derived tests. An identification of real
solution sets, or even of their tangent complexes, proves less.
The later Kuranishi argument is designed to establish precisely
this functorial identification.

\begin{exercise}
Let $f\colon U\to\R^r$ and let $\{U_i\}$ cover $U$.
Use the functor-of-points formula to verify that the corresponding
zero loci form an effective open atlas of $\mathcal Z_f$.
\end{exercise}

\begin{exercise}
For a nontrivial finite group $G$, calculate $*\times_{BG}*$ and
explain both the effectiveness of $*\to BG$ and its failure to be
an open inclusion.
\end{exercise}

All the descent and representability results used in this chapter are
contained in \cite[Section~2.1]{Steffens_Representability_Elliptic_Moduli}.

% !TEX root = ../main.tex

\chapter{Section stacks and derived solution stacks}
\label[chapter]{chap:Section_Stacks_Jets_Solution_Stacks}

A differential operator assigns to a function, section, or map an expression
in its derivatives. Its solutions often describe a geometric condition:
harmonicity, stationarity of an energy, or compatibility with a complex
structure. The equation can be defined for smooth families of unknowns.
The resulting functor of solutions is a stack, whose geometric
representability is a separate question.

Our running example is the nonlinear elliptic equation
\[
 P(u)=\Delta u+u^2=0
\]
on a closed connected Riemannian manifold. Its ordinary solution set is
a single point, but its infinitesimal solutions and obstructions are
both nontrivial. The eventual conclusion is concrete: near zero, its
derived solution object is the multiple point defined by $c^2=0$.
\Cref{chap:Elliptic_Representability_Derived_Charts} will prove this by
eliminating the mean-zero part of $u$, leaving its constant part $c$
as the remaining variable.

This chapter sets up the functor that the later analytic argument will
represent. Geometric examples lead to section stacks and their derived
fibres. Jets then express the finite-order condition, and linearization
gives the definition of ellipticity. Each construction is formulated
for families, so it can be used over a parameter stack.

\section{Differential equations from geometry}

The examples below identify the unknowns, their equations, and the
bundles in which the equations take values. We begin with scalar
functions and then consider maps, where the output bundle depends
on the unknown.

\begin{example}[Harmonic functions]
For a Riemannian manifold $N$, the Laplacian is an operator
\[
 \Delta\colon\Cinfty(N)\longrightarrow\Cinfty(N).
\]
We use the nonnegative convention $\Delta=d^*d$.
The equation $\Delta u=0$ says that $u$ is harmonic.
It occurs in potential theory and in extending boundary data into a domain.
On a closed connected manifold, integration by parts shows that harmonic
functions are constant:
\[
 0=\int_Nu\Delta u=\int_N|du|^2.
\]
Boundary-value problems require additional analytic hypotheses; the
representability theorem considered here concerns closed fibres.

Both source and target are sections of the trivial real line bundle
$N\times\R\to N$. The operator takes an entire section as its input
and returns another section.
\end{example}

\begin{example}[Critical points of energies]
Let $N$ be a closed Riemannian manifold and let $V\colon\R\to\R$ be smooth. Consider
\[
 \mathcal E(u)=\int_N\left(\tfrac12|du|^2+V(u)\right)\,d\mathrm{vol}.
\]
Differentiating in the direction $v$ and integrating by parts gives
\[
 d\mathcal E_u(v)=\int_N\bigl(\Delta u+V'(u)\bigr)v\,d\mathrm{vol}.
\]
Therefore the critical points are the solutions of
$P(u)=0$, where
\[
 P\colon\Cinfty(N)\longrightarrow\Cinfty(N),\qquad
 P(u)=\Delta u+V'(u).
\]
The energy is a real-valued functional; its Euler--Lagrange operator $P$
has a function as its output. The metric and volume form identify the
first variation with this function.

Taking $V(u)=u^3/3$ gives the running example $P(u)=\Delta u+u^2$.
\end{example}

\begin{example}[Pseudo-holomorphic maps]
Let $(\Sigma,j)$ be a fixed Riemann surface and $(Z,J)$ an almost complex manifold.
A smooth map $u\colon\Sigma\to Z$ is pseudo-holomorphic when
\[
 \bar\partial_Ju=\tfrac12(du+J(u)\circ du\circ j)=0.
\]
The unknown is a section of $\Sigma\times Z\to\Sigma$.
Its output lies in
\[
 \Gamma\bigl(\Sigma,
 \operatorname{Hom}^{0,1}_{\R}(T\Sigma,u^*TZ)\bigr),
\]
where the superscript denotes complex-antilinear real linear maps.
The output bundle depends on $u$.

Globally, this is a section of a bundle over the space of maps.
Near a fixed map, exponential coordinates on $Z$ and identifications
of nearby output bundles turn it into an operator
\[
 Q\subseteq\Gamma(\Sigma,F)\longrightarrow\Gamma(\Sigma,E)
\]
for fixed finite-rank bundles $F,E$.
\Cref{chap:Elliptic_Equations_Kuranishi_Models} explains the coordinate principle, and \Cref{chap:Pseudo_Holomorphic_Curves_Derived_Moduli} carries out
the Cauchy--Riemann application. Keeping track of these source and target
bundles is necessary before applying an ellipticity theorem.
\end{example}

\section{Section stacks}

To apply the functor-of-points viewpoint, we retain smooth families of
sections and their pullbacks. We first fix the domain manifold, then
allow it to vary over a base.

For a submersion $p\colon X\to N$, the ordinary section set is
\[
 \Gamma(N;X)=\{u\colon N\to X:\ p\circ u=\operatorname{id}_N\}.
\]
If $X$ is a vector bundle, this is a vector space.
A family parametrized by an ordinary manifold $T$ is a smooth map
$T\times N\to X$ over $N$. Smoothness is joint smoothness in the
parameter and in the point of $N$.

The same formula makes sense with derived test objects.
For $T$ derived, it uses maps of derived stacks, and so retains
infinitesimal and higher data in the parameter.

\begin{definition}[Section stack]
\label[definition]{def:Section_Stack}
The section stack of $X\to N$ is
\[
 \Sec_N(X)(T)=\Hom_{/N}(T\times N,X).
\]
\end{definition}

This is the mapping stack often written $\operatorname{Map}_N(N,X)$.
The notation $\Gamma(N;X)$ will be reserved for ordinary smooth
sections, and, for a vector bundle on compact $N$, their Fr\'echet
manifold. The section stack is a functor on parameter objects.

A family of manifolds over a smooth stack $S$
is a map $q\colon M\to S$ whose pullback along every ordinary manifold
$T\to S$ is a submersion of ordinary manifolds. The family is proper
if each such pullback is proper. For a manifold base this is the usual
notion of a submersion, or a proper submersion, respectively.
Suppose $q\colon M\to S$ is such a family and $X\to M$ is a family
of submersions.

\begin{definition}[Relative section stack]
\label[definition]{def:Relative_Section_Stack}
The relative section stack $\Sec_{M/S}(X)\to S$ is characterized by
\[
 \Hom_{/S}(T,\Sec_{M/S}(X))
 \cong\Hom_{/M}(T\times_S M,X).
\]
It is the Weil restriction $\operatorname{Res}_{M/S}(X)$.
\end{definition}

The terminology describes the right adjoint to pullback along $M\to S$.
Its existence here follows from the locally cartesian closed
$\infty$-category of stacks
\cite[Definition~3.2.1]{Steffens_Representability_Elliptic_Moduli}.
The formula also proves base-change compatibility:
\[
 \Sec_{M/S}(X)\times_S T
 \cong\Sec_{M_T/T}(X_T),
 \qquad M_T=M\times_S T.
\]
Over $s\in S$, its fibre is the section stack on $M_s$.

\section{Solution stacks as derived fibres}

Let $P\colon\Sec_N(X)\to\Sec_N(Y)$ be any morphism of stacks, and let
$b$ be a section of $Y$. Define
\begin{equation}\label{eq:Solution_Stack_Inhomogeneous}
 \Sol(P;b)=\Sec_N(X)\times_{\Sec_N(Y)}^{\mathrm{der}}*,
\end{equation}
where the second map classifies $b$.
The adjective derived emphasizes the category in which the pullback is
taken; all pullbacks in derived stacks are already homotopy pullbacks.

On a derived parameter $T$, a solution is a family $u$ together with a
homotopy from $P(u)$ to the constant family $b_T$.
Thus a solution stack can be defined without any finite-order assumption.
One does need a morphism of stacks, not merely a map between sets
of smooth sections.

For closed $N$, a smooth map between the corresponding convenient
section manifolds supplies such a morphism. This follows from the
comparison between section stacks and convenient manifolds explained
in \Cref{chap:Elliptic_Representability_Smooth_Derived_Bridge}. For example, on $S^1$ the map
\[
 u\longmapsto\left(x\longmapsto\int_{S^1}u\,d\theta\right)
\]
is a smooth linear map of section manifolds and defines a solution stack.
It is not a finite-order differential operator: changing $u$ away from
$x$ can change the output at $x$ while leaving every derivative there fixed.

For a parameter base $S$, a right-hand side is a section
$b\colon S\to\Sec_{M/S}(Y)$, and the same pullback is formed over $S$.
All these constructions commute with base change.
The finite-order differential condition introduced next provides a
geometric way to define $P$ and the analytic structure needed for
representability.

\section{Jets and finite-order operators}

Jets encode the dependence of an operator on finitely many derivatives.
Their relative construction also makes the induced map of section
stacks available on derived parameters.

A $k$-jet of a local section of $X\to N$ at $x\in N$ records its value
and its derivatives through order $k$ at $x$.
Two sections define the same jet when those derivatives agree in local
coordinates. The chain rule makes this independent of the coordinates.
These jets form the bundle $J_N^k(X)\to N$.

For example, for $X=N\times\R$, a first jet records a scalar and a covector.
A second jet records the second derivatives as well; a connection
identifies that last datum with a symmetric tensor.
The jet prolongation sends a section $u$ to $j^ku$.

\begin{definition}[Finite-order differential operator]
\label[definition]{def:Finite_Order_Differential_Operator}
A nonlinear differential operator of order at most $k$ from $X$ to $Y$
over $N$ is a smooth map
\[
 p\colon J_N^k(X)\longrightarrow Y
\]
over $N$. Its induced map on sections is
$P(u)=p\circ j^ku$.
\end{definition}

Using the lowercase $p$ distinguishes the finite-dimensional map on
jets from the induced map $P$ on sections.
The latter is determined at $x$ by finitely many derivatives of $u$ at
that same point. A general map on sections need not have this locality.

For $\Delta u+u^2$, $k=2$. In local coordinates the Laplacian depends
on derivatives of $u$ of order at most two, while $u^2$ uses only its value.
For the Cauchy--Riemann equation only first jets are needed.

For $M\to S$, relative jets $J^k_{M/S}(X)$ record derivatives in the
directions tangent to the fibres $M_s$. No differentiation in the
parameter $s$ is involved. The resulting map $P$ is therefore an
$S$-family of operators on the fibre manifolds.

The derived jet stack has a coordinate-independent construction.
Let $(M/S)^{(k)}$ be the order-$k$ infinitesimal neighbourhood of the
relative diagonal, with projections $\pi_1,\pi_2$ to $M$. Then
\[
 J^k_{M/S}(X)
 =\operatorname{Res}_{(M/S)^{(k)}/M}
       \bigl(X\times_{M,\pi_1}(M/S)^{(k)}\bigr),
\]
where the Weil restriction uses $\pi_2$.
It parametrizes sections on infinitesimal neighbourhoods.
For an ordinary submersion this recovers the classical jet bundle.
The construction, the prolongation map, and base-change compatibility are
proved in
\cite[Definition~3.3.7 and Proposition~3.3.9]
{Steffens_Representability_Elliptic_Moduli}.

In particular, a map $p\colon J^k_{M/S}(X)\to Y$ induces a morphism
of relative section stacks. This makes the preceding solution
construction available over all derived test objects.

\section{Differential moduli problems and ellipticity}

We now collect the geometric data of an equation into a differential
moduli problem. Its linearization gives the deformation-obstruction
complex, and its principal symbol determines whether elliptic analysis
can be applied.

A useful common formulation asks two operators to have the same output:
\[
 P_1(u_1)=P_2(u_2).
\]
It includes $P(u)=b$ by taking the second unknown bundle to be $N\to N$,
whose only section is the identity, and taking its operator to be $b$.

\begin{definition}
\label[definition]{def:Differential_Moduli_Problem}
Over a family $M\to S$, a differential moduli problem consists of
submersions $Y_1,Y_2,X\to M$ and maps
\[
 p_i\colon J^k_{M/S}(Y_i)\longrightarrow X.
\]
Its solution stack is
\[
 \Sol(\mathcal E)=
 \Sec_{M/S}(Y_1)
 \times_{\Sec_{M/S}(X)}
 \Sec_{M/S}(Y_2),
\]
with maps induced by $p_1,p_2$.
\end{definition}

This is \cite[Definitions~3.4.1--3.4.2]
{Steffens_Representability_Elliptic_Moduli}.
The order $k$ is a common order: an operator of lower order can be
composed with a jet truncation map. Ellipticity below is then a
condition on the symbol for that chosen common order.

When $S$ is a smooth stack, these data are specified on manifold
families over $S$ and are compatible with pullback. Relative jets and
section stacks commute with that pullback, so the equations and
solution stacks descend. This includes families with symmetries
encoded in the base stack.

Fix first $S=*$ and a solution $u$ of $P(u)=b$.
Variations of $u$ are sections of
$F=u^*T_{X/N}$; variations of the output at $b$ are sections of
$E=b^*T_{Y/N}$. Differentiating the finite jet formula gives a linear
differential operator
\[
 D_uP\colon\Gamma(N;F)\longrightarrow\Gamma(N;E).
\]
In coordinates it has the form
\[
 D_uP(v)=\sum_{|\alpha|\leq k}A_\alpha(x)\partial^\alpha v.
\]
Its coefficients depend on $u$ and its derivatives, but the expression
is linear in $v$.

For a matching problem and solution $(u_1,u_2)$ with common output $w$,
put
\[
 F=u_1^*T_{Y_1/N}\oplus u_2^*T_{Y_2/N},\qquad
 E=w^*T_{X/N}.
\]
The linearized matching equation is
\begin{equation}\label{eq:Tangent_Complex_Solution_Stack}
 D(v_1,v_2)=D_{u_1}P_1(v_1)-D_{u_2}P_2(v_2).
\end{equation}
Its deformation-obstruction complex is
$[\Gamma(N;F)\xrightarrow D\Gamma(N;E)]$ in degrees $0,-1$.
The kernel is the infinitesimal solution space.
The cokernel receives obstructions to solving inhomogeneous correction
equations, as in the finite-dimensional zero-locus calculation.
The representability proof will identify this complex with the
finite-dimensional tangent complex of a local derived chart.

Ellipticity is a condition on the highest-order part of the linearized
operator. In particular, it depends on the source and target bundles
of the full matching equation.

\begin{definition}[Principal symbol and ellipticity]
\label[definition]{def:Elliptic_Differential_Moduli_Problem}
For a linear operator $D$ of order $k$, its principal symbol is
\[
 \sigma_D(x,\xi)=\sum_{|\alpha|=k}A_\alpha(x)\xi^\alpha
 \colon F_x\longrightarrow E_x.
\]
It is elliptic if this map is invertible for every nonzero covector $\xi$.
A differential moduli problem is elliptic if its linearized matching
operator is elliptic at every solution on every fibre $M_s$.
\end{definition}

This real symbol convention omits the optional factors of $i$ used in
Fourier analysis; the invertibility condition is unchanged.
For the nonnegative Laplacian the real symbol with this convention is
$-|\xi|^2$, while the Fourier convention gives $|\xi|^2$.
Both are invertible away from the zero covector.

Ellipticity concerns the highest-order part. On a compact manifold
it leads to finite-dimensional kernel and cokernel after Sobolev
completion. It does not assert surjectivity. The latter failure is
exactly why a derived zero locus can be needed.

\section{Solutions and linearization of the running example}

Return to $P(u)=\Delta u+u^2$ on a closed connected manifold $N$.
If $P(u)=0$, integration gives
\[
 0=\int_N(\Delta u+u^2)=\int_Nu^2,
\]
so $u=0$. Its ordinary solution set is one point.

Differentiating gives
\[
 D_uP(v)=\Delta v+2uv.
\]
The zeroth-order term does not change the principal symbol, so this
operator is elliptic. At zero it is $\Delta$.
The standard solvability theorem for the Laplacian on a closed connected
manifold says that its image consists of the functions of integral zero.
Consequently,
\[
 \ker\Delta\cong\R,\qquad\operatorname{coker}\Delta\cong\R.
\]
Constants give representatives in both cases.

This explains the usefulness of the example. Its real zero set is
particularly simple, while its infinitesimal solutions and obstructions
are both nontrivial. \Cref{chap:Elliptic_Representability_Derived_Charts} will obtain its local Kuranishi equation
explicitly, recovering a multiple point.

\section{Further reading}

The jet tower makes the highest-order part of an operator explicit.
We describe it for vector bundles, then illustrate parameter dependence
with a family of metrics and give two further calculations.

For a vector bundle $F\to N$, the jet tower has exact sequences
\[
 0\longrightarrow\operatorname{Sym}^k(T^*N)\otimes F
 \longrightarrow J_N^k(F)\longrightarrow J_N^{k-1}(F)
 \longrightarrow0.
\]
The kernel consists of the highest derivatives when all lower derivatives
vanish. Connections can split these sequences, but the exact sequences
themselves are canonical. This is also the part of the jet data
used by the principal symbol.

A smooth family of metrics $g_s$ gives operators
$P_s(u)=\Delta_{g_s}u+u^2$ on the fibres. For a test family
$T\to S$, the equation is evaluated using the pulled-back metric.
The parameter affects the coefficients, while all derivatives of $u$
remain in the fibre directions.
The relative solution stack is defined before any assertion about
its projection to $S$ being a submersion or being representable.

\begin{exercise}
Compute the Euler--Lagrange operator for
$\int_N(\tfrac12|du|^2+\tfrac14u^4-\tfrac12u^2)$ and its linearization
at each constant solution. Check ellipticity.
\end{exercise}

\begin{exercise}
Let $P_1,P_2$ be first-order linear operators with a common target.
Write the principal symbol of the matching equation.
Explain why ellipticity of either operator individually does not
imply ellipticity of their difference on the direct sum of the sources.
\end{exercise}

The section and jet constructions are in
\cite[Sections~3.2--3.3 and Definitions~3.4.1--3.4.2]
{Steffens_Representability_Elliptic_Moduli}.
The remainder of the seminar proves the representability statement
for the elliptic problems defined here.

% This chapter is the complete working manuscript for the eleventh seminar talk.
% The final section contains supplementary material outside the default live
% route.

\chapter{Fredholm equations and finite-dimensional reduction}
\label[chapter]{chap:Elliptic_Equations_Kuranishi_Models}

Talk 10 constructed the intrinsic solution stack of a differential equation
and calculated its tangent complex. For
\[
  P(u)=\Delta u+u^2
\]
at its unique solution, that complex is
\[
  [\Cinfty(M)\xrightarrow{\Delta}\Cinfty(M)].
\]
Its terms are infinite-dimensional, but its deformation and obstruction
spaces are both one-dimensional.

This chapter asks how finite-dimensional cohomology becomes an actual
finite-dimensional nonlinear model. Sobolev completion makes Banach calculus
available. Ellipticity makes the linearization Fredholm. Adding a
finite-dimensional obstruction space makes the augmented equation
submersive, and the Banach implicit function theorem then produces a
finite-dimensional manifold carrying an obstruction section.

The obstruction-space construction is the only general reduction in the live
route. Lyapunov--Schmidt coordinates are used only to calculate the running
example and are developed in full in the supplement. The chapter ends with a
precise boundary: the resulting Kuranishi model describes the ordinary smooth
solution germ and its tangent complex, but has not yet been identified with an
open restriction of the intrinsic derived solution stack.

\section{From smooth sections to a Fredholm problem}
\label[section]{sec:Smooth_Sections_To_Fredholm_Problem}

\subsection{Why an analytic completion is needed}
\label[subsection]{sec:Why_Analytic_Completion_Needed}

Let \(E\to M\) and \(F\to M\) be smooth vector bundles over a closed
manifold, and let
\[
  P\colon\mathcal U\subseteq\Gamma(E)\longrightarrow\Gamma(F)
\]
be a nonlinear differential operator of order \(k\). Here \(\mathcal U\) is
an open subset of the smooth section space. Near a solution \(u_0\), its
linearization is a linear differential operator
\[
  D_{u_0}P\colon\Gamma(E)\longrightarrow\Gamma(F)
\]
of order at most \(k\).

The space \(\Gamma(E)\) is naturally a Fr\'echet space. One may measure a
smooth section and all of its derivatives, but no single norm defines this
topology. The ordinary inverse and implicit function theorems for Banach
spaces therefore do not apply directly. There are inverse function theorems
for suitable Fr\'echet spaces, most notably the Nash--Moser theorem, but they
would introduce much more analysis than is needed here.

The standard alternative for an elliptic equation is to work temporarily at
a fixed Sobolev regularity. Choose bundle metrics, a Riemannian metric, and a
connection. For a nonnegative integer \(r\), the Sobolev norm is schematically
\[
  \lVert u\rVert_{H^r}^2
  =\sum_{j=0}^r\lVert\nabla^ju\rVert_{L^2}^2.
\]
The completion of \(\Gamma(E)\) in this norm is a Hilbert space
\(H^r(E)\). Different choices of metrics and connections give equivalent
norms. Moreover,
\[
  \Gamma(E)=\bigcap_{r\geq0}H^r(E).
\]
Thus smooth sections can be recovered by retaining all Sobolev levels.

An operator of order \(k\) loses \(k\) derivatives. Its natural completion
therefore has the form
\begin{equation}
  \label{eq:Sobolev_Extension_Nonlinear_Operator}
  P_l\colon\mathcal U_{l+k}\subseteq H^{l+k}(E)
  \longrightarrow H^l(F),
\end{equation}
not \(H^l(E)\to H^l(F)\). For sufficiently large \(l\), Sobolev embedding
and multiplication theorems imply that \(P_l\) is smooth. In the geometric
setting of jet-defined nonlinear operators, this is part of the functorial
Sobolev-scale construction in Steffens's Lemma~3.4.5
\cite[Lemma~3.4.5]{Steffens_Representability_Elliptic_Moduli}.

\begin{warning}[Completion changes the ambient space]
  \label[warning]{warn:Sobolev_Completion_Changes_Ambient_Space}
  A Sobolev completion introduces nonsmooth sections. It is a device for
  applying Banach calculus, not a claim that the moduli problem should be
  redefined using sections of one arbitrarily chosen regularity. Elliptic
  regularity will be needed to return to smooth solutions.
\end{warning}

\subsection{The linear elliptic theorem}
\label[subsection]{sec:Linear_Elliptic_Fredholm_Theorem}

We recall the functional-analytic property supplied by ellipticity.

\begin{definition}[Fredholm operator]
  \label[definition]{def:Fredholm_Operator}
  A bounded linear operator \(L\colon X\to Y\) between Banach spaces is
  \emph{Fredholm} if
  \[
    \dim\ker L<\infty,
    \qquad
    \operatorname{im}L\subseteq Y\text{ is closed},
    \qquad
    \dim\operatorname{coker}L<\infty.
  \]
  Its \emph{index} is
  \[
    \operatorname{ind}L
    =\dim\ker L-\dim\operatorname{coker}L.
  \]
\end{definition}

\begin{theorem}[Elliptic operators are Fredholm]
  \label[theorem]{thm:Elliptic_Operators_Are_Fredholm}
  Let \(M\) be closed and let
  \[
    L\colon\Gamma(E)\longrightarrow\Gamma(F)
  \]
  be an elliptic linear differential operator of order \(k\). Its Sobolev
  extension
  \[
    L_l\colon H^{l+k}(E)\longrightarrow H^l(F)
  \]
  is Fredholm. Its kernel consists of smooth sections, and its cokernel can be
  represented by a finite-dimensional space of smooth sections of \(F\).
\end{theorem}

This is the principal input from linear elliptic theory. The symbol condition
called ellipticity in the preceding chapter becomes a finite-dimensionality
statement after Sobolev completion. In particular,
\[
  H^0[H^{l+k}(E)\xrightarrow{L_l}H^l(F)]
  =\ker L_l
\]
and
\[
  H^1[H^{l+k}(E)\xrightarrow{L_l}H^l(F)]
  =\operatorname{coker}L_l
\]
are finite-dimensional.

The theorem does not yet give a finite-dimensional nonlinear model. It says
only that the failure of the linearized equation to be invertible is
finite-dimensional. The implicit function theorem will solve all remaining
directions.

\subsection{The nonlinear return to smoothness}
\label[subsection]{sec:Nonlinear_Return_To_Smoothness}

The second quoted input is nonlinear elliptic regularity.

\begin{theorem}[Nonlinear elliptic regularity]
  \label[theorem]{thm:Nonlinear_Elliptic_Regularity}
  Let \(P\) be a nonlinear elliptic differential operator of order \(k\).
  Suppose that \(l\) is sufficiently large, that
  \(u\in H^{l+k}(E)\), and that
  \[
    P_l(u)=f
  \]
  for a smooth section \(f\in\Gamma(F)\). Then \(u\) is smooth.
\end{theorem}

The required lower bound on \(l\) depends on the operator and the dimension
of \(M\). Since we are free to begin at an arbitrarily high Sobolev level, the
precise optimal bound is irrelevant to the reduction. What matters is that a
sufficiently regular weak solution with smooth right-hand side is smooth.
Steffens records exactly this form in Lemma~3.4.6
\cite[Lemma~3.4.6]{Steffens_Representability_Elliptic_Moduli}.

This theorem will be applied not only to \(P_l(u)=0\), but also to an
augmented equation
\[
  P_l(u)+\iota(c)=0,
\]
where \(c\) belongs to a finite-dimensional space and \(\iota(c)\) is a
smooth section. Thus every point of the finite-dimensional manifold produced
below is represented by a smooth section \(u\).

The analytic division of labor is now clear.
\[
  \begin{gathered}
    \text{Sobolev completion}
    \quad\Longrightarrow\quad
    \text{Banach calculus},\\
    \text{ellipticity}
    \quad\Longrightarrow\quad
    \text{Fredholm linearization},\\
    \text{elliptic regularity}
    \quad\Longrightarrow\quad
    \text{smooth solutions}.
  \end{gathered}
\]

\section{The Fredholm--Kuranishi reduction}
\label[section]{sec:Fredholm_Kuranishi_Reduction}

We now isolate the nonlinear construction from its elliptic origin. Let
\(X\) and \(Y\) be Banach spaces, let \(U\subseteq X\) be open, and let
\[
  P\colon U\longrightarrow Y
\]
be smooth. Fix \(x_0\in U\) with \(P(x_0)=0\), and suppose that
\[
  L=D_{x_0}P\colon X\longrightarrow Y
\]
is Fredholm.

\subsection{Splitting the linearized equation}
\label[subsection]{sec:Splitting_Linearized_Equation}

Put \(K=\ker L\). Since \(K\) is finite-dimensional, it admits a closed
complement \(X'\subseteq X\). Since \(\operatorname{im}L\) is closed and has
finite codimension, choose a finite-dimensional subspace \(C\subseteq Y\)
which maps isomorphically to \(\operatorname{coker}L\). We then have
splittings
\begin{equation}
  \label{eq:Fredholm_Splittings}
  X=K\oplus X',
  \qquad
  Y=\operatorname{im}L\oplus C.
\end{equation}
Write
\[
  \iota\colon C\hookrightarrow Y
\]
for the chosen inclusion. The restriction
\begin{equation}
  \label{eq:Fredholm_Complement_Isomorphism}
  L|_{X'}\colon X'\iso\operatorname{im}L
\end{equation}
is an isomorphism of Banach spaces.

The kernel \(K\) consists of infinitesimal solutions. The space \(C\)
records the failure of the linearized equation to be surjective. Neither
splitting in \Cref{eq:Fredholm_Splittings} is canonical. The resulting chart
will accordingly be a presentation rather than an intrinsic object.

\subsection{Adding an obstruction space}
\label[subsection]{sec:Adding_Obstruction_Space}

Consider the augmented map
\begin{equation}
  \label{eq:Augmented_Fredholm_Map}
  \widetilde P\colon U\times C\longrightarrow Y,
  \qquad
  \widetilde P(x,c)=P(x)+\iota(c).
\end{equation}
Its derivative at \((x_0,0)\) is
\[
  D_{(x_0,0)}\widetilde P(v,e)=L(v)+\iota(e).
\]
This map is surjective by the second splitting in
\Cref{eq:Fredholm_Splittings}. Its kernel is
\[
  \ker D_{(x_0,0)}\widetilde P
  =K\times\{0\}.
\]

We use the following form of the Banach implicit function theorem.

\begin{theorem}[Banach submersion theorem]
  \label[theorem]{thm:Banach_Submersion_Theorem}
  Let \(f\colon U\to Y\) be a smooth map between Banach spaces. If
  \(D_xf\) is surjective and has complemented kernel, then, near \(x\), the
  level set \(f^{-1}(f(x))\) is a smooth Banach submanifold with tangent space
  \(\ker D_xf\).
\end{theorem}

In the present application, the kernel is finite-dimensional and hence
complemented. Surjectivity is open for the family of Fredholm operators, so
after shrinking \(U\times C\), we may assume that
\(D\widetilde P\) remains surjective throughout the neighborhood under
consideration.

\begin{theorem}[Fredholm--Kuranishi reduction]
  \label[theorem]{thm:Fredholm_Kuranishi_Reduction}
  Let \(P\colon U\to Y\) and \(x_0\) be as above. After choosing the
  splittings in \Cref{eq:Fredholm_Splittings}, there are a neighborhood
  \(W\subseteq U\times C\) of \((x_0,0)\), a finite-dimensional manifold
  \[
    V=W\cap\widetilde P^{-1}(0),
  \]
  and a smooth map
  \[
    s\colon V\longrightarrow C,
    \qquad
    s(x,c)=c,
  \]
  with the following properties.
  \begin{enumerate}
    \item One has \(\dim V=\dim K\).
    \item The zero locus \(s^{-1}(0)\) is naturally identified with the germ
      of \(P^{-1}(0)\) at \(x_0\).
    \item At every \(v=(x,0)\in s^{-1}(0)\), there is a natural
      quasi-isomorphism
      \begin{equation}
        \label{eq:Kuranishi_Fredholm_Tangent_Comparison}
        [T_vV\xrightarrow{D_vs}C]
        \longrightarrow
        [X\xrightarrow{D_xP}Y].
      \end{equation}
    \item At the base point, the left side of
      \Cref{eq:Kuranishi_Fredholm_Tangent_Comparison} is
      \[
        K\xrightarrow{0}C
      \]
  \end{enumerate}
\end{theorem}

\begin{proof}
  The derivative of \(\widetilde P\) at \((x_0,0)\) is a split surjection
  with kernel \(K\times\{0\}\). By
  \Cref{thm:Banach_Submersion_Theorem}, its zero set is therefore a smooth
  finite-dimensional manifold near \((x_0,0)\), with dimension \(\dim K\).
  This gives \(V\).

  The equation \(s(x,c)=0\) says that \(c=0\). On \(V\), the augmented
  equation also says
  \[
    P(x)+\iota(c)=0.
  \]
  Consequently,
  \[
    s^{-1}(0)
    =\{(x,0)\in W\mid P(x)=0\},
  \]
  which proves the second assertion.

  It remains to compare tangent complexes. At a zero \(v=(x,0)\) of \(s\),
  the tangent space of \(V\) is
  \[
    T_vV
    =\{(\xi,e)\in X\oplus C
      \mid D_xP(\xi)+\iota(e)=0\}.
  \]
  Projection to the two coordinates gives maps
  \[
    \operatorname{pr}_X\colon T_vV\longrightarrow X,
    \qquad
    D_vs=\operatorname{pr}_C\colon T_vV\longrightarrow C.
  \]
  The pair \((\operatorname{pr}_X,-\iota)\) defines the map of complexes in
  \Cref{eq:Kuranishi_Fredholm_Tangent_Comparison}, since
  \[
    D_xP(\operatorname{pr}_X(\xi,e))
    =-\iota(e)
    =-\iota(D_vs(\xi,e)).
  \]

  On degree-zero cohomology, projection to \(X\) gives an isomorphism
  \[
    \ker(D_vs)
    =\{(\xi,0)\mid D_xP(\xi)=0\}
    \iso
    \ker(D_xP).
  \]
  On degree-one cohomology, the induced map is
  \[
    C/\operatorname{im}(D_vs)
    \longrightarrow
    Y/\operatorname{im}(D_xP),
    \qquad
    [e]\longmapsto[-\iota(e)].
  \]
  Surjectivity of
  \[
    D_xP+\iota\colon X\oplus C\longrightarrow Y
  \]
  shows that this map is surjective. If \(\iota(e)=D_xP(\xi)\), then
  \((\xi,-e)\in T_vV\), so \(e\) belongs to
  \(\operatorname{im}(D_vs)\). It is therefore injective as well. The map of
  complexes is a quasi-isomorphism.

  At \((x_0,0)\), one has \(T_{(x_0,0)}V=K\times\{0\}\). Hence
  \(D_{(x_0,0)}s=0\), giving the final assertion.
\end{proof}

The trivial vector bundle \(V\times C\to V\), together with the section
\(s\), is a Kuranishi model in the sense of Chapter~8. Its virtual dimension
is
\begin{equation}
  \label{eq:Fredholm_Index_Virtual_Dimension}
  \dim V-\dim C
  =\dim\ker L-\dim\operatorname{coker}L
  =\operatorname{ind}L.
\end{equation}

\subsection{Returning to smooth elliptic equations}
\label[subsection]{sec:Returning_To_Smooth_Elliptic_Equations}

Apply \Cref{thm:Fredholm_Kuranishi_Reduction} to the Sobolev extension
\(P_l\) in \Cref{eq:Sobolev_Extension_Nonlinear_Operator}. By
\Cref{thm:Elliptic_Operators_Are_Fredholm}, the derivative at a smooth
solution is Fredholm. Choose \(C\) to be represented by smooth sections of
\(F\). Every point \((u,c)\in V\) satisfies
\[
  P_l(u)=-\iota(c),
\]
whose right-hand side is smooth. By
\Cref{thm:Nonlinear_Elliptic_Regularity}, the section \(u\) is smooth.
Consequently, the finite-dimensional manifold \(V\) and its zero locus do not
contain spurious nonsmooth solutions.

We have proved the promised analytic statement.

\begin{corollary}[Local Kuranishi model for an elliptic equation]
  \label[corollary]{cor:Local_Kuranishi_Model_Elliptic_Equation}
  Let \(P\) be a nonlinear elliptic differential operator on a closed
  manifold, and let \(u_0\) be a smooth solution. After choosing a sufficiently
  high Sobolev completion and a finite-dimensional obstruction space, there is
  a finite-dimensional Kuranishi model \((V,V\times C,s)\) such that
  \(s^{-1}(0)\) is the germ of smooth solutions near \(u_0\). At every nearby
  solution, its tangent complex is quasi-isomorphic to the Sobolev completion
  of the linearized differential operator.
\end{corollary}

Linear elliptic regularity identifies the kernel and cokernel of the Sobolev
linearization with those of the operator on smooth sections. Thus the
finite-dimensional tangent complex also recovers the smooth
deformation-obstruction complex from the solution stack in Chapter~10.

\section{The running equation}
\label[section]{sec:Laplacian_Square_Fredholm_Reduction}

We now carry out the reduction completely for the running example. Let \(M\)
be a closed connected Riemannian manifold and let
\[
  P(u)=\Delta u+u^2.
\]
Choose an integer \(l>\dim(M)/2\). Sobolev multiplication makes
\[
  P_l\colon H^{l+2}(M)\longrightarrow H^l(M)
\]
a smooth map. Its derivative at zero is
\[
  D_0P_l=\Delta\colon H^{l+2}(M)\longrightarrow H^l(M).
\]

\subsection{Constants and mean-zero functions}
\label[subsection]{sec:Constants_Mean_Zero_Functions}

Normalize the projection onto constant functions by
\begin{equation}
  \label{eq:Mean_Projection}
  \Pi(f)
  =\frac{1}{\operatorname{vol}(M)}\int_Mf.
\end{equation}
For every Sobolev index \(r\), put
\[
  H^r_0(M)=\ker\Pi.
\]
Then
\begin{equation}
  \label{eq:Constant_Mean_Zero_Splitting}
  H^r(M)=\R\oplus H^r_0(M).
\end{equation}
The kernel of \(\Delta\) is the space of constants, and its image is the
space of mean-zero functions. The restriction
\begin{equation}
  \label{eq:Laplacian_Mean_Zero_Isomorphism}
  \Delta\colon H^{l+2}_0(M)\iso H^l_0(M)
\end{equation}
is an isomorphism. Hence the Fredholm splittings are
\[
  K=\R,
  \qquad
  X'=H^{l+2}_0(M),
  \qquad
  \operatorname{im}\Delta=H^l_0(M),
  \qquad
  C=\R.
\]

Write
\[
  u=a+v,
  \qquad
  a\in\R,
  \quad
  v\in H^{l+2}_0(M).
\]
The mean-zero component of the equation is
\begin{equation}
  \label{eq:Laplacian_Square_Mean_Zero_Equation}
  (1-\Pi)\bigl(\Delta v+(a+v)^2\bigr)=0.
\end{equation}
The derivative of its left side with respect to \(v\) at \((a,v)=(0,0)\) is
the isomorphism in \Cref{eq:Laplacian_Mean_Zero_Isomorphism}. The implicit
function theorem therefore gives a unique local solution
\[
  v=\varphi(a)
\]
of \Cref{eq:Laplacian_Square_Mean_Zero_Equation}.

At this point the example simplifies more than a general
Lyapunov--Schmidt reduction. For every constant \(a\), the choice \(v=0\)
solves the mean-zero equation because
\[
  (1-\Pi)(a^2)=0.
\]
By local uniqueness,
\begin{equation}
  \label{eq:Laplacian_Square_Mean_Zero_Correction_Vanishes}
  \varphi(a)=0
\end{equation}
for all sufficiently small \(a\).

\subsection{The exact obstruction map}
\label[subsection]{sec:Obstruction_Map_Exactly_A_Squared}

The remaining constant component of the equation is
\[
  \kappa(a)
  =\Pi(P_l(a+\varphi(a))).
\]
Using \Cref{eq:Laplacian_Square_Mean_Zero_Correction_Vanishes}, we obtain
\begin{equation}
  \label{eq:Laplacian_Square_Exact_Obstruction_Map}
  \kappa(a)
  =\Pi(a^2)
  =a^2.
\end{equation}
Thus no further coordinate change is needed. The local finite-dimensional
reduction is literally
\[
  \R\longrightarrow\R,
  \qquad
  a\longmapsto a^2.
\]

The same calculation can be displayed using the augmented equation. Include
the constants into \(H^l(M)\) and set
\[
  \widetilde P_l(u,c)=\Delta u+u^2+c.
\]
The mean-zero component again gives \(v=0\). The constant component gives
\[
  a^2+c=0.
\]
Hence the finite-dimensional manifold cut out by the augmented equation is
\[
  V=\{(a,0,-a^2)\mid |a|\text{ is small}\}\cong\R.
\]
Projection to the obstruction coordinate is the section
\[
  s(a)=-a^2.
\]
This is exactly the obstruction-space chart from
\Cref{thm:Fredholm_Kuranishi_Reduction}.

\subsection{The local derived singularity}
\label[subsection]{sec:Laplacian_Square_Local_Derived_Singularity}

The ordinary zero locus of \(a^2\) is one point. Its derived zero locus is
\begin{equation}
  \label{eq:Laplacian_Square_Candidate_Derived_Model}
  \R\times_{\R}^{\mathbf R}\{0\},
  \qquad
  a\longmapsto a^2.
\end{equation}
Its animated \(\Cinfty\)-ring is represented by the Koszul algebra
\[
  \bigl(\Cinfty(\R)\otimes\Lambda(e),\partial e=a^2\bigr),
\]
which was calculated in Chapter~5. Its tangent complex at zero is
\begin{equation}
  \label{eq:Laplacian_Square_Finite_Tangent_Complex}
  [\R\xrightarrow{0}\R].
\end{equation}

The inclusions of constant functions give a map from
\Cref{eq:Laplacian_Square_Finite_Tangent_Complex} to the Sobolev Fredholm
complex
\begin{equation}
  \label{eq:Laplacian_Square_Sobolev_Fredholm_Complex}
  [H^{l+2}(M)\xrightarrow{\Delta}H^l(M)].
\end{equation}
It induces an isomorphism on the kernels and cokernels, hence is a
quasi-isomorphism. Linear elliptic regularity gives a second quasi-isomorphism
to the smooth complex
\[
  [C^\infty(M)\xrightarrow{\Delta}C^\infty(M)]
\]
calculated in Chapter~10.

We have therefore reached the exact finite-dimensional model anticipated at
the beginning of the seminar:
\[
  \begin{gathered}
    \Delta u+u^2=0
    \quad\rightsquigarrow\quad
    a^2=0,\\
    [C^\infty(M)\xrightarrow{\Delta}C^\infty(M)]
    \quad\simeq\quad
    [\R\xrightarrow{0}\R].
  \end{gathered}
\]
The arrow \(\rightsquigarrow\) denotes the classical finite-dimensional
reduction proved in this chapter. Replacing it by an identification of derived
solution stacks is the next task.

\subsection{What the reduction proves, and what remains}
\label[subsection]{sec:Fredholm_Reduction_What_Remains}

The construction in this chapter solves the local analytic problem. It turns
one nonlinear elliptic equation near one solution into a finite-dimensional
section and preserves the entire two-term deformation-obstruction complex.
It also explains the local origin of quasi-smoothness: the candidate derived
model is the zero locus of a section of a finite-rank vector bundle.

Several choices entered the construction:
\begin{enumerate}
  \item A Sobolev index \(l\).
  \item A complement to the kernel of the linearization.
  \item A finite-dimensional subspace representing its cokernel.
  \item A neighborhood on which the augmented operator remains submersive.
\end{enumerate}
Different choices need not give identical Kuranishi charts. This is expected
and was the central lesson of Chapter~8. The intrinsic object should not be
defined to be any one of these presentations.

Steffens instead begins with the solution stack from Chapter~10. The analytic
construction is used to prove that an open restriction of this stack is
represented by the derived zero locus of the obstruction section. Once that
comparison is established, the uniqueness of representing objects supplies
the compatibility among the different local reductions, and the local
representability theorem from Chapter~9 glues them.

The remaining questions are therefore precise.

\begin{enumerate}
  \item Why does the Sobolev reduction compute a pullback in derived smooth
    stacks, rather than only the ordinary solution set?
  \item Why is the result independent of fixing one Sobolev regularity?
  \item How does the construction behave in smooth families of equations?
  \item How do the resulting local representatives glue over an arbitrary
    smooth-stack base?
\end{enumerate}

These questions are the proof architecture of the elliptic representability
theorem.

\section{Summary and supplementary materials}
\label[section]{sec:Elliptic_Kuranishi_Summary_Supplementary_Materials}

\subsection{Takeaways}
\label[subsection]{sec:Fredholm_Reduction_Takeaways}

The live route can be compressed into ten statements.

\begin{enumerate}
  \item Smooth section spaces are Fr\'echet, so the ordinary Banach implicit
    function theorem does not apply directly.
  \item An order-\(k\) differential operator extends at high regularity as
    \(H^{l+k}\to H^l\).
  \item The Sobolev extension of an elliptic linear operator on a closed
    manifold is Fredholm.
  \item Fredholmness says that the failure of invertibility is
    finite-dimensional.
  \item Adding a finite-dimensional representative of the cokernel makes the
    augmented nonlinear equation submersive.
  \item Its zero set is a finite-dimensional manifold carrying an obstruction
    section.
  \item The zero locus of that section is the original solution germ.
  \item The tangent complex of the Kuranishi chart is quasi-isomorphic to the
    Fredholm deformation-obstruction complex.
  \item Elliptic regularity shows that the solutions found after Sobolev
    completion are smooth.
  \item For \(P(u)=\Delta u+u^2\), the obstruction map is exactly
    \(a\mapsto a^2\).
\end{enumerate}

\subsection{Supplement: Lyapunov--Schmidt coordinates}
\label[subsection]{sec:Lyapunov_Schmidt_Derived_Model}

The obstruction-space construction is the form which appears in Steffens's
proof. A second description makes the finite-dimensional obstruction map more
explicit. It also gives the most efficient route through the running example.

\subsubsection{Solving the image equation}
\label[subsubsection]{sec:Solving_Image_Equation}

Retain the splittings
\[
  X=K\oplus X',
  \qquad
  Y=\operatorname{im}L\oplus C
\]
from \Cref{eq:Fredholm_Splittings}. Let
\[
  q\colon Y\longrightarrow\operatorname{im}L,
  \qquad
  \pi_C\colon Y\longrightarrow C
\]
be the two projections. Write a point near \(x_0\) uniquely as
\[
  x=x_0+k+w,
  \qquad
  k\in K,
  \quad
  w\in X'.
\]
The equation \(P(x)=0\) has two components:
\begin{align}
  qP(x_0+k+w)&=0,
    \label{eq:Lyapunov_Schmidt_Image_Equation}\\
  \pi_CP(x_0+k+w)&=0.
    \label{eq:Lyapunov_Schmidt_Obstruction_Equation}
\end{align}

The derivative with respect to \(w\) of the left side of
\Cref{eq:Lyapunov_Schmidt_Image_Equation} at \((k,w)=(0,0)\) is
\[
  qL|_{X'}=L|_{X'}
  \colon X'\iso\operatorname{im}L.
\]
The Banach implicit function theorem therefore gives a unique smooth map
\[
  \varphi\colon K_0\longrightarrow X'
\]
on a neighborhood \(K_0\subseteq K\) of zero such that
\[
  qP(x_0+k+\varphi(k))=0.
\]
The values of \(\varphi\) are the infinite-dimensional directions which have
been solved away.

The remaining equation is finite-dimensional.

\begin{definition}[Lyapunov--Schmidt obstruction map]
  \label[definition]{def:Lyapunov_Schmidt_Obstruction_Map}
  The \emph{Lyapunov--Schmidt obstruction map} is
  \begin{equation}
    \label{eq:Lyapunov_Schmidt_Obstruction_Map}
    \kappa\colon K_0\longrightarrow C,
    \qquad
    \kappa(k)=\pi_CP(x_0+k+\varphi(k)).
  \end{equation}
\end{definition}

By construction, there is an isomorphism of germs of ordinary zero loci
\begin{equation}
  \label{eq:Lyapunov_Schmidt_Zero_Locus_Identification}
  P^{-1}(0)_{x_0}
  \cong
  \kappa^{-1}(0)_0,
  \qquad
  x_0+k+\varphi(k)\longleftrightarrow k.
\end{equation}

The linear term of \(\kappa\) vanishes. Indeed, differentiating the image
equation at zero gives
\[
  L(D_0\varphi(k))=0.
\]
The left side belongs to \(\operatorname{im}L\), while
\(D_0\varphi(k)\in X'\), so
\[
  D_0\varphi=0.
\]
It follows that
\[
  D_0\kappa(k)
  =\pi_CL(k+D_0\varphi(k))
  =0.
\]
Thus all kernel directions are infinitesimal solutions at the base point. The
nonlinear terms of \(\kappa\) determine which of them extend to actual nearby
solutions.

\subsubsection{The two constructions agree}
\label[subsubsection]{sec:Two_Fredholm_Reductions_Agree}

The augmented equation
\[
  P(x)+\iota(c)=0
\]
first implies \(qP(x)=0\). Hence every point of the manifold \(V\) from
\Cref{thm:Fredholm_Kuranishi_Reduction} has the form
\[
  x=x_0+k+\varphi(k).
\]
Its \(C\)-component then says
\[
  \kappa(k)+c=0.
\]
Consequently, the map
\[
  k\longmapsto
  \bigl(x_0+k+\varphi(k),-\kappa(k)\bigr)
\]
is a local parametrization of \(V\), and the obstruction section becomes
\[
  s(k)=-\kappa(k).
\]

Thus the obstruction-space construction and Lyapunov--Schmidt reduction are
not competing methods. They are two descriptions of the same local chart.
The first is coordinate-free once the obstruction space has been chosen. The
second uses kernel and image splittings to write the chart as a map between
finite-dimensional vector spaces.

The sign is harmless. Multiplication by \(-1\) is an automorphism of the
obstruction space, so the derived zero loci of \(\kappa\) and \(-\kappa\) are
isomorphic.

\subsubsection{The tangent complex and virtual dimension}
\label[subsubsection]{sec:Lyapunov_Schmidt_Tangent_Complex}

The Kuranishi chart in Lyapunov--Schmidt coordinates is the section
\[
  \kappa\colon K_0\longrightarrow K_0\times C
\]
of the trivial vector bundle. Its candidate derived zero locus is
\begin{equation}
  \label{eq:Candidate_Derived_Lyapunov_Schmidt_Zero_Locus}
  Z^{\mathrm{der}}(\kappa)
  =K_0\times_C^{\mathbf R}\{0\}.
\end{equation}
At a zero \(k\), its tangent complex is
\[
  [K\xrightarrow{D_k\kappa}C].
\]
By \Cref{thm:Fredholm_Kuranishi_Reduction}, this complex is
quasi-isomorphic to
\[
  [X\xrightarrow{D_{x_k}P}Y],
  \qquad
  x_k=x_0+k+\varphi(k).
\]
At the base point, \(D_0\kappa=0\), so the complex is
\begin{equation}
  \label{eq:Kernel_Cokernel_Tangent_Complex}
  [K\xrightarrow{0}C].
\end{equation}

The quasi-isomorphism is more important than an equality of dimensions. It
identifies the deformation and obstruction spaces and behaves correctly at
nearby solutions where \(D_k\kappa\) need not vanish. Its Euler
characteristic is the Fredholm index by
\Cref{eq:Fredholm_Index_Virtual_Dimension}.

\begin{warning}[What has and has not been proved]
  \label[warning]{warn:Candidate_Derived_Local_Model}
  We have proved that the ordinary solution germ is the zero locus of
  \(\kappa\), and that the tangent complex of its derived zero locus agrees
  with the Fredholm deformation-obstruction complex. These facts strongly
  identify \Cref{eq:Candidate_Derived_Lyapunov_Schmidt_Zero_Locus} as the
  correct local derived model. They do not by themselves prove that this
  derived zero locus represents an open substack of the solution stack from
  Chapter~10.
\end{warning}

The missing assertion concerns the full functor of derived families, not only
its ordinary points and tangent complex. Steffens proves it by showing that
the augmented map is submersive in the derived stack setting. The Sobolev
scale enters that proof as a whole rather than through one fixed completion.


\subsection{Supplement: reduction with parameters}
\label[subsection]{sec:Supplement_Fredholm_Reduction_With_Parameters}

The relative theorem will use a parameterized version of the construction.
Let \(B\) be a finite-dimensional manifold and let
\[
  P\colon B\times U\longrightarrow Y
\]
be smooth. Suppose that \(P(b_0,x_0)=0\) and that the vertical derivative
\[
  L=D_xP(b_0,x_0)\colon X\longrightarrow Y
\]
is Fredholm. Choose \(C\subseteq Y\) representing its cokernel and consider
\[
  \widetilde P(b,x,c)=P(b,x)+\iota(c).
\]
Its derivative in the \((x,c)\)-directions is surjective at
\((b_0,x_0,0)\). Hence the full derivative is surjective, and the zero set
\[
  V=\widetilde P^{-1}(0)
\]
is a finite-dimensional manifold near that point.

The projection \(V\to B\) is a submersion. Indeed, given
\(\dot b\in T_bB\), surjectivity of
\[
  D_xP+\iota\colon X\oplus C\longrightarrow Y
\]
allows us to choose \((\dot x,\dot c)\) satisfying
\[
  D_bP(\dot b)+D_xP(\dot x)+\iota(\dot c)=0.
\]
Then \((\dot b,\dot x,\dot c)\in T_vV\) lifts \(\dot b\).

Projection to \(C\) again gives an obstruction section
\[
  s\colon V\longrightarrow C,
\]
and \(s^{-1}(0)\) is the family of original solutions. At a zero, the
relative tangent complex over \(B\) is quasi-isomorphic to the vertical
Fredholm complex \([X\xrightarrow{D_xP}Y]\).

This finite-dimensional parameterized statement is the local analytic model
behind Steffens's relative theorem. The actual proof must first reduce a
family of section stacks to this form and must then compare the resulting
pullback with the one in derived stacks.

\subsection{Supplement: the Sobolev tower}
\label[subsection]{sec:Supplement_Sobolev_Tower}

A fixed Sobolev completion is enough for the classical reduction, but it is
not an intrinsic replacement for the smooth section stack. Steffens therefore
uses the full tower
\[
  \Gamma(E)\longrightarrow\cdots\longrightarrow
  H^{l+1}(E)\longrightarrow H^l(E)\longrightarrow\cdots.
\]
The transition maps are compact and have dense image, and the smooth section
space is their inverse limit in convenient manifolds. An order-\(k\)
differential operator gives compatible maps
\[
  P_l\colon H^{l+k}(E)\longrightarrow H^l(F).
\]

After adding the same smooth obstruction space at every level, the Banach
implicit function theorem produces finite-dimensional augmented solution
manifolds \(V_l\). Nonlinear elliptic regularity implies that their points are
smooth at sufficiently high regularity. Linear elliptic regularity identifies
their tangent spaces across successive levels. The tower of manifolds
\(V_l\) is therefore eventually locally constant.

This statement contains the analytic heart of Steffens's argument, but it is
not yet the derived conclusion. The Sobolev tower is a limit diagram of smooth
stacks, not a limit diagram of derived \(\Cinfty\)-stacks. The next chapter
will explain how the stacky-submersion criterion bypasses this obstruction.

\subsection{Supplement: a parametrix explanation of Fredholmness}
\label[subsection]{sec:Supplement_Parametrix_Fredholmness}

The proof of \Cref{thm:Elliptic_Operators_Are_Fredholm} belongs to elliptic
analysis, but its mechanism can be summarized. Ellipticity permits the
construction of a pseudodifferential parametrix \(Q\) satisfying
\[
  QL=1-R_1,
  \qquad
  LQ=1-R_2,
\]
where \(R_1\) and \(R_2\) are smoothing operators. On Sobolev spaces over a
closed manifold, smoothing operators are compact. Thus \(L\) is invertible
modulo compact operators, which implies that it is Fredholm.

The same identities improve regularity. If \(Lu=f\) is smooth, then
\[
  u=Qf+R_1u.
\]
Both terms on the right are more regular than the original Sobolev estimate
alone suggests. Iteration gives smoothness. Nonlinear regularity applies this
idea to the linearization along a solution, with estimates controlling the
nonlinear remainder.

This sketch explains why Fredholmness and regularity arise from the same
elliptic mechanism. It is not a proof of the pseudodifferential parametrix
theorem.

\subsection{Supplement: dependence on choices}
\label[subsection]{sec:Supplement_Fredholm_Reduction_Choices}

The finite-dimensional vector spaces \(K\) and \(C\) are determined up to
isomorphism at the base point, but their embeddings and complements are
chosen. Changing a complement changes the solved map \(\varphi\). Enlarging
the obstruction space changes both the dimension of \(V\) and the section
\(s\). Shrinking the neighborhood changes the domain of the chart.

The operations from Chapter~8 already predict the appropriate comparisons.
Shrinking does not change the germ. Adding a redundant obstruction direction
is accompanied by a stabilization of the chart. Changes of complements give
local coordinate changes. These observations explain why different
reductions should present the same derived germ.

They do not replace the representability proof. Constructing comparisons
between every pair of choices, and then constructing all higher
compatibilities, would reproduce the Kuranishi-atlas problem. The intrinsic
solution stack avoids this task: once each reduction is shown to represent an
open restriction of the same stack, their comparisons and higher coherence
follow from uniqueness.

\subsection{Related literature}
\label[subsection]{sec:Fredholm_Reduction_Related_Literature}

The primary source for the chapter is Steffens,
\emph{Representability of Elliptic Moduli Problems in Derived
\(C^\infty\)-Geometry}. Section~1.0.1 explains the passage from Fr\'echet
section spaces to Sobolev completions and separates the analytic problem from
the derived-pullback problem. Lemma~3.4.5 constructs the functorial Sobolev
scale, and Lemma~3.4.6 records the nonlinear regularity theorem
\cite[Section~1.0.1 and Lemmas~3.4.5--3.4.6]{Steffens_Representability_Elliptic_Moduli}.
The proof of Theorem~3.4.3 adds a cokernel representative, applies the Banach
implicit function theorem at every Sobolev level, and then performs the
stack-theoretic comparison which we have postponed
\cite[Proof of Theorem~3.4.3 and
Remark~3.4.7]{Steffens_Representability_Elliptic_Moduli}.

The Banach implicit function theorem may be found in Lang's
\emph{Differential Manifolds}
\cite[Chapter~I]{Lang_Differential_Manifolds}. For the nonlinear elliptic
regularity input, Steffens points to Taylor's
\emph{Pseudodifferential Operators and Nonlinear PDE}, specifically
Theorem~1.2.D
\cite[Theorem~1.2.D]{Taylor_Pseudodifferential_Nonlinear_PDE}.

Pardon begins from the same linear principle: an elliptic complex on a compact
manifold is quasi-isomorphic to a finite-dimensional complex. Section~2 also
explains why this formulation behaves better in families than taking kernels
and cokernels separately
\cite[Section~2]{Pardon_Representability_Elliptic_Fredholm}. In the first step
of the proof sketch of Theorem~5.1, Newton--Picard iteration represents the
regular locus, where the linearization is surjective. The later
finite-dimensional parameter thickening plays the role of adding obstruction
directions
\cite[Proof sketch of Theorem~5.1]{Pardon_Representability_Elliptic_Fredholm}.

% This chapter is the complete working manuscript for the twelfth seminar
% talk. The final section contains supplementary material outside the default
% live route.

\chapter{Elliptic representability I: the smooth-to-derived bridge}
\label[chapter]{chap:Elliptic_Representability_Smooth_Derived_Bridge}

Talk 10 defined the solution stack of a differential moduli problem before
making any analytic choices. Talk 11 began with one elliptic equation near one
solution, chose a Sobolev completion and a finite-dimensional obstruction
space, and produced a Kuranishi model
\[
  \kappa\colon V\longrightarrow C.
\]
Its ordinary zero locus is the local solution set, and its tangent complex is
quasi-isomorphic to the elliptic deformation-obstruction complex.

These two constructions do not yet coincide. The intrinsic solution stack
classifies smooth families over arbitrary test objects. The Kuranishi model
was obtained by applying the Banach implicit function theorem at a chosen
Sobolev regularity. Agreement on ordinary solutions and tangent complexes
does not prove that the model represents the same derived moduli functor.

The purpose of this chapter is to bridge that gap. We localize the intrinsic
stack until it is the derived zero locus of one nonlinear operator on smooth
section spaces. We then use a functorial tower of Sobolev completions to prove
one geometric assertion: the augmented smooth operator is submersive as a
morphism of derived $\Cinfty$-stacks. The Sobolev tower is used inside the
proof, but it is not used to define the moduli object.

This is the technical midpoint of Steffens's representability argument. The
next chapter will take two finite-dimensional pullbacks, construct the local
derived chart, and globalize it.

\section{The theorem and the missing bridge}
\label[section]{sec:Elliptic_Representability_Theorem_Gap}

\subsection{The intrinsic solution stack}
\label[subsection]{sec:Intrinsic_Solution_Stack_Recalled_Twelve}

Let $S$ be a smooth stack and let $M\to S$ be an $S$-family of manifolds. An
elliptic differential moduli problem
\[
  \mathcal E=(X\to M,Y_1\to M,Y_2\to M,P_1,P_2)
\]
is as in \Cref{def:Elliptic_Differential_Moduli_Problem}. The differential
operators induce maps of section stacks
\[
  P_i\colon
  \operatorname{Sect}_{M/S}(Y_i)
  \longrightarrow
  \operatorname{Sect}_{M/S}(X).
\]
Its solution stack is the pullback
\begin{equation}
  \label{eq:Elliptic_Solution_Stack_Twelve}
  \begin{tikzcd}
    \Sol(\mathcal E) \rar \dar \drar[pullback]
      & \operatorname{Sect}_{M/S}(Y_1) \dar["P_1"] \\
    \operatorname{Sect}_{M/S}(Y_2) \rar["P_2"']
      & \operatorname{Sect}_{M/S}(X).
  \end{tikzcd}
\end{equation}
This definition requires neither properness nor ellipticity. It is an
intrinsic construction in the $\infty$-category of derived
$\Cinfty$-stacks.

At a solution $t=(\sigma_1,\sigma_2,\eta)$ over $s\in S$, the relative
tangent complex is
\begin{equation}
  \label{eq:Elliptic_Relative_Tangent_Complex_Twelve}
  \left[
    \begin{matrix}
      \Gamma(M_s,\sigma_1^*T_{Y_{1,s}/M_s})\\
      {}\oplus
      \Gamma(M_s,\sigma_2^*T_{Y_{2,s}/M_s})
    \end{matrix}
    \xrightarrow{D_{\sigma_1}P_1-D_{\sigma_2}P_2}
    \Gamma(M_s,\sigma_X^*T_{X_s/M_s})
  \right]
\end{equation}
in cohomological degrees $0$ and $1$, by
\Cref{thm:Tangent_Complex_Solution_Stack}. Ellipticity says that the displayed
differential is elliptic. If $M_s$ is compact, its kernel and cokernel are
finite-dimensional, although the two chain groups are not.

\subsection{The destination}
\label[subsection]{sec:Relative_Elliptic_Representability_Destination}

We isolate the geometric hypothesis on the domains.

\begin{definition}[Proper family of manifolds]
  \label[definition]{def:Proper_Family_Manifolds_Twelve}
  Let $S$ be a smooth stack. A morphism $M\to S$ is an
  \emph{$S$-family of manifolds} if it is submersive as a morphism of smooth
  stacks. It is a \emph{proper $S$-family of manifolds} if it is also proper.
\end{definition}

For manifolds this is a proper submersion. Ehresmann's theorem then makes it
locally a smooth fiber bundle with compact fiber. Steffens proves the
corresponding local statement over a smooth-stack base
\cite[Proposition~2.1.61]{Steffens_Representability_Elliptic_Moduli}.

\begin{theorem}[Relative elliptic representability]
  \label[theorem]{thm:Relative_Elliptic_Representability}
  Let $S$ be a smooth stack, let $M\to S$ be a proper $S$-family of
  manifolds, and let $\mathcal E$ be an elliptic differential moduli problem
  over $M$. Then
  \[
    \Sol(\mathcal E)\longrightarrow S
  \]
  is representable by quasi-smooth derived $\Cinfty$-schemes locally of
  finite presentation.
\end{theorem}

This is Steffens's Theorem~3.4.3
\cite[Theorem~3.4.3]{Steffens_Representability_Elliptic_Moduli}. We state the
theorem now so that every local construction has a visible destination. Its
proof will be completed in the next chapter.

The conclusion is relative. For every representable map $T\to S$, the base
change
\[
  T\times_S\Sol(\mathcal E)
\]
is represented by a quasi-smooth derived $\Cinfty$-scheme locally of finite
presentation over $T$.

The theorem does not assert compactness, add broken or nodal domains, orient
the determinant complex, or construct a virtual fundamental class. It proves
that the solution functor already defined in
\Cref{eq:Elliptic_Solution_Stack_Twelve} is geometric.

\subsection{Why the Fredholm reduction is not enough}
\label[subsection]{sec:Fredholm_Reduction_Not_Enough_Twelve}

Suppose temporarily that $S$ is a manifold, that $M=S\times N$ with $N$
compact, and that the problem has been placed in vector-bundle coordinates.
It is then governed by a family
\[
  P\colon S\times Q\longrightarrow S\times\Gamma(E;N),
\]
where $Q\subseteq\Gamma(F;N)$ is open. At a solution $(s,\sigma)$, choose a
finite-dimensional space $C$ representing the cokernel of the vertical
linearization and a map
\[
  \iota\colon C\longrightarrow\Gamma(E;N)
\]
with smooth image. The augmented operator is
\begin{equation}
  \label{eq:Augmented_Operator_Informal_Twelve}
  \widetilde P(s',\tau,c)
  =\bigl(s',P_{s'}(\tau)+\iota(c)\bigr).
\end{equation}

Talk 11 applied the implicit function theorem to a Sobolev extension of this
map. That produces an ordinary finite-dimensional augmented solution
manifold, but it leaves three gaps:
\begin{enumerate}
  \item The intrinsic stack uses smooth sections, while the analytic
    construction uses Sobolev sections.
  \item A construction near ordinary points need not classify families over
    derived test objects.
  \item The Kuranishi model depends on auxiliary choices until it is
    identified with an open part of the intrinsic stack.
\end{enumerate}

Smooth sections are the inverse limit of their Sobolev completions in smooth
geometry. The inclusion of smooth stacks into derived $\Cinfty$-stacks need
not preserve such countable limits. Consequently, it would be invalid simply
to take a tower of Sobolev-level derived zero loci and declare its inverse
limit to be the intrinsic solution stack.

Steffens avoids this inverse limit. The tower is used to prove that
\Cref{eq:Augmented_Operator_Informal_Twelve} is a submersion in smooth stacks.
Submersiveness can then be transported into derived geometry, where only
finite pullbacks remain.

\section{Localizing the intrinsic problem}
\label[section]{sec:Localizing_Intrinsic_Elliptic_Problem_Twelve}

\subsection{A fixed compact domain}
\label[subsection]{sec:Fixed_Compact_Domain_Twelve}

Relative representability is local on $S$. Pull back along a smooth atlas,
then shrink around the image of the chosen solution. We may suppose that $S$
is a manifold. By properness and local triviality, we may further identify
\begin{equation}
  \label{eq:Proper_Family_Local_Product_Twelve}
  M\iso S\times N,
\end{equation}
where $N$ is compact.

The two parts of this reduction have different roles. The product description
gives one fixed space of sections and one functorial Sobolev scale throughout
the local family. Compactness makes the Sobolev extensions of elliptic
operators Fredholm.

Fix a solution
\[
  t=(s,\sigma_1,\sigma_2,\eta)
  \in\Sol(\mathcal E).
\]
The path $\eta$ identifies the two output sections and their tangent data. We
write the common output as $\sigma_X$.

\subsection{Section-stack charts}
\label[subsection]{sec:Section_Stack_Charts_Twelve}

A local addition on a submersion $Y\to S\times N$ identifies a neighborhood
of a section $\sigma$ with a neighborhood of the zero section in the vertical
tangent bundle $\sigma^*T_{Y/(S\times N)}$. Since $N$ is compact, this
pointwise construction induces an open chart on the stack of sections.

\begin{proposition}[Local chart on a section stack]
  \label[proposition]{prop:Local_Section_Stack_Chart_Twelve}
  Let $S$ be a manifold, let $N$ be compact, and let
  $Y\to S\times N$ be an $S$-family of submersions. Near a section over
  $s\in S$, there are an open neighborhood $S'\subseteq S$, a vector bundle
  $F\to N$, and an open substack inclusion
  \[
    S'\times Q\longrightarrow
    \operatorname{Sect}_{S'\times N/S'}(Y|_{S'\times N}),
  \]
  where $Q\subseteq\Gamma(F;N)$ is open and contains the zero section.
\end{proposition}

We quote this from Steffens
\cite[Proposition~3.2.22]{Steffens_Representability_Elliptic_Moduli}. The
geometric slogan is that nearby sections are vertical vector fields along the
chosen section. The stack-level assertion says that this description is
compatible with arbitrary parameters and defines an open substack.

Apply the proposition to $Y_1$, $Y_2$, and $X$. Relative jets commute with
the required base changes, by Steffens's Proposition~3.3.9
\cite[Proposition~3.3.9]{Steffens_Representability_Elliptic_Moduli}. The
finite-order operators therefore restrict to maps between the local section
charts.

We will also use the following finite-limit observation.

\begin{lemma}[Open charts pass to pullbacks]
  \label[lemma]{lem:Open_Charts_Pass_To_Pullbacks_Twelve}
  Consider a morphism between two cospans in an $\infty$-category of stacks
  with finite limits:
  \[
    \begin{tikzcd}[column sep=large]
      A \rar \dar[hook] & C \dar[hook] & B \lar \dar[hook] \\
      A' \rar & C' & B' \lar.
    \end{tikzcd}
  \]
  If the three vertical maps are open substack inclusions, then
  \[
    A\times_C B\longrightarrow A'\times_{C'}B'
  \]
  is an open substack inclusion.
\end{lemma}

\begin{proof}
  Open substack inclusions are monomorphisms and are stable under products and
  base change. The induced square
  \[
    \begin{tikzcd}
      A\times_C B \rar \dar \drar[pullback]
        & A\times B \dar \\
      A'\times_{C'}B' \rar
        & A'\times B'
    \end{tikzcd}
  \]
  is a pullback because two points of $A$ and $B$ with equal images in $C'$
  already have equal images in $C$. The result is therefore a base change of
  an open inclusion.
\end{proof}

This is the open-inclusion case of Steffens's Lemma~3.4.4
\cite[Lemma~3.4.4]{Steffens_Representability_Elliptic_Moduli}. Applying it to
the three section charts gives an open substack
\begin{equation}
  \label{eq:Local_Solution_Open_Substack_Twelve}
  \Sol(\mathcal E)_t\longrightarrow\Sol(\mathcal E)
\end{equation}
containing $t$.

\subsection{One nonlinear PDE}
\label[subsection]{sec:One_Nonlinear_PDE_Twelve}

In vector-bundle coordinates, the equality $P_1=P_2$ becomes the zero
equation for their difference. Put
\[
  F=F_1\times_NF_2,
  \qquad
  E=\sigma_X^*T_{X/M}.
\]
The finite-order condition gives an open vector-bundle neighborhood
$W\subseteq J_N^k(F)$ on which the local operator is defined. Let
\[
  Q=\{\tau\in\Gamma(F;N):j^k\tau(N)\subseteq W\}.
\]
Compactness of $N$ makes $Q$ open in the smooth section space. We obtain one
smooth map over $S$,
\begin{equation}
  \label{eq:Local_Nonlinear_PDE_Twelve}
  P\colon S\times Q\longrightarrow S\times\Gamma(E;N),
  \qquad
  (s',\tau)\longmapsto(s',P_{s'}(\tau)).
\end{equation}

The derived zero locus of \Cref{eq:Local_Nonlinear_PDE_Twelve} is precisely
the open stack in \Cref{eq:Local_Solution_Open_Substack_Twelve}. Ellipticity
says that at each solution in this chart the vertical derivative
\begin{equation}
  \label{eq:Vertical_Elliptic_Linearization_Twelve}
  D_\tau P_{s'}\colon
  \Gamma(F;N)\longrightarrow\Gamma(E;N)
\end{equation}
is elliptic.

The local equation has not replaced the intrinsic stack by an analytic
model. It describes an open restriction of that stack before any Sobolev
completion has been chosen.

\section{The Sobolev tower and obstruction augmentation}
\label[section]{sec:Sobolev_Tower_Obstruction_Augmentation_Twelve}

\subsection{The functorial tower}
\label[subsection]{sec:Functorial_Sobolev_Tower_Twelve}

For sufficiently large $l$, the order-$k$ operator in
\Cref{eq:Local_Nonlinear_PDE_Twelve} extends to a smooth map of Banach
manifolds
\begin{equation}
  \label{eq:Sobolev_Level_Operator_Twelve}
  P_l\colon S\times Q_{k+l}
  \longrightarrow S\times H^l(E;N),
\end{equation}
where $Q_{k+l}\subseteq H^{k+l}(F;N)$ is the open subset satisfying the
corresponding jet condition. The shift by $k$ records the loss of derivatives
of an order-$k$ operator.

Steffens's functorial Sobolev-scale lemma supplies four properties which are
used in the proof:
\begin{enumerate}
  \item Smooth sections are the inverse limit of the Sobolev spaces.
  \item The transition maps are compact and have dense image.
  \item Open jet conditions determine compatible open subsets $Q_{k+l}$.
  \item Jet prolongation and the nonlinear operator extend compatibly through
    the tower.
\end{enumerate}
The precise statement is Lemma~3.4.5
\cite[Lemma~3.4.5]{Steffens_Representability_Elliptic_Moduli}. In particular,
we have a commuting tower
\[
  \begin{tikzcd}[column sep=small]
    S\times Q \rar \dar["P"']
      & \cdots \rar
      & S\times Q_{k+l+1} \rar \dar["P_{l+1}"']
      & S\times Q_{k+l} \rar \dar["P_l"']
      & \cdots \\
    S\times\Gamma(E;N) \rar
      & \cdots \rar
      & S\times H^{l+1}(E;N) \rar
      & S\times H^l(E;N) \rar
      & \cdots.
  \end{tikzcd}
\]
Both rows are limit diagrams in convenient manifolds and hence in smooth
stacks.

We use convenient manifolds only as a setting in which smooth section spaces
and the displayed limits exist. No general convenient inverse function
theorem enters the proof.

\begin{warning}[Do not take the derived inverse limit]
  \label[warning]{warn:No_Derived_Sobolev_Inverse_Limit_Twelve}
  The two rows above need not remain limit diagrams after passage from smooth
  stacks to derived $\Cinfty$-stacks. In particular, one may not identify the
  intrinsic solution stack with the inverse limit of the Sobolev-level
  derived zero loci.
\end{warning}

The warning also has a geometric explanation. At one fixed Sobolev level,
reparametrization groups generally fail to act smoothly. A fixed completion
is therefore poorly suited to the global family structure of a moduli
problem. The intrinsic stack remains a stack of smooth sections.

\subsection{A compatible obstruction augmentation}
\label[subsection]{sec:Compatible_Obstruction_Augmentation_Twelve}

Fix the solution $(s,\sigma)$. The Sobolev extension of
$D_\sigma P_s$ is Fredholm. Choose a finite-dimensional space
\[
  C\cong\operatorname{coker}(D_\sigma P_s)
\]
and represent it by smooth sections through a map
\[
  \iota\colon C\longrightarrow\Gamma(E;N).
\]
Define
\[
  \widetilde P_l(s',\tau,c)
  =\bigl(s',P_{l,s'}(\tau)+\iota(c)\bigr).
\]

At $(s,\sigma,0)$, the vertical derivative of $\widetilde P_l$ is
surjective. Openness of the Fredholm locus and upper semicontinuity of
cokernel dimension give compatible neighborhoods
\[
  \widetilde U_l\subseteq S\times Q_{k+l}\times C
\]
on which every $\widetilde P_l$ is a Fredholm submersion. Their smooth cone
point is an open neighborhood
\[
  \widetilde U\subseteq S\times Q\times C.
\]
It carries the smooth augmented operator
\begin{equation}
  \label{eq:Augmented_Smooth_Operator_Twelve}
  \widetilde P\colon
  \widetilde U\longrightarrow S\times\Gamma(E;N),
  \qquad
  (s',\tau,c)\longmapsto
  \bigl(s',P_{s'}(\tau)+\iota(c)\bigr).
\end{equation}

The Banach implicit function theorem applies at each Sobolev level. The
remaining question is whether the smooth map
\Cref{eq:Augmented_Smooth_Operator_Twelve} is submersive in the sense needed
for derived pullbacks.

\section{Testing submersiveness on smooth families}
\label[section]{sec:Testing_Submersiveness_Smooth_Families_Twelve}

\subsection{The categorical criterion}
\label[subsection]{sec:Smooth_Derived_Submersion_Criterion_Twelve}

The comparison we need is a consequence of Steffens's
Proposition~2.2.18(3) and Corollary~2.2.19.

\begin{proposition}[Smooth-to-derived submersion criterion]
  \label[proposition]{prop:Smooth_To_Derived_Submersion_Criterion_Twelve}
  Let $f\colon\mathcal X\to\mathcal Y$ be a morphism of smooth stacks. Suppose
  that for every manifold $Z$ and every map $Z\to\mathcal Y$, the pullback
  formed in smooth stacks
  \[
    Z\times_{\mathcal Y}\mathcal X\longrightarrow Z
  \]
  is represented locally by a submersion of ordinary manifolds. Then the map
  obtained from $f$ in derived $\Cinfty$-stacks is submersive.
\end{proposition}

The result is not merely the assertion that ordinary submersions remain
submersions. Proposition~2.2.18 supplies the representable testing criterion,
and Corollary~2.2.19 transports the resulting notion from smooth to derived
stacks
\cite[Proposition~2.2.18 and Corollary~2.2.19]{Steffens_Representability_Elliptic_Moduli}.

Geometrically, the criterion retains exactly the property needed for finite
pullbacks. A pullback along a submersion is transverse, so its ordinary smooth
pullback already computes the corresponding derived pullback.

\subsection{Sobolev test pullbacks}
\label[subsection]{sec:Sobolev_Test_Pullbacks_Twelve}

\begin{proposition}[Derived submersiveness of the augmented operator]
  \label[proposition]{prop:Augmented_Operator_Derived_Submersive_Twelve}
  After shrinking $\widetilde U$ around $(s,\sigma,0)$, the map
  \Cref{eq:Augmented_Smooth_Operator_Twelve} is submersive as a morphism of
  derived $\Cinfty$-stacks.
\end{proposition}

\begin{proof}
  Let $Z$ be a manifold and let
  \[
    g\colon Z\longrightarrow S\times\Gamma(E;N)
  \]
  be a map of smooth stacks. Compose $g$ with the map to each Sobolev level
  and form
  \begin{equation}
    \label{eq:Sobolev_Test_Pullback_Twelve}
    Z_l
    =Z\times_{S\times H^l(E;N)}\widetilde U_l.
  \end{equation}
  Since $\widetilde P_l$ is a Fredholm submersion, the Banach implicit
  function theorem shows that $Z_l$ is a finite-dimensional manifold and
  $Z_l\to Z$ is a submersion. The spaces $Z_l$ form an inverse tower.

  We show that this tower is eventually locally constant. First consider
  points. A point of $Z_l$ is represented by $(z,\tau,c)$ satisfying
  \[
    P_l(\tau)+\iota(c)=g(z).
  \]
  The right-hand side and $\iota(c)$ are smooth sections. Nonlinear elliptic
  regularity therefore makes $\tau$ smooth. Thus the tower eventually stops
  acquiring new points locally.

  Next consider tangent spaces relative to $Z$. At a smooth point, the
  vertical tangent space of $Z_l\to Z$ is the kernel of the linearized
  augmented operator
  \[
    D\widetilde P_l\colon
    H^{k+l}(F;N)\oplus C\longrightarrow H^l(E;N).
  \]
  Linear elliptic regularity makes every element of this kernel smooth. Hence
  the transition maps eventually induce isomorphisms on vertical tangent
  spaces. Since both stages are submersive over $Z$, their total tangent
  spaces are then isomorphic as well, so the transition map is a local
  diffeomorphism.

  Point stabilization and tangent stabilization together make the inverse
  tower eventually locally constant. Its limit in smooth stacks is therefore
  represented locally by one of the finite-dimensional manifolds $Z_l$, and
  its map to $Z$ is a submersion. Functoriality of the Sobolev tower identifies
  that limit with
  \[
    Z\times_{S\times\Gamma(E;N)}\widetilde U.
  \]
  The criterion in
  \Cref{prop:Smooth_To_Derived_Submersion_Criterion_Twelve} proves the claim.
\end{proof}

The source compresses one analytic point in this proof. Pointwise elliptic
regularity alone does not give a neighborhood on which one Sobolev stage
works uniformly. Steffens's compatible choice of the neighborhoods
$\widetilde U_l$, together with the regularizing property of Lemma~3.4.6,
supplies the required local uniformity
\cite[Proof of Theorem~3.4.3, pp.~69--70]{Steffens_Representability_Elliptic_Moduli}.

The argument has not taken a limit in derived stacks. It took a limit in
smooth stacks, extracted the finite-limit property of submersiveness, and
transported only that property into derived geometry. This is the
smooth-to-derived bridge.

\subsection{The exact handoff}
\label[subsection]{sec:Exact_Handoff_Derived_Chart_Thirteen}

We now know that
\[
  \widetilde P\colon
  \widetilde U\longrightarrow S\times\Gamma(E;N)
\]
is derived-submersive. Consequently, its derived pullback along the zero
section is represented by the corresponding transverse smooth pullback. That
pullback is an ordinary finite-dimensional manifold $V$.

The original equation is not yet obtained by this first pullback, because
$V$ still carries the obstruction coordinate $c$. The next chapter will
project $V$ to $S\times C$ and impose $c=0$ by a second, finite-dimensional
derived pullback. The first pullback removes the infinite-dimensional
analysis. The second retains the finite-dimensional obstruction theory.

\section{Summary and supplementary materials}
\label[section]{sec:Smooth_Derived_Bridge_Summary_Supplementary}

\subsection{Takeaways}
\label[subsection]{sec:Smooth_Derived_Bridge_Takeaways}

The live route can be compressed into eight statements.
\begin{enumerate}
  \item The solution stack is defined intrinsically before choosing a
    completion or obstruction space.
  \item Properness reduces the domain family locally to $S\times N$ with
    $N$ compact.
  \item Local additions and relative jets reduce an open part of the
    intrinsic stack to the derived zero locus of one smooth family $P$.
  \item A functorial Sobolev tower supplies Banach manifolds without
    redefining the moduli problem at finite regularity.
  \item A finite-dimensional obstruction space makes every Sobolev-level
    augmented operator a Fredholm submersion after shrinking.
  \item Smooth test pullbacks are finite-dimensional submersions at every
    level.
  \item Nonlinear and linear elliptic regularity make their towers eventually
    locally constant.
  \item The testing criterion turns this smooth statement into derived
    submersiveness of the augmented operator.
\end{enumerate}

\subsection{Supplement: the general open-pullback lemma}
\label[subsection]{sec:Supplement_General_Open_Pullback_Twelve}

Steffens's Lemma~3.4.4 is more general than
\Cref{lem:Open_Charts_Pass_To_Pullbacks_Twelve}. It considers a morphism of
cospans
\[
  \begin{tikzcd}[column sep=large]
    X_0 \rar \dar["f_0"'] & X_1 \dar["f_1"] & X_2 \lar \dar["f_2"] \\
    Y_0 \rar & Y_1 & Y_2 \lar
  \end{tikzcd}
\]
and a geometric class of morphisms stable under products and base change. If
the $f_i$ lie in that class and the natural comparison conditions hold, then
the induced map
\[
  X_0\times_{X_1}X_2
  \longrightarrow
  Y_0\times_{Y_1}Y_2
\]
lies in the same class. Open inclusions are the case used in the live proof.

\subsection{Supplement: why finite regularity is not the moduli object}
\label[subsection]{sec:Supplement_Finite_Regularity_Not_Moduli_Twelve}

One might try to avoid the Sobolev tower by defining the moduli problem at one
large Sobolev index. This would make the implicit function theorem immediate,
but it would make the global object depend on an arbitrary analytic choice.
It also behaves badly under reparametrization: diffeomorphism groups of the
domain generally do not act smoothly on a fixed Sobolev completion.

Steffens's Remark~3.4.7 emphasizes the alternative. Define the moduli functor
using smooth families, use Sobolev spaces only to prove local
submersiveness, and return to the intrinsic smooth stack before taking the
derived pullback
\cite[Remark~3.4.7]{Steffens_Representability_Elliptic_Moduli}.

\subsection{Related literature}
\label[subsection]{sec:Smooth_Derived_Bridge_Related_Literature}

The primary source is Steffens's proof of Theorem~3.4.3. Proposition~3.2.22
provides section-stack charts, Proposition~3.3.9 controls relative jets,
Lemmas~3.4.4--3.4.6 provide the open-pullback and Sobolev inputs, and
Proposition~2.2.18 together with Corollary~2.2.19 supplies the categorical
bridge
\cite[Proposition~2.2.18, Corollary~2.2.19, Proposition~3.2.22,
Proposition~3.3.9, and Lemmas~3.4.4--3.4.6]{Steffens_Representability_Elliptic_Moduli}.
Pardon's published survey separates the ordinary regularity calculation from
the categorical extension to derived tests by a different method. That
comparison belongs to the end of the next chapter
\cite[Section~5]{Pardon_Representability_Elliptic_Fredholm}.

% This chapter is the complete working manuscript for the thirteenth seminar
% talk. The final section contains supplementary source material outside the
% default live route.

\chapter{Elliptic representability II: derived charts and globalization}
\label[chapter]{chap:Elliptic_Representability_Derived_Charts}

The previous chapter ended with one precise statement. Near a chosen solution
of an elliptic differential moduli problem, the augmented operator
\[
  \widetilde P\colon
  \widetilde U\longrightarrow S\times\Gamma(E;N)
\]
is submersive as a morphism of derived $\Cinfty$-stacks. All
infinite-dimensional analysis was concentrated in the proof of that
statement.

This chapter turns it into geometry. We first take the zero fiber of
$\widetilde P$. Submersiveness makes this an ordinary finite-dimensional
manifold $V$. We then project $V$ to the finite-dimensional obstruction space
and take one further derived zero fiber. Pullback pasting identifies the
result with an open restriction of the intrinsic solution stack. Varying the
solution produces an effective open atlas and completes the proof of
\Cref{thm:Relative_Elliptic_Representability}.

The final part compares this mechanism with Pardon's organization of the same
Fredholm geometry. Steffens proves derived submersiveness of an augmented
operator through a Sobolev tower. Pardon first represents the ordinary
regular locus, extends proper section stacks from smooth to derived test
objects, and makes an arbitrary solution regular by enlarging the parameter
space. The comparison is structural. Pardon's published article gives a
proof sketch for pseudo-holomorphic section problems, not a second complete
proof of Steffens's theorem in its full generality.

\section{The finite-dimensional derived chart}
\label[section]{sec:Finite_Dimensional_Derived_Chart_Thirteen}

\subsection{The augmented zero fiber is ordinary}
\label[subsection]{sec:Augmented_Zero_Fiber_Ordinary_Thirteen}

Retain the local data from the previous chapter. Thus $S$ is a manifold,
$N$ is compact,
\[
  P\colon S\times Q\longrightarrow S\times\Gamma(E;N)
\]
is the local family of elliptic equations, and
\[
  \widetilde U\subseteq S\times Q\times C
\]
is an augmented neighborhood on which
\begin{equation}
  \label{eq:Augmented_Operator_Thirteen}
  \widetilde P(s',\tau,c)
  =\bigl(s',P_{s'}(\tau)+\iota(c)\bigr)
\end{equation}
is derived-submersive.

Let $0_E\colon S\to S\times\Gamma(E;N)$ be the zero section and form the
derived pullback
\begin{equation}
  \label{eq:Augmented_Zero_Fiber_Thirteen}
  \begin{tikzcd}
    V \rar \dar \drar[pullback]
      & \widetilde U \dar["\widetilde P"] \\
    S \rar["0_E"']
      & S\times\Gamma(E;N).
  \end{tikzcd}
\end{equation}
By \Cref{prop:Augmented_Operator_Derived_Submersive_Twelve}, the right-hand
map is submersive. The square is therefore transverse, so its derived
pullback is represented by the ordinary smooth pullback
\begin{equation}
  \label{eq:Augmented_Solution_Manifold_Thirteen}
  V
  =\{(s',\tau,c)\in\widetilde U:
       P_{s'}(\tau)+\iota(c)=0\}.
\end{equation}
This is a finite-dimensional manifold, and $V\to S$ is a submersion.

The first pullback has removed the infinite-dimensional ambient section
spaces. It has not yet imposed the original equation, because a point of $V$
may have $c\neq0$.

\subsection{The obstruction section}
\label[subsection]{sec:Relative_Obstruction_Section_Thirteen}

Projection to the obstruction coordinate restricts to a map over $S$,
\begin{equation}
  \label{eq:Obstruction_Section_Thirteen}
  \kappa\colon V\longrightarrow S\times C,
  \qquad
  (s',\tau,c)\longmapsto(s',c).
\end{equation}
Let $0_C\colon S\to S\times C$ be the zero section. Define
\begin{equation}
  \label{eq:Derived_Obstruction_Zero_Locus_Thirteen}
  \begin{tikzcd}
    \mathcal Z_t \rar \dar \drar[pullback]
      & V \dar["\kappa"] \\
    S \rar["0_C"']
      & S\times C.
  \end{tikzcd}
\end{equation}
On ordinary points, $c=0$ reduces
$P_{s'}(\tau)+\iota(c)=0$ to $P_{s'}(\tau)=0$. The derived statement is
stronger and follows by pasting pullback squares.

To make the pasting precise, let
\[
  U_0
  =\widetilde U\times_{S\times Q\times C}(S\times Q),
\]
where $S\times Q\to S\times Q\times C$ is the zero section. The map
$U_0\to S\times Q$ is an open inclusion and
$\widetilde P|_{U_0}=P|_{U_0}$. Pasting
\Cref{eq:Augmented_Zero_Fiber_Thirteen} and
\Cref{eq:Derived_Obstruction_Zero_Locus_Thirteen} gives an isomorphism
\begin{equation}
  \label{eq:Derived_Chart_Local_Solution_Isomorphism_Thirteen}
  \mathcal Z_t
  \iso
  U_0\times_{S\times\Gamma(E;N)}S.
\end{equation}
The right side is the derived zero locus of $P$ restricted to $U_0$. By
\Cref{eq:Local_Solution_Open_Substack_Twelve} and
\Cref{lem:Open_Charts_Pass_To_Pullbacks_Twelve}, it is an open restriction
of the intrinsic stack $\Sol(\mathcal E)$ containing $t$.

This argument is the decisive comparison. Equality of ordinary zero sets and
agreement of tangent complexes would not identify the derived moduli
functors. The pasted universal properties do.

\subsection{Quasi-smoothness and the tangent complex}
\label[subsection]{sec:Quasi_Smoothness_Tangent_Thirteen}

Both $V$ and $S\times C$ are finite-dimensional manifolds over $S$. After
shrinking around $t$, they are affine $\Cinfty$-schemes of finite
presentation. The derived zero locus in
\Cref{eq:Derived_Obstruction_Zero_Locus_Thirteen} is therefore a quasi-smooth
derived $\Cinfty$-scheme locally of finite presentation.

At $v=(s,\sigma,0)\in V$, its relative tangent complex is
\begin{equation}
  \label{eq:Kuranishi_Tangent_Complex_Thirteen}
  \left[
    T_v(V/S)\xrightarrow{D_v\kappa}C
  \right]
\end{equation}
in cohomological degrees $0$ and $1$. We compare it directly with the
elliptic complex. Put
\[
  L=D_\sigma P_s\colon
  \Gamma(F;N)\longrightarrow\Gamma(E;N).
\]
The derivative of the augmented equation is
\[
  A=(L,\iota)\colon
  \Gamma(F;N)\oplus C
  \longrightarrow
  \Gamma(E;N),
\]
and it is surjective. Moreover,
\[
  T_v(V/S)=\ker A,
  \qquad
  D_v\kappa=\operatorname{pr}_C|_{\ker A}.
\]
There is therefore a chain map
\begin{equation}
  \label{eq:Kuranishi_Elliptic_Chain_Map_Thirteen}
  \begin{tikzcd}[column sep=large]
    \ker A \rar["\operatorname{pr}_C"] \dar["\operatorname{pr}_{\Gamma(F)}"']
      & C \dar["-\iota"] \\
    \Gamma(F;N) \rar["L"']
      & \Gamma(E;N).
  \end{tikzcd}
\end{equation}
It commutes because $L(\xi)+\iota(c)=0$ for $(\xi,c)\in\ker A$.

\begin{proposition}[Finite and elliptic tangent complexes agree]
  \label[proposition]{prop:Finite_Elliptic_Tangent_Complexes_Agree}
  The chain map in
  \Cref{eq:Kuranishi_Elliptic_Chain_Map_Thirteen} is a quasi-isomorphism.
  Consequently,
  \[
    \mathbb T_{\mathcal Z_t/S,v}
    \simeq
    [\Gamma(F;N)\xrightarrow{L}\Gamma(E;N)].
  \]
\end{proposition}

\begin{proof}
  On degree-zero cohomology, the map identifies
  \[
    \ker(D_v\kappa)
    =\{(\xi,0):L\xi=0\}
  \]
  with $\ker L$.

  For degree one, the induced map is
  \[
    C/\operatorname{im}(D_v\kappa)
    \longrightarrow
    \operatorname{coker}L,
    \qquad
    [c]\longmapsto[-\iota(c)].
  \]
  It is surjective because $A=(L,\iota)$ is surjective. Its kernel consists
  of those $c$ for which $\iota(c)=-L\xi$ for some $\xi$. Equivalently,
  $(\xi,c)\in\ker A$, so $c$ lies in the image of $D_v\kappa$. Hence the map
  is an isomorphism.
\end{proof}

The two pullbacks have complementary roles:
\[
  \begin{gathered}
    \widetilde P\text{ is submersive}
    \quad\Longrightarrow\quad
    V\text{ is an ordinary manifold},\\
    \kappa\text{ need not be submersive}
    \quad\Longrightarrow\quad
    \mathcal Z_t\text{ may carry derived structure}.
  \end{gathered}
\]
The augmentation removes the infinite-dimensional failure of transversality.
The derived zero locus retains the remaining finite-dimensional obstruction.

\section{Globalization and the representability theorem}
\label[section]{sec:Globalization_Representability_Thirteen}

\subsection{An effective open atlas}
\label[subsection]{sec:Effective_Open_Atlas_Thirteen}

The construction above applies near every solution $t$. For each such point,
we obtain a quasi-smooth derived $\Cinfty$-scheme $\mathcal Z_t$ locally of
finite presentation and an open morphism
\[
  \mathcal Z_t\longrightarrow\Sol(\mathcal E)
\]
whose image contains $t$.

These maps cover the solution stack. More precisely, every family of
solutions over a test manifold is locally on that manifold contained in one
of the chosen section-stack charts and augmented neighborhoods. Thus
\begin{equation}
  \label{eq:Effective_Open_Atlas_Thirteen}
  \coprod_t\mathcal Z_t
  \longrightarrow
  \Sol(\mathcal E)
\end{equation}
is an effective epimorphism and each summand is an open substack.

There is no independent coordinate-change problem to solve. On an overlap,
both charts represent open restrictions of the same intrinsic stack. Their
intersection is the corresponding pullback over $\Sol(\mathcal E)$, and its
coherent transition data are inherited from that stack.

\subsection{Completion of the proof}
\label[subsection]{sec:Completion_Relative_Representability_Thirteen}

\begin{proof}[Proof of \Cref{thm:Relative_Elliptic_Representability}]
  Work first after a representable base change to a manifold chart of $S$.
  Around every solution, the localization in Chapter~12 gives one smooth
  family of elliptic equations over a compact manifold. The Sobolev-scale
  argument and
  \Cref{prop:Augmented_Operator_Derived_Submersive_Twelve} make its augmented
  operator derived-submersive.

  The two pullbacks in
  \Cref{eq:Augmented_Zero_Fiber_Thirteen,eq:Derived_Obstruction_Zero_Locus_Thirteen}
  produce a quasi-smooth derived $\Cinfty$-scheme locally of finite
  presentation. Pullback pasting identifies it with an open substack of the
  intrinsic solution stack. Varying the solution yields the effective open
  atlas in \Cref{eq:Effective_Open_Atlas_Thirteen}.

  Apply the open-atlas criterion
  \Cref{thm:Open_Local_Representability}. It follows that the solution stack
  after base change to the manifold chart is represented. Finally, use the
  relative consequence
  \Cref{cor:Local_Representability_Relative} to descend this conclusion over
  the original smooth-stack base $S$. Quasi-smoothness and local finite
  presentation are local on the representing scheme, so they descend as
  well.
\end{proof}

This final globalization is deliberately short. The technical work was the
construction and identification of one open chart. Once those charts are
open substacks of a single intrinsic stack, Proposition~2.1.48 and
Corollary~2.1.49 perform the gluing
\cite[Proposition~2.1.48 and Corollary~2.1.49]{Steffens_Representability_Elliptic_Moduli}.

\subsection{What the hypotheses and choices do}
\label[subsection]{sec:Hypotheses_Choices_Thirteen}

The proof assigns a distinct role to every hypothesis.
\begin{enumerate}
  \item The smooth-stack base permits descent to manifold charts while
    retaining families and automorphisms.
  \item Properness gives compact fibers and local triviality of the family of
    domains.
  \item Finite order makes the nonlinear equation functorial through relative
    jets and compatible with the Sobolev tower.
  \item Ellipticity gives Fredholm linearizations and the regularity which
    returns Sobolev solutions to smooth sections.
  \item Finite-dimensional obstruction spaces turn the local problem into a
    quasi-smooth finite-presentation derived zero locus.
\end{enumerate}

Local trivializations, local additions, Sobolev indices, complements, and
representatives of the cokernel all occur in the construction. They do not
occur in the definition of $\Sol(\mathcal E)$. Every resulting chart is
identified with an open restriction of that same intrinsic stack. This is the
precise sense in which the representing object is independent of the
analytic choices.

\section{The proof pattern in finite dimensions}
\label[section]{sec:Proof_Pattern_Finite_Dimensions_Thirteen}

\subsection{The complete route}
\label[subsection]{sec:Complete_Representability_Route_Thirteen}

The proof can be read as the following chain of constructions:
\[
  \begin{aligned}
    &\text{Intrinsic elliptic solution stack}\\
    &\quad\rightsquigarrow
      \text{local smooth nonlinear operator }P\\
    &\quad\rightsquigarrow
      \text{derived-submersive augmentation }\widetilde P\\
    &\quad\rightsquigarrow
      \text{ordinary manifold }V\xrightarrow{\kappa}S\times C\\
    &\quad\rightsquigarrow
      V\times_{S\times C}^{\mathbf R}S\\
    &\quad\rightsquigarrow
      \text{effective open atlas}.
  \end{aligned}
\]
The first two arrows use differential geometry and jets. The third uses
Fredholm analysis, the Sobolev tower, and the smooth-to-derived criterion.
The fourth is a transverse pullback. The fifth is finite-dimensional derived
geometry. The last uses descent.

\subsection{The model calculation}
\label[subsection]{sec:X_Squared_Regular_Enlargement_Thirteen}

The whole route is visible in the function
\[
  f\colon\R\longrightarrow\R,
  \qquad
  f(x)=x^2.
\]
Its ordinary zero locus is one point, but its derived zero locus has tangent
complex
\[
  [\R\xrightarrow{0}\R]
\]
at the origin. Add an obstruction coordinate and define
\[
  \widetilde f\colon\R^2\longrightarrow\R,
  \qquad
  \widetilde f(x,c)=x^2+c.
\]
This map is a submersion. Its zero fiber
\[
  V=\{(x,c):x^2+c=0\}
\]
is an ordinary one-dimensional manifold. Restrict the second coordinate to
obtain
\[
  \kappa\colon V\longrightarrow\R,
  \qquad
  \kappa(x,-x^2)=-x^2.
\]
Pullback pasting gives
\begin{equation}
  \label{eq:X_Squared_Regular_Enlargement_Thirteen}
  \R\times_{\R}^{\mathbf R}\{0\}
  \iso
  V\times_{\R}^{\mathbf R}\{0\},
\end{equation}
where the map on the left is $f$ and the map on the right is $\kappa$.

The left side presents the singular equation directly. The right side first
solves a regular enlarged equation and then imposes the obstruction parameter
by a derived fiber. Steffens's proof is the elliptic family version of this
calculation.

\subsection{Representability is not a virtual class}
\label[subsection]{sec:Representability_Not_Virtual_Class_Thirteen}

The theorem proves that the moduli functor is represented and that its local
derived geometry is quasi-smooth. It does not automatically produce an
enumerative invariant. A virtual fundamental class additionally requires a
suitable global setting, typically including compactness or properness of the
moduli object and orientation data for its determinant complex.

The bordism interlude explained what a compact oriented derived manifold can
contribute. The present theorem supplies a represented quasi-smooth object,
but it does not establish those further hypotheses. This distinction remains
essential in applications.

\section{Comparison with Pardon's regular-locus engine}
\label[section]{sec:Pardon_Regular_Locus_Comparison_Thirteen}

\subsection{The regular locus on smooth tests}
\label[subsection]{sec:Pardon_Regular_Locus_Smooth_Tests_Thirteen}

Pardon's published argument begins with a pseudo-holomorphic section problem
over a proper family of smooth compact curves. Schematically, one has an
equation morphism
\[
  \Phi\colon
  \operatorname{Sect}_{C/B}(W)
  \longrightarrow
  \operatorname{Sect}_{C/B}(H)
\]
and the solution stack is its fiber over the zero section.

The \emph{regular locus} consists of the solutions where the total
linearization of $\Phi$ is surjective. Ordinary parametric Fredholm theory
represents this locus, on smooth test manifolds, by an ordinary smooth
manifold. The new question is not analytic:
\begin{quote}
  Why does the same manifold classify regular families over derived test
  objects?
\end{quote}

\subsection{Proper section stacks and derived tests}
\label[subsection]{sec:Pardon_Proper_Section_Stacks_Thirteen}

Let
\[
  i\colon\mathrm{Sm}\longrightarrow\mathrm{Der}
\]
denote the inclusion of manifolds into derived smooth manifolds. Restriction
of stacks has a left adjoint $i_!$, the sheaf left Kan extension. For a stack
$F$ on manifolds, $i_!F$ is the universal extension of $F$ to derived tests
which respects the required colimits and descent.

Pardon's Proposition~5.2 says, in the form needed here, that the section stack
of a submersion over a proper smooth source is obtained by this universal
extension:
\begin{equation}
  \label{eq:Pardon_Proper_Section_Stack_Thirteen}
  i_!\operatorname{Sect}_{C/B}(W)_{\mathrm{Sm}}
  \iso
  \operatorname{Sect}_{C/B}(W)_{\mathrm{Der}}.
\end{equation}
Properness and the softness supplied by partitions of unity are the
geometric inputs. The statement does not say that derived families are
ordinary families. It says that their coherent behavior on derived tests is
determined by the ordinary smooth functor.

Left adjoints do not preserve arbitrary pullbacks. Pardon's Lemma~5.3 gives
the additional fact used in the proof: sheaf left Kan extension preserves the
pullbacks in which the distinguished morphism is submersive. On the regular
locus, the equation morphism is submersive. Applying the lemma to its zero
fiber shows that
\begin{equation}
  \label{eq:Pardon_Regular_Locus_Extension_Thirteen}
  i_!\operatorname{Sol}(\Phi)^{\mathrm{reg}}_{\mathrm{Sm}}
  \iso
  \operatorname{Sol}(\Phi)^{\mathrm{reg}}_{\mathrm{Der}}.
\end{equation}
Thus the ordinary manifold which represents the regular locus on smooth tests
also represents it on derived tests
\cite[Proposition~5.2 and Lemma~5.3]{Pardon_Representability_Elliptic_Fredholm}.

\subsection{Parameter thickening}
\label[subsection]{sec:Pardon_Parameter_Thickening_Thirteen}

At a nonregular solution, choose a finite-dimensional space $C$ mapping
surjectively to the cokernel of the linearization. Enlarge the parameter base
by $C$ and add the corresponding inhomogeneous terms to the equation. The
chosen solution is regular in this enlarged family. The preceding argument
represents the enlarged regular locus by an ordinary manifold. The original
problem is recovered by taking its derived fiber over the zero value of the
new parameter.

This is the same enlargement-and-derived-fiber pattern as
\Cref{eq:X_Squared_Regular_Enlargement_Thirteen}. The two proofs organize its
middle step differently:
\begin{enumerate}
  \item Steffens adds an obstruction coordinate to the operator and uses the
    functorial Sobolev tower to prove that the augmented smooth operator is
    derived-submersive.
  \item Pardon enlarges the parameter base, represents the ordinary regular
    locus, uses proper section stacks to extend it to derived tests, and
    returns by a derived parameter fiber.
\end{enumerate}

Both methods obtain local representatives as open substacks of an intrinsic
solution functor. Consequently, both avoid constructing a strict cocycle of
independently chosen Kuranishi charts.

\subsection{Scope and source status}
\label[subsection]{sec:Pardon_Scope_Source_Status_Thirteen}

Pardon's published Theorem~5.1 concerns pseudo-holomorphic section problems
over proper families of smooth compact curves. Its conclusion and hypotheses
are not identical to Steffens's general theorem on elliptic differential
moduli problems. The article explicitly presents the argument as a brief
proof sketch
\cite[Theorem~5.1 and its proof sketch]{Pardon_Representability_Elliptic_Fredholm}.

Pardon's July 2026 manuscript develops a much broader logarithmic and
elliptic framework and gives a longer account of the regular-locus strategy.
It also explicitly identifies itself as unfinished work in progress and
contains unresolved internal placeholders. In particular, the proposed
logarithmic derived-base reduction is not a completed proof source. We use
the manuscript only to clarify the architecture of the published sketch
\cite[Section~N.6.4]{Pardon_Logarithmic_Derived_Moduli}.

\section{Summary and supplementary materials}
\label[section]{sec:Derived_Charts_Globalization_Summary_Supplementary}

\subsection{Takeaways}
\label[subsection]{sec:Derived_Charts_Globalization_Takeaways}

The live route has nine conclusions.
\begin{enumerate}
  \item The zero fiber of the derived-submersive augmented operator is an
    ordinary finite-dimensional manifold $V$.
  \item Projection to the obstruction coordinate gives a finite-dimensional
    map $\kappa\colon V\to S\times C$.
  \item Its derived zero locus recovers the original local equation.
  \item Pullback pasting identifies this derived zero locus with an open
    restriction of the intrinsic solution stack.
  \item Its tangent complex is quasi-isomorphic to the elliptic
    deformation-obstruction complex.
  \item The local chart is quasi-smooth and locally of finite presentation.
  \item The local charts form an effective open atlas, so the open-atlas
    theorem proves relative representability.
  \item The analytic choices affect presentations, not the intrinsic
    represented stack.
  \item Pardon's regular-locus proof uses the same enlargement-and-derived-
    fiber pattern with a different smooth-to-derived engine.
\end{enumerate}

\subsection{Supplement: the precise published Pardon endpoint}
\label[subsection]{sec:Supplement_Pardon_Published_Endpoint_Thirteen}

Pardon's published Derived Regularity Theorem states representability for a
pseudo-holomorphic section problem
\[
  W\longrightarrow C\longrightarrow B
\]
over a derived smooth stack, with $C\to B$ a proper family of real
two-dimensional domains and with the relevant complex structures. It also
identifies the derived-to-topological comparison for the solution stack
\cite[Theorem~5.1]{Pardon_Representability_Elliptic_Fredholm}.

The printed statement says that $B\to C$ is proper of relative dimension two,
although the displayed composable maps and the proof require $C\to B$. This
is best treated as a typographical reversal, not as a new mathematical
hypothesis.

\subsection{Related literature}
\label[subsection]{sec:Derived_Charts_Globalization_Related_Literature}

Steffens's Theorem~3.4.3 is the primary source for the completed
representability proof. The local chart and globalization occupy the final
part of its proof, while Proposition~2.1.48 and Corollary~2.1.49 supply the
open-atlas step
\cite[Theorem~3.4.3, Proposition~2.1.48, and
Corollary~2.1.49]{Steffens_Representability_Elliptic_Moduli}.

Pardon's published proceedings article supplies the comparison through
Theorem~5.1, Proposition~5.2, and Lemma~5.3
\cite[Section~5]{Pardon_Representability_Elliptic_Fredholm}. The July 2026
manuscript is useful for terminology and proof order, but its work-in-progress
status must be retained whenever it is cited
\cite[Sections~P.5 and N.6.4]{Pardon_Logarithmic_Derived_Moduli}.

% !TEX root = ../main.tex

\chapter{Pseudo-holomorphic curves as a derived moduli problem}
\label[chapter]{chap:Pseudo_Holomorphic_Curves_Derived_Moduli}

Let $(\Sigma,j)$ be a closed Riemann surface and $(Z,J)$ an almost complex
manifold. A map $u\colon\Sigma\to Z$ is pseudo-holomorphic when its derivative
is complex linear. This first-order equation provides an application of the
elliptic representability argument of the preceding three chapters.

We first write the equation as a section of its bundle of
complex-antilinear forms and put it into the local coordinates required
by the representability proof. The main calculation is its principal
symbol. Ellipticity then gives a quasi-smooth derived scheme of maps
from a fixed domain, with tangent complex given by the linearized
Cauchy--Riemann equation. Constant maps provide an explicit example
of its deformation and obstruction groups.

Allowing the domain to vary produces a derived moduli stack over the
smooth stack of Riemann surfaces. This relative formulation keeps
deformations of the map distinct from variations and automorphisms of
its domain. We conclude by explaining what further input would be
needed for nodal compactifications and enumerative invariants.

\section{The Cauchy--Riemann equation}
\label[section]{sec:Cauchy_Riemann_Equation_Fourteen}

The condition that a derivative be complex linear is expressed by
the vanishing of its complex-antilinear part. We recall this
decomposition and apply it pointwise to a smooth map.

An almost complex structure on a smooth manifold $Z$ is a vector-bundle
endomorphism
\[
  J\colon TZ\longrightarrow TZ
\]
satisfying $J^2=-\operatorname{id}_{TZ}$. Each tangent space $T_xZ$ thereby
becomes a complex vector space, with multiplication by $i$ given by $J_x$.
The adjective \emph{almost} records that these pointwise complex structures need
not arise from complex coordinates on $Z$.

A Riemann surface can be described similarly as a smooth real surface
$\Sigma$ equipped with an endomorphism $j\colon T\Sigma\to T\Sigma$ satisfying
$j^2=-1$. In real dimension two every almost complex structure is integrable,
but this fact will not be needed. The tangent-bundle formulation puts the
differential equation directly in view.

We begin with the relevant linear algebra. Let $(V,j)$ and $(W,J)$ be complex
vector spaces regarded as real vector spaces with chosen complex structures.
For a real-linear map $L\colon V\to W$, define
\begin{equation}
  \label{eq:Linear_Antilinear_Decomposition_Fourteen}
  L^{1,0}=\frac12(L-JLj),
  \qquad
  L^{0,1}=\frac12(L+JLj).
\end{equation}

\begin{lemma}[Complex-linear decomposition]
  \label[lemma]{lem:Complex_Linear_Decomposition_Fourteen}
  The map $L^{1,0}$ is complex linear, the map $L^{0,1}$ is complex
  antilinear, and
  \[
    L=L^{1,0}+L^{0,1}.
  \]
  Consequently, $L$ is complex linear if and only if $L^{0,1}=0$.
\end{lemma}

\begin{proof}
  Using $j^2=J^2=-1$ gives
  \[
    L^{1,0}j=\tfrac12(Lj+JL)=JL^{1,0},\qquad
    L^{0,1}j=\tfrac12(Lj-JL)=-JL^{0,1}.
  \]
  The sum of the two parts is $L$. Thus $L^{0,1}=0$ implies that
  $L$ is complex linear; conversely, $Lj=JL$ gives $JLj=-L$ and
  hence $L^{0,1}=0$.
\end{proof}

For $p\in\Sigma$ and $x\in Z$, write
\[
  \Hom^{0,1}_{j,J}(T_p\Sigma,T_xZ)
\]
for the real-linear maps $A\colon T_p\Sigma\to T_xZ$ satisfying
$A\circ j=-J\circ A$. These vector spaces form a vector bundle over
$\Sigma\times Z$. Restriction along the graph of a map $u$ gives the bundle
of complex-antilinear one-forms with values in $u^*TZ$, whose sections are
denoted
\[
  \Omega^{0,1}(\Sigma,u^*TZ)
  =
  \Gamma\bigl(
    \Hom^{0,1}_{j,J}(T\Sigma,u^*TZ)
  \bigr).
\]

\begin{definition}[Nonlinear Cauchy--Riemann operator]
  \label[definition]{def:Nonlinear_Cauchy_Riemann_Operator_Fourteen}
  For a smooth map $u\colon\Sigma\to Z$, define
  \begin{equation}
    \label{eq:Nonlinear_Cauchy_Riemann_Operator_Fourteen}
    \bar\partial_Ju
    =
    \frac12\bigl(du+J\circ du\circ j\bigr)
    \in\Omega^{0,1}(\Sigma,u^*TZ).
  \end{equation}
  The map $u$ is \emph{$J$-pseudo-holomorphic}, or
  \emph{$J$-holomorphic}, if $\bar\partial_Ju=0$.
\end{definition}

By \Cref{lem:Complex_Linear_Decomposition_Fourteen}, the equation is
equivalent to
\[
  du\circ j=J\circ du.
\]
When $Z=\mathbb C^n$ with its standard complex structure, this is the ordinary
Cauchy--Riemann equation in local coordinates.

The dependence on the value $u(p)$ makes this a nonlinear equation in
general: varying $u$ changes the endomorphism $J_{u(p)}$ as well as $du$.

The equation is invariant under reparametrization. If
$\varphi\colon(\Sigma',j')\to(\Sigma,j)$ is biholomorphic, then
\[
  \bar\partial_J(u\circ\varphi)=0
\]
whenever $\bar\partial_Ju=0$. Indeed,
$d\varphi\circ j'=j\circ d\varphi$, and substitution in
\Cref{eq:Nonlinear_Cauchy_Riemann_Operator_Fourteen} proves the claim. This
elementary invariance is the first indication that a moduli object for varying
domains must retain automorphisms.

Steffens formulates the application for a symplectic manifold $(Z,\omega)$
equipped with an $\omega$-tame almost complex structure, meaning
\[
  \omega(v,Jv)>0
\]
for every nonzero tangent vector $v$. The symplectic form is fundamental for
energy estimates and compactness, but the equation, the symbol calculation,
and the local representability argument below use only $J$. We will state the
representability theorem first for an almost complex target and return to the
role of $\omega$ at the end.

\section{The fixed-domain solution stack}
\label[section]{sec:Cauchy_Riemann_Local_Coordinates}

The output of the Cauchy--Riemann operator lies in a bundle depending
on the map. First jets define the global section and its derived zero
locus. Near any chosen map, trivializing the equation bundle gives a
single operator between sections of fixed vector bundles.

The equation depends only on the value and first derivative of $u$. A first
jet at $p\in\Sigma$ can be represented by a pair
\[
  (x,L),
  \qquad
  x\in Z,
  \quad
  L\colon T_p\Sigma\longrightarrow T_xZ.
\]
The complex-antilinear projection defines a smooth map on first jets,
\begin{equation}
  \label{eq:Cauchy_Riemann_Jet_Map_Fourteen}
  (x,L)
  \longmapsto
  \left(x,\frac12(L+J_xLj_p)\right).
\end{equation}
Thus pseudo-holomorphicity is a first-order differential equation in the sense
of \Cref{chap:Section_Stacks_Jets_Solution_Stacks}.

For the moment keep the domain fixed. The smooth mapping object
\[
  \mathcal B=C^\infty(\Sigma,Z)
\]
has tangent space
\[
  T_u\mathcal B=\Gamma(\Sigma,u^*TZ)
\]
at $u$. Over $\mathcal B$ there is a bundle $\mathcal F\to\mathcal B$ whose
fiber is
\[
  \mathcal F_u=\Omega^{0,1}(\Sigma,u^*TZ).
\]
The nonlinear Cauchy--Riemann operator is a section
\[
  \bar\partial_J\colon\mathcal B\longrightarrow\mathcal F.
\]
To take its derived zero locus, regard these as derived section stacks:
\[
 \mathcal B=\Sec_\Sigma(\Sigma\times Z),\qquad
 \mathcal F=\Sec_\Sigma(H^{0,1}),
\]
where $H^{0,1}\to\Sigma\times Z$ is the bundle defined above and the second
section stack uses its composite map to $\Sigma$.
The comparison in \Cref{chap:Elliptic_Representability_Smooth_Derived_Bridge} identifies their smooth-stack restrictions
with the indicated convenient manifolds.
The derived solution object is
$\mathcal B\times_{\mathcal F}^{\mathrm{der}}\mathcal B$, using
$\bar\partial_J$ and the zero section.

Fix a map $u$. Exponential coordinates identify nearby maps with a
neighbourhood $Q$ of zero in $\Gamma(\Sigma;F)$, where $F=u^*TZ$.
Choose a connection preserving $J$. Parallel transport along the short
exponential paths gives complex-linear isomorphisms
$\tau_v\colon u^*TZ\to(\exp_u v)^*TZ$.
These identify the nearby bundles of complex-antilinear forms with
\[
 E=\Hom^{0,1}_{j,J}(T\Sigma,u^*TZ).
\]
The transformed equation is
\begin{equation}\label{eq:Local_Cauchy_Riemann_Operator_Fourteen}
 P(v)=\tau_v^{-1}\bar\partial_J(\exp_u v)
 \colon Q\longrightarrow\Gamma(\Sigma;E).
\end{equation}
All coordinate changes depend pointwise on $v$. Thus $P$ is still a
first-order differential operator between sections of fixed vector bundles.
The local section-stack comparison identifies its derived zero fibre with
an open part of the original solution stack.

This is the form needed for the argument of \Cref{chap:Elliptic_Equations_Kuranishi_Models,chap:Elliptic_Representability_Smooth_Derived_Bridge,chap:Elliptic_Representability_Derived_Charts}.
The first-jet presentation in Steffens's Construction~4.0.1 requires a
correction to its output bundle and a local reduction before applying the
fixed-order ellipticity criterion. The correction and the local reduction are discussed in
\Cref{sec:Pseudoholomorphic_Summary_Supplements_Fourteen}.

\section{Linearization and ellipticity}
\label[section]{sec:Cauchy_Riemann_Linearization_Fourteen}

We differentiate at a solution and calculate the principal symbol of
the resulting operator. The symbol is invertible on every nonzero
covector, so the elliptic theory of the preceding chapter applies.
The index formula and the constant-map example describe the resulting
finite-dimensional linear data.

Let $u$ be pseudo-holomorphic. Because $\bar\partial_Ju=0$, the vertical
derivative of the section at $u$ is intrinsically defined:
\[
  D_u\colon
  \Gamma(\Sigma,u^*TZ)
  \longrightarrow
  \Omega^{0,1}(\Sigma,u^*TZ).
\]
Choose a torsion-free connection $\nabla$ on $TZ$ in order to write it.

\begin{proposition}[Linearized Cauchy--Riemann operator]
  \label[proposition]{prop:Linearized_Cauchy_Riemann_Operator_Fourteen}
  For $\xi\in\Gamma(\Sigma,u^*TZ)$, the linearization is
  \begin{equation}
    \label{eq:Linearized_Cauchy_Riemann_Operator_Fourteen}
    D_u\xi
    =
    \frac12\bigl(
      \nabla\xi+J\circ\nabla\xi\circ j
    \bigr)
    +
    \frac12(\nabla_\xi J)\circ du\circ j.
  \end{equation}
  In particular, $D_u$ is a real Cauchy--Riemann type operator: its
  first-order part is a Cauchy--Riemann operator, while the second summand is
  zeroth order.
\end{proposition}

\begin{proof}
  Choose a smooth family $u_t$ with $u_0=u$ and
  $\frac{d}{dt}|_{t=0}u_t=\xi$. Covariant differentiation gives
  \[
    \left.\frac{D}{dt}\right|_{t=0}du_t=\nabla\xi.
  \]
  Differentiating the target almost complex structure gives
  \[
    \left.\frac{D}{dt}\right|_{t=0}J(u_t)=\nabla_\xi J.
  \]
  The complex structure $j$ on the fixed domain does not vary. Applying the
  product rule to
  \Cref{eq:Nonlinear_Cauchy_Riemann_Operator_Fourteen} proves
  \Cref{eq:Linearized_Cauchy_Riemann_Operator_Fourteen}.

  The connection is used only to compare the bundles $u_t^*TZ$ for varying
  $t$. The vertical derivative of a bundle section at a zero is independent
  of this choice, so $D_u$ is intrinsic.
\end{proof}

The first summand differentiates $\xi$. The term involving $\nabla_\xi J$
contains no derivative of $\xi$, and hence contributes nothing to the
principal symbol. Consequently the principal symbol is the same as that of the first summand.

For a first-order differential operator, the principal symbol is obtained by
retaining only the coefficient of the first derivative. Fix $p\in\Sigma$,
$0\neq\zeta\in T_p^*\Sigma$, and $v\in T_{u(p)}Z$.

\begin{theorem}[Ellipticity of the Cauchy--Riemann operator]
  \label[theorem]{thm:Cauchy_Riemann_Ellipticity_Fourteen}
  The principal symbol of $D_u$ is
  \begin{equation}
    \label{eq:Cauchy_Riemann_Principal_Symbol_Fourteen}
    \sigma_\zeta(D_u)(v)(Y)
    =
    \frac12\bigl(
      \zeta(Y)v+\zeta(jY)Jv
    \bigr).
  \end{equation}
  For every nonzero $\zeta$, this is an isomorphism
  \[
    T_{u(p)}Z
    \iso
    \Hom^{0,1}_{j,J}(T_p\Sigma,T_{u(p)}Z)
  \]
  of real vector spaces. Consequently, $D_u$ is elliptic.
\end{theorem}

\begin{proof}
  The first-order part of $D_u$ is
  \[
    \xi\longmapsto
    \frac12(\nabla\xi+J\circ\nabla\xi\circ j).
  \]
  Replacing $\nabla_Y\xi$ by $\zeta(Y)v$ gives
  \Cref{eq:Cauchy_Riemann_Principal_Symbol_Fourteen}.

  First observe that the resulting map is complex antilinear in $Y$. Indeed,
  \begin{align*}
    \sigma_\zeta(D_u)(v)(jY)
    &=\frac12\bigl(
       \zeta(jY)v-\zeta(Y)Jv
      \bigr)\\
    &=-J\sigma_\zeta(D_u)(v)(Y).
  \end{align*}

  Choose a $j$-compatible inner product on $T_p\Sigma$ and let $Y$ be the
  metric dual of $\zeta$. Then
  \[
    \zeta(Y)>0,
    \qquad
    \zeta(jY)=0.
  \]
  If $\sigma_\zeta(D_u)(v)=0$, evaluation at $Y$ gives
  \[
    0=\frac12\zeta(Y)v,
  \]
  and therefore $v=0$. The symbol is injective. Since $T_p\Sigma$ has complex
  dimension one, a complex-antilinear map out of $T_p\Sigma$ is determined by
  its value on one nonzero real direction. Hence the source and target of the
  symbol have the same real dimension, so the symbol is an isomorphism.
\end{proof}

This proof is pointwise. It uses neither a symplectic form on $Z$ nor
integrability of $J$. It also works uniformly when the domain and the map vary
in smooth families.

Since $\Sigma$ is closed, after completing in Sobolev norms,
the elliptic Fredholm theorem applies.

\begin{corollary}[Fredholm deformation operator]
  \label[corollary]{cor:Cauchy_Riemann_Fredholm_Fourteen}
  For sufficiently large $k$, the extension
  \[
    D_u\colon
    H^{k+1}(\Sigma,u^*TZ)
    \longrightarrow
    H^k\Omega^{0,1}(\Sigma,u^*TZ)
  \]
  is Fredholm. Its kernel and cokernel are finite-dimensional and are
  represented by smooth sections.
\end{corollary}

\begin{proof}
  Fredholmness follows from
  \Cref{thm:Cauchy_Riemann_Ellipticity_Fourteen} and the elliptic Fredholm
  theorem used in
  \Cref{thm:Elliptic_Fredholm_Regularity}.
  Elliptic regularity identifies the Sobolev kernel and cokernel with their
  smooth counterparts.
\end{proof}

If $Z$ has complex dimension $n$, Riemann--Roch gives the real index
\begin{equation}
  \label{eq:Cauchy_Riemann_Index_Fourteen}
  \operatorname{ind}_{\mathbb R}(D_u)
  =
  n\chi(\Sigma)
  +2\left\langle c_1(u^*TZ),[\Sigma]\right\rangle.
\end{equation}
This is the Riemann--Roch formula
\cite[Theorem~3.22]{Wendl_Lectures_Holomorphic_Curves}.
It is the virtual dimension of the moduli problem
relative to the fixed domain. For an unparametrized problem with varying
domains, the deformation theory of the domain must also be included.

\begin{example}[Constant maps]
Fix $x\in Z$ and let $u_x\colon\Sigma\to Z$ be the constant map with value
$x$. Since $du_x=0$, it is pseudo-holomorphic, and the zeroth-order term in
\Cref{eq:Linearized_Cauchy_Riemann_Operator_Fourteen} vanishes. Thus
$D_{u_x}$ is the Dolbeault operator on the trivial complex vector bundle with
fiber $T_xZ$:
\[
  \bar\partial\colon
  \Omega^0(\Sigma)\otimes_{\mathbb C}T_xZ
  \longrightarrow
  \Omega^{0,1}(\Sigma)\otimes_{\mathbb C}T_xZ.
\]
For connected $\Sigma$ this gives
\begin{equation}
  \label{eq:Constant_Map_Deformation_Obstruction_Fourteen}
  \ker D_{u_x}\cong T_xZ,
  \qquad
  \operatorname{coker}D_{u_x}
  \cong
  H^{0,1}(\Sigma)\otimes_{\mathbb C}T_xZ.
\end{equation}

Indeed, holomorphic functions on a closed connected Riemann surface are
constant, giving the kernel. Dolbeault cohomology is, by definition, the
cokernel of the scalar $\bar\partial$ operator on a surface, giving the
second formula. If $\Sigma$ has genus $g$ and $\dim_{\mathbb C}Z=n$, these
groups have real dimensions $2n$ and $2ng$. The index is $2n(1-g)$.

For a sphere, the cokernel vanishes. For a positive-genus domain,
constant maps have nonzero obstruction space, although the visible family of
constant maps is the ordinary manifold $Z$.
The calculation concerns the full mapping problem at a constant map;
it does not identify the entire moduli object with the constant-map locus.
\end{example}

These kernel and cokernel calculations give the linear data.
The Kuranishi construction identifies their two-term complex with
the intrinsic derived tangent complex, as we now explain.

\section{Deformation theory with the domain fixed}
\label[section]{sec:Pseudo_Holomorphic_Deformation_Theory_Fourteen}

The local Kuranishi construction turns the elliptic deformation
complex into an intrinsic tangent complex. We use this identification
to interpret regularity, obstruction groups, and virtual dimension
while keeping the domain fixed.

Let
\[
  \mathcal M_{(\Sigma,j)}(Z,J)
\]
denote the derived solution stack for the fixed domain $(\Sigma,j)$. On
ordinary test manifolds, a $T$-point is a smooth map
\[
  u\colon T\times\Sigma\longrightarrow Z
\]
whose restriction to every slice is pseudo-holomorphic. On a derived test
object, the equation is imposed by a derived zero locus and hence retains its
infinitesimal intersection data.

The analytic deformation complex is the two-term complex with differential
$D_u$. Its identification with a derived tangent complex follows from the
finite-dimensional Kuranishi reduction.

\begin{theorem}[Fixed-domain tangent complex]
  \label[theorem]{thm:Fixed_Domain_Tangent_Complex_Fourteen}
  At a pseudo-holomorphic map $u\colon(\Sigma,j)\to(Z,J)$, the tangent complex
  of the fixed-domain derived solution stack is
  \begin{equation}
    \label{eq:Fixed_Domain_Tangent_Complex_Fourteen}
    \mathbb T_u\mathcal M_{(\Sigma,j)}(Z,J)
    \simeq
    \left[
      \Gamma(\Sigma,u^*TZ)
      \xrightarrow{D_u}
      \Omega^{0,1}(\Sigma,u^*TZ)
    \right]
  \end{equation}
  in homological degrees $0$ and $-1$.
\end{theorem}

\begin{proof}
  Use the local operator of
  \Cref{sec:Cauchy_Riemann_Local_Coordinates}.
  Its derivative at the chosen zero is $D_u$, which is elliptic.
  The construction of \Cref{chap:Elliptic_Representability_Derived_Charts} gives a finite-dimensional chart
  $\kappa\colon Z_U\to V$ and a quasi-isomorphism
  \[
    [T_uZ_U\longrightarrow V]\longrightarrow
    [\Gamma(\Sigma,u^*TZ)\xrightarrow{D_u}
      \Omega^{0,1}(\Sigma,u^*TZ)].
  \]
  The left side is the tangent complex of the finite-dimensional derived
  zero locus, by the calculation in \Cref{chap:Stable_Linearization_And_Tangent_Complexes}. This proves the formula
  without applying a finite-dimensional cotangent theorem directly to
  the infinite-dimensional mapping object.
\end{proof}

Consequently,
\begin{align*}
  H_0\mathbb T_u
  &\cong\ker D_u,\\
  H_{-1}\mathbb T_u
  &\cong\operatorname{coker}D_u.
\end{align*}
The first group is the space of infinitesimal deformations of the map with the
domain fixed. The second is its obstruction space. Both are finite-dimensional
by \Cref{cor:Cauchy_Riemann_Fredholm_Fourteen}.

\begin{definition}[Regular pseudo-holomorphic map]
  \label[definition]{def:Regular_Pseudo_Holomorphic_Map_Fourteen}
  A pseudo-holomorphic map $u$ is \emph{regular} if $D_u$ is surjective.
\end{definition}

At a regular map the obstruction group vanishes. The Banach implicit function
theorem then identifies the nearby ordinary solution space with a smooth
manifold of dimension $\operatorname{ind}_{\mathbb R}(D_u)$. Derivedly, the
tangent complex is concentrated in degree $0$.

If the cokernel is nonzero, it receives obstructions to extending an
infinitesimal deformation to an actual family. A particular obstruction
class can still vanish. The Kuranishi map describes these nonlinear
conditions; its derivative alone need not determine them.

The virtual dimension at $u$ is
\[
  \dim H_0\mathbb T_u-\dim H_{-1}\mathbb T_u
  =\operatorname{ind}_{\mathbb R}(D_u),
\]
which is given by \Cref{eq:Cauchy_Riemann_Index_Fourteen}.

\section{Varying domains and relative representability}
\label[section]{sec:Pseudo_Holomorphic_Families_Representability_Fourteen}

We now retain the complex structure and automorphisms of the domain
as part of the moduli problem. The stack of compact Riemann surfaces
supplies a universal proper family. The relative Cauchy--Riemann
section defines the solution stack over this base, and the local
elliptic argument proves its relative representability.

The fixed-domain problem ignores reparametrizations and variations of the
complex structure. To retain both, let
\[
  S=\operatorname{Surf}_{\mathbb C}
\]
denote the smooth stack of compact Riemann surfaces. For a test manifold $T$,
an object of $S(T)$ is a proper submersion
\[
  \pi\colon C\longrightarrow T
\]
with compact surface fibers, together with a smooth complex structure on the
vertical tangent bundle $T(C/T)$. Morphisms are Cartesian diagrams whose maps
on fibers are biholomorphic.

The stack retains smooth dependence on parameters and the automorphisms of
each surface. For example, the Riemann sphere has automorphism group
$\operatorname{PSL}_2(\mathbb C)$, which appears as isotropy at its point of $S$.

The stack carries a universal family
\begin{equation}
  \label{eq:Universal_Riemann_Surface_Family_Fourteen}
  M=\widetilde{\operatorname{Surf}}_{\mathbb C}
  \longrightarrow
  S=\operatorname{Surf}_{\mathbb C}.
\end{equation}
Its fiber over $(\Sigma,j)$ is the surface $\Sigma$ itself. In particular,
$M\to S$ is a proper family of smooth compact manifolds, which is exactly the
domain hypothesis in the elliptic representability theorem.

Kodaira--Spencer theory implies that $S$ is a smooth Artin
$\Cinfty$-stack. We quote this result, as Steffens does in
\cite[Remark~4.0.3]{Steffens_Representability_Elliptic_Moduli}. Its role will
be to turn relative representability over $S$ into an absolute derived Artin
statement.

The Artin terminology allows smooth submersive atlases, with representable
overlaps, in place of the open atlases for schemes. Such a stack can have
positive-dimensional groups of automorphisms. In the derived version the
atlas objects are derived schemes. Pulling a smooth atlas of $S$ back to a
relatively representable moduli stack gives an atlas of this kind.
Here a derived Artin stack is called quasi-smooth when it admits a smooth
atlas by quasi-smooth derived schemes.

Put $Y=Z\times M$ over $M$. A section of $Y$ over a family $C\to T$ is a
smooth family of maps $C\to Z$. The corresponding section stack
\[
  \mathcal B
  =\Sec_{M/S}(Y)
\]
parametrizes the unknown maps.

Over $Y$ there is a finite-rank vector bundle $H^{0,1}\to Y$ whose fiber at
$(c,x)$ is
\[
  \Hom^{0,1}_{j_c,J_x}
  \bigl(T_c(M/S),T_xZ\bigr).
\]
Pulling this bundle back along a map-section and taking relative sections
produces a bundle $\mathcal F\to\mathcal B$. Its fiber at a family $u$ is the
family of spaces
\[
  \Omega^{0,1}(C_t,u_t^*TZ).
\]

The complex-antilinear projection of the relative first jet defines a section
\[
  \operatorname{CR}\colon\mathcal B\longrightarrow\mathcal F,
  \qquad
  u\longmapsto\bar\partial_Ju.
\]
The construction uses only the vertical tangent bundle, the fiberwise complex
structure, and $J$. It is therefore natural under Cartesian base change of
families of surfaces.

\begin{definition}[Derived moduli stack of pseudo-holomorphic maps]
  \label[definition]{def:Derived_Moduli_Pseudoholomorphic_Maps_Fourteen}
  The \emph{derived moduli stack of pseudo-holomorphic maps}, denoted
  $\mathcal M(Z,J)$, is the derived zero locus of the Cauchy--Riemann section
  over $\mathcal B$.
\end{definition}

It comes with a natural morphism
\[
  \mathcal M(Z,J)
  \longrightarrow
  \operatorname{Surf}_{\mathbb C}.
\]
Over an ordinary test manifold $T\to S$ classifying $C\to T$, its ordinary
points are precisely smooth maps $u\colon C\to Z$ whose restrictions to the
fibers are pseudo-holomorphic. The derived zero locus retains the higher
intersection data of this equation.

\begin{theorem}[Pseudo-holomorphic representability]
  \label[theorem]{thm:Pseudoholomorphic_Representability_Fourteen}
  Let $(Z,J)$ be an almost complex manifold. The morphism
  \begin{equation}
    \label{eq:Pseudoholomorphic_Representability_Morphism_Fourteen}
    \mathcal M(Z,J)
    \longrightarrow
    \operatorname{Surf}_{\mathbb C}
  \end{equation}
  is representable by quasi-smooth derived $\Cinfty$-schemes locally of finite
  presentation. Consequently, $\mathcal M(Z,J)$ is a quasi-smooth derived
  Artin $\Cinfty$-stack.
\end{theorem}

\begin{proof}
  Representability is local on the target. Pull back along a map
  $T\to\operatorname{Surf}_{\mathbb C}$ from a manifold. It classifies a
  smooth proper family of Riemann surfaces $C\to T$. Locally on $T$, Ehresmann
  triviality identifies the underlying family with
  $T'\times\Sigma\to T'$, although the fiberwise complex structure may still
  vary with the parameter.

  Fix a solution $u$ in this family. The local section-stack construction from
  \Cref{sec:Cauchy_Riemann_Local_Coordinates} identifies a neighborhood of $u$ with
  $T'\times Q$, where $Q$ is an open subset of a smooth section space. After
  trivializing the nearby equation bundles, the family Cauchy--Riemann section
  becomes a smooth family of nonlinear first-order operators
  \begin{equation}
    \label{eq:Family_Local_Cauchy_Riemann_Operator_Fourteen}
    P\colon
    T'\times Q
    \longrightarrow
    T'\times\Gamma(\Sigma;E).
  \end{equation}
  Its vertical linearization at $u$ is $D_u$, which is elliptic by
  \Cref{thm:Cauchy_Riemann_Ellipticity_Fourteen}.

  The local Kuranishi construction of \Cref{chap:Elliptic_Representability_Smooth_Derived_Bridge,chap:Elliptic_Representability_Derived_Charts} applies to
  this operator. It produces a quasi-smooth derived
  $\Cinfty$-scheme locally of finite presentation representing a neighborhood
  of $u$ in the solution stack. These neighborhoods cover the pullback over
  $T$, and the open-atlas criterion glues them. This proves relative
  representability.

  Finally, $\operatorname{Surf}_{\mathbb C}$ is a smooth Artin
  $\Cinfty$-stack. Relative quasi-smooth representability over this base makes
  the total solution stack a quasi-smooth derived Artin $\Cinfty$-stack.
\end{proof}

Steffens states this conclusion in the symplectic tame setting in
\cite[Construction~4.0.1 and Remark~4.0.3]{Steffens_Representability_Elliptic_Moduli}.
The proof above uses the local zero-section form of the representability
argument, for the reason explained in
\Cref{warn:Steffens_Five_Tuple_Fourteen}. It also shows that the local
representability statement depends only on $J$.

In particular, given a manifold $T$ and a
specified family $C\to T$ of Riemann surfaces, the pullback
\[
  T\times_{\operatorname{Surf}_{\mathbb C}}\mathcal M(Z,J)
\]
is represented by a quasi-smooth derived $\Cinfty$-scheme locally of finite
presentation over $T$. Its relative tangent complex at a map $u$ is
\begin{equation}
  \label{eq:Relative_Tangent_Complex_Final_Fourteen}
  \mathbb T_u\bigl(
    \mathcal M(Z,J)/\operatorname{Surf}_{\mathbb C}
  \bigr)
  \simeq
  \left[
    \Gamma(\Sigma,u^*TZ)
    \xrightarrow{D_u}
    \Omega^{0,1}(\Sigma,u^*TZ)
  \right].
\end{equation}
The absolute tangent complex also contains deformations and infinitesimal
automorphisms of the domain. The corresponding fiber sequence is recorded in
\Cref{eq:Absolute_Tangent_Fiber_Sequence_Fourteen}.

\section{Compactification and enumerative questions}
\label[section]{sec:Pseudoholomorphic_Boundary_Fourteen}

Properness of the domain family ensures compact fibres on which elliptic
analysis applies. Compactness of the solution object is another issue.
Even with compact target, sequences of maps can develop bubbles, and smooth
domains can degenerate to nodal curves. These limits require an enlarged
moduli problem, typically stable nodal maps with fixed genus, marked points,
and homology class.

Representing that compactification requires analysis near its nodal objects,
including gluing. An enumerative application then needs a properness theorem
and appropriate orientations and virtual classes. \Cref{chap:Derived_Bordism_And_Fundamental_Classes} supplies a
bordism class for a compact quasi-smooth derived manifold. The total
moduli stack here may have isotropy, so a virtual-class construction for
stacks requires additional theory.

When $J$ is tamed by a symplectic form $\omega$, the area
$\int_\Sigma u^*\omega$ is positive for a nonconstant pseudo-holomorphic map
from a connected domain, and depends only on its homology class.
With the associated $J$-invariant metric, it equals the energy of the map.
Thus fixing the class gives an
energy bound, an essential input to compactness. The bound allows bubbling
and explains why nodal maps enter the compactification.

\section{Further reading}
\label[section]{sec:Pseudoholomorphic_Summary_Supplements_Fourteen}

The following discussion compares the local equation used in the proof
with the symmetric presentation in the source, records the absolute
tangent complex when the domain varies, and describes further choices
of geometric data.

\begin{warning}[The printed five-tuple]
  \label[warning]{warn:Steffens_Five_Tuple_Fourteen}
  In Steffens's Construction~4.0.1, the common output of the displayed
  Cauchy--Riemann and zero maps is the full relative first-jet bundle. Taken
  literally, its linearized matching complex has complementary
  complex-linear jet directions in degree minus one, in addition to the
  Cauchy--Riemann Fredholm complex.

  Replacing the common output by the bundle of complex-antilinear maps
  removes those extra obstruction directions. The resulting symmetric matching operator still
  has a zeroth-order component, equating the underlying map values, and a
  first-order Cauchy--Riemann component. It is therefore not elliptic in the
  fixed-order sense of \cite[Definition~3.4.2]{Steffens_Representability_Elliptic_Moduli}.
  The general theorem cannot be invoked verbatim from that five-tuple.
\end{warning}

The derived intersection of the Cauchy--Riemann section with the zero section
can be displayed as
\begin{equation}
  \label{eq:Symmetric_Cauchy_Riemann_Pullback_Fourteen}
  \begin{tikzcd}[column sep=large]
    \mathcal M(Z,J)
      \rar \dar \drar[pullback]
      & \mathcal B \dar["\operatorname{CR}"] \\
    \mathcal B
      \rar["0"']
      & \mathcal F.
  \end{tikzcd}
\end{equation}
At a solution $u$, put
\[
  A=\Gamma(\Sigma,u^*TZ),
  \qquad
  B=\Omega^{0,1}(\Sigma,u^*TZ).
\]
Linearizing this pullback presentation gives
\begin{equation}
  \label{eq:Symmetric_Tangent_Complex_Fourteen}
  \left[
    A\oplus A
    \xrightarrow{\delta}
    A\oplus B
  \right],
  \qquad
  \delta(\xi_1,\xi_2)
  =
  (\xi_1-\xi_2,D_u\xi_1).
\end{equation}
The first component differentiates equality of the underlying maps, and the
second differentiates the Cauchy--Riemann equation.

\begin{lemma}[Cancellation of the repeated field]
  \label[lemma]{lem:Repeated_Field_Cancellation_Fourteen}
  There is a chain isomorphism
  \[
    [A\oplus A\xrightarrow{\delta}A\oplus B]
    \iso
    [A\xrightarrow{D_u}B]
    \oplus
    [A\xrightarrow{\operatorname{id}}A].
  \]
  Consequently,
  \Cref{eq:Symmetric_Tangent_Complex_Fourteen} is quasi-isomorphic to the
  Cauchy--Riemann deformation complex.
\end{lemma}

\begin{proof}
  On the source use the coordinates
  \[
    (\xi_1,\xi_2)
    \longmapsto
    (\xi_1,\xi_1-\xi_2).
  \]
  Reorder the target $A\oplus B$ as $B\oplus A$. In these coordinates the
  differential is
  \[
    (\xi,r)\longmapsto(D_u\xi,r),
  \]
  which is the displayed direct sum. The identity complex is contractible.
\end{proof}

The cancellation recovers the Cauchy--Riemann deformation complex.
Its identification with the intrinsic tangent complex uses the local
Kuranishi reduction proved above. The first-order symbol of the displayed
matching operator is still $(\xi_1,\xi_2)\mapsto(0,\sigma(D_u)\xi_1)$,
which is not invertible. Passing to the single local equation removes the
repeated map variable before applying elliptic analysis.

The complex in \Cref{eq:Relative_Tangent_Complex_Final_Fourteen} holds the
domain and its complex structure fixed. The absolute tangent complex also
contains the deformation theory of $(\Sigma,j)$. It fits into the relative
tangent fiber sequence
\begin{equation}
  \label{eq:Absolute_Tangent_Fiber_Sequence_Fourteen}
  \mathbb T_u\bigl(
    \mathcal M(Z,J)/\operatorname{Surf}_{\mathbb C}
  \bigr)
  \longrightarrow
  \mathbb T_u\mathcal M(Z,J)
  \longrightarrow
  \mathbb T_{(\Sigma,j)}\operatorname{Surf}_{\mathbb C}.
\end{equation}
The last term records deformations and infinitesimal automorphisms of the
domain. For a smooth Riemann surface, these are measured by
$H^1(\Sigma,T\Sigma)$ and $H^0(\Sigma,T\Sigma)$, respectively,
placed in tangent degrees $0$ and $1$. Thus passing from the relative
complex to the absolute one includes both changes of complex structure
and reparametrizations.
In this example the absolute tangent complex has amplitude in $[-1,1]$;
the degree $1$ term accounts for infinitesimal automorphisms. This is
compatible with quasi-smoothness of the Artin stack as defined through
its smooth atlases.

This distinction prevents an ambiguity in expected-dimension formulas. The
index in \Cref{eq:Cauchy_Riemann_Index_Fourteen} is the relative virtual
dimension for a fixed domain. An unparametrized moduli problem with varying
domains must also account for
$\mathbb T_{(\Sigma,j)}\operatorname{Surf}_{\mathbb C}$.

For connected genus-$g$ domains and no marked points, Riemann--Roch for
$T\Sigma$ gives the real virtual dimension $6g-6$ of the domain stack.
Adding this to the relative index gives
\[
 \operatorname{vdim}_{\R}\mathcal M(Z,J)
 =(n-3)(2-2g)+2\langle c_1(u^*TZ),[\Sigma]\rangle.
\]
In genus zero the contribution $-6$ accounts for the six-dimensional
reparametrization group. This is an expected dimension for the stack,
including its infinitesimal automorphisms.

Several variations fit the same representability argument.
\begin{enumerate}
  \item One may let $J$ vary in a smooth family by enlarging the base stack.
  \item One may allow a smooth family of almost complex target manifolds.
  \item One may add marked points by adding sections of the universal curve to
    the objects classified by the domain stack.
  \item One may formulate equivariant problems over quotient stacks, thereby
    retaining isotropy rather than passing to a naive orbit set.
\end{enumerate}
The first three variants are noted in
\cite[Remark~4.0.2]{Steffens_Representability_Elliptic_Moduli}. All still
concern smooth compact domains.

\begin{exercise}
At a constant map from a genus-$g$ surface to a complex $n$-fold, compute
the dimensions of the relative deformation and obstruction spaces.
Check the Riemann--Roch formula.
\end{exercise}

\begin{exercise}
Write the first-order principal symbol of the symmetric matching operator
in \Cref{eq:Symmetric_Tangent_Complex_Fourteen}.
Show why it fails to be invertible even though its complex is
quasi-isomorphic to the elliptic Cauchy--Riemann complex.
\end{exercise}

The primary derived-geometric source is
\citet{Steffens_Representability_Elliptic_Moduli}.
The construction in
\cite[Construction~4.0.1]{Steffens_Representability_Elliptic_Moduli}
defines the stack of smooth compact
Riemann surfaces, its universal family, the jet-level Cauchy--Riemann map, and
the solution stack. Varying data, marked points, and the derived Artin
conclusion are discussed in
\cite[Remarks~4.0.2--4.0.3]{Steffens_Representability_Elliptic_Moduli}.
The source correction needed for the printed five-tuple is recorded in
\Cref{warn:Steffens_Five_Tuple_Fourteen}.

For the analytic background, Wendl's
\emph{Lectures on Holomorphic Curves in Symplectic and Contact Geometry}
develops almost complex structures and pseudo-holomorphic maps in Section~2.1,
real Cauchy--Riemann type operators and their linearization in
Sections~2.3--2.4, regularity in Section~2.5, and Fredholm and index theory in
Sections~3.3--3.4
\cite{Wendl_Lectures_Holomorphic_Curves}. These sections support the
linearization, symbol, Fredholm, and Riemann--Roch statements used here.


\backmatter

\printbibliography

\end{document}
